% Additional notation transported from the night manuscript.
\providecommand{\LL}{\mathcal L}
\providecommand{\cT}{\mathcal T}
\providecommand{\HP}{\mathbb H\mathbb P}
\providecommand{\tr}{\operatorname{tr}}
\providecommand{\sgn}{\operatorname{sgn}}
\providecommand{\vir}{\mathrm{vir}}

\section{Geometric and characteristic conventions}\label{sec:V-known-inputs}
The index and signature polynomials used below are
\cite[equations (1.2)--(1.4), pp.~3--4]{ffEsoleKang}; the
rank-at-most-four integral classification is
\cite[Theorem 2.2.1, p.~61]{ffPetersSterk}.  In the almost
complex application, the Todd polynomial and the cotangent-twisted
polynomial are the corresponding elliptic-symbol indices, so their
integrality does not require an integrable complex structure.
For $c_1=c_3=0$ the substitution gives
\begin{equation}
\operatorname{Td}=\frac{3c_2^2-c_4}{720},\qquad
\chi^1=\frac{3c_2^2-31c_4}{180},\qquad
\sigma=\frac{3c_2^2+14c_4}{45}.
\label{maps:eq:F.1}
\end{equation}
In particular $\chi^1-4\operatorname{Td}=-c_4/6$, and
$\operatorname{Td}=(3\sigma-c_4)/48$.  For the rank-four middle
group and $c_4=6$ this forces $\sigma=2$, since
$\sigma\equiv2\pmod {16}$ and $|\sigma|\leq4$.
The cited lattice theorem then gives $I_{3,1}$ and the displayed
polynomials give the complete numerical record
\begin{equation}
 H^4(X;\mathbb Z)\cong I_{3,1},\qquad c_4=6,\qquad
 c_2^2=2,\qquad\sigma=2,\qquad\operatorname{Td}=0.
 \label{maps:eq:F.2}
\end{equation}
These substitutions are the form of the known inputs used in
the construction; the longer credited re-derivation remains
[S6\_09, \nolinkurl{maps:src:F.1.2}].

Mandal's connected-sum criterion is used only in its stated
orientation: $\ell\mathbb{HP}^2\#m(-\mathbb{HP}^2)$ has an
orientation-compatible almost complex structure exactly when
$\ell=2m+1$ and $m$ is odd
\cite[Corollary 1.4, p.~4, and its proof in \S5]{ffMandal}.
Its normalized tangent classes are those of
\cite[Lemma 2.1, p.~5, and Theorem 2.4, p.~7]{ffMandal}.
The credited application retained in the compendium is
[S6\_09, \nolinkurl{maps:src:F.1.3}].


\section{The marked cusp and preliminary lattice restrictions}



\begin{theorem}[Euler number of a Kuhn cusp]\label{maps:src:F.1.1}
\StatementRecord{maps:src:F.1.1}{proved}{Historical events: Claude Fable 5.2 (Anthropic), effort high, 16 November 2026: carried as background without fault found. GPT-6 Astra (OpenAI), ChatGPT Pro mode, effort Pro, S6\_09\_PAPERS, 17 November 2026: targeted proof check. Also checked in this writing. Pre-continuation writing events: GPT-6 Astra (OpenAI), ChatGPT Pro mode, effort Pro, S6\_10\_PAPERS, 18 November 2026 02:02:02.394029+00:00: targeted proof check. Continuation mathematical and source checks: GPT-6 Astra (OpenAI), ChatGPT Pro mode, effort Pro, S6\_10\_PAPERS continuation, 19 November 2026: computation or formal certificate; independent re-proof; used and rechecked in this run. Continuation metadata events: GPT-6 Astra (OpenAI), ChatGPT Pro mode, effort Pro, S6\_10\_PAPERS continuation, 19 November 2026: used and rechecked in this run; status-vocabulary reconciliation; metadata reconciliation only; no new mathematical check. \emph{Pre-continuation recorded proof scope.} GPT-6 Astra (OpenAI, ChatGPT Pro mode), effort Pro, 2026-11-17; GPT-6 Astra (OpenAI, ChatGPT Pro mode), S6\_06\_BLIND, effort Pro, 2026-11-17; Claude Fable 5.2 (Anthropic), effort high, 16 November 2026; also checked in this writing by GPT-6 Astra (OpenAI), ChatGPT Pro mode, effort Pro, 18 November 2026}{the hypotheses, notation and definitions stated here; the explicit constructions and calculations in the proof}


Let $g:M^8\to S^2$ be a proper triangulable map, a $T^6$-bundle
off finitely many values.  Suppose its unique non-quotient special
fibre is the rank-three Kuhn toroidal fibre with period-lattice
index $d$, and all its other special reduced fibres are free
finite quotients of $T^6$.  Then $\chi(M)=6d$.

\end{theorem}


\begin{proof}
\begin{enumerate}[label=\arabic*.]
\item A fundamental cube has six unimodular tetrahedra.  Taking the
quotient by an index-$d$ sublattice multiplies every cell count
by $d$, so the quotient cell counts are $d(1,7,12,6)$.
A cell of dimension $r$ corresponds to a torus stratum
$(\mathbb C^*)^{3-r}$.  Only the $6d$ point strata contribute
to Euler characteristic.  Euler integration over the regular
base gives zero since $\chi(T^6)=0$, and a finite free torus
quotient has Euler characteristic $0/|G|=0$.  Additivity now
gives the formula.  This checks the rank-three instance of the
localization mechanism in \cref{cdp:st:C2a} and \cref{cdp:st:C2c}.  It requires the specified
toroidal filling, not just its unipotent monodromy matrix.

\end{enumerate}
\end{proof}




\begin{proposition}[the diagonal obstruction]\label{maps:src:F.1.4}
\StatementRecord{maps:src:F.1.4}{proved}{Historical events: Claude Fable 5.2 (Anthropic), effort high, 16 November 2026: carried as background without fault found. GPT-6 Astra (OpenAI), ChatGPT Pro mode, effort Pro, S6\_09\_PAPERS, 17 November 2026: targeted proof check. Also checked in this writing. Pre-continuation writing events: GPT-6 Astra (OpenAI), ChatGPT Pro mode, effort Pro, S6\_10\_PAPERS, 18 November 2026 02:02:02.394029+00:00: targeted proof check. Continuation mathematical and source checks: GPT-6 Astra (OpenAI), ChatGPT Pro mode, effort Pro, S6\_10\_PAPERS continuation, 19 November 2026: independent re-proof; used and rechecked in this run. Continuation metadata events: GPT-6 Astra (OpenAI), ChatGPT Pro mode, effort Pro, S6\_10\_PAPERS continuation, 19 November 2026: used and rechecked in this run; status-vocabulary reconciliation; metadata reconciliation only; no new mathematical check. \emph{Pre-continuation recorded proof scope.} GPT-6 Astra (OpenAI, ChatGPT Pro mode), effort Pro, 2026-11-17; Claude Fable 5.2 (Anthropic), effort high, 16 November 2026; also checked in this writing by GPT-6 Astra (OpenAI), ChatGPT Pro mode, effort Pro, 18 November 2026}{the hypotheses, notation and definitions stated here; the explicit constructions and calculations in the proof}


If $A\in\operatorname{GL}_n(\mathbb Z)$ has no eigenvalue one,
every affine map $x\mapsto Ax+b$ of $\mathbb R^n/\mathbb Z^n$
has a fixed point.  Thus the diagonal elliptic finite monodromies
in Lead 7 admit no free affine lifts.

\end{proposition}


\begin{proof}
\begin{enumerate}[label=\arabic*.]
\item The real map $I-A$ is invertible.  A real solution of
$(I-A)x=b$ maps to a fixed point in the torus.  The diagonal
copies of the order-four and order-six elliptic rotations have
no eigenvalue one, so the conclusion applies to every translation
part, not only to integral or invariant choices.

\end{enumerate}
\end{proof}




\begin{proposition}[the cusp normal form]\label{maps:src:F.2.1}
\StatementRecord{maps:src:F.2.1}{proved}{Historical events: Claude Fable 5.2 (Anthropic), effort high, 16 November 2026: carried as background without fault found. GPT-6 Astra (OpenAI), ChatGPT Pro mode, effort Pro, S6\_09\_PAPERS, 17 November 2026: targeted proof check. Also checked in this writing. Pre-continuation writing events: GPT-6 Astra (OpenAI), ChatGPT Pro mode, effort Pro, S6\_10\_PAPERS, 18 November 2026 02:02:02.394029+00:00: targeted proof check. Continuation mathematical and source checks: GPT-6 Astra (OpenAI), ChatGPT Pro mode, effort Pro, S6\_10\_PAPERS continuation, 19 November 2026: independent re-proof; used and rechecked in this run. Continuation metadata events: GPT-6 Astra (OpenAI), ChatGPT Pro mode, effort Pro, S6\_10\_PAPERS continuation, 19 November 2026: used and rechecked in this run; status-vocabulary reconciliation; metadata reconciliation only; no new mathematical check. \emph{Pre-continuation recorded proof scope.} GPT-6 Astra (OpenAI, ChatGPT Pro mode), effort Pro, 2026-11-17; GPT-6 Astra (OpenAI, ChatGPT Pro mode), S6\_06\_BLIND, effort Pro, 2026-11-17; Claude Fable 5.2 (Anthropic), effort high, 16 November 2026; also checked in this writing by GPT-6 Astra (OpenAI), ChatGPT Pro mode, effort Pro, 18 November 2026}{the hypotheses, notation and definitions stated here; the explicit constructions and calculations in the proof}


Under the integral index-one condition in F.2, the cusp matrix is
conjugate in $\operatorname{GL}_6(\mathbb Z)$ to
\[
 U=\begin{pmatrix}I_3&I_3\\0&I_3\end{pmatrix}.
\]
Its integral centralizer consists exactly of
$\left(\begin{smallmatrix}P&Q\\0&P\end{smallmatrix}\right)$
with $P\in\operatorname{GL}_3(\mathbb Z)$ and
$Q\in\operatorname{Mat}_{3\times3}(\mathbb Z)$.

\end{proposition}


\begin{proof}
\begin{enumerate}[label=\arabic*.]
\item Let $N=M_0-I$.  The index-one condition says that the image
is the saturated kernel lattice and that
$\Lambda/\ker N\to\operatorname{im}N$ is an isomorphism.
Choose a basis $a_1,a_2,a_3$ of this image and lifts $b_i$ with
$Nb_i=a_i$.  Saturation shows that the six vectors form an
integral basis.  It gives the displayed matrix.  Commuting a
block matrix with $N$ forces its lower-left block to vanish
and its two diagonal blocks to agree.  Invertibility over
$\mathbb Z$ is then equivalent to unimodularity of that block.

\end{enumerate}
\end{proof}




\begin{proposition}[a complete finite-order filter]\label{maps:src:F.2.2}
\StatementRecord{maps:src:F.2.2}{proved}{Historical events: Claude Fable 5.2 (Anthropic), effort high, 16 November 2026: carried as background without fault found. GPT-6 Astra (OpenAI), ChatGPT Pro mode, effort Pro, S6\_09\_PAPERS, 17 November 2026: targeted proof check. Also checked in this writing. Pre-continuation writing events: GPT-6 Astra (OpenAI), ChatGPT Pro mode, effort Pro, S6\_10\_PAPERS, 18 November 2026 02:02:02.394029+00:00: targeted proof check. Continuation mathematical and source checks: GPT-6 Astra (OpenAI), ChatGPT Pro mode, effort Pro, S6\_10\_PAPERS continuation, 19 November 2026: computation or formal certificate; independent re-proof; used and rechecked in this run. Continuation metadata events: GPT-6 Astra (OpenAI), ChatGPT Pro mode, effort Pro, S6\_10\_PAPERS continuation, 19 November 2026: used and rechecked in this run; status-vocabulary reconciliation; metadata reconciliation only; no new mathematical check. \emph{Pre-continuation recorded proof scope.} GPT-6 Astra (OpenAI, ChatGPT Pro mode), S6\_06\_BLIND, effort Pro, 2026-11-17; Claude Fable 5.2 (Anthropic), effort high, 16 November 2026; also checked in this writing by GPT-6 Astra (OpenAI), ChatGPT Pro mode, effort Pro, 18 November 2026}{the hypotheses, notation and definitions stated here; \cref{maps:src:F.1.4}}


The linear order of a finite matrix in $\operatorname{GL}_6(\mathbb Z)$
which admits a free affine lift belongs to
\[
 \{1,2,3,4,5,6,8,10,12\}.
 \tag{F.3}\label{maps:eq:F.3}
\]
There are 38 possible degree-six cyclotomic characteristic
polynomials with an eigenvalue one, and 19 of them commute
with some real complex structure.  These are necessary types,
not 38 or 19 conjugacy classes of admissible tuples.

\end{proposition}


\begin{proof}
\begin{enumerate}[label=\arabic*.]
\item Proposition \ref{maps:src:F.1.4} requires a nonzero one-eigenspace.  The other
cyclotomic factors have total degree at most five.  The elementary
identity
\[
 \frac{\phi(n)^2}{n}=
 \prod_{p^a\parallel n}p^{a-2}(p-1)^2\geq\frac12
\]
implies $n\leq50$ when $\phi(n)\leq5$.
The complete enumeration through 50 gives
$n=1,2,3,4,5,6,8,10,12$.  Enumerating their multiplicities
with total degree six and at least one factor $T-1$ gives 38
polynomials.  A commuting complex structure exists on a real
finite-order module exactly when the multiplicities of the
real eigenvalues $1$ and $-1$ are even; nonreal pairs already
carry such a structure.  Imposing those two parity conditions
leaves 19.  Taking least common multiples of their root orders
gives exactly \eqref{maps:eq:F.3}, also for the 38-polynomial list.
The code prints every polynomial and every exterior invariant
rank.  This is a bounded exhaustive enumeration for the
stated filter.  It does not bound affine multiplicities when
the affine action has larger order than its linear part.

\end{enumerate}
\end{proof}




\begin{proposition}[Counterexample: determination by peripheral data alone]\label{maps:src:F.2.3}
\StatementRecord{maps:src:F.2.3}{proved}{Historical events: Claude Fable 5.2 (Anthropic), effort high, 16 November 2026: carried as background without fault found. GPT-6 Astra (OpenAI), ChatGPT Pro mode, effort Pro, S6\_09\_PAPERS, 17 November 2026: targeted proof check. Also checked in this writing. Pre-continuation writing events: GPT-6 Astra (OpenAI), ChatGPT Pro mode, effort Pro, S6\_10\_PAPERS, 18 November 2026 02:02:02.394029+00:00: targeted proof check. Continuation mathematical and source checks: GPT-6 Astra (OpenAI), ChatGPT Pro mode, effort Pro, S6\_10\_PAPERS continuation, 19 November 2026: computation or formal certificate; independent re-proof; targeted proof check; used and rechecked in this run. Continuation metadata events: GPT-6 Astra (OpenAI), ChatGPT Pro mode, effort Pro, S6\_10\_PAPERS continuation, 19 November 2026: used and rechecked in this run; status-vocabulary reconciliation; metadata reconciliation only; no new mathematical check. \emph{Pre-continuation recorded proof scope.} Claude Fable 5.2 (Anthropic), effort high, 16 November 2026; also checked in this writing by GPT-6 Astra (OpenAI), ChatGPT Pro mode, effort Pro, 18 November 2026}{the hypotheses, notation and definitions stated here; the explicit constructions and calculations in the proof}


Monodromy matrices and affine twist vectors alone do not
determine the homology of a compact glued manifold.  A cusp
filling, or an equivalent relative topological specification,
is additional data.

\end{proposition}


\begin{proof}
\begin{enumerate}[label=\arabic*.]
\item Use the constructed fourfold $Z$ and blow up a smooth point
of the smooth locus of its cusp fibre.  The operation takes
place entirely inside that filling.  It leaves every regular
fibre, every finite quotient model, every peripheral monodromy
and every affine twist unchanged.  The total space remains
smooth and compact, and its map to the same base remains
holomorphic.  Blowing up a point in complex dimension four
replaces a ball by the disc bundle of $\mathcal O_{\mathbb P^3}(-1)$.
Mayer--Vietoris along their common $S^7$ boundary adds a copy
of $\mathbb Z$ in degrees two, four and six.  Thus the Euler
number increases by three and the homology changes.  The
peripheral index-one unipotent matrix is unaffected.
The stronger specification of a particular height-one Kuhn
filling excludes this example.  With that extra specification
the question remains a substantive classification problem;
we do not claim to refute it.

\end{enumerate}
\end{proof}




\begin{theorem}[the self-dual parity obstruction]\label{maps:src:F.3.1}
\StatementRecord{maps:src:F.3.1}{proved}{Historical events: Claude Fable 5.2 (Anthropic), effort high, 16 November 2026: carried as background without fault found. GPT-6 Astra (OpenAI), ChatGPT Pro mode, effort Pro, S6\_09\_PAPERS, 17 November 2026: targeted proof check. Also checked in this writing. Pre-continuation writing events: GPT-6 Astra (OpenAI), ChatGPT Pro mode, effort Pro, S6\_10\_PAPERS, 18 November 2026 02:02:02.394029+00:00: targeted proof check. Continuation mathematical and source checks: GPT-6 Astra (OpenAI), ChatGPT Pro mode, effort Pro, S6\_10\_PAPERS continuation, 19 November 2026: independent re-proof; used and rechecked in this run. Continuation metadata events: GPT-6 Astra (OpenAI), ChatGPT Pro mode, effort Pro, S6\_10\_PAPERS continuation, 19 November 2026: used and rechecked in this run; status-vocabulary reconciliation; metadata reconciliation only; no new mathematical check. \emph{Pre-continuation recorded proof scope.} Claude Fable 5.2 (Anthropic), effort high, 16 November 2026; also checked in this writing by GPT-6 Astra (OpenAI), ChatGPT Pro mode, effort Pro, 18 November 2026}{the hypotheses, notation and definitions stated here; \cref{maps:src:F.1.1}}


Suppose a compact smooth eight-manifold has the torus fibration
and the specified Kuhn cusp of \ref{maps:src:F.1.1} with index one, and its
other reduced fibres are free quotients of complex three-tori
by finite holomorphic affine actions.  If its rational
$H^1$ monodromy representation is self-dual, then it cannot
satisfy $b_1=b_2=0$.  In particular no symplectic-monodromy
datum of F.2 can give the three-connected fourfold required
in F.4, for either $k=2$ or $k=3$ (or any number $k$ of finite
points).

\end{theorem}


\begin{proof}
\begin{enumerate}[label=\arabic*.]
\item Let $\mathcal V$ be the rank-six cohomology local system on
the punctured sphere and let $j$ be its inclusion.  On a finite
quotient fibre, averaging gives $R^1g_*\mathbb Q=(j_*\mathcal V)$
at that point.  At the Kuhn cusp, its normalization components
are simply connected and the degree-one branch complex is
the cochain complex of its dual $T^3$.  Specialization takes
its three deck classes to the covectors annihilating the
angular lattice, exactly the three-dimensional monodromy
invariant subspace.  Hence the same equality holds at the cusp.
It holds everywhere.

\item Put $a=\dim H^0(B,j_*\mathcal V)$ and
$b=\dim H^2(B,j_*\mathcal V)$.  The only differential that
can kill a class in $H^0(R^1)$ or in $H^2(R^0)$ is
$d_2:H^0(R^1)\to H^2(R^0)=\mathbb Q$.  If $b_1=b_2=0$,
it must be an isomorphism, so $a=1$.  Moreover
$H^1(B,j_*\mathcal V)$ must vanish because its position
$(1,1)$ has no possible incoming or outgoing differential.
Self-duality and compactly supported duality give $b=a=1$.

\item At the $i$th finite point, let $r_i$ be the dimension of the
fixed subspace in $H^1$.  The derivative is complex linear,
so its one-eigenspace has even real dimension.  Each $r_i$
is therefore even.  The cusp fixed dimension is three.
Euler additivity gives
\[
 \chi(B,j_*\mathcal V)=6(1-k)+3+\sum_{i=1}^k r_i,
\]
an odd integer.  Consequently
\[
 \dim H^1(B,j_*\mathcal V)=a+b-\chi(B,j_*\mathcal V)
\]
is odd and cannot vanish.  This is the required contradiction.
A preserved nondegenerate alternating form identifies a
representation with its dual, so symplectic monodromy is
included.  The explicit form
$(6\gamma\wedge\delta+u\wedge w)\oplus\omega_E$
checks this hypothesis for the example of Part E.

\end{enumerate}
\end{proof}




\begin{proposition}[the remaining rank patterns]\label{maps:src:F.3.2}
\StatementRecord{maps:src:F.3.2}{proved}{Historical events: Claude Fable 5.2 (Anthropic), effort high, 16 November 2026: carried as background without fault found. GPT-6 Astra (OpenAI), ChatGPT Pro mode, effort Pro, S6\_09\_PAPERS, 17 November 2026: targeted proof check. Also checked in this writing. Pre-continuation writing events: GPT-6 Astra (OpenAI), ChatGPT Pro mode, effort Pro, S6\_10\_PAPERS, 18 November 2026 02:02:02.394029+00:00: targeted proof check. Continuation mathematical and source checks: GPT-6 Astra (OpenAI), ChatGPT Pro mode, effort Pro, S6\_10\_PAPERS continuation, 19 November 2026: computation or formal certificate; independent re-proof; used and rechecked in this run. Continuation metadata events: GPT-6 Astra (OpenAI), ChatGPT Pro mode, effort Pro, S6\_10\_PAPERS continuation, 19 November 2026: used and rechecked in this run; status-vocabulary reconciliation; metadata reconciliation only; no new mathematical check. \emph{Pre-continuation recorded proof scope.} Claude Fable 5.2 (Anthropic), effort high, 16 November 2026; also checked in this writing by GPT-6 Astra (OpenAI), ChatGPT Pro mode, effort Pro, 18 November 2026}{the hypotheses, notation and definitions stated here; the explicit constructions and calculations in the proof}


If a possibly non-self-dual holomorphic datum with the same
local models had $b_1=b_2=0$, put
$b=\dim H^0(B,\mathcal V^*)$.  For $k=2$, the only possible
unordered pairs $(r_1,r_2;b)$ are
\[
 (2,2;0),\qquad(2,4;2).
\]
For $k=3$, they are
\[
 (2,2,6;0),\ (2,4,4;0),\ (2,4,6;2),\ (4,4,4;2).
 \tag{F.4}\label{maps:eq:F.4}
\]

\end{proposition}


\begin{proof}
\begin{enumerate}[label=\arabic*.]
\item The preceding Leray argument still gives $a=1$ and vanishing
$H^1(B,j_*\mathcal V)$.  Euler additivity therefore gives
$b=\sum r_i+3-1-6(k-1)$.  Each $r_i$ belongs to
$\{2,4,6\}$ because a free affine lift requires an eigenvalue
one and complex linearity makes its multiplicity even.
Also $0\leq b\leq\min(3,r_1,\ldots,r_k)$, by inclusion in
each local invariant subspace.  Enumerating these three-value
lists gives exactly \eqref{maps:eq:F.4}, as independently printed by the
complete bounded filter.  No existence of a tuple or of a
period family with these ranks is asserted.

\end{enumerate}
\end{proof}




\begin{proposition}[resolution of the classification problem]\label{maps:src:F.2.4}
\StatementRecord{maps:src:F.2.4}{proved}{Historical events: Claude Fable 5.2 (Anthropic), effort high, 16 November 2026: carried as background without fault found. GPT-6 Astra (OpenAI), ChatGPT Pro mode, effort Pro, S6\_09\_PAPERS, 17 November 2026: targeted proof check. Also checked in this writing. Pre-continuation writing events: GPT-6 Astra (OpenAI), ChatGPT Pro mode, effort Pro, S6\_10\_PAPERS, 18 November 2026 02:43:58.020152+00:00: targeted proof check. Continuation mathematical and source checks: GPT-6 Astra (OpenAI), ChatGPT Pro mode, effort Pro, S6\_10\_PAPERS continuation, 19 November 2026: independent re-proof; used and rechecked in this run. Continuation metadata events: GPT-6 Astra (OpenAI), ChatGPT Pro mode, effort Pro, S6\_10\_PAPERS continuation, 19 November 2026: used and rechecked in this run; status-vocabulary reconciliation; metadata reconciliation only; no new mathematical check. \emph{Pre-continuation recorded proof scope.} Claude Fable 5.2 (Anthropic), effort high, 16 November 2026; also checked in this writing by GPT-6 Astra (OpenAI), ChatGPT Pro mode, effort Pro, 18 November 2026}{\cref{maps:src:F.2.16,maps:src:F.2.30,maps:src:F.2.31,maps:src:F.2.35,maps:src:F.2.36,maps:src:F.2.37,maps:src:F.2.38,maps:src:F.2.39}}


The previously incomplete linear classification and all-twist
integral topology computation are completed by \ref{maps:src:F.2.31} and
\ref{maps:src:F.2.35}--39, with the fixed cusp filling and full marking
specified as required by \ref{maps:src:F.2.3}.  Every linear tuple has an
exact canonical representative and a returned unimodular
conjugator; the finite bound depends on the supplied input.
The unbounded trace family proves that no fixed global
height cutoff can replace this procedure.  All actual affine
orders, numerators, clutch vectors and cusp markings are
included in the general fundamental-group and integral
homology algorithms.  The complete constructive survivor
list, its integral tests and the common oriented smooth
type are given in \ref{maps:src:F.2.30}.

This conclusion does not claim an executed enumeration of
the infinite general class set.  The exhaustive finite
involution catalogue, the general comparison tests, the
three full topology integrations and the unified wrapper
checks have the exact executed scopes stated in \ref{maps:src:F.2.35}--39.
The earlier obstruction record is retained verbatim in
\path{records/fourfolds_F2_4_superseded.tex}; its finite-order
filter and homology-level obstructions were superseded by
the complete canonical and chain constructions.


\end{proposition}
\begin{proof}
\begin{enumerate}[label=\arabic*.]
\item The integral intertwiner module and its centralizer order in
\cref{maps:src:F.2.36} decide simultaneous integral conjugacy. Enumerating
tuples in increasing height and then lexicographic order, and applying
that decision to each candidate, gives the canonical tuple and the
returned unimodular intertwiner of \cref{maps:src:F.2.37}. The supplied
tuple itself occurs by its own height, proving termination for this input.
\item The affine power vectors, integral normal closure and marked
extension presentation in \cref{maps:src:F.2.31} give the full fundamental
group. The finite-cap, cusp and complement maps of
\cref{maps:src:F.2.35} form the actual gluing cone, whose finite integer
differentials determine every integral homology group by Smith reduction.
All orders, affine numerators, the clutch vector and the cusp basis enter
those maps explicitly.
\item The unbounded conjugacy invariant of \cref{maps:src:F.2.16}
rules out any uniform finite height box: a finite box contains only
finitely many tuples and hence only finitely many values of that trace.
The survivor criterion and its smooth classification are exactly
\cref{maps:src:F.2.30}; the involution catalogue and complete algorithmic
assembly are \cref{maps:src:F.2.38,maps:src:F.2.39}.
\end{enumerate}
\end{proof}



\section{The integral eight-manifold}



\begin{theorem}[a non-self-dual topological candidate]\label{maps:src:F.2.5}
\StatementRecord{maps:src:F.2.5}{proved}{Historical events: Claude Fable 5.2 (Anthropic), effort high, 16 November 2026: carried as background without fault found; targeted proof check. GPT-6 Astra (OpenAI), ChatGPT Pro mode, effort Pro, S6\_09\_PAPERS, 17 November 2026: targeted proof check. Also checked in this writing. Pre-continuation writing events: GPT-6 Astra (OpenAI), ChatGPT Pro mode, effort Pro, S6\_10\_PAPERS, 18 November 2026 02:02:02.394029+00:00: computation or formal certificate; targeted proof check. Continuation mathematical and source checks: GPT-6 Astra (OpenAI), ChatGPT Pro mode, effort Pro, S6\_10\_PAPERS continuation, 19 November 2026: computation or formal certificate; independent re-proof; targeted proof check; used and rechecked in this run. Continuation metadata events: GPT-6 Astra (OpenAI), ChatGPT Pro mode, effort Pro, S6\_10\_PAPERS continuation, 19 November 2026: used and rechecked in this run; status-vocabulary reconciliation; metadata reconciliation only; no new mathematical check. \emph{Continuation recorded proof scope.} Complete current five-step proof re-derived before the corresponding target readings: actual marked caps and meridians, exact integral normal closure, rational restriction complex, relative almost complex extension, Todd-signature constraint and integral lattice. Root and VI frozen scopes are aggregated; all current numerical claims have fresh exact certificates. The former additional Hermite, full-table and homogeneous-space identification routes are preserved with their former scopes.}{\cref{maps:src:F.2.1,night:P1.1,night:P1.9,night:P1.11,maps:src:F.1.2}; \cite{ffAchingerOgus20,StacksProject,HirzebruchTopologicalMethods,ffPetersSterk}}


There is a closed smooth oriented eight-manifold with the prescribed
index-one Kuhn cusp, two free affine fillings of multiplicities $5,4$,
trivial fundamental group, and rational Betti numbers
\[
 (1,0,0,0,4,0,0,0,1). \tag{F.5}\label{maps:eq:F.5}
\]
It admits an almost complex structure.  Its signature is two and
its middle intersection lattice modulo torsion is
$3\langle1\rangle\oplus\langle-1\rangle$.
Integral $H_2,H_3$ and a holomorphic period realization are not
asserted here.

\end{theorem}


\begin{proof}
\begin{enumerate}[label=\arabic*.]
\item We use the following matrices on the marked dual lattice:
\[
A=\begin{pmatrix}
1&0&0&1&-1&1\\0&1&-1&-1&1&0\\0&1&0&-1&0&0\\
0&2&-1&-1&0&1\\0&2&0&-1&0&0\\0&1&1&-1&0&0
\end{pmatrix},\quad
B=\begin{pmatrix}
1&0&0&0&0&0\\0&0&-1&0&1&0\\0&0&-1&0&0&1\\
0&0&-1&0&0&0\\0&1&-1&0&0&0\\0&0&-1&1&0&0
\end{pmatrix}.
\]
Put $M_1=A^t$, $M_2=B^t$, $M_0=(M_1M_2)^{-1}$ and
\[
 v_1=(1,3,-2,-1,2,0)^t,\qquad v_2=(-1,0,0,0,0,0)^t.
\]
Theorem~\ref{night:P1.1} gives their orders, the primitive cusp
marking, the coinvariant homomorphism $\ell(x)=x_0$, and freeness of
every nonidentity power of the specified affine actions. Consequently
\[
 (T^6\times\Delta_u)/
 \langle(x,u)\mapsto(M_ix+v_i/m_i,\zeta_{m_i}u)\rangle,
 \qquad (m_1,m_2)=(5,4),
\]
are smooth oriented finite caps. On the universal angular cover
$u=r e^{2\pi ia}$ their actual marked collar maps are
\[
 \alpha_i(x,u)=x-av_i,
 \qquad
 \alpha_i(M_ix+v_i/m_i,\zeta_{m_i}u)=M_i\alpha_i(x,u).
\]
At the cusp we take the actual index-one Kuhn cap of
Theorem~\ref{night:P1.9}, in the marking of
Proposition~\ref{maps:src:F.2.1}. The three displayed boundary bundles
therefore glue to the marked torus bundle over the intervening pair
of pants, with zero additional clutch translation. Every boundary
is paired. The result is a closed connected smooth eight-manifold,
oriented by the cusp complex structure and the continued fibre/base
orientation.

\item The small cusp cap retracts onto its central fibre $C$ through
the actual mapping-cylinder homeomorphism
\cite[Theorem 4.1.1 and Lemmas 4.1.2--4.1.3, pp.~480--482]{ffAchingerOgus20}.
Indeed its divisorial log chart is
$\mathbb N\to\mathbb N^r$, $1\mapsto(1,\ldots,1)$; it is log smooth,
saturated and vertical, and the cap is proper by
Theorem~\ref{night:P1.9}. In the resulting cylinder the homotopy
$[x,r]\mapsto[x,(1-t)r]$ fixes $C$ and retracts onto it.

Before the period quotient, the components of $C$ and their
nonempty intersections are smooth complete toric strata. Each is
simply connected: a loop moves into its dense torus by general
position, and the one-parameter loops along the rays of a smooth
maximal cone bound the coordinate discs; those rays are an integral
basis. Van Kampen sends component-crossing paths to edges and triple
intersections to triangles of the Kuhn nerve in $\mathbb R^3$.
Thus the covering core has trivial fundamental group. Its period
quotient has deck group $\mathbb Z^3$, so $\pi_1(C)=\mathbb Z^3$.
The nearby-torus map is the quotient of $\Lambda$ by the saturated
image of $M_0-I$, by the actual $H_1$ collapse in
Theorem~\ref{night:P1.11}. It is onto, and the cusp meridian dies
because its zero section extends.

We also record the invariant normal closure needed below. Let
$K$ be the smallest $\langle M_1,M_2\rangle$-invariant subgroup
containing $(M_0-I)\Lambda$, and let
$J=(M_1-I)\Lambda+(M_2-I)\Lambda=\ker\ell$.
The equality $\ell M_i=\ell$ gives $K\subseteq\ker\ell$.
On the actual abelian quotient $\Lambda/K$, the induced maps $a,b$
satisfy $ab=1$, $a^5=1$ and $b^4=1$. Thus $a^4=a^5=1$, so
$a=b=1$ and $J\subseteq K$. Consequently $K=J=\ker\ell$ as
integral subgroups; the quotient argument uses no saturation.


Let $a_i$ be the two finite peripheral lifts. The universal covering
deck transformation of the $i$th cap has $m_i$th power equal to
translation by $v_i$, since $M_iv_i=v_i$. The actual collar map in
Step~1 sends $a_i$ to that transformation. Van Kampen therefore gives
\[
 a_i\lambda a_i^{-1}=M_i\lambda,
 \qquad a_1^5=v_1,\quad a_2^4=v_2,\quad a_1a_2=1.
\]
Write $a_1=t$ and $a_2=t^{-1}$. Then
$t=t^5t^{-4}=v_1+v_2$ lies in the abelian fibre subgroup. Conjugation
is consequently trivial, and that subgroup is a quotient of
$\Lambda_{\langle M_1,M_2\rangle}=\mathbb Z$, with coordinate $\ell$.
Now $\ell(v_1+v_2)=0$, so $t=0$, while $t^5=v_1$ kills its generator
$\ell(v_1)=1$. Every generator vanishes, proving $\pi_1=1$.
Equivalently, after this abelian reduction, the relations on
$(\ell,a_1,a_2)$ have matrix
\[
 \begin{pmatrix}-1&5&0\\1&0&4\\0&1&1\end{pmatrix},
 \qquad\det=-1.
\]

\item We calculate rational cohomology directly from the marked
restriction maps. For $1\leq q\leq4$ put
\[
 L_q=\bigwedge^q\Lambda^\vee\otimes\mathbb Q,
 \quad A_q=\bigwedge^qA,\quad B_q=\bigwedge^qB,
 \quad C_q=(B_qA_q)^{-1}.
\]
The finite restriction images are
$J_{1,q}=\ker(A_q-I)$ and $J_{2,q}=\ker(B_q-I)$: transfer followed
by restriction is multiplication by the finite covering degree,
which is invertible over $\mathbb Q$. The marked collapse in
Theorem~\ref{night:P1.11} gives
$J_{0,q}=\ker(C_q-I)$ at the cusp, with one additional rational
kernel in degree four. Proper base change identifies these with the
actual stalk maps
\cite[Tag 09V5, Lemma 20.18.1; chapter PDF p.~30]{StacksProject}.

Cutting the sphere along the two based peripheral arcs gives the
following cellular complex for the sheaf of restriction images:
\[
 J_{1,q}\oplus J_{2,q}\oplus J_{0,q}
 \xrightarrow{d_0}L_q^2\xrightarrow{\Delta_q}L_q,
\]
\[
 d_0(j_1,j_2,j_0)=(j_1-j_0,j_2-A_qj_0),
 \qquad
 \Delta_q(x,y)=(A_q-I)x+(B_q-I)y.
\]
The cusp value is transported across the first arc by $A_q$ before
the second subtraction. The relative disc boundary is $\Delta_q$;
their composition is zero because $B_qA_qj_0=j_0$.
The exterior matrices are their minors in lexicographically ordered
exterior bases. Rational elimination gives
\[
\begin{array}{c|r|rrr|rr}
q&\dim L_q&\dim J_{1,q}&\dim J_{2,q}&\dim J_{0,q}
 &\operatorname{rank}d_0&\operatorname{rank}\Delta_q\\\hline
1&6&2&2&3&6&6\\
2&15&3&3&9&15&15\\
3&20&4&4&9&17&20\\
4&15&3&3&9&15&15
\end{array}
\]
Here are explicit nonzero minors for the two rank calculations.
For each $J_{i,q}$ use the rational nullspace basis obtained from
reduced row echelon form: order the free columns increasingly, set
one free coordinate to one and the others to zero, and solve the
pivot coordinates. Number rows and columns from zero.
For $\Delta_q$ use all rows and the following columns:
\[
\begin{array}{c|l|r}
q&\text{columns}&\det\\\hline
1&1,2,3,5,7,8&2\\
2&0,1,2,4,\ldots,12,16,20,24&-6\\
3&0,\ldots,7,10,\ldots,17,20,24,30,36&-6\\
4&0,\ldots,7,10,11,12,13,15,21,25&6
\end{array}
\]
For $d_0$, use columns $(0,1,2,3,5,6)$ when $q=1$ and all
columns otherwise, with the following row sets:
\[
\begin{array}{c|l|r}
q&\text{rows}&\det\\\hline
1&0,1,3,4,6,7&1\\
2&0,\ldots,8,10,12,13,15,20,21&-2/3\\
3&0,\ldots,10,13,16,20,21,30,31&-4/3\\
4&0,\ldots,11,15,16,25&1/3
\end{array}
\]
All are determinants of the matrices specified above. They give
the stated full row ranks for $\Delta_q$ and full column ranks for
$d_0$ in degrees two through four. In degree one the common fixed
covector $e_0$ gives the nonzero kernel vector $(e_0,e_0,e_0)$ of
$d_0$, so its rank is at most six; the displayed minor gives equality.
The finite invariant dimensions count exterior products of the
specified eigenvalues with product one; the cusp dimensions are
those of Theorem~\ref{night:P1.11}.

These ranks determine the cohomology: $H^0$ of the image sheaf is
$\mathbb Q$ in coefficient degree one, $H^1$ is $\mathbb Q^3$ in
coefficient degree three, and every other entry in these four
coefficient degrees is zero. The cusp skyscraper kernel adds
$\mathbb Q$ to $H^0(R^4\pi_*\mathbb Q)$ and no higher cohomology.
Connected fibres give $R^0\pi_*\mathbb Q=\mathbb Q$.

In the Leray sequence
\cite[Tag 01F2, Lemma 20.13.4; chapter PDF p.~22]{StacksProject},
the only remaining low-degree differential is
$d_2^{0,1}:\mathbb Q\to\mathbb Q$. Its kernel is $H^1(M;\mathbb Q)=0$
by Step~2, so it is an isomorphism. Hence
$b_1=b_2=b_3=0$ and $b_4=3+1=4$. Poincar\'e duality determines all
remaining Betti numbers and gives \eqref{maps:eq:F.5}. Also, the
smooth and finite quotient fibres have Euler number zero, while the
Kuhn cusp has Euler number six by Theorem~\ref{night:P1.9}. Additivity
over the base strata therefore gives $\chi(M)=6$ directly.

\item Each finite monodromy admits an invariant real linear complex
structure: its $1$ and $-1$ real eigenspaces have even dimension, and
each remaining conjugate eigenpair gives a real two-plane. Choose
a complex structure on each summand. Reversing it on one fixed
two-plane selects the required orientation. These structures commute
with $M_i$, so their direct sums with the root-disc structure descend
to the actual affine caps. The cusp cap already has its complex
structure. Their vertical complex structures on the boundary collars
are translation-invariant in the marked torus fibres.

The space $\mathcal J=\operatorname{GL}_6^+(\mathbb R)/
\operatorname{GL}_3(\mathbb C)$ of allowable vertical structures
retracts to $\operatorname{SO}(6)/\operatorname{U}(3)$ by the explicit
conjugation
\[
 J\longmapsto P_tJP_t^{-1},
 \qquad P_t=\bigl((1-t)I+t(I+J^tJ)\bigr)^{1/2}.
\]
At $t=1$ the structure is orthogonal, since $J$ preserves the metric
$I+J^tJ$; an orthogonal $J$ is fixed throughout this homotopy.
The inclusion $\operatorname{U}(3)\to\operatorname{SO}(6)$ sends
$\operatorname{diag}(e^{i\theta},1,1)$ to one-plane rotation and
induces the surjection $\mathbb Z\to\mathbb Z/2$ on fundamental
groups. Indeed the determinant fibration and
$\operatorname{SU}(2)=S^3\to\operatorname{SU}(3)\to S^5$ give the
first group; the quaternion cover of $\operatorname{SO}(3)$ and
the sphere fibrations for $\operatorname{SO}(n)$ give the second,
with a full rotation lifting from $1$ to $-1$. The homogeneous-space
exact sequence therefore gives $\pi_0\mathcal J=\pi_1\mathcal J=0$.

Over the two-dimensional pair of pants the prescribed boundary
sections of the associated $\mathcal J$-bundle extend over relative
one-cells by connectedness and over relative two-cells by simple
connectivity. Smooth approximation relative to the collars gives
a smooth vertical structure. Horizontal splittings form an affine
space, so a partition of unity extends the prescribed complex
splittings from those collars. Combining the vertical structure with
the pulled-back base structure and gluing the caps gives the required
almost complex structure on $M$.

\item We have $c_1=c_3=0$ rationally because $H^2(M;\mathbb Q)$ and
$H^6(M;\mathbb Q)$ vanish, and $c_4[M]=\chi(M)=6$. The characteristic
number formulas in Theorem~\ref{maps:src:F.1.2}, with their source
\cite[equations (1.2)--(1.4), pp.~3--4]{ffEsoleKang}, give
\[
 \operatorname{Td}[M]=\frac{3c_2^2[M]-6}{720},
 \qquad
 \sigma(M)=\frac{3c_2^2[M]+84}{45}
           =16\operatorname{Td}[M]+2.
\]
The Todd genus of this almost complex manifold is an integer
\cite[Theorem 24.5.4 and corollary, p.~184]{HirzebruchTopologicalMethods}.
Since $|\sigma(M)|\leq b_4=4$, the displayed integer relation forces
$\operatorname{Td}[M]=0$ and $\sigma(M)=2$. Thus the real middle
pairing has positive and negative indices $3$ and $1$.

Integral Poincar\'e duality makes the middle pairing modulo torsion
unimodular. An even unimodular lattice has signature divisible by
eight, so this lattice is odd. The classification of indefinite odd
unimodular lattices
\cite[Theorem 2.2.1, p.~61]{ffPetersSterk} gives its integral isometry
class $3\langle1\rangle\oplus\langle-1\rangle$.
\end{enumerate}
\end{proof}


\label{sec:fourfold}

We use column vectors.  The matrices $A,B$ in the supplied data act on
cohomology; $M_1=A^t$, $M_2=B^t$ act on
$\Lambda=\mathbb Z^6=H_1(T^6;\mathbb Z)$, and
$M_0=(M_1M_2)^{-1}$.  Positive based meridians satisfy
$g_1g_2g_0=1$.  The exterior-square basis of $\bigwedge^2\mathbb Z^4$
is $(01,02,03,12,13,23)$.  All assertions about an integral matrix below
are checked over $\mathbb Z$ or $\mathbb Q$, without floating-point
arithmetic.  The files in \texttt{code/part1/} contain the complete
matrices, including the changes of basis in every Smith computation.

\begin{theorem}\label{night:P1.1}
\StatementRecord{night:P1.1}{proved}{Historical events: GPT-6 Astra (OpenAI), ChatGPT Pro mode, effort Pro, S6\_06\_BLIND, 16--18 November 2026: targeted proof check. Also checked in this writing. Pre-continuation writing events: GPT-6 Astra (OpenAI), ChatGPT Pro mode, effort Pro, S6\_10\_PAPERS, 18 November 2026 02:02:02.394029+00:00: computation or formal certificate; targeted proof check. Continuation mathematical and source checks: GPT-6 Astra (OpenAI), ChatGPT Pro mode, effort Pro, S6\_10\_PAPERS continuation, 19 November 2026: computation or formal certificate; independent re-proof; used and rechecked in this run. Continuation metadata events: GPT-6 Astra (OpenAI), ChatGPT Pro mode, effort Pro, S6\_10\_PAPERS continuation, 19 November 2026: used and rechecked in this run; status-vocabulary reconciliation; metadata reconciliation only; no new mathematical check. }{Definitions and the displayed constructions}
The finite homology monodromies have orders $5,4$, respectively, and
characteristic polynomials
$(T-1)^2\Phi _5(T)$ and $(T-1)^2(T+1)^2(T^2+1)$.
For $N=M_0-I$ we have $N^2=0$, $\operatorname{rank}N=3$, and
nonzero Smith factors $1,1,1$.  Moreover
\[
 \Lambda^{\langle M_1,M_2\rangle}=0,\qquad
 \Lambda_{\langle M_1,M_2\rangle}=\mathbb Z,
\]
where the quotient homomorphism is $\ell(x)=x_0$.
On the cohomology lattice the corresponding groups are $\mathbb Ze_0$
and $\mathbb Z/2$.  Every nonidentity power of each specified affine
action $\tau_i$ is free.
\end{theorem}
\begin{proof}
\begin{enumerate}[label=\arabic*.]
\item Exact multiplication gives
\[
N=\begin{pmatrix}
0&0&0&0&0&0\\0&0&0&1&0&-1\\0&0&0&0&1&1\\
-1&0&0&0&0&0\\1&0&0&0&0&0\\-1&0&0&0&0&0
\end{pmatrix},\qquad
S=\begin{pmatrix}
0&0&0&1&0&0\\0&1&0&0&0&0\\0&0&1&0&0&0\\
-1&0&0&0&1&0\\1&0&0&0&0&1\\-1&0&0&0&0&0
\end{pmatrix}.
\]
The matrix $S$ is unimodular and
$S^{-1}M_0S=\left(\begin{smallmatrix}I_3&I_3\\0&I_3\end{smallmatrix}\right)$.
The columns of $S$ are henceforth denoted
$(a_1,a_2,a_3,b_1,b_2,b_3)$, so that $Nb_j=a_j$.
This is also an integral certificate of the Smith assertion.
The Hermite basis of
$(M_1-I)\Lambda+(M_2-I)\Lambda$ is $(e_1,\ldots,e_5)$;
the stacked matrix of the two fixed-vector equations has rank six.
Transposing the matrices gives a fixed lattice $\mathbb Ze_0$ and
coinvariant Smith factors $(1,1,1,1,1,2)$.  The pairing
$\Lambda^\vee\times\Lambda\longrightarrow\mathbb Z$ identifies the
fixed cohomology vector $e_0$ with the covector $\ell$.

\item We have $M_iv_i=v_i$, $\ell(v_1)=1$, and $\ell(v_2)=-1$.
Thus $\tau_i^k(x)=M_i^kx+kv_i/m_i$, and $\tau_i^{m_i}=1$ on the
torus.  A fixed point of $\tau_i^k$ would give
$k\ell(v_i)/m_i=0$ in $\mathbb R/\mathbb Z$.  For
$0<k<m_i$ this is impossible.  The order and characteristic-polynomial
computations, and the full fixed-lattice calculations, are reproduced by
\texttt{tuple\_check.py}; its output records every power tested.

\end{enumerate}
\end{proof}




\begin{theorem}\label{night:P1.3}
\StatementRecord{night:P1.3}{proved}{Historical events: GPT-6 Astra (OpenAI), ChatGPT Pro mode, effort Pro, S6\_06\_BLIND, 16--18 November 2026: targeted proof check. Also checked in this writing. Pre-continuation writing events: GPT-6 Astra (OpenAI), ChatGPT Pro mode, effort Pro, S6\_10\_PAPERS, 18 November 2026 02:02:02.394029+00:00: computation or formal certificate; targeted proof check. Continuation mathematical and source checks: GPT-6 Astra (OpenAI), ChatGPT Pro mode, effort Pro, S6\_10\_PAPERS continuation, 19 November 2026: computation or formal certificate; independent re-proof; used and rechecked in this run. Continuation metadata events: GPT-6 Astra (OpenAI), ChatGPT Pro mode, effort Pro, S6\_10\_PAPERS continuation, 19 November 2026: used and rechecked in this run; status-vocabulary reconciliation; metadata reconciliation only; no new mathematical check. }{Definitions and the displayed constructions}
The two finite filling pieces have first integral homology
$\mathbb Z^2$ and $\mathbb Z^2\oplus\mathbb Z/2$.
\end{theorem}
\begin{proof}
\begin{enumerate}[label=\arabic*.]
\item For $K_i=T^6/\langle\tau_i\rangle$, the universal covering deck group
has commuting translations $x^u$, $u\in\Lambda$, and a generator $a_i$,
with the complete relations
\[
a_ix^ua_i^{-1}=x^{M_iu},\qquad a_i^{m_i}=x^{v_i}.
\]
Indeed the affine lifts satisfy these relations, and they reduce every
word uniquely to $x^ua_i^r$, $0\leq r<m_i$.
The abelianization relation matrix has columns
$((M_i-I)e_j,0)$ and $(-v_i,m_i)$.
Smith reduction of these explicit seven-by-seven matrices gives the
groups above.  The disk quotient
$(T^6\times\Delta)/\langle(\tau_i,\zeta_{m_i})\rangle$
retracts onto $K_i$ by equivariant radial contraction.  Freeness supplies
ordinary smooth quotient charts by disjoint translates of small balls.

\end{enumerate}
\end{proof}


\begin{lemma}[The actual order-four torsion transgression]
\label{mapsaudit:order-four-torsion-transgression}
\StatementRecord{mapsaudit:order-four-torsion-transgression}{proved}{No earlier separately identified check is recorded. Checked only in this writing. Pre-continuation writing events: GPT-6 Astra (OpenAI), ChatGPT Pro mode, effort Pro, S6\_10\_PAPERS, 18 November 2026 02:02:02.394029+00:00: targeted proof check. Continuation mathematical and source checks: GPT-6 Astra (OpenAI), ChatGPT Pro mode, effort Pro, S6\_10\_PAPERS continuation, 19 November 2026: computation or formal certificate; independent re-proof; targeted proof check; used and rechecked in this run. Continuation metadata events: GPT-6 Astra (OpenAI), ChatGPT Pro mode, effort Pro, S6\_10\_PAPERS continuation, 19 November 2026: used and rechecked in this run; status-vocabulary reconciliation; metadata reconciliation only; no new mathematical check. \emph{Continuation recorded proof scope.} The complete current cellular proof was written independently before the target proof was opened. The actual finite covering restrictions, integral sheaf complex, primitive minors, and unit torsion differential are explicit. Fresh exact programs verify these calculations. The complete former proof is preserved with its separate review scope; the cusp-collapse input has its own review.}{\cref{maps:src:F.2.5}; the integral quotient and explicit certificates in \cref{maps:src:F.2.6}; \cref{maps:eq:F.6}; the displayed cellular and boundary maps}

Use the order-four filling and the marked matrices of \cref{maps:src:F.2.5}.
Let $T=\bigwedge^{q-1}D^t$ on the integral vertical cochains of the
five-torus. If $w$ represents a torsion element of
$\operatorname{coker}(T-I)$, its point-supported degree-$q$ class has
Leray differential equal, up to the one boundary-orientation sign, to
the image of $w$ in the global degree-$(q-1)$ coinvariants. There is
no multiplicity-four factor. For each of $q=2,3,4$, the unique nonzero local torsion class maps
to the nonzero class in the global $\mathbb Z/2$ quotient.
\end{lemma}


\begin{proof}
\begin{enumerate}[label=\arabic*.]
\item Use the explicit matrix $C_2$ in the proof of
\cref{maps:src:F.2.6}.  In the original vertical marking $D$ is the
bottom-right block of $M_2$; multiplication gives
$DC_2=C_2D_2$, where $D_2$ is the four-cycle followed by the sign
on its fifth coordinate.  The change of cochain coordinates is
$\bigwedge^kC_2^t$, and the inverse is $\bigwedge^kC_2^{-t}$.
Thus the product cell structure of the signed-permutation torus,
transported by this integral torus isomorphism, has zero cellular
differential and the stated action $T=\bigwedge^{q-1}D^t$.
The fibre cells and their prisms give the actual mapping-torus
cochain complex.  The degree-$q$ suspension classes are consequently
the coinvariants of $T$ on vertical degree-$q-1$ cochains.

\item Write $u=-x_0$ in the original marking.  The affine action is
$(u,y,z)\mapsto(u+1/4,Dy,iz)$, and
$W=e^{-2\pi i u}z$ is invariant.  With core angle $a=4u\bmod1$
and normal angle $k=\arg W/(2\pi)$, the coarse angle is $a+4k$.
Let $g,h$ denote the two corresponding loops.  The integral basis
\[
 v=g^4h^{-1},\qquad t=g,
 \qquad [v,t]_{[g,h]}=\begin{pmatrix}4&1\\-1&0\end{pmatrix}
\]
has determinant one.  Its coarse degrees are zero and one; its
vertical monodromies are $T^4=I$ and $T$ respectively.
It specifies the actual generic six-torus fibre and the primitive
base loop in the boundary product.

\item A suspension cochain representing $[w]$ has values $(w,0)$
on $(g,h)$.  Following the four consecutive core prisms transports
their values by $I,T,T^2,T^3$, giving
\[
 (w,0)_{(g,h)}\longmapsto(N_Tw,w)_{(v,t)},
 \qquad N_T=I+T+T^2+T^3.
 \tag{FT.1}\label{maps:eq:FT.1}
\]
For a torsion coinvariant there are an integer $n\ne0$ and an
integral cochain $b$ with $nw=(T-I)b$.  Therefore
$nN_Tw=N_T(T-I)b=0$, and $N_Tw=0$ in the free cochain lattice.
The fibre restriction vanishes; its base-circle component is $[w]$.
Replacing $w$ by $w+(T-I)b$ changes that component by its cellular
coboundary.  The formula is consequently a well-defined map of the
stated integral coinvariant groups.

\item The base cover by a disc at the finite point and its
complement gives the relative \v{C}ech connecting map.  On their
annular intersection the degree-zero fibre restriction is zero and
the next component is the one-cocycle just computed.
The connecting map sends it to $\pm[w]$ in the global coinvariants;
the two ordinary peripheral loops impose the relations
$\bigwedge^{q-1}M_1^t-I$ and $\bigwedge^{q-1}M_2^t-I$.
These are the two successive boundaries in the same base--fibre
complex, so their composition is the Leray $d_2$ of
\eqref{maps:eq:F.6}.  The one boundary-orientation choice accounts
for the sign.  Its coefficient is one, as witnessed by the
degree-one loop $t$ above.  Point-supported kernels or cokernels in
the comparison of restriction sheaves have no positive cohomology,
so they preserve this degree-two base quotient.

\item In signed-permutation coordinates the negative orbits in
vertical degrees $1,2,3$ are represented respectively by
$e_4$, $e_0\wedge e_2$ and $e_0\wedge e_2\wedge e_4$.
The orbit relations give exactly one local $\mathbb Z/2$ in each
case.  Transport by $\bigwedge^kC_2^{-t}$ and then include vertically
in the six-lattice.  This is the same map as the vertical inclusion
followed by $\bigwedge^kQ_2$ used in
\cref{maps:src:F.2.6}, because $P_2=(-1)\oplus C_2$.
The explicit $w_k,\rho_k,E_k,z_k$ there satisfy
\[
 \rho_k\Delta_k=0\pmod2,\qquad
 \rho_kw_k=1\pmod2,\qquad E_kz_k=2w_k,
 \qquad |\det E_k|=2.
\]
Thus the unique nonzero local class maps to the unique nonzero
global class in each of the three degrees.  Formula
\eqref{maps:eq:FT.1} sends the zero representative to zero as well.
\end{enumerate}
\end{proof}




\begin{theorem}[integral three-connectivity of the candidate]\label{maps:src:F.2.6}
\StatementRecord{maps:src:F.2.6}{proved}{Historical events: Claude Fable 5.2 (Anthropic), effort high, 16 November 2026: carried as background without fault found. Claude Fable 5.2 (Anthropic), effort high, 17 November 2026: carried as background without fault found. GPT-6 Astra (OpenAI), ChatGPT Pro mode, effort Pro, S6\_09\_PAPERS, 17 November 2026: targeted proof check. Also checked in this writing. Pre-continuation writing events: GPT-6 Astra (OpenAI), ChatGPT Pro mode, effort Pro, S6\_10\_PAPERS, 18 November 2026 02:02:02.394029+00:00: targeted proof check. Continuation mathematical and source checks: GPT-6 Astra (OpenAI), ChatGPT Pro mode, effort Pro, S6\_10\_PAPERS continuation, 19 November 2026: computation or formal certificate; independent re-proof; used and rechecked in this run. Continuation metadata events: GPT-6 Astra (OpenAI), ChatGPT Pro mode, effort Pro, S6\_10\_PAPERS continuation, 19 November 2026: used and rechecked in this run; status-vocabulary reconciliation; metadata reconciliation only; no new mathematical check. \emph{Continuation recorded proof scope.} The complete current cellular proof was written independently before the target proof was opened. The actual finite covering restrictions, integral sheaf complex, primitive minors, and unit torsion differential are explicit. Fresh exact programs verify these calculations. The complete former proof is preserved with its separate review scope; the cusp-collapse input has its own review.}{\cref{maps:src:F.2.5}; \cref{night:P1.11}; \cref{night:P1.1}; \cref{maps:src:F.2.1}; \cref{maps:eq:F.5}; the displayed integral cellular maps; \cite[Theorem 3.2, p.~195; Corollary 3.3, p.~196; Theorem 3.30, p.~241; Theorem 4.32, p.~366]{Hatcher2002}}


The manifold of Theorem \ref{maps:src:F.2.5} has free integral homology with
Betti numbers \eqref{maps:eq:F.5}.  In particular $H_2=H_3=0$, and it is
three-connected.

\end{theorem}


\begin{proof}
\begin{enumerate}[label=\arabic*.]
\item We use the marked construction of \cref{maps:src:F.2.5} and the
integral cusp-collapse assertion of \cref{night:P1.11}.
Explicitly, $M_1=A^t$, $M_2=B^t$ and the vectors $v_1,v_2$ are
those printed in item~1 of the proof of \cref{maps:src:F.2.5};
$S$ is the displayed integral marking in item~1 of the proof of
\cref{night:P1.1}.  Write
$m_1=5$, $m_2=4$, and put
\[
\begin{split}
C_1&=\begin{pmatrix}
-1&-1&0&-1&0\\0&1&0&1&1\\0&1&0&0&0\\
0&-1&-1&0&-1\\0&0&1&0&0
\end{pmatrix},\\[2pt]
C_2&=\begin{pmatrix}
-1&0&-1&0&-1\\0&1&0&0&0\\0&0&0&1&0\\
0&-1&0&-1&1\\0&0&1&0&0
\end{pmatrix}.
\end{split}
\]
For the specified translation vectors $v_i$, the matrices
$P_i=[v_i,(0;C_i)]$ are unimodular.  Direct multiplication gives
\[
 P_i^{-1}M_iP_i=1\oplus D_i,\qquad Q_i=P_i^{-t},
\]
where $D_1$ sends $e_j$ to $e_{j+1\bmod5}$, and $D_2$ sends
$e_j$ to $e_{j+1\bmod4}$ for $0\le j<4$ and $e_4$ to $-e_4$.
These equalities give the actual integral coordinate maps.  In these
coordinates the affine action is
\[
 (u,y,z)\longmapsto
 (u+1/m_i,D_i y,e^{2\pi i/m_i}z).
\]
Thus its core $F_i$ is the mapping torus of $D_i:T^5\to T^5$.
The function $W=e^{-2\pi i u}z$ trivializes its normal disc bundle.
With $\theta=m_i u\bmod1$, the coarse-base function on that bundle is
\[
                  q=W^{m_i}e^{2\pi i\theta}.
\]
In particular the boundary is $F_i\times S^1_W$ with this specified
map to the base circle.

\item We compute the integral restriction images on this model.
Let $V=\mathbb Z^5$ be the vertical cohomology lattice, let
$T_{i,k}=\bigwedge^kD_i^t$, and let
$N_{i,k}=\sum_{a=0}^{m_i-1}T_{i,k}^{a}$.
The signed permutations preserve the product cell structure with one
vertex and one edge on each circle.  Its cellular differential is zero.
Fibre cells and their prisms give the mapping-torus complex
\[
 C^q(F_i)=\bigwedge^qV\oplus\bigwedge^{q-1}V,
 \qquad d(a,b)=(0,(T_{i,q}-I)a).
\]
The actual $m_i$-fold covering by the generic six-torus pulls back
$(a,b)$ to $(a,N_{i,q-1}b)$: the horizontal cell is the union of the
$m_i$ consecutive prisms, with their transported coefficients.
Consequently its image in $L_q=\bigwedge^q\mathbb Z^6$ is
\[
 J_{i,q}=\bigwedge^q Q_i
 \left(\operatorname{Inv}T_{i,q}\ \oplus\
 e_0\wedge N_{i,q-1}\bigwedge^{q-1}V\right).
\]
Here vertical coordinates occupy slots $1,\ldots,5$ in the six-lattice.
This formula specifies both the lattice and its inclusion.

We fix its column basis as follows.  Exterior coordinates are in
increasing lexicographic order, with indices starting at zero.
For each signed-permutation orbit, start at its least coordinate
and sum its successive signed images until it returns there.
A positive orbit of length $l$ gives a primitive invariant column;
the norm image has the column $m_i/l$ times that column.
A negative orbit has zero invariant lattice and one $\mathbb Z/2$
in its coinvariant group.  Indeed the successive relations identify
all its coordinates with signed copies of the first, and the last
relation is twice that first coordinate.  Order the positive orbits
by their least coordinates, with vertical columns before horizontal
columns.  For $D_1$ there is no torsion in vertical degrees $1$ through
$4$.  For $D_2$ there is one negative orbit in each such degree; the
positive-orbit counts are $1,2,2,1$.

\item At the cusp, \cref{night:P1.11} supplies a primitive collapse
image and free homology, with the only collapse cokernel equal to
$\mathbb Z$ in degree four.  Dualizing those integral maps therefore
gives
\[
 J_{0,q}=\operatorname{Inv}\bigl(\bigwedge^q M_0^t\bigr),
\]
with cohomological restriction kernel zero for $q=1,2,3$ and
$\mathbb Z$ for $q=4$.
Here are explicit primitive bases, which also fix the column order
in the calculation below.  In the supplied cusp marking $S$, write
the cohomology basis as $x_0,x_1,x_2,y_0,y_1,y_2$; its action sends
$x_i$ to $x_i+y_i$ and fixes $y_i$.
In degree one take $(y_0,y_1,y_2)$.  In degree two take, in order,
\[
 y_i\wedge y_j\ (i<j),\qquad
 x_i\wedge y_i\ (0\le i\le2),\qquad
 x_i\wedge y_j+x_j\wedge y_i\ (i<j).
\]
In degree three take $y_0\wedge y_1\wedge y_2$, then
$x_i\wedge y_j\wedge y_k$ in lexicographic $(j,k)$ order with
$j<k$ and $i=j,k$, and then
\[
\begin{split}
 &x_0\wedge y_1\wedge y_2+x_1\wedge y_0\wedge y_2,\\
 &x_0\wedge y_1\wedge y_2-x_2\wedge y_0\wedge y_1.
\end{split}
\]
Their coordinate equations show that these bases are primitive and
have ranks $3,9,9$.  For degree four set
$J=\left(\begin{smallmatrix}0&I_3\\-I_3&0\end{smallmatrix}\right)$,
and let $H(e_I)$ be the complementary exterior basis vector, with
the sign of the permutation $(I,I^c)$.  Apply
$H\bigwedge^2J$ to the displayed degree-two basis.
The identity $JU=U^{-t}J$ gives the invariant equation, and the
signed-permutation maps $H$ and $\bigwedge^2J$ preserve primitivity.
Transport all these columns by $\bigwedge^qS^{-t}$ to obtain
$J_{0,q}$ in the original marking.  Thus its degree-four rank is nine.

\item Let $\mathcal R_q=R^qf_*\mathbb Z$ on the coarse base $S^2$,
let $j$ include its thrice-punctured locus, and let $\mathcal S_q$
be the image of $\mathcal R_q\to j_*L_q$.
Since $M$ is compact and the base is Hausdorff, the proper-base-change
map identifies its stalks with fibre cohomology
\cite[Tag 09V5, Lemma 20.18.1; chapter PDF p.~30]{StacksProject}.
Its three stalk images are exactly the lattices above.  The kernel
$\mathcal K_q$ is supported at those points: for $q=2,3$ it is the
order-four $\mathbb Z/2$, and for $q=4$ it is that $\mathbb Z/2$
together with the cusp $\mathbb Z$.
Point-supported sheaves have no positive cohomology, so
$H^p(\mathcal R_q)=H^p(\mathcal S_q)$ for $p>0$.

Put $A_q=\bigwedge^qM_1^t$ and $B_q=\bigwedge^qM_2^t$.
The supplied based marking gives
$\bigwedge^qM_0^t=(B_qA_q)^{-1}$.
Oriented cellular duality on the punctured sphere and its two-loop
spine gives the compact-support complex $L_q^2\to L_q$ in degrees
one and two.  Attaching the three stalk lattices gives the full complex
\[
\begin{gathered}
 J_{1,q}\oplus J_{2,q}\oplus J_{0,q}
 \xrightarrow{d_0}L_q^2\xrightarrow{\Delta_q}L_q,\\
 d_0(j_1,j_2,j_0)=(j_1-j_0,j_2-A_qj_0),\\
 \Delta_q(a,b)=(A_q-I)a+(B_q-I)b.
\end{gathered}
\]
For the cusp column, the two consecutive peripheral paths carry
$j_0$ and $A_qj_0$; their boundary is
$(B_qA_q-I)j_0=0$.  The reverse cusp orientation gives their negatives.
The two finite columns are the corresponding single peripheral
loops.  This is the relative cellular sequence for extension by zero
and the three stalk inclusions, so it retains their integral indices.
It also verifies $\Delta_qd_0=0$ with the displayed transports.

Here are finite determinant certificates for its cohomology.  They
use precisely the ordered bases just given, so every matrix is
specified by the marked $M_i,S,C_i$ and exterior minors.
For each $q$ select the following columns $\mathcal P_q$ of
$\Delta_q$ and call their square matrix $E_q$:
\[
\begin{aligned}
\mathcal P_1&=(1,2,5,8,9,10),\\
\mathcal P_2&=(0,1,2,8,9,11,12,13,16,17,23,24,25,26,28),\\
\mathcal P_3&=(1,2,3,5,6,10,11,12,13,15,16,18,19,\\
             &\hspace{22mm}26,29,31,32,34,35,37),\\
\mathcal P_4&=(0,1,5,7,9,10,11,13,16,17,20,23,26,28,29).
\end{aligned}
\]
Direct integer multiplication gives
\[
 \det E_q=(2,2,2,-2)_q,\qquad
 E_q^{-1}\Delta_q\in\operatorname{Mat}(\mathbb Z).
\]
Hence $E_q\mathbb Z^{\binom6q}$ is the entire image of
$\Delta_q$ and its cokernel is $\mathbb Z/2$.
Let $\mathcal F_q$ be the increasing complement of $\mathcal P_q$
in $\{0,\ldots,2\binom6q-1\}$.
These free-coordinate rows identify $\ker\Delta_q$ integrally:
the pivot coordinates of a kernel vector are
$-E_q^{-1}(\Delta_q)_{\mathcal F_q}$ times its free coordinates.
The relation matrix is therefore $R_q=(d_0)_{\mathcal F_q}$.
Its sizes, ranks and primitive minors are
\[
\begin{array}{c|c|c|l}
q&\text{size of }R_q&\operatorname{rank}R_q&
\text{maximal minors}\\ \hline
1&6\times7&6&-20,\ -5,\ 4\\
2&15\times15&15&\det R_2=1\\
3&20\times17&17&-1\\
4&15\times15&15&\det R_4=-1
\end{array}
\]
For $q=1$ the three column sets are respectively
$(0,1,2,3,5,6)$, $(0,1,2,4,5,6)$, $(0,2,3,4,5,6)$,
using all rows.  For $q=3$ use all columns and rows
$(0,1,\ldots,13,15,16,17)$.
These determinants have greatest common divisor one, so the
quotient by the relations is free.  The stalk column counts are
$(2,2,3)$, $(3,3,9)$, $(4,4,9)$ and $(3,3,9)$.
Thus the kernel dimensions and the quotient ranks give exactly
\[
\begin{array}{c|ccc}
q&H^0(\mathcal S_q)&H^1(\mathcal S_q)&H^2(\mathcal S_q)\\ \hline
1&\mathbb Z&0&\mathbb Z/2\\
2&0&0&\mathbb Z/2\\
3&0&\mathbb Z^3&\mathbb Z/2\\
4&0&0&\mathbb Z/2.
\end{array}
\]

\item We compute the torsion differential on the actual boundary.
Let $g$ be the core-circle loop of $F_2$ and $h$ the $W$-circle loop.
Their coarse-base degrees are one and four.  The basis
\[
 v=g^4h^{-1},\qquad t=g,
 \qquad [v,t]_{[g,h]}=\begin{pmatrix}4&1\\-1&0\end{pmatrix}
\]
has determinant one; $v$ lies in the generic torus fibre and $t$
projects once around the base.
For $k=q-1$ let $T=\bigwedge^kD_2^t$ and $N=I+T+T^2+T^3$.
A torsion class $[w]\in\operatorname{coker}(T-I)$ has a mapping-torus
prism representative with values $w$ on $g$ and zero on $h$.
Subdivision into successive transported prisms gives values $Nw$
on $v$ and $w$ on $t$.  An integral equation $nw=(T-I)a$ with
$n\ne0$ implies $nNw=0$, hence $Nw=0$ in the free lattice.
Its fibre restriction is therefore zero and its first boundary
Wang component is exactly $[w]$.
The relative-base connecting map takes that component to its global
coinvariant class.  In the cellular fibre and base \v{C}ech double
complex these are precisely the two successive boundaries defining
the Leray differential.  Thus, up to the single boundary-orientation
sign,
\[
 d_2[w]=\pm[w]\in L_{q-1}/\operatorname{im}\Delta_{q-1}.
 \tag{F.6}\label{maps:eq:F.6}
\]
Changing $w$ by $(T-I)a$ changes a boundary representative by a
coboundary.  The coefficient is one because the loop $t$ has degree
one over the coarse-base circle.

For $k=1,2,3$ the negative signed orbits have representatives
$e_4$, $e_0\wedge e_2$, $e_0\wedge e_2\wedge e_4$.
Map them by the vertical inclusion followed by $\bigwedge^kQ_2$;
call the resulting columns $w_k\in L_k$.
The following rows $\rho_k$, read modulo two, annihilate
$\Delta_k$:
\[
\begin{aligned}
\rho_1&=(0,1,1,1,1,1),\\
\rho_2&=(1,1,1,1,1,1,1,0,0,1,0,1,0,1,1),\\
\rho_3&=(0,0,0,0,0,0,0,0,0,0,1,1,0,1,0,1,1,1,0,0).
\end{aligned}
\]
The explicit representatives and integer witnesses are
\[
\begin{aligned}
w_1&=(0,0,1,1,1,0)^t,\\
w_2&=(0,0,0,0,0,0,0,0,-1,0,0,-1,0,-1,-1)^t,\\
w_3&=(0,0,0,0,0,0,0,0,0,0,0,0,1,0,1,1,0,0,0,0)^t,\\
z_1&=(0,0,0,-1,-1,-1)^t,\\
z_2&=(-1,0,1,-2,-2,-2,4,2,-1,-3,2,0,-2,0,2)^t,\\
z_3&=(-2,-2,-3,0,0,-1,-1,-1,2,2,-1,3,-1,4,3,\\
   &\hspace{23mm}-3,-2,1,-4,1)^t.
\end{aligned}
\]
They satisfy $E_kz_k=2w_k$ and $\rho_kw_k=1\bmod2$.
The nonzero local torsion class therefore maps to the nonzero global
class in each degree $q=2,3,4$.

\item Since $H^0(\mathcal S_q)=0$ in these three degrees, the
point-supported kernels give
\[
 H^0(\mathcal R_2)=\mathbb Z/2,\quad
 H^0(\mathcal R_3)=\mathbb Z/2,\quad
 H^0(\mathcal R_4)=\mathbb Z\oplus\mathbb Z/2.
\]
The Leray sequence
\cite[Tag 01F2, Lemma 20.13.4; chapter PDF p.~22]{StacksProject}
has only base columns zero, one and two.
The maps from $(0,2)$ to $(2,1)$ and from $(0,3)$ to $(2,2)$ are
the isomorphisms just proved.  The restriction of the map from
$(0,4)$ to $(2,3)$ to its torsion summand is also an isomorphism.
Writing the image of a free cusp generator as $\epsilon\in\mathbb Z/2$,
its kernel is the free group generated by $(1,\epsilon)$.
All total-degree-three terms consequently vanish.  In total degree
four the remaining terms are $\mathbb Z$ and $\mathbb Z^3$.
Their extension is free, so
\[
                   H^3(M;\mathbb Z)=0,\qquad
                   H^4(M;\mathbb Z)=\mathbb Z^4.
\]
By \cref{maps:src:F.2.5}, $\pi_1(M)=0$ and the rational homology
in degrees two and three vanishes.  Universal coefficients
\cite[Theorem 3.2, p.~195, and Corollary 3.3, p.~196]{Hatcher2002}
identify the torsion of $H_2$ and $H_3$ with the corresponding
$\operatorname{Ext}$ subgroups of $H^3$ and $H^4$; both are zero.
These finitely generated groups have rational rank zero, hence
$H_2=H_3=0$.  Integral Poincar\'e duality
\cite[Theorem 3.30, p.~241]{Hatcher2002} determines the remaining
groups as $\mathbb Z,\mathbb Z^4,\mathbb Z$ in degrees $0,4,8$ and
zero otherwise.  Finally, the Hurewicz homomorphisms
\cite[Theorem 4.32, p.~366]{Hatcher2002} give successively
$\pi_2=H_2=0$ and $\pi_3=H_3=0$.
\end{enumerate}
\end{proof}




\begin{theorem}\label{night:P1.10}
\StatementRecord{night:P1.10}{proved}{Historical events: GPT-6 Astra (OpenAI), ChatGPT Pro mode, effort Pro, S6\_06\_BLIND, 16--18 November 2026: targeted proof check. Also checked in this writing. Pre-continuation writing events: GPT-6 Astra (OpenAI), ChatGPT Pro mode, effort Pro, S6\_10\_PAPERS, 18 November 2026 02:02:02.394029+00:00: targeted proof check. Continuation mathematical and source checks: GPT-6 Astra (OpenAI), ChatGPT Pro mode, effort Pro, S6\_10\_PAPERS continuation, 19 November 2026: computation or formal certificate; independent re-proof; used and rechecked in this run. Continuation metadata events: GPT-6 Astra (OpenAI), ChatGPT Pro mode, effort Pro, S6\_10\_PAPERS continuation, 19 November 2026: used and rechecked in this run; status-vocabulary reconciliation; metadata reconciliation only; no new mathematical check. }{\cref{night:P1.11}, \cref{night:P1.9}}
For every real parameter $c>-h_{\min}$ the corrected construction gives
a compact smooth complex fourfold $M_c$ and a proper holomorphic map
to $\mathbb P^1$ with the two specified multiple fibres and the
index-one Kuhn cusp.  With its complex orientation,
$\pi_1(M_c)=1$ and $\chi(M_c)=6$.
\end{theorem}
\begin{proof}
\begin{enumerate}[label=\arabic*.]
\item Evaluation of $\Lambda$ on $F_c$ maps $\Lambda_{\mathbb R}$
isomorphically onto the underlying real bundle of $F_c^\vee$:
the kernel is zero exactly by $F_c\cap\bar F_c=0$.
The image is consequently a full discrete lattice.  Quotient by its
holomorphic period sections gives a holomorphic torus bundle on the
nonspecial base.

\item On a finite root chart $u^{m_i}=t_i$, divide the pullback torus bundle by
\[
(x,u)\longmapsto(M_ix+v_i/m_i,\zeta_{m_i}u).
\]
Freeness proves smoothness, including the central fibre.  The special
divisor is $m_i$ times the reduced free quotient.  On the punctured
root disk translation by
$-(\log u/(2\pi i))v_i$, interpreted by its period section, identifies
this affine quotient with the original Jacobian bundle.  A change of
logarithm gives an integral lattice period; replacing $u$ by
$\zeta_{m_i}u$ cancels exactly the affine term $v_i/m_i$.
This proves the holomorphic gluing and its sign.
At the cusp use P1.9 and the exponential period map.  The zero section
$\zeta=0$ extends to $z=1$ in the dense torus of the component marked
by the vertex zero.  The resulting map is proper by a finite compact
cover of the base by the three local pieces and a compact remainder
with torus fibres.  All total-space charts and their overlaps have
already been specified, so this produces a smooth compact complex
fourfold, and in particular an almost complex structure.

\item For completeness we record every van Kampen relation.  The main
bundle has its zero section and fundamental group
$\Lambda\rtimes\langle g_1,g_2\rangle$.
The finite pieces impose
$g_i^{m_i}=v_i$ and $g_ix^ug_i^{-1}=x^{M_iu}$; the logarithmic gluing
above supplies these signs.  The cusp zero section kills
$g_0=(g_1g_2)^{-1}$ and the cusp collapse kills $\operatorname{im}N$.
The latter assertion follows directly from the toric moment model in
P1.11: angular cycles $a_j$ collapse and the base cycles $b_j$ survive.
Its fundamental group is $\mathbb Z^3$, since the simply connected
complete toric components and their connected toric intersections have
dual complex $\mathbb R^3/\mathbb Z^3$; alternatively, van Kampen on
the moment cells kills the angular loops generated by the primitive
facet normals and retains exactly the three base loops.
Thus the full presentation is
\[
\left\langle \Lambda,g_1,g_2\ \middle|\
 [\Lambda,\Lambda]=1,\quad
 g_ix^ug_i^{-1}=x^{M_iu},\quad
 g_1^5=v_1,\ g_2^4=v_2,\ g_1g_2=1,\ \operatorname{im}N=0
\right\rangle.
\]
The last relation is redundant after the preceding ones but records
the actual cusp map.  Since $g_2=g_1^{-1}$, the two induced actions are
inverse; their orders divide both five and four and are therefore one.
The remaining fibre group is $\mathbb Z$ through $\ell$.  With its
generator $z$, the additive relations are
$5g_1-z=0$, $4g_2+z=0$, $g_1+g_2=0$.
Their matrix
\[
\begin{pmatrix}5&0&-1\\0&4&1\\1&1&0\end{pmatrix}
\]
has determinant $-1$, proving triviality.
Euler additivity gives $\chi(M_c)=6$: smooth tori and finite free
torus quotients have Euler number zero, and P1.9 gives the remaining
contribution.

\end{enumerate}
\end{proof}




\begin{theorem}\label{night:P1.11}
\StatementRecord{night:P1.11}{proved}{Historical events: Claude Opus 5.5 (Anthropic), effort max, sweep 1, scope s6, 18 November 2026: computation or formal certificate. GPT-6 Astra (OpenAI), ChatGPT Pro mode, effort Pro, S6\_06\_BLIND, 16--18 November 2026: targeted proof check. Also checked in this writing. Pre-continuation writing events: GPT-6 Astra (OpenAI), ChatGPT Pro mode, effort Pro, S6\_10\_PAPERS, 18 November 2026 02:02:02.394029+00:00: targeted proof check. Continuation mathematical and source checks: GPT-6 Astra (OpenAI), ChatGPT Pro mode, effort Pro, S6\_10\_PAPERS continuation, 19 November 2026: computation or formal certificate; targeted proof check; used and rechecked in this run. Continuation metadata events: GPT-6 Astra (OpenAI), ChatGPT Pro mode, effort Pro, S6\_10\_PAPERS continuation, 19 November 2026: used and rechecked in this run; status-vocabulary reconciliation; metadata reconciliation only; no new mathematical check. \emph{Continuation recorded proof scope.} The complete current eight-step proof is covered by the separately written prereading derivations and subsequent exact comparisons of the V, VI and root reviewers. They check the actual parameter transports, marked Voronoi--Kuhn fan, compatible common-character moment maps, deck and angular section, mapping-cylinder collapse, integral homology, marked duality maps, primitive image and saturated kernel. Fresh independent programs verify the integer matrices and finite face data. The former extra routes remain preserved with their individual checking scopes.}{\cref{night:P1.9}}
The index-one central cusp fibre $C$ has torsion-free integral homology
with ranks $(1,3,9,9,10,3,1)$.
The collapse $\kappa:T^6\to C$ has primitive image in every degree;
it is onto except in degree four, where its cokernel is $\mathbb Z$.
Its image ranks are $(1,3,9,9,9,3,1)$.
In every degree its kernel is the saturation of
$(\bigwedge^kM_0-I)\bigwedge^k\Lambda$.
\end{theorem}
\begin{proof}
\begin{enumerate}[label=\arabic*.]
\item We first remove the unit matrix $C(q)$ topologically, by an explicit
deformation, without changing the integral marking.  Use the family
$C_u(q)=(1-u)C(q)+u iI_3$, followed by the reverse of the same
family starting at $C(q)=0$.  The two flows compose to a marked
comparison with $C(q)=0$.  The action and properness estimates
of P1.9 are uniform on both parameter intervals.  Local quotient charts form a smooth
family of pairs, with the local equation $q=x_0x_1x_2x_3$ independent
of $u$.  Local lifts of $\partial/\partial u$, tangent to every
coordinate stratum and annihilating $q$, can be patched by a partition
of unity.  Over a smaller compact disk their flow exists for the whole
interval.  This gives a stratified diffeomorphism of the central
fibres and a diffeomorphism of the nearby fibres, preserving the
period marking.  The same construction on the real oriented blowup
preserves the collapse map.  We may therefore calculate at $C(q)=0$.

\item Here is a concrete topological model.  Put $G=4I_3-\mathbf1\mathbf1^t$.
Its positive eigenvalues are $1,4,4$.  The Voronoi polytope in dual
coordinates is cut out by
\[
|v^ty|\leq \tfrac12v^tGv
 \quad (0\ne v\in\{0,1\}^3).
\]
The value on a binary vector with $k$ nonzero entries is $k(4-k)$.
These fourteen primitive facet normals give the star of a vertex in
the Kuhn triangulation.  One checks its Delaunay tetrahedra by the
successive chains
$0,e_i,e_i+e_j,e_1+e_2+e_3$ for the six permutations $(i,j,k)$;
solving their three perpendicular-bisector equations and the fourteen
inequalities verifies precisely their common faces.  The translates
of this Voronoi polytope by $G\mathbb Z^3$ tile $\mathbb R^3$.

\item For a simplex $\sigma$ with vertices $v_0,\ldots,v_r$, put
$L_\sigma=\sum_i\mathbb Z(v_i-v_0)$ and
$D_\sigma=\bigcap_{v\in\sigma}P_v$, where $P_v=P+Gv$.
The component torus has lattice $A=\mathbb Z^3$; the common-stratum
marking is the explicit map
\[
 (A\oplus\mathbb Z)/\sum_i\mathbb Z(v_i,1)
 \longrightarrow A/L_\sigma,\qquad
 [(a,k)]\longmapsto a-kv_0.
\]
The Kuhn unimodularity makes $L_\sigma$ primitive.

We choose compatible moment maps as follows.  The polytope $2P_v$
is integral, since its vertex facet normals are unimodular and its
facet constants are integral.  Embed the component $X_v$ using the
characters $\chi^{-m}$ for
$m\in\mathcal A_v=2P_v\cap A^\vee$, and define
\[
 \mu_v([Z_m])=
 \frac{\sum_{m\in\mathcal A_v}m|Z_m|^2}
 {2\sum_{m\in\mathcal A_v}|Z_m|^2}.
\]
The sign is specified: the ray $w-v$ tends to the face on which
$\langle w-v,-m\rangle$ is minimal, hence to the upper common
Voronoi face of $P_v$ and $P_w$.  Thus the character polytope is
$-2P_v$, and $\mu_v$ is minus one-half its usual moment map.
The monomial moment theorem
\cite[equations (12.2.3)--(12.2.4) and Theorem 12.2.5, p.~571]{CoxLittleSchenck2011}
identifies the nonnegative toric part with $P_v$.
For the squared weights here, the intervening map on that part is
$[a_m]\mapsto[a_m^2]$, with inverse $[b_m]\mapsto[\sqrt{b_m}]$:
both preserve the homogeneous binomial relations and positive common
projective rescaling.  Write $s_v:P_v\to X_v$ for the resulting
continuous positive section.

On a shared face put
$\mathcal A_\sigma=2D_\sigma\cap A^\vee$ and choose
$m_0\in\mathcal A_\sigma$.
For $m\in\mathcal A_\sigma$, the difference $m-m_0$ annihilates
$L_\sigma$.  The ratio $Z_m/Z_{m_0}$ is consequently the restriction
of the same character
\[
 \bigl(-(m-m_0),\langle m-m_0,v\rangle\bigr)
       \in(A\oplus\mathbb Z)^\vee
\]
for every vertex $v$ of $\sigma$.  At $C=0$ all deck unit multipliers are one, and the toric gluing
identifies these character functions exactly.  On either component
the restricted moment map is therefore the identical expression
\[
 \frac{\sum_{m\in\mathcal A_\sigma}
       m|\chi^{-(m-m_0)}|^2}
 {2\sum_{m\in\mathcal A_\sigma}|\chi^{-(m-m_0)}|^2}.
\]
The equality persists on every smaller face by restricting the index
set.  Positive character values also agree, so the sections $s_v$
glue on the entire shared orbit closure.  The angular maps are the
displayed quotient homomorphisms $A\to A/L_\sigma$.

The period deck map is $(a,k)\mapsto(a+k\nu,k)$, sending $v$ to
$v+\nu$.  Its component marking is unchanged because
$a+k\nu-k(v+\nu)=a-kv$.
Moreover $\mathcal A_{v+\nu}=\mathcal A_v+2G\nu$; the common
character factor cancels in projective ratios.  Hence
\[
 \mu_{v+\nu}(T_\nu p)=\mu_v(p)+G\nu,
 \qquad s_{v+\nu}(y+G\nu)=T_\nu s_v(y).
\]
These give an explicit well-defined map
\[
 \begin{gathered}
 \mathcal Q=
 \bigl((\mathbb R^3/G\mathbb Z^3)\times
       (\mathbb R^3/\mathbb Z^3)\bigr)/\!\sim
 \longrightarrow C,\\
 [y,\theta]\longmapsto e^{2\pi i\theta}\,s_v(y)
 \quad(y\in P_v).
 \end{gathered}
\]
where over $\operatorname{relint}D_\sigma$ the relation collapses
the subtorus $L_\sigma\otimes(\mathbb R/\mathbb Z)$.
On each stratum the inverse takes its moment value and its angular
character class in $T_A/T_{L_\sigma}$; the latter is obtained by
$z\mapsto z/|z|$ on its nonzero characters.
The map is a continuous bijection, and its source is compact and
$C$ Hausdorff, so it is a homeomorphism.
The angular identity $[y,0]$ gives the continuous section, with
the same face and deck compatibilities just proved.

\item The actual nearby-fibre collapse is the quotient map in this model.
Give the semistable family its divisorial logarithmic structure.
The local equations $q=x_0x_1x_2x_3$ and their face restrictions
make it fine, saturated, log smooth and vertical over the standard
logarithmic disc.  The family is proper by P1.9.
In its logarithmic central fibre, retain all arguments of the $x_j$
with their sum fixed by $\arg q$.  The integral angular lattice over
a face with vertices $v_0,\ldots,v_r$ is the kernel of the height
map on $\mathbb Z^3\oplus\mathbb Z$.  Forgetting the arguments of
the zero coordinates induces exactly
\[
 \mathbb Z^3\longrightarrow
 (\mathbb Z^3\oplus\mathbb Z)/
       \sum_i\mathbb Z(v_i,1)
 \ \cong\ \mathbb Z^3/L_\sigma.
\]
The last isomorphism sends the class of $(a,k)$ to
$a-kv_0\pmod {L_\sigma}$; its inverse sends $a$ to $(a,0)$.
Thus no finite index enters this angular quotient.
The polar maps agree under the specified integral monomial overlaps.
At $C(q)=0$ the deck multipliers have positive real values, so the
logarithmic fibre at angle one is the product of the base and angular
three-tori in item~3, with their given markings.

The mapping-cylinder comparison for a proper vertical log-smooth
family identifies this logarithmic forgetful map with the actual
nearby-fibre collapse, compatibly with the central strata
\cite[formula (3-3-1), p.~478; Theorem 4.1.1 and Lemmas
4.1.2--4.1.3, pp.~480--482]{ffAchingerOgus20}.
Its angular cycles are $a_j$ and its base cycles are $b_j$.
In these column-basis conventions monodromy fixes each $a_j$
and sends $b_j$ to $b_j+a_j$, as specified by $Nb_j=a_j$.
The stratified deformation in item~1 transports this actual map
back to the original bounded period matrix.

\item Filter $\mathcal Q$ by the dimension $d$ of its dual base cells.
Over the interior of the dual cell of a $p$-simplex, $d=3-p$ and the
angular quotient has lattice
$Q_\sigma=\mathbb Z^3/L_\sigma$ of rank $d$.
Excision for the cell and its boundary gives the relative chain
groups of the resulting homology spectral sequence:
\[
E^1_{d,q}=\bigoplus_{\dim\sigma=3-d}\bigwedge^qQ_\sigma.
\]
Orient dual cells by the orientation of the base and the ordered
simplex.  The differential is the signed quotient map for each
codimension-one coface.  More explicitly, with integral projection
matrices $Q_\sigma$ and chosen right inverses $S_\sigma$,
\[
D_{p,q}[\tau,\sigma]=
 \sum_{\sigma=\partial_j\tau\ {\rm mod}\ \mathbb Z^3}
 (-1)^j\bigwedge^q(Q_\tau S_\sigma).
\tag{1.11.1}
\]
All $L_\sigma$ are primitive because the tetrahedra are unimodular.
Thus these projections and right inverses are integral.  Formula
(1.11.1) specifies every integer entry, including multiple incidences
under translation identifications.  The program
\texttt{cusp\_complex.py} enumerates the $26$ simplex classes, all
$74$ codimension-one incidences, and exports every matrix, basis, and
Smith certificate to \texttt{cusp\_complex.json}.

\item The nonzero $E^2_{d,q}$ ranks are
\[
\begin{array}{c|rrrr}
q\backslash d&0&1&2&3\\\hline
0&1&3&3&1\\
1&0&6&8&3\\
2&0&0&7&3\\
3&0&0&0&1
\end{array}
\]
and every group is free.  Support leaves only the possible higher
differential $d_2:E^2_{3,0}\to E^2_{1,1}$.
The angular identity section splits the base projection as a map of
filtered spaces, so the base fundamental class survives.  This kills
that differential.  All subsequent differentials have zero source or
target.  The filtration extensions are extensions of free abelian
groups and therefore split, proving the asserted integral homology.

\item The product-to-quotient map induces on $E^1$ the exterior powers of
the quotient projections.  The newly authored program
\path{P1_11_collapse_fresh.py} in \path{code/continuation_V/}
computes the induced maps on integral cycle and quotient bases.
Their images are primitive, and the only nonzero graded cokernel
is $\mathbb Z$ at $(d,q)=(2,2)$.  The summed image ranks are
\[
             r_k=(1,3,9,9,9,3,1).
\]
The same exact exterior-power calculation gives these ranks for
$\ker(\bigwedge^kM_0-I)$.
The actual collapse satisfies
$\kappa_*(\bigwedge^kM_0-I)=0$, because a monodromy circuit
has the same retraction to the central fibre in the mapping-cylinder
comparison.  Therefore its rank is at most $r_k$.
A filtered map has rank at least the sum of the ranks of its
graded maps; hence the actual rank is exactly $r_k$.

Choose integral bases compatible with the source and target
filtrations; their graded groups are free.  In each graded block
choose the unit Smith pivots of the computed map.  The corresponding
$r_k$ rows and columns of the actual filtered map form a block
triangular submatrix with unimodular diagonal pivot blocks.
Its determinant is $\pm1$.  Since the full map has rank $r_k$,
these unit pivots prove that its actual image is primitive.
Thus the actual map is onto in every degree except four, where its
cokernel is the free group of rank $10-9=1$.
This rank argument also rules out additional lower-filtration images
of source classes in the kernels of the graded maps.

\item For a fully marked certificate we represent a base cycle $b_J$ by
Poincar\'e duality from the Alexander--Whitney cup cocycle on the
complementary coordinates.  On an ordered $p$-simplex its coefficient
is
\[
\operatorname{sgn}(J,J^c)
 \prod_{r=1}^p(v_r-v_{r-1})_{(J^c)_r}.
\]
Tensoring with an angular exterior vector and applying (1.11.1)
gives the matrices exported by \texttt{cusp\_marked\_collapse.py}.
They annihilate $\bigwedge^kM_0-I$ exactly.  The actual collapse also
annihilates that operator, since turning the smoothing parameter
does not change its image in the central fibre.  The computed ranks
equal the ranks of the monodromy coinvariants; primitivity therefore
identifies both kernels with the saturation of its image.
Consequently the exported matrices are matrices of the actual
collapse in suitable integral target bases.  This last argument also
eliminates any ambiguity from choosing splittings of the filtration.

\end{enumerate}
\end{proof}




\begin{theorem}\label{night:P1.12}
\StatementRecord{night:P1.12}{proved}{Historical events: GPT-6 Astra (OpenAI), ChatGPT Pro mode, effort Pro, S6\_06\_BLIND, 16--18 November 2026: targeted proof check. Also checked in this writing. Pre-continuation writing events: GPT-6 Astra (OpenAI), ChatGPT Pro mode, effort Pro, S6\_10\_PAPERS, 18 November 2026 02:02:02.394029+00:00: targeted proof check. Continuation mathematical and source checks: GPT-6 Astra (OpenAI), ChatGPT Pro mode, effort Pro, S6\_10\_PAPERS continuation, 19 November 2026: computation or formal certificate; independent re-proof; targeted proof check; used and rechecked in this run. Continuation metadata events: GPT-6 Astra (OpenAI), ChatGPT Pro mode, effort Pro, S6\_10\_PAPERS continuation, 19 November 2026: used and rechecked in this run; status-vocabulary reconciliation; metadata reconciliation only; no new mathematical check. \emph{Continuation recorded proof scope.} The complete current ten-step proof has renewed prereading coverage by an explicit union: root\textquotesingle s actual evaluation, transported complex structure and finite logarithmic collar isotopy were frozen at 08:15:08 UTC; the complement, exact boundary kernels, marked collapse and Mayer--Vietoris argument were independently derived by the ledger reviewer before the old proof at 08:48:08 UTC; VI separately rederived the based mapping-space and $\operatorname{Ext}(H_4,\mathbb Z)$ routes before their isolated targets. The Euler substitution is in the ledger reviewer\textquotesingle s frozen proof. The new complement program passed 444 exact checks before target reading; 177 cone-morphism and 145 exported-basis checks then compared every retained displayed matrix, including two-torsion. Root\textquotesingle s fresh program verifies the retained evaluation and collar maps. The old complement steps and their data are retained exactly for P1.19. No value is assumed for a higher cusp suspension map. The renewed source comparison remains at the stated constructed-family scope.}{\cref{night:P1.1,night:P1.9,night:P1.10,night:P1.11}; the actual marked evaluation and finite logarithmic collar maps; \cite[Theorem 4.1.1 and Lemmas 4.1.2--4.1.3, pp.~480--482]{ffAchingerOgus20}; \cite[Theorem 3.2, p.~195; Theorem 3.30, p.~241; Theorem 4.32, p.~366]{Hatcher2002}}
The fourfold $M_c$ has torsion-free integral homology
\[
H_k(M_c;\mathbb Z)=
\begin{cases}
\mathbb Z,&k=0,8,\\
\mathbb Z^4,&k=4,\\
0,&k=1,2,3,5,6,7.
\end{cases}
\]
In particular it is three-connected.
\end{theorem}
\begin{proof}
\begin{enumerate}[label=\arabic*.]
\item On the punctured base the evaluation map is
\[
 E_c:\Lambda_{\mathbb R}\longrightarrow F_c^\vee,
 \qquad E_c(x)(\omega)=\omega(x).
\]
Its kernel is zero: a real vector annihilating $F_c$ also annihilates
$\overline{F_c}$, and these two spaces span the complexified cohomology
space. Both real dimensions are six. If $\Omega$ is the evaluation
matrix of a local frame, its inverse is explicitly
\[
 E_c^{-1}(z)=R^{-1}
       \begin{pmatrix}\operatorname{Re}z\\\operatorname{Im}z\end{pmatrix},
 \qquad
 R=\begin{pmatrix}\operatorname{Re}\Omega\\\operatorname{Im}\Omega\end{pmatrix}.
\]
Under a frame change $\Omega'=G\Omega$, both $R$ and the real coordinate
column of $z$ are multiplied by the real matrix of $G$. Thus this
inverse glues. Evaluation sends each element of $\Lambda$ to its
specified period section and induces the marked real torus bundle
isomorphism, preserving the zero section. For a peripheral transport
$\gamma$ with matrix $M$, the actual evaluation square is
\[
 E_{c,\gamma p}(Mx)=\Phi_\gamma E_{c,p}(x),
\]
where $\Phi_\gamma$ is the transport in the period quotient. Consequently
\[
 J_c=E_c^{-1}iE_c=R^{-1}J_{\rm std}R,
 \qquad
 J_{\rm std}=\begin{pmatrix}0&-I_3\\I_3&0\end{pmatrix},
 \qquad J_c^2=-I_6,
\]
and $J_c(\gamma p)=M J_c(p)M^{-1}$.

\item At a finite site put $u=r e^{2\pi ia}$ on the universal angular
cover and $L=\log r/(2\pi)$. The prescribed attaching translation in
Theorem~\ref{night:P1.10} becomes
\[
 E_c^{-1}\!\left(-\frac{\log u}{2\pi i}E_c(v)\right)
       =-av+LJ_c(u)v.
\]
The affine action is $(x,u)\mapsto(Mx+v/m,\zeta_m u)$, and $Mv=v$ by
Theorem~\ref{night:P1.1}. Write
\[
 \alpha_0(x,u)=x-av,
 \qquad h_c(u)=LJ_c(u)v,
 \qquad \alpha_c=\alpha_0+h_c.
\]
The angular shift is $a\mapsto a+1/m$, and the displayed transport of
$J_c$ gives the exact identities
\[
 h_c(\zeta_m u)=M h_c(u),\qquad
 \alpha_c(Mx+v/m,\zeta_m u)=M\alpha_c(x,u).
\]
Hence $h_c$ is a section of the real local-system vector bundle on
the attaching annulus. Choose a scalar cutoff $\chi$ which is one at
the inner boundary of an attaching collar and zero at its outer
boundary. Translation by $-\chi h_c$ on that collar has inverse
translation by $\chi h_c$; the translations by $-s\chi h_c$,
$0\leq s\leq1$, give an equivariant isotopy. At the inner boundary
it carries $\alpha_c$ to $\alpha_0$ and at the outer boundary it is
the identity. The two finite collars are disjoint.

\item Write $L_k=\bigwedge^k\Lambda$ and $W_{i,k}=\bigwedge^kM_i$.
The cellular model of the main torus bundle over a graph with two
oriented edges has
\[
G_k=L_k\oplus L_{k-1}\oplus L_{k-1},\qquad
d_G(f,x_1,x_2)=((W_{1,k-1}-I)x_1+(W_{2,k-1}-I)x_2,0,0).
\tag{1.12.1}
\]
The boundary mapping torus at $i$ has
$B_{i,k}=L_k\oplus L_{k-1}$ and
\[
d_{B_i}(f,x)=((W_{i,k-1}-I)x,0).
\tag{1.12.2}
\]
The finite boundary maps to $G$ by identity on $L_k$ and inclusion
of the suspension into its respective edge.  The cusp boundary,
corresponding to $(g_1g_2)^{-1}$, maps by
\[
j_{0,G}(f,x)=(f,-W_{2,k-1}W_{0,k-1}x,-W_{0,k-1}x).
\tag{1.12.3}
\]
This follows by traversing the second edge backwards and then the
first backwards, transporting the fibre chain to each initial end.
Indeed its boundary is $W_{0,k-1}-I$ because
\[
-(W_1-I)W_2W_0-(W_2-I)W_0=W_0-I.
\]

\item Here are full models for the finite pieces and their attaching maps.
Put $\varepsilon_i=\ell(v_i)$ and
$P_i=(\varepsilon_iv_i,e_1,\ldots,e_5)$.
This is unimodular.  It splits
$\Lambda=\mathbb Z(\varepsilon_iv_i)\oplus K$, $K=\ker\ell$,
and $P_i^{-1}M_iP_i=1\oplus R_i$.
The cap retracts onto the mapping torus of $R_i$ on $T_K$.
Its chains are
\[
E_{i,k}=\bigwedge^kK\oplus\bigwedge^{k-1}K,
\qquad d_{E_i}(f,x)=((\bigwedge^{k-1}R_i-I)x,0).
\tag{1.12.4}
\]
In these split coordinates, with $v=\varepsilon_iv_i$, the finite
boundary-to-cap chain map is specified on all summands by
\[
\begin{array}{c|c}
\text{boundary chain}&\text{cap chain}\\\hline
x\in\bigwedge^kK&(x,0)\\
v\wedge x&(0,\varepsilon_i\sum_{j=0}^{m_i-1}
                  (\bigwedge^{k-1}R_i)^jx)\\
g_i\wedge x&(0,x)\\
g_i\wedge v\wedge x&0.
\end{array}
\tag{1.12.5}
\]
Changing back uses
$\operatorname{diag}(\bigwedge^kP_i^{-1},
\bigwedge^{k-1}P_i^{-1})$.
Geometrically this map is
$(t,y,\theta)\mapsto[y,\theta+\varepsilon_im_it]$ into the mapping
torus.  The $t$-circle sweeps $m_i$ strips, giving the norm sum; the
$g_i$-circle sweeps one strip; their product maps to a single
base direction, giving zero on the last row.  This proves the map
including its sign and the possible finite-cap torsion.

\item We justify that these are integral cellular models, rather than a
formal use of invariant rational forms.  In the bases with columns
\[
K_1=\begin{pmatrix}
-1&-1&0&-1&0\\0&1&0&1&1\\0&1&0&0&0\\
0&-1&-1&0&-1\\0&0&1&0&0
\end{pmatrix},\qquad
K_2=\begin{pmatrix}
-1&0&-1&0&-1\\-1&1&0&1&-1\\0&1&-1&1&-1\\
0&-1&0&-1&1\\1&-1&1&0&1
\end{pmatrix},
\]
both unimodular, $R_1$ is the permutation of five coordinates in a
cycle, and $R_2$ is a four-cycle together with multiplication by $-1$
on the fifth coordinate.  The minimal product cell structure of
circles is preserved by these signed permutations; its cellular
boundary is zero and its action on cells is exactly the exterior
power.  Formula (1.12.5) is obtained by subdividing its swept
base interval into the indicated strips.  The common torus can be
marked into these two structures by based cellular approximations.
Their cellular maps are the exterior-power matrices because the
cellular differentials vanish.  Based maps between these aspherical
tori with a fixed map on $\pi_1$ have a contractible space of
homotopies: lifting to the universal covers gives equivariant maps
to a contractible vector space and linear interpolation relative to
the base point.  Thus the choices introduce no extra attaching class
on the graph.  Composing the two edge cylinders gives exactly
(1.12.3).  The explicit conjugators were found in the bounded search
$[-1,1]^5$ and then verified by exact multiplication in
\texttt{finite\_permutation\_models.py}.

\item Let $F_k:B_{1,k}\oplus B_{2,k}\to G_k\oplus E_{1,k}\oplus E_{2,k}$
have the main boundary maps as its first row, and the negatives of
(1.12.5) in the corresponding two cap rows.  The chain complex of
$Y$ is the mapping cone
\[
C_k(Y)=G_k\oplus E_{1,k}\oplus E_{2,k}
       \oplus B_{1,k-1}\oplus B_{2,k-1},\qquad
d_Y=\begin{pmatrix}d_G\oplus d_{E_1}\oplus d_{E_2}&F_{k-1}\\
0&-d_{B_1}\oplus-d_{B_2}\end{pmatrix}.
\tag{1.12.6}
\]
Equations (1.12.1)--(1.12.6), minors for exterior powers, and the input
matrices specify every entry of every differential.
The full arrays are exported to \texttt{complement\_export.json}; its
program verifies $d^2=0$ in all degrees and computes Smith
decompositions.  Chain ranks in degrees zero through eight are
$(3,22,71,132,155,118,57,16,2)$.
The resulting homology of $Y$ is free, with ranks
$(1,1,3,9,12,9,3,0,0)$.

\item The boundary $B_0$ has
\[
H_1(B_0)=\mathbb Z^4,\qquad H_2(B_0)=\mathbb Z^{12},\qquad
H_3(B_0)=\mathbb Z^{18}\oplus\mathbb Z/2.
\]
The homomorphism $j_Y:H_k(B_0)\to H_k(Y)$ from (1.12.3) is
surjective for $k=1,2,3$.  In degrees two and three, the program
\texttt{complement\_kernel\_check.py} supplies integral cycle bases
of its kernel whose entire suspension block in (1.12.2) is zero.
This is an exact zero matrix, not a vanishing after tensoring with
$\mathbb Q$.  Therefore these kernels are represented by cycles of
the nearby torus itself.  The restriction of the actual map to the
cusp cap on these kernels is consequently the already proved
collapse of P1.11, regardless of its value on other boundary
suspension cycles.

\item For transparency, in the Smith bases exported with these cycles the
combined map in degree one is
\[
\begin{pmatrix}20&1&0&0\\1&0&-1&-1\\1&0&0&0\\0&0&2&1\end{pmatrix},
\qquad\det=1.
\tag{1.12.7}
\]
The cusp meridian maps to zero here because the zero section extends.
In degree two $\ker j_Y=\mathbb Z^9$ and its map to $H_2(C)$ is
\[
\begin{pmatrix}
0&0&1&-1&0&-1&0&0&0\\
0&0&1&0&0&-1&0&1&0\\
0&0&1&-1&0&-2&0&1&0\\
0&0&0&0&-1&-1&0&0&-1\\
1&1&-1&0&0&1&1&0&0\\
1&2&0&0&1&0&0&0&0\\
-1&-1&0&0&0&0&0&0&0\\
0&0&0&0&1&0&0&0&0\\
0&0&0&0&0&-1&0&0&0
\end{pmatrix},\qquad\det=1.
\tag{1.12.8}
\]
In degree three $\ker j_Y=\mathbb Z/2\oplus\mathbb Z^9$.
Its collapse matrix, with the torsion generator first, is
\[
\begin{pmatrix}
0&1&0&1&0&0&0&0&0&1\\
0&1&-1&0&0&0&-1&0&-1&0\\
0&0&0&0&0&0&1&0&1&0\\
0&0&0&-1&1&-1&0&-1&0&0\\
0&0&0&0&0&0&0&1&0&0\\
0&0&0&0&0&0&1&0&0&0\\
0&0&0&0&0&1&0&1&0&0\\
0&0&1&0&1&-2&1&-1&2&0\\
0&1&0&-1&1&-2&0&-1&0&0
\end{pmatrix}.
\tag{1.12.9}
\]
The determinant of its last nine columns is one.  In particular the
torsion of the boundary and the second finite cap has been included
throughout the calculation.

\item We now apply Mayer--Vietoris to $M_c=Y\cup C_{\rm cap}$.
Because $j_Y$ is onto, subtracting a preimage of the $Y$-coordinate
identifies the cokernel of $(j_Y,j_C)$ with
$H_k(C)/j_C(\ker j_Y)$; its kernel is
$\ker j_Y\cap\ker j_C$.
Equations (1.12.7)--(1.12.9) therefore say that the combined map is an
isomorphism in degrees one and two and is onto in degree three.
The exact sequence yields $H_2(M_c)=H_3(M_c)=0$.
We already proved $\pi_1(M_c)=0$, and successive applications of
Hurewicz now give $\pi_2(M_c)=\pi_3(M_c)=0$.

\item Integral Poincar\'e duality and the universal coefficient theorem
give $H_5=H_6=H_7=0$.  They also show that $H_4$ is free: its possible
torsion would be $\operatorname{Ext}(H_4,\mathbb Z)$ in $H^5$, which
vanishes by duality with $H_3=0$.
Finally $\chi=6$ gives $\operatorname{rank}H_4=4$.
This proves all claimed integral homology groups without making an
assumption about the higher-degree cusp suspension chain term.

\end{enumerate}
\end{proof}




\begin{theorem}\label{night:P1.19}
\StatementRecord{night:P1.19}{proved}{Historical events: Claude Opus 5.5 (Anthropic), effort max, sweep 1, scope s6, 18 November 2026: computation or formal certificate. GPT-6 Astra (OpenAI), ChatGPT Pro mode, effort Pro, S6\_06\_BLIND, 16--18 November 2026: targeted proof check. Also checked in this writing. Pre-continuation writing events: GPT-6 Astra (OpenAI), ChatGPT Pro mode, effort Pro, S6\_10\_PAPERS, 18 November 2026 02:50:20.091270+00:00: computation or formal certificate; targeted proof check. Continuation metadata events: GPT-6 Astra (OpenAI), ChatGPT Pro mode, effort Pro, S6\_10\_PAPERS continuation, 19 November 2026: used and rechecked in this run; status-vocabulary reconciliation; metadata reconciliation only; no new mathematical check. }{\cref{night:P1.11}, \cref{night:P1.12}, \cref{night:P1.9}}
There is an independent computation from the full global cellular
complex.  Its chain ranks in degrees zero through eight are
\[
 (4,26,87,162,200,156,79,23,3),
\]
and its differentials, after the explicit integral chain equivalences
below, are the matrices of case one in
\path{code/part1/global_ext_cases.json}.
Their integral Smith calculation gives
\[
 H_k(M;\mathbb Z)=
 \begin{cases}
 \mathbb Z,&k=0,8,\\
 \mathbb Z^4,&k=4,\\
 0,&\text{otherwise}.
 \end{cases}
\]
This calculation uses local Poincar\'e--Lefschetz duality of the
smooth cusp cap to determine its secondary attaching class.  It does
not use global Poincar\'e duality or the homology conclusion of P1.12.
In particular, the map obtained by setting the whole cusp suspension
chain term to zero is not the actual relative cusp model; its precise
missing term is accounted for here.
\end{theorem}
\begin{proof}
\begin{enumerate}[label=\arabic*.]
\item We retain the actual complement complex and finite-cap maps of P1.12.
The issue left by that argument was the full cusp boundary map.
We first determine its homology maps, then its secondary chain term,
and finally give the explicit integral reduction used by the program.

\item Write the cusp lattice as $A\oplus B$, with ordered generators
$a_1,a_2,a_3,b_1,b_2,b_3$ and $Nb_i=a_i$.
The boundary is the mapping torus of
\[
 (b,a)\longmapsto(b,a+b)
\]
on $T_B^3\times T_A^3$.
For the pure toric model, its retraction to the central moment model
has the following explicit description.  On the dual cell of a Kuhn
simplex $\sigma$, choose any vertex $v$ of $\sigma$ and put
\[
 (b,a,\theta)\longmapsto
   [b,a+\theta(b-v)],\qquad0\leq\theta\leq1.
\tag{1.19.1}
\]
The brackets denote the angular quotient by $L_\sigma$.
Changing $v$ changes the displayed angular coordinate by a real
multiple of a vector in $L_\sigma$, so the map is unchanged.
Replacing $b$ by $b+n$ also replaces $v$ by $v+n$; hence the formula
is periodic.  At $\theta=1$ the integer translation $v$ disappears
on the angular torus, and the map agrees with the cusp monodromy
identification.  These observations prove that (1.19.1) is well-defined
on the boundary mapping torus.  In a monomial chart of P1.9 it is
exactly the operation of forgetting the arguments of the vanishing
coordinates, with the argument of $q$ recorded by $\theta$.
Thus it is the actual boundary retraction.  The deformation removing
the bounded period units in P1.11 transports this map by a homotopy
and preserves its homology class.

\item We prove that the suspension classes have zero image in cusp homology.
Let $D$ be the derivation of $\bigwedge^\bullet(A\oplus B)$ extending
$Nb_i=a_i$, $Na_i=0$.  Then $\bigwedge M_0=\exp D$ and
$\ker(\bigwedge M_0-I)=\ker D$ over $\mathbb Q$.
An integral vector in this kernel has the same property.
The summand with $p$ base generators and $q$ angular generators is
sent by $D$ to the summand with $(p-1,q+1)$; consequently each
homogeneous summand of an invariant vector is separately killed.
Let $E$ replace $a_i$ by $b_i$. In the orthonormal exterior basis
$E=D^t$ and $[D,E]=(q-p)I$ on that summand.  For $Dv=0$,
\[
 (q-p)\|v\|^2
 =\langle[D,E]v,v\rangle=\|Ev\|^2\geq0.
\]
Thus an invariant $k$-cycle has base degree at most $\lfloor k/2\rfloor$.

\item This assertion about cycles can be realized geometrically, rather
than only in a rational associated graded group.  Give $T_B^3$ its
minimal product CW structure and put
$F_p=T_A^3\times(T_B^3)^{(p)}$.
Its cellular differential is zero, so
$H_*(F_p;\mathbb Z)\to H_*(T^6;\mathbb Z)$ is injective.
The shear preserves $F_p$ and admits cellular approximation through
maps preserving this filtration.  The corresponding mapping-cylinder
cells give the integral complex
$\operatorname{Cone}(\bigwedge M_0-I)$.
An invariant cellular vector of base degree at most $p$ therefore
has a suspension lift whose entire projection to $T_B^3$ has dimension
at most $p$.  Equivalently, its closing boundary already bounds in
$F_p$, by the just-proved injectivity.  This construction fixes the
suspension representatives used in the Wang complex and is compatible
with the based product-cell comparison in P1.12.

\item Formula (1.19.1) preserves $b$, so the image of such a lift also has
base projection of dimension at most $p$.  It represents filtration
at most $p$ in the dual-cell filtration of P1.11.  Here is a direct
cellular justification of this last step.  In every dual cell of
dimension greater than $p$, choose an interior point outside that
projection.  Radial retraction from this point pushes the projected
cycle to the cell boundary.  On the central moment model it lifts
by keeping the angular coordinate fixed; on the boundary the angular
quotient is the canonical further quotient.  The lift is continuous
and fixes the boundary.  Performing these retractions in descending
cell dimension gives a homologous cycle on the dual $p$-skeleton.
There are finitely many cells modulo periods, so this operation
terminates.

\item The integral $E_2=E_\infty$ table of P1.11 has nonzero entries only
at angular degree $q\leq d$, where $d$ is the base-cell degree.
For a suspension of a $k$-cycle, the total degree is $k+1$, whereas
$d\leq p\leq\lfloor k/2\rfloor$. Its angular degree would therefore
be $k+1-d>d$, where that table is zero.  The suspension image
vanishes.  On the fibre summand the map is precisely the marked
collapse $\kappa_k$ already exported in
\texttt{cusp\_marked\_collapse.json}.  Thus the homology maps agree
with those of the chain map $f_0=(\kappa,0)$, in every degree.

\item Agreement on homology does not identify these two chain maps.
We now determine the remaining difference.  Integral Smith reduction
splits a finite free complex into free homology generators, contractible
unit pairs, and two-term pairs with differential equal to a torsion
invariant factor.  The cusp homology is free.  The boundary homology
has just one torsion group, $\mathbb Z/2$ in degree three.
Consequently, after maps on homology have been fixed, the difference
of two maps from the boundary complex to the cusp complex is
represented by one vector
\[
 \varepsilon\in H_4(C;\mathbb Z)/2H_4(C;\mathbb Z)
                 \cong(\mathbb Z/2)^{10}.
\tag{1.19.2}
\]
This follows directly by applying a map to the upper generator of
the two-term pair $\mathbb Z\xrightarrow{2}\mathbb Z$; changing a
chain homotopy changes its image by twice an integral vector.
There are no other secondary terms.  Thus no rational formality or
unproved equivariant formality is used at this point.

\item In the marked cusp degree-four basis, the primitive functional
\[
 \psi=(0,0,0,1,-1,1,1,-1,-1,1)
\tag{1.19.3}
\]
annihilates $\kappa_4$ and identifies its cokernel with $\mathbb Z$.
The exact multiplication $\psi\kappa_4=0$ and the primitive Smith
factors are exported.  Setting $\varepsilon=0$ gives relative
cusp homology $H_4(C,\partial)=\mathbb Z^{10}\oplus\mathbb Z/2$,
as computed in \texttt{cusp\_relative\_candidate.json}.
More intrinsically, the relative exact sequence in this degree is
an extension of $\mathbb Z^9\oplus\mathbb Z/2$ by the cokernel
$\mathbb Z$ detected by $\psi$.  Its torsion extension class is
exactly $\psi\varepsilon$ modulo two.

\item The actual cusp neighborhood is a compact oriented smooth
8-manifold with boundary: take the inverse image of a sufficiently
small closed disc, whose boundary circle avoids the critical value.
It retracts onto $C$.  Poincar\'e--Lefschetz duality and the already
computed free cusp homology give
\[
 H_4(\text{cusp cap},\partial;\mathbb Z)
       \cong H^4(C;\mathbb Z)=\mathbb Z^{10}.
\]
The preceding extension must therefore be nonsplit on its torsion
summand.  This determines the required secondary class for the global
calculation:
\[
                 \psi\varepsilon=1\pmod2.
\tag{1.19.4}
\]
This also proves the stated failure of the wholly zero suspension
chain map as a relative cusp model.

\item We next give the integral equivalence which reduces every class
satisfying (1.19.4) to one displayed matrix.  Write the degree-four
homology gluing map in the plus-sign convention as
\[
 j_4:H_4(\partial)\longrightarrow H_4(Y)\oplus H_4(C).
\]
Multiplication by $-1$ on the cusp chain summand is the chain
isomorphism from the minus-sign convention in the exported whole
diagram to this convention.  The matrix $j_4$ is $22$ by $18$,
has rank eighteen, and all its nonzero Smith factors are one.
Its cokernel is free of rank four. In its exported cokernel basis,
the map from $H_4(C)$ is
\[
 \begin{pmatrix}0\\0\\0\\\psi\end{pmatrix}.
\tag{1.19.5}
\]
This is a direct integer matrix calculation, independent of the
homology of $M$.

\item Let $e_4$ denote the fourth coordinate vector of $\mathbb Z^{10}$,
so $\psi e_4=1$.  The nine columns
$e_i-\psi_i e_4$, $i\ne4$, form an integral basis $K$ of $\ker\psi$;
$[e_4,K]$ is unimodular.  The file
\texttt{cusp\_ext\_reduction.json} gives integral lifts $W$ satisfying
\[
 j_4W=\begin{pmatrix}0\\K\end{pmatrix}.
\]
It also gives actual boundary cycles $X$ representing these lifts,
and integral complement chains $T$ with
\[
 d_YT=j_YX,\qquad f_0X=K.
\tag{1.19.6}
\]
Every entry is checked by multiplication in the program.

\item For clarity, the chain automorphisms implementing the reduction
are equally explicit.  Let $r$ be the exported integral degree-four
cocycle representing the torsion generator in $H^4(\partial;\mathbb Z)$.
In the ordered 35-dimensional boundary chain group it is
$r=-e_{17}^{t}-e_{18}^{t}$.  It satisfies
\[
 r d_{\partial,5}=0,\qquad rX=0,
 \qquad 2r=t\,d_{\partial,4},
\tag{1.19.7}
\]
where the integral row $t$ is exported as well.
For a column $v\in\mathbb Z^9$, the endomorphism
$\eta=(Xv)r$ of degree-four boundary chains, zero in all other
degrees, is a chain map and has $\eta^2=0$.
Thus $I+\eta$ is an integral chain automorphism with inverse
$I-\eta$. Precomposing the attaching map by it changes the
secondary vector by $Kv$, up to the explicitly given homotopy
$(Tv)r$ in the complement, by (1.19.6).
The last equation in (1.19.7) removes an even multiple of $e_4$
by the homotopy $e_4t$.  Decomposing $\varepsilon$ in the
unimodular basis $[e_4,K]$, equation (1.19.4) therefore reduces
the actual global attaching map to the canonical correction
\[
 f_4=f_{0,4}+e_4r,
 \qquad f_k=f_{0,k}\quad(k\ne4).
\tag{1.19.8}
\]
These formulas prove the chain equivalence; they do not merely
compare the ranks of the two complexes.

\item The program \texttt{global\_ext\_cases.py} inserts (1.19.8) into
the complete main-piece/three-cap mapping cone, retaining every
finite-cap differential and every boundary map of P1.12.
It checks $d^2=0$ over $\mathbb Z$ and exports every differential.
Its odd case is the actual global model by the preceding argument.
The even case is also computed as a diagnostic, and happens to give
the same global homology; it is not used to identify the local cap.
The independent program \texttt{global\_cell\_certificate.py}
exports the Smith left and right matrices, kernel bases, incoming
matrices in those bases, and their Smith left and right matrices
for the odd case.  The differential ranks in degrees one through
eight are
\[
                 (3,23,64,98,98,58,21,2).
\]
All nonzero incoming Smith factors in the cycle lattices are one.
Subtracting incoming and outgoing ranks from the displayed chain
ranks gives $(1,0,0,0,4,0,0,0,1)$, with no torsion.
This completes the full cellular comparison without using the
global duality conclusion of P1.12.

\end{enumerate}
\end{proof}


\section{The genus-two pencil and the lifted period planes}



\begin{theorem}\label{night:P1.4}
\StatementRecord{night:P1.4}{proved}{Historical events: GPT-6 Astra (OpenAI), ChatGPT Pro mode, effort Pro, S6\_06\_BLIND, 16--18 November 2026: targeted proof check. Also checked in this writing. Pre-continuation writing events: GPT-6 Astra (OpenAI), ChatGPT Pro mode, effort Pro, S6\_10\_PAPERS, 18 November 2026 02:02:02.394029+00:00: computation or formal certificate; targeted proof check. Continuation mathematical and source checks: GPT-6 Astra (OpenAI), ChatGPT Pro mode, effort Pro, S6\_10\_PAPERS continuation, 19 November 2026: computation or formal certificate; targeted proof check. Continuation metadata events: GPT-6 Astra (OpenAI), ChatGPT Pro mode, effort Pro, S6\_10\_PAPERS continuation, 19 November 2026: used and rechecked in this run; status-vocabulary reconciliation; metadata reconciliation only; no new mathematical check. }{Definitions and the displayed constructions}
The genus-two pencil has one nodal fibre away from $0,\infty$, at
$\lambda_*=-256/3125$.  Good reduction at $0,\infty$ requires root
orders $8,5$.  On this root orbifold its Hodge bundle is
$\mathcal H=T^3\oplus T$, where
\[
T=\mathcal O\bigl((2/5)[\infty]-(3/8)[0]\bigr),\qquad
\deg T=1/40,
\]
and $T^3$ is the line generated meromorphically by $dx/y$.
\end{theorem}
\begin{proof}
\begin{enumerate}[label=\arabic*.]
\item The discriminant of $x^5+\lambda^3x+\lambda^4$ is
$\lambda^{15}(3125\lambda+256)$.
At $\lambda_*$ the double root is $64/625$;
the second $x$-derivative and first $\lambda$-derivative there are
nonzero.  The remaining roots are simple.  The odd-degree
hyperelliptic compactification is smooth at its branch point at
infinity.  This gives one node, a smooth total surface at that node,
and genus-one normalization.

\item At zero set $\lambda=u^8$, $x=u^6X$, $y=u^{15}Y$; at infinity set
$\lambda=u^{-5}$, $x=u^{-4}X$, $y=u^{-10}Y$.
The resulting smooth equations and differentials are
\[
\begin{array}{c|c|c|c}
 &\text{equation}&dx/y&x\,dx/y\\\hline
0&Y^2=X^5+X+u^2&u^{-9}dX/Y&u^{-3}X\,dX/Y\\
\infty&Y^2=X^5+uX+1&u^6dX/Y&u^2X\,dX/Y.
\end{array}
\]
Their central polynomials are square-free.  The deck maps at zero and
infinity are $(X,Y)\mapsto(iX,\zeta_8Y)$ and
$(X,Y)\mapsto(\zeta_5^{-1}X,Y)$, respectively.
The eigenvalues on $H^{1,0}$ are $(\zeta_8,\zeta_8^3)$ and
$(\zeta_5^{-1},\zeta_5^{-2})$.  Including conjugates gives characteristic
polynomials $T^4+1$ and $\Phi_5$, proving minimality of the root orders.
The fourth power at zero is the hyperelliptic involution, acting as
$-I$ on $H^1$ and as $I$ on its exterior square.

\item The two meromorphic differential generators have orbifold divisors
$(6/5)[\infty]-(9/8)[0]$ and
$(2/5)[\infty]-(3/8)[0]$.
At the node they are independent sections of the dualizing Hodge
extension: one has opposite nonzero residues at the two points of the
normalization, and the linear combination vanishing at the node is a
regular nonzero elliptic differential.  Thus there is no additional
divisor at $\lambda_*$.  These computations prove the asserted direct
sum and its degrees.  The discriminant and all substitutions are
checked in \texttt{pencil\_charts.py/.out}.

\end{enumerate}
\end{proof}




\begin{theorem}\label{night:P1.5}
\StatementRecord{night:P1.5}{proved}{Historical events: GPT-6 Astra (OpenAI), ChatGPT Pro mode, effort Pro, S6\_06\_BLIND, 16--18 November 2026: targeted proof check. Also checked in this writing. Pre-continuation writing events: GPT-6 Astra (OpenAI), ChatGPT Pro mode, effort Pro, S6\_10\_PAPERS, 18 November 2026 02:02:02.394029+00:00: targeted proof check. Continuation mathematical and source checks: GPT-6 Astra (OpenAI), ChatGPT Pro mode, effort Pro, S6\_10\_PAPERS continuation, 19 November 2026: independent re-proof; used and rechecked in this run. Continuation metadata events: GPT-6 Astra (OpenAI), ChatGPT Pro mode, effort Pro, S6\_10\_PAPERS continuation, 19 November 2026: used and rechecked in this run; status-vocabulary reconciliation; metadata reconciliation only; no new mathematical check. }{Definitions and the displayed constructions}
The formula in the prompt with $P$ the image of
$\ell_0\wedge\mathcal H$, and $\mathcal H$ the rank-two Hodge bundle,
cannot produce a rank-three subbundle.  The required rank-three image
is instead the image of $\ell_0\wedge\mathcal H_{\mathrm{dR}}$, where
$\mathcal H_{\mathrm{dR}}$ is the full rank-four de Rham bundle.
\end{theorem}
\begin{proof}
\begin{enumerate}[label=\arabic*.]
\item The kernel of $\ell_0\otimes\mathcal H\to\bigwedge^2\mathcal H$
is $\ell_0\otimes\ell_0$, so its image has rank one.  Projection and
addition of a map into the invariant line cannot increase its domain
rank.  For the corrected domain there is an injective holomorphic map
\[
\ell_0\otimes(\mathcal H_{\mathrm{dR}}/\ell_0)
 \longrightarrow\bigwedge^2\mathcal H_{\mathrm{dR}},\quad
u\otimes[v]\longmapsto u\wedge v.
\]
Every nonzero image vector is decomposable and has wedge square zero;
$\omega^2\ne0$.  Projection modulo $\mathcal O\omega$ is therefore
injective on its rank-three image.  This specifies the comparison map
and the corrected bundle used below.

\end{enumerate}
\end{proof}




\begin{theorem}[the geometric rank-five quotient]\label{maps:src:F.3.5}
\StatementRecord{maps:src:F.3.5}{proved}{Historical events: Claude Fable 5.2 (Anthropic), effort high, 16 November 2026: targeted proof check. GPT-6 Astra (OpenAI), ChatGPT Pro mode, effort Pro, S6\_09\_PAPERS, 17 November 2026: targeted proof check. Also checked in this writing. Pre-continuation writing events: GPT-6 Astra (OpenAI), ChatGPT Pro mode, effort Pro, S6\_10\_PAPERS, 18 November 2026 02:02:02.394029+00:00: targeted proof check. Continuation mathematical and source checks: GPT-6 Astra (OpenAI), ChatGPT Pro mode, effort Pro, S6\_10\_PAPERS continuation, 19 November 2026: computation or formal certificate; targeted proof check; used and rechecked in this run. Continuation metadata events: GPT-6 Astra (OpenAI), ChatGPT Pro mode, effort Pro, S6\_10\_PAPERS continuation, 19 November 2026: used and rechecked in this run; status-vocabulary reconciliation; metadata reconciliation only; no new mathematical check. \emph{Pre-continuation recorded proof scope.} GPT-6 Astra (OpenAI, ChatGPT Pro mode), effort Pro, 2026-11-17; GPT-6 Astra (OpenAI, ChatGPT Pro mode), S6\_06\_BLIND, effort Pro, 2026-11-17; Claude Fable 5.2 (Anthropic), effort high, 2026-11-16; also checked in this writing by GPT-6 Astra (OpenAI), ChatGPT Pro mode, effort Pro, 18 November 2026}{the hypotheses, notation and definitions stated here; \cref{maps:src:F.2.5}}


Put $\lambda_*=-256/3125$ and
$U=\mathbb P^1\setminus\{0,\infty,\lambda_*\}$.  The smooth projective
models of
\[
 C_\lambda:\quad y^2=x^5+\lambda^3x+\lambda^4
 \tag{F.7}\label{maps:eq:F.7}
\]
give a polarized rank-four rational cohomology system $\mathcal H$ on
$U$.  Its primitive exterior square is rationally isomorphic to
the rank-five quotient of the system with matrices $A,B$ in
Theorem \ref{maps:src:F.2.5} by its invariant first coordinate line.
On the orbifold $\mathcal B'$ obtained by taking roots of orders
five at infinity and eight at zero, the extended Hodge bundle is
\[
 \mathcal E=\mathcal T^3\oplus\mathcal T,\qquad
 \mathcal T=\mathcal O_{\mathcal B'}(2D_\infty-3D_0),\qquad
 \deg\mathcal T=\frac1{40}.
 \tag{F.8}\label{maps:eq:F.8}
\]
Here $D_\infty,D_0$ are the reduced root divisors; the extension at
$\lambda_*$ is the nilpotent canonical extension.

\end{theorem}


\begin{proof}
\begin{enumerate}[label=\arabic*.]
\item The discriminant of the degree-five polynomial in \eqref{maps:eq:F.7} is exactly
$\lambda^{15}(3125\lambda+256)$.  At $\lambda_*$ its unique double
root is $r=64/625$ and the coefficient of $(x-r)^2$ is
$10r^3\ne0$.  Its other roots are simple.  Thus the fibre has one
ordinary node and irreducible genus-one normalization.  The local
model is a smoothing of a nonseparating node; the vanishing circle
is a primitive nonseparating homology class.  Cutting and regluing
an annular neighborhood proves that its monodromy is a single
Dehn twist, and hence a primitive rank-one transvection on $H^1$.
In particular its logarithm $N_H$ has square zero and rank one.

\item At zero, the substitutions $\lambda=s^8$, $x=s^6X$, $y=s^{15}Y$
give
\[
 Y^2=X^5+X+s^2.
\]
This is smooth at $s=0$.  The generator $s\mapsto\zeta_8s$ at
fixed $(x,y)$ acts by $(X,Y)\mapsto(iX,\zeta_8Y)$, so its
eigenvalues on the two holomorphic differentials are
$\zeta_8,\zeta_8^3$.  The full rational characteristic polynomial
is $g(T)=T^4+1$.
At infinity put $\lambda=s^{-5}$, $x=s^{-4}X$, $y=s^{-10}Y$.
The resulting smooth family is
\[
 Y^2=X^5+sX+1.
\]
The deck generator acts by $(X,Y)\mapsto(\zeta_5^4X,Y)$,
with holomorphic eigenvalues $\zeta_5^4,\zeta_5^3$.  Its full
characteristic polynomial is $f(T)=T^4+T^3+T^2+T+1$.
Inverses have the same two rational characteristic polynomials.

\item We identify the representation, rather than merely its spectra.
Choose positive peripheral generators $T_1,T_2$ at infinity and
zero, in that order, with $T_1T_2$ inverse to nodal monodromy.
Set $R=T_1^{-1}$ and $S=T_2$.  Then $R-S$ has rank one and their
characteristic polynomials are $f,g$.  The companion presentation
is Levelt's theorem in \cite[Theorem 3.5, pp.~331--332;
Corollary 3.6, p.~332]{BeukersHeckman1989}.  Its generators are
related to the present pair by
\[
 h_\infty=R=T_1^{-1},\qquad h_0=S^{-1}=T_2^{-1},\qquad
 h_1=R^{-1}S=T_1T_2,\qquad h_\infty h_1h_0=I.
\]
The following calculation supplies the actual rational change
of basis for this presentation.
If $u$ spans $\operatorname{im}(R-S)$, it is cyclic for $R$:
the span of its iterates is stable under both $R$ and $S$, and
the irreducibility of $f$ makes that span all of $\mathbb Q^4$.
In the rational basis $u,Ru,R^2u,R^3u$, the matrix of $R$ is
the companion matrix of $f$, and $S=R-u\varphi$ for a covector
$\varphi$.  The determinant identity
\[
 g(T)=f(T)+\varphi\operatorname{adj}(TI-R)u
 \tag{F.9}\label{maps:eq:F.9}
\]
determines $\varphi$ uniquely: the four coordinates of
$\operatorname{adj}(TI-R)u$ are a basis of the polynomials of
degree at most three, as is seen from their successively distinct
leading degrees in the cyclic basis.  It follows that any two
such pairs are simultaneously rationally conjugate.  Apply the
lemma to the explicit companion pair
\[
 C_f=\begin{pmatrix}
 0&0&0&-1\\1&0&0&-1\\0&1&0&-1\\0&0&1&-1
 \end{pmatrix},\qquad
 C_g=\begin{pmatrix}
 0&0&0&-1\\1&0&0&0\\0&1&0&0\\0&0&1&0
 \end{pmatrix}.
 \tag{F.10}\label{maps:eq:F.10}
\]
Their difference has rank one and their characteristic polynomials
are $f,g$.  More explicitly, put
\[
 \begin{aligned}
 \mathsf B&=(u,Ru,R^2u,R^3u),\\
 d&=(C_f-C_g)e_4=(0,-1,-1,-1)^t,\\
 \mathsf B_c&=(d,C_fd,C_f^2d,C_f^3d).
 \end{aligned}
\]
Here $\det\mathsf B_c=1$.  Uniqueness of the covector in
\eqref{maps:eq:F.9} gives the actual intertwiner
$\mathsf P=\mathsf B\mathsf B_c^{-1}$, with
$R\mathsf P=\mathsf P C_f$ and $S\mathsf P=\mathsf P C_g$.
Thus $(T_1,T_2)$ is rationally conjugate to
$(C_f^{-1},C_g)$, with the chosen peripheral conventions.

\item Use the lexicographic basis in $\bigwedge^2\mathbb Q^4$.
Its common invariant vector is $(1,1,1,1,1,1)^t$.
Let $Q$ be the identity matrix of size six with first column
replaced by this vector.  It is unimodular.  Direct multiplication
gives the exact identities
\[
 \left(Q^{-1}\bigwedge^2(C_f^{-1})Q\right)_{2:6,2:6}
   =A_{2:6,2:6},\qquad
 \left(Q^{-1}\bigwedge^2(C_g)Q\right)_{2:6,2:6}
   =B_{2:6,2:6}.
 \tag{F.11}\label{maps:eq:F.11}
\]
They are checked over $\mathbb Z$ in
\texttt{fourfolds\_periods.py}, not inferred from eigenvalues.
The invariant line is the inverse of the polarization, up to a
nonzero rational scalar.  Projection along that line identifies
the quotient with the primitive exterior square $\mathcal W$.

\item The differentials $\omega_1=dx/y$ and $\omega_2=x\,dx/y$ span
the extended Hodge bundle everywhere.  At the node they are
independent dualizing differentials: $\omega_1$ has nonzero
opposite residues on the normalization and
$(x-r)\omega_1$ is a nonzero holomorphic differential there.
Here the extension assertion can be checked on the nodal
annulus, without invoking an additional extension theorem.
In coordinates $uv=t$ a relative dualizing differential has
a convergent Laurent expansion
$\sum_n a_n(t)u^n\,du/u$.  Integration on the vanishing circle
is $2\pi i a_0(t)$.  Integration across the annulus has the
single logarithmic term $a_0(t)\log t$; integration of every
nonconstant Laurent monomial and of the complementary fixed
arcs gives a holomorphic remainder.  Thus the period vector
has the form $\exp(sN_H)$ times a holomorphic vector,
$s=\log t/(2\pi i)$.  Untwisting by $\exp(-sN_H)$ is exactly
the nilpotent canonical extension used here.  The two independent
dualizing limits remain a basis of its Hodge subbundle.
Their isotropic pairing, zero away from zero, remains zero
at the limit by holomorphic continuation, so that limit is
Lagrangian.
At zero their expressions in the root coordinates are
$s^{-9}dX/Y$ and $s^{-3}X\,dX/Y$; at infinity they are
$s^6dX/Y$ and $s^2X\,dX/Y$.
Thus their divisor lines are
$\mathcal O(6D_\infty-9D_0)$ and
$\mathcal O(2D_\infty-3D_0)$, proving \eqref{maps:eq:F.8}.
The nonzero residue of $\omega_1$ also shows
\[
 N_H\ell_\infty\ne0,
 \qquad \ell=\mathcal T^3\subset\mathcal E,
 \tag{F.12}\label{maps:eq:F.12}
\]
where $\ell_\infty$ denotes its limit at the nodal cusp.
Indeed integration on the vanishing circle tends to
$2\pi i$ times that residue; the transvection sends the cohomology
class of a differential to that period times its vanishing
cohomology vector.

\end{enumerate}
\end{proof}




\begin{theorem}\label{night:P1.2}
\StatementRecord{night:P1.2}{proved}{Historical events: GPT-6 Astra (OpenAI), ChatGPT Pro mode, effort Pro, S6\_06\_BLIND, 16--18 November 2026: targeted proof check. Also checked in this writing. Pre-continuation writing events: GPT-6 Astra (OpenAI), ChatGPT Pro mode, effort Pro, S6\_10\_PAPERS, 18 November 2026 02:43:58.020152+00:00: targeted proof check. Continuation mathematical and source checks: GPT-6 Astra (OpenAI), ChatGPT Pro mode, effort Pro, S6\_10\_PAPERS continuation, 19 November 2026: computation or formal certificate; independent re-proof; used and rechecked in this run. Continuation metadata events: GPT-6 Astra (OpenAI), ChatGPT Pro mode, effort Pro, S6\_10\_PAPERS continuation, 19 November 2026: used and rechecked in this run; status-vocabulary reconciliation; metadata reconciliation only; no new mathematical check. }{Definitions and the displayed constructions}
The lower blocks of $A,B$ are precisely the quotient representations of
$\bigwedge^2(C_f^{-1}),\bigwedge^2C_g$ in the supplied basis $Q$.
The integral quotient is the dual of the primitive exterior-square
lattice.  The natural map from the primitive lattice into that quotient
has cokernel $\mathbb Z/2$; these two integral representations are not
isomorphic.  Their rationalizations are isomorphic by that same map.
\end{theorem}
\begin{proof}
\begin{enumerate}[label=\arabic*.]
\item The two matrices obtained from their two-by-two minors are
\[
Q^{-1}\bigwedge^2(C_f^{-1})Q=
\begin{pmatrix}
1&-1&0&1&0&0\\0&1&-1&-1&1&0\\0&1&0&-1&0&0\\
0&2&-1&-1&0&1\\0&2&0&-1&0&0\\0&1&1&-1&0&0
\end{pmatrix},
\]
\[
Q^{-1}\bigwedge^2C_gQ=
\begin{pmatrix}
1&0&1&0&0&0\\0&0&-1&0&1&0\\0&0&-1&0&0&1\\
0&0&-1&0&0&0\\0&1&-1&0&0&0\\0&0&-1&1&0&0
\end{pmatrix}.
\]
Their lower five-by-five blocks are the asserted blocks, entry by entry.
Put $L=\bigwedge^2\mathbb Z^4$, let $(\ ,\ )$ be its unimodular wedge
pairing, and set $\omega=(1,1,1,1,1,1)^t$.
Then $(\omega,\omega)=2$ and the primitive sublattice
$P_{\mathbb Z}=\omega^\perp$ is the kernel of
$(1,-1,1,1,-1,1)$.  In the basis with columns
\[
\begin{pmatrix}
1&-1&-1&1&-1\\1&0&0&0&0\\0&1&0&0&0\\
0&0&1&0&0\\0&0&0&1&0\\0&0&0&0&1
\end{pmatrix},
\]
the inclusion followed by projection
$j:P_{\mathbb Z}\longrightarrow L/\mathbb Z\omega$ has matrix
\[
j=\begin{pmatrix}
0&1&1&-1&1\\-1&2&1&-1&1\\-1&1&2&-1&1\\
-1&1&1&0&1\\-1&1&1&-1&2
\end{pmatrix},\qquad \det j=2.
\]
Its Smith factors are $(1,1,1,1,2)$.  The unimodular pairing on $L$
identifies $L/\mathbb Z\omega$ with $P_{\mathbb Z}^\vee$;
under this identification $j$ is the pairing map.

\item We rule out every possible integral intertwiner.  Denote the lower blocks above by $A_5,B_5$.
The Krylov matrix $(e_3,A_5e_3,A_5^2e_3,A_5^3e_3,A_5^4e_3)$
has determinant $-1$, so $A_5$ is cyclic.  Any rational endomorphism
commuting with $A_5$ therefore has the form
$u_0I+u_1A_5+\cdots+u_4A_5^4$.
The entries $(1,2),(1,4),(2,4),(3,1)$ of its commutator with $B_5$
give, on $(u_1,u_2,u_3,u_4)^t$, the coefficient matrix
\[
\begin{pmatrix}
0&1&0&1\\1&-1&0&0\\1&-1&-1&1\\1&1&0&1
\end{pmatrix},\qquad \det=-1.
\]
Thus $u_1=u_2=u_3=u_4=0$, proving that the simultaneous rational
centralizer is exactly $\mathbb QI$.
The natural map $j$ above is equivariant and rationally invertible.
For any rational equivariant map $f:P_{\mathbb Q}\to
(L/\mathbb Z\omega)_{\mathbb Q}$, the composite $fj^{-1}$ lies
in this centralizer.  Consequently $f=cj$ for some $c\in\mathbb Q$.
An integral isomorphism would have determinant $\pm1$, whereas
$\det(cj)=2c^5$.  Its $2$-adic valuation would give
$1+5v_2(c)=0$, impossible because $v_2(c)\in\mathbb Z$.
This proves the asserted failure of integral isomorphism and retains
$j$ itself as the explicit rational isomorphism.

\end{enumerate}
\end{proof}




\begin{theorem}\label{night:P1.7}
\StatementRecord{night:P1.7}{proved}{Historical events: GPT-6 Astra (OpenAI), ChatGPT Pro mode, effort Pro, S6\_06\_BLIND, 16--18 November 2026: targeted proof check. Also checked in this writing. Pre-continuation writing events: GPT-6 Astra (OpenAI), ChatGPT Pro mode, effort Pro, S6\_10\_PAPERS, 18 November 2026 02:02:02.394029+00:00: targeted proof check. Continuation mathematical and source checks: GPT-6 Astra (OpenAI), ChatGPT Pro mode, effort Pro, S6\_10\_PAPERS continuation, 19 November 2026: computation or formal certificate; independent re-proof; used and rechecked in this run. Continuation metadata events: GPT-6 Astra (OpenAI), ChatGPT Pro mode, effort Pro, S6\_10\_PAPERS continuation, 19 November 2026: used and rechecked in this run; status-vocabulary reconciliation; metadata reconciliation only; no new mathematical check. \emph{Continuation recorded proof scope.} Complete current proof derived before the printed target was opened: actual primitive vanishing orbit, unimodular cup pairing, explicit integral marking and inverse, simultaneous companion reversal with pairing sign, based cusp conjugation and the exact exterior quotient. The former rational-kernel route is preserved separately and replaced here by the already renewed in-paper rational input. The fresh program \texttt{code/continuation\_root/P1\_7\_integral\_marking\_fresh.py} passes all 24 exact checks. The source comparison and replacement hashes are preserved in the continuation records.}{\cref{maps:src:F.3.5,night:P1.2}}
The integral $H^1$ monodromy of the pencil has a companion marking with
finite pair $(C_f^{-1},C_g)$, and cusp product $C_fC_g^{-1}$.
Its exterior-square quotient is consequently the integral quotient
representation appearing in P1.2.
\end{theorem}
\begin{proof}
\begin{enumerate}[label=\arabic*.]
\item Let $L=H^1(C_\lambda,\mathbb Z)$ with its integral cup pairing,
and write $T_1,T_2$ for the finite generators of
\cref{maps:src:F.3.5}. Put $A=C_f$ and $B=C_g$.
The rational intertwiner proved there is an actual map
\[
 q:L_{\mathbb Q}\longrightarrow\mathbb Q^4,
 \qquad qT_1q^{-1}=A^{-1},\quad qT_2q^{-1}=B.
\]
Its relation to the Levelt presentation is
$(h_\infty,h_1,h_0)=(T_1^{-1},T_1T_2,T_2^{-1})$;
the ordered product is one. The known rigidity input and its explicit
rational change of basis are recorded in \cref{maps:src:F.3.5}.
The nodal peripheral is $c=(T_1T_2)^{-1}$. Its primitive
Picard--Lefschetz class $\delta\in L$ and coefficient one are the
geometric node input proved there: the nonseparating vanishing circle
has a transverse integral cycle meeting it once, and the annulus
$xy=t$ turns once. Thus $T_1T_2$ acts by
$v\mapsto v+\epsilon\delta\langle\delta,v\rangle$,
where $\epsilon\in\{1,-1\}$ records the fixed inverse and pairing
conventions.

\item A $B$-invariant alternating form has upper entries
$(u,v,u,u,v,u)$. Invariance under $A$, evaluated on the first and
fourth basis vectors, gives $v=-u$. The common invariant forms are
therefore exactly the multiples of
\[
 J=\begin{pmatrix}
 0&1&-1&1\\-1&0&1&-1\\1&-1&0&1\\-1&1&-1&0
 \end{pmatrix},\qquad
 d=\begin{pmatrix}1\\1\\1\\0\end{pmatrix}.
\]
The Pfaffian of $J$ is $1$, and direct multiplication gives
\[
 A^tJA=B^tJB=J,\qquad
 A^{-1}B-I=d e_4^t,\qquad d^tJ=e_4^t.
\]
Here $e_4$ is the fourth standard vector. Write the transported cup
form as $sJ$ and $q(\delta)=k d$, with $s,k\in\mathbb Q^*$.
Comparison with the actual Picard--Lefschetz formula gives
$\epsilon s k^2=1$.

\item The integral classes
$\delta_i=T_1^{-i}\delta$, $0\leq i\leq3$, have rational images
equal to the columns of $kD$, where
\[
 D=(d,Ad,A^2d,A^3d)=
 \begin{pmatrix}1&0&-1&0\\1&1&-1&-1\\1&1&0&-1\\0&1&0&0\end{pmatrix},
 \qquad \det D=1.
\]
Their actual pairing is $\epsilon G$, with
\[
 G=D^tJD=
 \begin{pmatrix}0&1&0&0\\-1&0&1&0\\0&-1&0&1\\0&0&-1&0\end{pmatrix},
 \qquad \det G=1.
\]
These four classes are a basis of $L_{\mathbb Q}$. For any $v\in L$,
write $v=\sum_i x_i\delta_i$ with $x_i\in\mathbb Q$.
The vector $(\langle\delta_i,v\rangle)_i=\epsilon Gx$ is integral.
Since $G^{-1}$ is integral, so is $x$. Thus the four classes are an
integral basis of $L$.

\item The required integral marking is the precise map
$\mu=k^{-1}q:L\longrightarrow\mathbb Z^4$.
Indeed $q(L)=kD\mathbb Z^4=k\mathbb Z^4$. Its inverse is specified
on standard vectors by
\[
 \mu^{-1}(e_j)=\sum_{i=0}^3(D^{-1})_{i+1,j}\delta_i,
 \qquad
 D^{-1}=\begin{pmatrix}1&-1&1&0\\0&0&0&1\\0&-1&1&0\\1&-1&0&1\end{pmatrix}.
\]
This supplies both integral arrows and retains the cup form
$\epsilon J$. In this marking the finite pair is $(A^{-1},B)$.

\item The exact simultaneous companion reversal is
\[
 W=\begin{pmatrix}0&0&0&1\\0&0&1&0\\0&1&0&0\\1&0&0&0\end{pmatrix}.
\]
It satisfies $W^2=I$, $\det W=1$, $WAW=A^{-1}$,
$WBW=B^{-1}$ and $W^tJW=-J$. These identities give the actual
integral comparison of the two companion conventions, including
the pairing sign. For the based cusp used in the statement take
$c'=T_1^{-1}cT_1=(T_2T_1)^{-1}$. Then
\[
 \mu c\mu^{-1}=B^{-1}A,\qquad
 \mu c'\mu^{-1}=A(B^{-1}A)A^{-1}=AB^{-1}.
\]
The homology matrices are the inverse transposes under the actual
evaluation pairing with $L$.

\item Apply $\bigwedge^2\mu$ to the integral wedge lattice. In the
lexicographic basis of \cref{night:P1.2}, the invariant vector is
$\omega=(1,1,1,1,1,1)^t$. Let $Q$ have first column $\omega$ and
remaining columns the last five standard vectors, as in that result.
The induced quotient is the actual map
\[
 [u\wedge v]\longmapsto
 \operatorname{pr}_{2,\ldots,6}
 Q^{-1}\bigl(\mu(u)\wedge\mu(v)\bigr).
\]
Its kernel before passage to the quotient is precisely the invariant
line. Its inverse is induced by $(\bigwedge^2\mu^{-1})Q$ on the last
five standard vectors. The finite quotient matrices are the lower
five-by-five blocks of
\[
 Q^{-1}(\bigwedge^2 A^{-1})Q,\qquad
 Q^{-1}(\bigwedge^2 B)Q.
\]
They are exactly the integral quotient matrices of
\cref{night:P1.2}, and the cusp follows by the same based word
$(T_2T_1)^{-1}$. The invariant forms, integral inverses, peripheral
words and exterior minors are independently reproduced in
\texttt{P1\_7\_integral\_marking\_fresh.py} in the directory 
\texttt{code/continuation\_root/}.
\end{enumerate}
\end{proof}




\begin{theorem}[a global nonpolarized period family]\label{maps:src:F.3.6}
\StatementRecord{maps:src:F.3.6}{proved}{Historical events: Claude Fable 5.2 (Anthropic), effort high, 16 November 2026: targeted proof check. GPT-6 Astra (OpenAI), ChatGPT Pro mode, effort Pro, S6\_09\_PAPERS, 17 November 2026: targeted proof check. Also checked in this writing. Pre-continuation writing events: GPT-6 Astra (OpenAI), ChatGPT Pro mode, effort Pro, S6\_10\_PAPERS, 18 November 2026 02:02:02.394029+00:00: targeted proof check. Continuation mathematical and source checks: GPT-6 Astra (OpenAI), ChatGPT Pro mode, effort Pro, S6\_10\_PAPERS continuation, 19 November 2026: computation or formal certificate; targeted proof check; used and rechecked in this run. Continuation metadata events: GPT-6 Astra (OpenAI), ChatGPT Pro mode, effort Pro, S6\_10\_PAPERS continuation, 19 November 2026: used and rechecked in this run; status-vocabulary reconciliation; metadata reconciliation only; no new mathematical check. \emph{Pre-continuation recorded proof scope.} GPT-6 Astra (OpenAI, ChatGPT Pro mode), effort Pro, 2026-11-17; GPT-6 Astra (OpenAI, ChatGPT Pro mode), S6\_06\_BLIND, effort Pro, 2026-11-17; Claude Fable 5.2 (Anthropic), effort high, 2026-11-16; also checked in this writing by GPT-6 Astra (OpenAI), ChatGPT Pro mode, effort Pro, 18 November 2026}{the hypotheses, notation and definitions stated here; \cref{maps:eq:F.11}; \cref{maps:eq:F.14}; \cref{maps:eq:F.16}; \cref{maps:eq:F.19}; \cref{maps:eq:F.12}; \cref{maps:src:F.2.1}; \cref{maps:eq:F.22}; \cref{maps:eq:F.20}; \cref{maps:eq:F.21}}


The integral rank-six system with cohomology monodromies $A,B$
admits a holomorphic subbundle $\mathcal F$ of rank three over
$U$ satisfying $\mathcal F\cap\overline{\mathcal F}=0$.
It extends holomorphically on the finite root charts.  At the
cusp an integral marking puts its periods in the form
\[
 Z(t)=\frac{\log t}{2\pi i}I_3+C(t),
 \tag{F.13}\label{maps:eq:F.13}
\]
where $C$ is holomorphic at zero and bounded there.

\end{theorem}


\begin{proof}
\begin{enumerate}[label=\arabic*.]
\item Write the exact sequence supplied by \eqref{maps:eq:F.11} as
\[
 0\longrightarrow\mathbb Q e\longrightarrow\mathcal V
   \longrightarrow\mathcal W\longrightarrow0.
 \tag{F.14}\label{maps:eq:F.14}
\]
Here $e$ is the first coordinate vector.  The nilpotent canonical
extension at the cusp and the ordinary extensions on finite root
charts preserve this sequence.  This follows directly by using
the frames $\exp(-sN)$: each residue has only eigenvalue zero,
and the same exponential construction respects the inclusion and
the quotient.  At finite roots the statement is ordinary exactness
of bundles with trivialized local monodromy.  Denote these extended
bundles by $\mathcal H_D,\mathcal V_D,\mathcal W_D$.

\item Let $\omega\in\bigwedge^2\mathcal H_D$ be the real, flat
nondegenerate bivector inverse to the polarization, and let
$p:\bigwedge^2\mathcal H_D\to\mathcal W_D$ be primitive projection.
Define the holomorphic rank-three bundle
\[
 \mathcal P=p(\ell\wedge\mathcal H_D).
 \tag{F.15}\label{maps:eq:F.15}
\]
Projection is injective on $\ell\wedge\mathcal H_D$: every
nonzero bivector there is decomposable whereas $\omega$ is not.
This argument applies also at the cusp.  The polarization gives
$\mathcal H_D/\mathcal E\simeq\mathcal E^*$, including there,
since $\mathcal E$ is Lagrangian.  Hence
\[
 0\longrightarrow\mathcal T^4\longrightarrow\mathcal P
 \longrightarrow\mathcal O\oplus\mathcal T^2\longrightarrow0.
 \tag{F.16}\label{maps:eq:F.16}
\]
After dualizing, the end terms are
$\mathcal O\oplus\mathcal T^{-2}$ and $\mathcal T^{-4}$.

\item We compute the obstruction to lifting $\mathcal P$ in \eqref{maps:eq:F.14}.
For the coarse projection $\rho:\mathcal B'\to\mathbb P^1$,
the elementary root-chart calculation of invariant sections gives
\[
 \rho_*\mathcal T^k=
 \mathcal O_{\mathbb P^1}
 \left(\left\lfloor\frac{2k}{5}\right\rfloor+
       \left\lfloor\frac{-3k}{8}\right\rfloor\right).
 \tag{F.17}\label{maps:eq:F.17}
\]
Higher direct images vanish, because taking invariants under a
finite group in characteristic zero is exact.  For $k=0,-2,-4$
the displayed degrees are respectively $0,-1,-1$.  Their first
cohomology vanishes.  The dual of \eqref{maps:eq:F.16} therefore gives
$H^1(\mathcal B',\mathcal P^*)=0$.
The pullback of \eqref{maps:eq:F.14} to $\mathcal P$ consequently splits as
holomorphic bundles: choose a holomorphic lift
$\sigma:\mathcal P\to\mathcal V_D$ of its inclusion in
$\mathcal W_D$.  This is an actual extension-class calculation
on a compact orbifold, not an approximation or an Oka principle.
To specify the choice concretely, take a finite root-coordinate
cover with local bundle lifts $\sigma_i$ and write their
differences as $\xi_{ij}e$, where
$\xi_{ij}$ is a Cech cocycle in $\mathcal P^*$.
Choose any solution $b_i-b_j=\xi_{ij}$ of this explicitly
vanishing cohomology class and glue the maps
$\sigma_i-b_ie$.  Every such solution is allowed in the
construction that follows.

\item There is also a canonical holomorphic functional
$\alpha:\mathcal P\to\mathcal O$.
For $v\in\mathcal P$, its unique preimage in
$\ell\wedge\mathcal H_D$ is
\[
 j(v)=v+\alpha(v)\omega
 \tag{F.18}\label{maps:eq:F.18}
\]
under the primitive decomposition.
At a smooth point, the positive Hodge metric implies that
$\ell$ and $\bar\ell$ form a nondegenerate symplectic pair.
Their common wedge line is $\ell\wedge\bar\ell$.
The bivector $\omega$ does not belong to
$\ell\wedge\mathcal H+\bar\ell\wedge\mathcal H$, since its
component on the two-dimensional symplectic orthogonal complement
is nonzero.  Projection $p$ therefore preserves the intersection,
and
\[
 \mathcal P\cap\bar{\mathcal P}
      =p(\ell\wedge\bar\ell).
 \tag{F.19}\label{maps:eq:F.19}
\]
This is the complexification of a real line.  If $u$ spans
$\ell$, then $i u\wedge\bar u$ is real and its coefficient on
$\omega$ is nonzero, by positivity.  Thus $\alpha$ is real and
nonzero on that real intersection line.  There is a unique smooth
real section $n$ of it such that $\alpha(n)=1$.

\item Since the quotient of $\sigma(n)$ is real, there is a unique
real smooth function $h$ on $U$ such that
\[
 \sigma(n)-\overline{\sigma(n)}=2ih e.
 \tag{F.20}\label{maps:eq:F.20}
\]
For $c\in\mathbb R$ set
\[
 \mathcal F_c=(\sigma+ic\alpha e)(\mathcal P).
 \tag{F.21}\label{maps:eq:F.21}
\]
This is a holomorphic rank-three subbundle projecting
isomorphically to $\mathcal P$.  An intersection with its complex
conjugate must project to \eqref{maps:eq:F.19}.  On its real generator the two
lifts differ by $2i(h+c)e$.  It follows, including arbitrary
complex multiples of that generator, that
\[
 \mathcal F_c\cap\overline{\mathcal F_c}=0
       \quad\Longleftrightarrow\quad h+c\ne0.
 \tag{F.22}\label{maps:eq:F.22}
\]

\item We verify the required cusp property before choosing $c$.
In a complex symplectic basis $a,b,d_1,d_2$ of the limiting
$\mathcal H_D$, relation \eqref{maps:eq:F.12} permits us to arrange
$N_Hb=a$, $N_Ha=N_Hd_1=N_Hd_2=0$ and $\ell_\infty=\mathbb Cb$.
Then
\[
 \mathcal P_\infty=
 \langle k=p(b\wedge a),\ b\wedge d_1,\ b\wedge d_2\rangle,
 \quad
 N_W(b\wedge d_i)=a\wedge d_i,
 \quad N_Wk=0.
 \tag{F.23}\label{maps:eq:F.23}
\]
The line generated by $k$ is
$\ker N_W/\operatorname{im}N_W$; all other kernel directions
come from the displayed image.  The exact matrices give
$N_V^2=0$, $\operatorname{rank}N_V=3$ and
$\operatorname{rank}N_W=2$.
Every lift of an image vector $a\wedge d_i$ is killed by $N_V$,
because that vector has a lift in $\operatorname{im}N_V$ and
$N_Ve=0$.  Rank three now forces $N_V$ to take any lift of $k$
to a nonzero multiple of $e$.  Therefore
$N_V:\mathcal F_{c,\infty}\to\operatorname{im}N_V$ is an
isomorphism for every $c$.  In particular its domain is a
complement to $\ker N_V$.

\item The primitive index-one cusp has a real integral basis
$a_i,b_i$ with $N_Vb_i=a_i$, by Proposition \ref{maps:src:F.2.1}, also for the
dual or inverse convention.  In this basis the untwisted
subbundle near zero is the graph $b+C_c(t)a$, with $C_c$
holomorphic and bounded.  Restoring the flat frame by
$\exp(sN_V)$, $s=\log t/(2\pi i)$, gives \eqref{maps:eq:F.13}.
Thus $\operatorname{Im}Z=\operatorname{Im}s\,I_3+
\operatorname{Im}C_c$ is nonsingular on a sufficiently small
punctured disc.  Apply this first with $c=0$.
Equation \eqref{maps:eq:F.22} shows that $h$ is nonzero and has constant sign
on that connected punctured disc.

\item At the two finite root points, all ingredients in \eqref{maps:eq:F.19}--\eqref{maps:eq:F.20}
are smooth and $\alpha(n)=1$ remains valid, so $h$ extends and
is bounded.  On the compact complement of the cusp disc it has
a finite bound $K$.  Choose $c$ with the same sign as $h$ near
the cusp by the exact rule
$c=\operatorname{sign}(h(t_0))(K+2)$, where $t_0$ is any point
of that punctured disc and $K$ is the supremum of $|h|$ on
the compact complement.  Then $h+c$ is nonzero both on that
disc and on its complement.  Equations \eqref{maps:eq:F.21}--\eqref{maps:eq:F.22} prove
global real transversality, and the preceding cusp calculation
proves the bounded correction for this same $c$.

\item Finally, the fourth power of the order-eight action on
$\mathcal H$ is $-I$.  It acts trivially on
$\mathcal W,\mathcal P,\alpha$ and $\mathcal V$.
Hence these bundles, the lift and the constructed periods descend
from $\mathcal B'$ to the orbifold $\mathcal B$ with orders five
and four.  Indeed any equivariant bundle map between bundles with
trivial action of the extra order-two stabilizer descends on its
root chart.  This gives exactly the finite orders required for
the rank-six datum.

\end{enumerate}
\end{proof}




\begin{theorem}\label{night:P1.8}
\StatementRecord{night:P1.8}{proved}{Historical events: GPT-6 Astra (OpenAI), ChatGPT Pro mode, effort Pro, S6\_06\_BLIND, 16--18 November 2026: targeted proof check. Also checked in this writing. Pre-continuation writing events: GPT-6 Astra (OpenAI), ChatGPT Pro mode, effort Pro, S6\_10\_PAPERS, 18 November 2026 02:02:02.394029+00:00: computation or formal certificate; targeted proof check. Continuation mathematical and source checks: GPT-6 Astra (OpenAI), ChatGPT Pro mode, effort Pro, S6\_10\_PAPERS continuation, 19 November 2026: computation or formal certificate; targeted proof check; used and rechecked in this run. Continuation metadata events: GPT-6 Astra (OpenAI), ChatGPT Pro mode, effort Pro, S6\_10\_PAPERS continuation, 19 November 2026: used and rechecked in this run; status-vocabulary reconciliation; metadata reconciliation only; no new mathematical check. }{\cref{night:P1.5}}
With the corrected rank-three bundle $P$ of P1.5, the supplied rank-six
system admits a holomorphic lifted family
$F_c=(\sigma+ic\alpha e)(P)$ on the punctured base.
For a specified splitting $\sigma$ there is a real analytic function
$h$ with finite attained minimum $h_{\min}$, and the exact domain of
real parameters giving $F_c\cap\overline{F_c}=0$ everywhere is
$c>-h_{\min}$.
At the cusp its periods have the form
$Z(q)=(\log q/(2\pi i))I_3+C(q)$ with $C$ holomorphic and bounded.
The construction extends to the two finite root charts and descends to
orders $5,4$ on the rank-six system.
\end{theorem}
\begin{proof}
\begin{enumerate}[label=\arabic*.]
\item Work first on the root orbifold with orders $5,8$, leaving the nodal
point as a logarithmic point.  The symplectic pairing and the Hodge
subbundle give an exact sequence
$0\to\mathcal H\to\mathcal H_{\mathrm{dR}}\to\mathcal H^\vee\to0$.
This remains exact at the node by the residue pairing in its plumbing
chart.  With $\ell_0=T^3$ the corrected bundle fits into
\[
0\longrightarrow T^4\longrightarrow P
 \longrightarrow\mathcal O\oplus T^2\longrightarrow0.
\]
Contraction in $\ell_0\wedge\mathcal H_{\mathrm{dR}}$ defines
$\alpha:P\to\mathcal O$, using the unique lift supplied by P1.5;
it is surjective.  Thus it is not an asserted contraction on the whole
quotient by $\omega$, where that formula would not be well defined.

\item For $a=2,4,-2,-4$ the invariant coarse direct image of $T^a$ has degree
$-1$, by taking the floors of the two coefficients of its rational
divisor.  Finite-stabilizer invariants are exact in characteristic zero
by averaging.  Consequently $H^0(T^a)=H^1(T^a)=0$ for these exponents.
The two-chart Laurent-series computation for $\mathcal O(-1)$ proves
this vanishing directly.  The extension therefore splits, and
\[
P\cong T^4\oplus T^2\oplus\mathcal O,\qquad
H^1(P^\vee)=0,\quad H^0(P^\vee)=\mathbb C\alpha,
\quad H^0(P)=\mathbb Cg,
\]
where $\alpha(g)=1$.

\item Let $V$ be the logarithmic extension of the supplied rank-six flat
bundle.  At the cusp use its nilpotent exponential frame, and on each
finite root chart use the trivial flat extension with the deck action.
These charts give the exact bundle sequence
$0\to\mathcal Oe\to V\to W\to0$.
Its pullback to $P$ splits because $H^1(P^\vee)=0$.
All holomorphic splittings are $\sigma+z\alpha e$, $z\in\mathbb C$.
To remove this scalar choice we require the $e$-coordinate of
$\sigma(g)$ to be zero in the fixed flat frame at $z=1/2$.

\item Choose a local generator $u$ of $\ell_0$ and put
$r=i\pi(u\wedge\bar u)\in W_{\mathbb R}$, $A=\alpha(r)>0$;
the sign is that of the polarization.  The symplectic orthogonal
decomposition into $\langle u,\bar u\rangle$ and its two-dimensional
complement proves $P\cap\bar P=\mathbb Cr$.
Indeed $u\wedge v-\bar u\wedge w=t\omega$ has zero component in the
exterior square of that complement on the left and a nonzero such
component on the right unless $t=0$.  With $t=0$ its intersection is
the line generated by $u\wedge\bar u$.
Define $h$ by the exact equality
\[
\sigma(r)-\overline{\sigma(r)}=2iAh e.
\]
Both $r$ and $A$ scale by $|a|^2$ when $u$ scales by $a$, so this is
well defined.  Comparing the unique lifts of a multiple of $r$ proves
\[
F_c\cap\overline{F_c}=0\quad\Longleftrightarrow\quad h+c\ne0.
\]

\item The sign and growth at the cusp follow from the actual extension
matrices.  Put $D=C_fC_g^{-1}-I$ and $N_V=(BA)^{-1}-I$.
Only the zeroth column of $D$ is nonzero, and it is $(0,1,1,1)^t$;
the top row of $N_V$ is $(0,0,0,-1,1,-1)$.
For arbitrary $a,b\in\mathbb C^4$, two-by-two minor expansion gives
\[
(N_V)_{0,*}(0,\pi(a\wedge Db))
 =\operatorname{contr}(a\wedge Db)=a_0b_0,
 \qquad N_W\pi(a\wedge Db)=0.
\]
These polynomial identities, with arbitrary symbols for the entries of
$a,b$, are checked in \texttt{cusp\_lift\_asymptotic.py}.
Write $u=\exp(sD)a(q)$, $s=\log q/(2\pi i)$ and $y=\operatorname{Im}s$.
The vector $a$ is holomorphic, and $a_0(0)\ne0$ because $dx/y$ has
nonzero nodal residue.  In the logarithmic frame we obtain
\[
\exp(-sN_W)r=i\pi(a\wedge\bar a)+2y\pi(a\wedge D\bar a),
\quad
A=i\operatorname{contr}(a\wedge\bar a)+2y|a_0|^2.
\]
The splitting is bounded in that frame.  On returning by
$\exp(sN_V)$, the only quadratic contribution to the imaginary
$e$-coordinate is $2|a_0|^2y^2$.  All other terms are $O(y)$,
uniformly for $\operatorname{Re}s$ in a fixed unit interval.
Dividing by $A=2|a_0|^2y+O(1)$ proves
\[
h(q)=y+O(1)=-\frac{\log|q|}{2\pi}+O(1).
\]
The function extends continuously over both finite root charts and
tends to $+\infty$ at the cusp.  Compactness away from that cusp gives
an attained minimum.  Connectedness and the intermediate value
theorem give $h(U)=[h_{\min},\infty)$, with finite-root points included
when testing the completed family.  Thus the domain is precisely the
displayed half-line.  Replacing $\sigma$ by $\sigma+ib\alpha e$ replaces
$h$ by $h+b$ and sends $F_c$ to $F_{c+b}$, proving the stated
comparison between choices.

\item At $q=0$ the map $N_V:F_{\log}\to\ker N_V$ is an isomorphism.
On the quotient its rank is two and its kernel is generated by
$\pi(a\wedge Da)$; the displayed scalar identity sends a lift of this
generator to $a_0^2e\ne0$.  In the integral Jordan marking, the
logarithmic three-plane is therefore a graph of a holomorphic matrix
$C(q)$; applying $\exp(sN_V)$ adds $sI_3$.  Evaluation on the homology
lattice gives the period matrix $(I_3,Z(q))$.
The fourth power of the order-eight deck transformation acts trivially
on $V$ and $P$, so the graph descends to the order-four chart.
The other chart has order five.  The affine translations are
holomorphic translations on the resulting tori, and their linear
parts preserve this graph by construction.

\end{enumerate}
\end{proof}


\paragraph{Scope of P1.8.}
The definition of $h$ is exact and specifies its normalization, its
range, and its cusp growth.  We have not evaluated $h_{\min}$ or
supplied the two independent thirty-digit period calculations requested
in item 1.4.  These numerical requirements remain separately recorded;
they are not consequences of the existence proof.

\begin{theorem}\label{night:P1.9}
\StatementRecord{night:P1.9}{proved}{Historical events: GPT-6 Astra (OpenAI), ChatGPT Pro mode, effort Pro, S6\_06\_BLIND, 16--18 November 2026: targeted proof check. Also checked in this writing. Pre-continuation writing events: GPT-6 Astra (OpenAI), ChatGPT Pro mode, effort Pro, S6\_10\_PAPERS, 18 November 2026 02:02:02.394029+00:00: targeted proof check. Continuation mathematical and source checks: GPT-6 Astra (OpenAI), ChatGPT Pro mode, effort Pro, S6\_10\_PAPERS continuation, 19 November 2026: independent re-proof; used and rechecked in this run. Continuation metadata events: GPT-6 Astra (OpenAI), ChatGPT Pro mode, effort Pro, S6\_10\_PAPERS continuation, 19 November 2026: used and rechecked in this run; status-vocabulary reconciliation; metadata reconciliation only; no new mathematical check. }{Definitions and the displayed constructions}
The periodic Kuhn fan, in the marking $S$, gives a smooth proper
holomorphic filling of the periods $Z(q)=sI_3+C(q)$ for a sufficiently
small disk.  Its central fibre is reduced with normal crossings.
Modulo the index-one period lattice it has respectively $1,7,12,6$
strata of dimensions $3,2,1,0$, each isomorphic to the corresponding
complex torus $(\mathbb C^*)^d$, and Euler number $6$.
An index-$d$ period sublattice gives $d$ times these stratum counts and
Euler number $6d$.
\end{theorem}
\begin{proof}
\begin{enumerate}[label=\arabic*.]
\item Let $\Sigma$ consist of the cones over all simplices of the periodic
Kuhn triangulation at height one, together with their faces.  For an
ordered maximal simplex $v_0,\ldots,v_3$, the columns $(v_j,1)$ form a
unimodular four-by-four matrix.  If $u_j$ are its dual integral basis,
the corresponding chart has coordinates $x_j=\chi^{u_j}$ and is
$\mathbb C^4$.  The base character is the sum of the four $u_j$, so
\[
q=x_0x_1x_2x_3.
\]
On a common face we invert the coordinates of omitted rays and change
dual bases; these are mutually inverse integral monomial maps.  These
explicit charts and overlaps construct a smooth toric space over the
disk with a reduced normal-crossings zero fibre.

\item The period group acts by
\[
(z,q)\longmapsto
\bigl(q^\lambda\exp(2\pi iC(q)\lambda)z,q\bigr),
\qquad \lambda\in\mathbb Z^3.
\]
The monomial part takes the cone over $\sigma$ to the cone over
$\sigma+\lambda$; the unit part is torus multiplication holomorphic in
$q$, which extends over each chart.  These maps compose according to
addition in $\mathbb Z^3$.  On the zero fibre a fixed point would
require a nonzero translation to stabilize a bounded simplex, which
is impossible.  For small nonzero $q$, the imaginary part of the
period matrix is $\operatorname{Im}(s)I_3+\operatorname{Im}C(q)$ and
is invertible; hence the action is free there too.

\item We give a compactness check for the infinite-chart construction.  Put
$u=\log|z|/\log|q|$.  A period changes $u$ by
$(I+E(q))\lambda$, where $E(q)=O(1/|\log|q||)$ uniformly.
On a small disk $(I+E)^{-1}$ has norm at most two.  Reduction of
$(I+E)^{-1}u$ modulo $\mathbb Z^3$ therefore gives a representative
with $u$ in a fixed bounded region.  Only finitely many simplices
meet that region.  In a simplex containing $u$, the numbers
$\log|x_j|/\log|q|$ are its barycentric coordinates; thus
$|x_j|\leq1$.  Finitely many closed polydisks in these charts give
compact representatives above a smaller closed base disk.  Every
central point is included in such representatives by translating its
finite star into the same bounded region and then using its complete
toric component.

\item The action is proper.  For compact coordinate neighborhoods approaching
the central fibre, only finitely many cones and their faces occur;
their translates can meet only when the translation is bounded.
Equivalently, the valuation estimate just given and the uniform bound
on $(I+E)^{-1}$ bound the translations for a sequence with $q\to0$.
On a closed annulus away from zero, properness is the usual
properness of a full lattice action.  These two checks give finitely
many intersecting translates of any pair of compact neighborhoods.
The quotient is therefore Hausdorff, is smooth by freeness, and is
proper over the disk by the compact representatives.  The map is flat
because $q$ is a nonzero divisor in each smooth chart.  Exponentiation
$z_j=\exp(2\pi i\zeta_j)$ identifies the nonzero fibres with
$\mathbb C^3/(\mathbb Z^3+Z(q)\mathbb Z^3)$.

\item A $p$-simplex supplies a stratum $(\mathbb C^*)^{3-p}$.
Enumeration of all faces of the six Kuhn tetrahedra, retaining their
translation identifications, gives $(1,7,12,6)$.
Additivity of compactly supported Euler characteristic and
$\chi_c(\mathbb C^*)=0$ leave six point contributions.
Inclusion of an index-$d$ period sublattice induces a finite covering
of these toric quotients.  Each simplex orbit has exactly $d$ lifts,
which proves the stated index dependence by an actual comparison map.

\end{enumerate}
\end{proof}




\begin{theorem}[compact holomorphic filling]\label{maps:src:F.3.7}
\StatementRecord{maps:src:F.3.7}{proved}{Historical events: Claude Fable 5.2 (Anthropic), effort high, 16 November 2026: carried as background without fault found. GPT-6 Astra (OpenAI), ChatGPT Pro mode, effort Pro, S6\_09\_PAPERS, 17 November 2026: targeted proof check. Also checked in this writing. Pre-continuation writing events: GPT-6 Astra (OpenAI), ChatGPT Pro mode, effort Pro, S6\_10\_PAPERS, 18 November 2026 02:02:02.394029+00:00: targeted proof check. Continuation mathematical and source checks: GPT-6 Astra (OpenAI), ChatGPT Pro mode, effort Pro, S6\_10\_PAPERS continuation, 19 November 2026: computation or formal certificate; targeted proof check; used and rechecked in this run. Continuation metadata events: GPT-6 Astra (OpenAI), ChatGPT Pro mode, effort Pro, S6\_10\_PAPERS continuation, 19 November 2026: used and rechecked in this run; status-vocabulary reconciliation; metadata reconciliation only; no new mathematical check. \emph{Pre-continuation recorded proof scope.} GPT-6 Astra (OpenAI, ChatGPT Pro mode), effort Pro, 2026-11-17; GPT-6 Astra (OpenAI, ChatGPT Pro mode), S6\_06\_BLIND, effort Pro, 2026-11-17; Claude Fable 5.2 (Anthropic), effort high, 16 November 2026; also checked in this writing by GPT-6 Astra (OpenAI), ChatGPT Pro mode, effort Pro, 18 November 2026}{the hypotheses, notation and definitions stated here; \cref{maps:src:F.2.6}; \cref{maps:src:F.2.5}; \cref{maps:eq:F.24}; \cref{maps:eq:F.13}; \cref{maps:eq:E.4}; \cref{maps:eq:E.5}; \cref{maps:src:E.1.2}; \cref{maps:src:E.1.3}}


The smooth eight-manifold of Theorem \ref{maps:src:F.2.6} carries an integrable
complex structure.  There is a proper holomorphic map to
$\mathbb P^1$ whose regular fibres are complex three-tori,
whose two finite reduced fibres are the free affine quotients
of orders five and four specified in Theorem \ref{maps:src:F.2.5}, and whose
remaining fibre is the index-one Kuhn degeneration.

\end{theorem}


\begin{proof}
\begin{enumerate}[label=\arabic*.]
\item We first fix the dual convention.  The preceding construction
has cohomology monodromies $A,B$; Theorem \ref{maps:src:F.2.5} uses homology
monodromies $A^t,B^t$.  All three special values are real.
Choose a positive real basepoint and paths along the positive
real circle to infinity and to zero.  Complex conjugation fixes
these paths and reverses the respective peripheral meridians.
Pull back the preceding period family by conjugation of the
base and conjugate its complex fibres.  Its Hodge subbundle is
holomorphic, with local expression
$\overline{\mathcal F(\bar\lambda)}$, and its cohomology
monodromies are $A^{-1},B^{-1}$.  Its homology monodromies are
therefore exactly $A^t,B^t$.
Put $M_1=A^t$, $M_2=B^t$, $M_o=M_2M_1$, $M_0=(M_1M_2)^{-1}$ and
$J=\operatorname{diag}(I_3,-I_3)$.  The based peripheral word is
$\kappa(g_0)=g_1^{-1}g_0^{-1}g_1$, so
$M_0=M_1M_o^{-1}M_1^{-1}$.  For the specified target marking $S$,
direct integral multiplication gives
\[
 S^{-1}M_0S=\begin{pmatrix}I_3&I_3\\0&I_3\end{pmatrix},\qquad
 S_o=M_1^{-1}SJ,\qquad
 S_o^{-1}M_oS_o=\begin{pmatrix}I_3&I_3\\0&I_3\end{pmatrix}.
\]
Any two markings of this positive standard form differ by
$G=\left(\begin{smallmatrix}P&Q\\0&P\end{smallmatrix}\right)$,
where $P\in\operatorname{GL}_3(\mathbb Z)$ and $Q$ is integral:
this follows by multiplying $G$ with the displayed unipotent
matrix.  The actual period change is
\[
 (I_3,Z)G=(P,Q+ZP),\qquad
 P^{-1}(P,Q+ZP)
  =(I_3,sI_3+P^{-1}Q+P^{-1}CP).
\]
Thus the bounded holomorphic logarithmic form holds in $S_o$.
In the real-normalized cusp coordinate,
$q_o=\bar q$ and $s_o=-\bar s$.
Base and fibre conjugation therefore give
$(I_3,-sI_3+C^\kappa(q))$, with
$C^\kappa(q)=\overline{C(\bar q)}$.
The period-cycle change $J$ gives the target marking
$M_1S_oJ=S$ and the positive periods
$(I_3,sI_3-C^\kappa(q))$.


\item For a holomorphic Hodge subbundle $\mathcal F$ with
$\mathcal F\cap\bar{\mathcal F}=0$, restriction of integral
homology functionals embeds the homology lattice discretely in
$\mathcal F^*$.  Real transversality makes its real span the
whole real vector space.  The quotient
$\mathcal F^*/\mathcal V_{\mathbb Z}^*$ is consequently a
holomorphic family of compact complex three-tori over $U$,
with its zero section.  This also proves local properness,
by a smoothly varying fundamental parallelepiped.

\item On the root disc at the finite point of order $m_i$, the family
extends smoothly and its deck map has linear part $M_i$.
The invariant vector $v_i$ of Theorem \ref{maps:src:F.2.5} gives a holomorphic
torsion section $v_i/m_i$.  Use the diagonal affine deck action
\[
 (z,s)\longmapsto(M_i z+v_i/m_i,\zeta_{m_i}s).
 \tag{F.24}\label{maps:eq:F.24}
\]
Here $M_i z$ denotes the induced complex-linear map
$\mathcal F_s^*\to\mathcal F_{\zeta_{m_i}s}^*$ from the marked
lattice action, not multiplication of a complex three-vector
by a six-by-six matrix.  Likewise $v_i/m_i$ means its evaluation
as a point of the target torus.
It has order $m_i$, since the $m_i$th translation is integral.
Every nonidentity power is fixed-point-free on the central
torus: the invariant integral covector $\ell$ takes $v_i$ to
$1$ or $-1$, so a fixed point would require the noninteger
$\pm j/m_i$ to vanish in $\mathbb R/\mathbb Z$.
Away from the central fibre the base rotation is already free.
The quotient is therefore smooth and proper over the coarse
disc, with precisely the required multiple fibre.

\item Its marked punctured family is holomorphically identified with
the original torus family by logarithmic translation.  On the
root cover use $r(s)=(\log s/(2\pi i))v_i$.  Then
$r(\zeta_{m_i}s)-M_i r(s)=v_i/m_i$, modulo the integral lattice;
the change of branch of $\log s$ is itself an integral vector.
Thus translation by $r$ conjugates the linear deck action to
\eqref{maps:eq:F.24}.  This fixes the affine clutching, including its sign.

\item At the cusp use the integral marking of \eqref{maps:eq:F.13}, the fan \eqref{maps:eq:E.4}
and the action \eqref{maps:eq:E.5} with $C_3=C$.
Its maximal charts are $\mathbb C^4$ with map
$t=z_0z_1z_2z_3$, and its transition maps are the toric changes
of coordinates multiplied by the holomorphic units
$\exp(2\pi i C(t)\lambda)$.
The actual multiplicative comparison is
\[
 q=\bar q_o,\qquad x=(\bar x_o)^{-1},\qquad
 \nu=-\nu_o,
 \qquad C_{\rm new}(q)=-C^\kappa(q).
\]
It carries the old period action to
$x\mapsto\exp(2\pi iC_{\rm new}(q)\nu)q^\nu x$.
On cocharacters it is
$\Phi=\operatorname{diag}(-I_3,1)$ together with conjugation.
The Kuhn simplex at $v$ with ordering $\sigma$ is carried to
the simplex at $-v-(1,1,1)$ with reversed ordering.
If $N_o$ and $N$ are their height-one generator matrices and
$R$ reverses four columns, then
$N=\Phi N_oR$ and $N^{-1}\Phi=RN_o^{-1}$.
Consequently the full affine-chart map is
$z_j=\overline{z_{o,3-j}}$.  These character identities also
give equality on every monomial overlap; the product of the
four coordinates is the conjugate base coordinate.
The already displayed period-action equality makes the map
descend to the analytic period quotient.  Hence the fan and
marking in this comparison are precisely the prescribed ones.

The proof of Theorem \ref{maps:src:E.1.2} used only boundedness of $C$,
integrality of the standard cusp, and unimodularity of that
fan.  It therefore proves smoothness, proper discontinuity
and properness for this model without a symmetry condition
on $C$.  Glue by the displayed integral marking and use no
additional cusp translation.

\item The three punctured overlap maps are holomorphic isomorphisms
with the same marked torus family, so their cocycle condition
is literal.  Proper local models and a compact complementary
base give a Hausdorff compact total space.  Every chart is
smooth.  The affine vectors, cusp fan and clutching are exactly
those in Theorems \ref{maps:src:F.2.5}--\ref{maps:src:F.2.6}.  Varying their bounded corrections
and varying the fibre complex structures along a smooth path
does not change these marked smooth fillings: on the finite
root models the underlying real torus family is the fixed
local-system bundle; on the cusp use the stratum-preserving
parameter vector fields and uniform properness estimate of
Theorem \ref{maps:src:E.1.3}.  Thus the resulting underlying smooth manifold
is the one whose integral topology was computed there.

\end{enumerate}
\end{proof}


\section{The oriented diffeomorphism type}

\paragraph{The retained diagnostic complex.}
The earlier file \texttt{global\_candidate\_complex.py} sets the entire
cusp suspension map to zero and has the same free Betti numbers.
That file and its output are retained as a diagnostic record.  They
are not used in P1.12 to identify the actual higher-degree gluing.
P1.19 supplies the requested independent full global cell computation.
It proves that the local secondary cusp term is nonzero, gives the
explicit integral reduction to the corrected global matrix, and
recovers every homology group without global Poincar\'e duality.

\begin{theorem}\label{night:P1.13}
\StatementRecord{night:P1.13}{proved}{Historical events: GPT-6 Astra (OpenAI), ChatGPT Pro mode, effort Pro, S6\_06\_BLIND, 16--18 November 2026: targeted proof check. Also checked in this writing. Pre-continuation writing events: GPT-6 Astra (OpenAI), ChatGPT Pro mode, effort Pro, S6\_10\_PAPERS, 18 November 2026 02:02:02.394029+00:00: targeted proof check. Continuation mathematical and source checks: GPT-6 Astra (OpenAI), ChatGPT Pro mode, effort Pro, S6\_10\_PAPERS continuation, 19 November 2026: computation or formal certificate; targeted proof check; used and rechecked in this run. Continuation metadata events: GPT-6 Astra (OpenAI), ChatGPT Pro mode, effort Pro, S6\_10\_PAPERS continuation, 19 November 2026: used and rechecked in this run; status-vocabulary reconciliation; metadata reconciliation only; no new mathematical check. }{\cref{night:P1.12}; \S\ref{sec:V-known-inputs}; \cite[equations (1.2)--(1.4), pp.~3--4]{ffEsoleKang}; \cite[Theorem 2.2.1, p.~61]{ffPetersSterk}}
With the complex orientation the fourfold has
\[
\sigma(M_c)=2,\qquad H^4(M_c;\mathbb Z)\cong I_{3,1},\qquad
w_2(M_c)=0,
\]
and there is a diagonal integral basis in which
$p_1/2=(1,1,1,1)$.
Its Chern numbers are
$c_4=6$, $c_2^2=2$, and
$c_1^4=c_1^2c_2=c_1c_3=0$; its Pontryagin numbers are
$p_1^2=8$, $p_2=14$.
\end{theorem}
\begin{proof}
\begin{enumerate}[label=\arabic*.]
\item P1.12 implies $c_1=c_3=0$ as integral cohomology classes;
$c_4[M_c]=\chi(M_c)=6$.  The Chern-root identities give
$p_1=-2c_2$, $p_2=c_2^2+2c_4$.
The symbol-index interpretation in \S\ref{sec:V-known-inputs},
with the opened polynomial \cite[equation (1.2), p.~3]{ffEsoleKang},
gives integrality of the Todd genus. Expanding
$\prod x_j/(1-e^{-x_j})$ to degree four gives
\[
\frac{3c_2^2-6}{720}\in\mathbb Z,
\quad\text{hence}\quad c_2^2=2+240k\quad(k\in\mathbb Z).
\]
The signature polynomial \cite[equation (1.4), p.~4]{ffEsoleKang}, applied to
this closed oriented smooth eight-manifold gives
\[
\sigma=\frac{7p_2-p_1^2}{45}
      =\frac{3c_2^2+84}{45}=2+16k.
\]
The rank-four middle form has $|\sigma|\leq4$, so $k=0$.
This proves the displayed characteristic numbers.

\item Duality makes the middle form unimodular.  Its signature is two and
its positive and negative indices are three and one.  The rank-at-most-four integral classification
\cite[Theorem 2.2.1, p.~61]{ffPetersSterk} applied to these
positive and negative indices gives the odd form \(I_{3,1}\).
Since $H^2(M_c;\mathbb Z/2)=0$, its $w_2$ is zero.
For the complex tangent bundle,
$w_4=c_2\bmod2=(p_1/2)\bmod2$.
The Wu identity in degree four, with $w_1=w_2=0$, says that $w_4$ is
characteristic for the middle pairing.  Thus the vector $p_1/2$ is
characteristic and has square $2$.

\item We give the required integral isometry of this pair explicitly as
a terminating procedure.  In a diagonal basis a characteristic vector
has four odd coordinates.  Change their signs to write it as
$(a,b,c,d)$ with all positive and
$a^2+b^2+c^2=d^2+2$.
For $d>1$,
\[
d<a+b+c<3d.
\]
The first inequality follows by squaring and using the positive cross
terms; the second follows from Cauchy--Schwarz and $d>1$.
Reflection in the integral norm-two vector $(1,1,1,1)$ sends this
vector to itself minus $(a+b+c-d)(1,1,1,1)$.
Its fourth coordinate is $2d-a-b-c$, whose absolute value is strictly
smaller than $d$.  The reflection preserves the integral lattice and
characteristic parity.  Repeating, with sign changes at each step,
terminates at $d=1$, when the norm equation forces $a=b=c=1$.
This supplies an integral isometry taking $p_1/2$ to $(1,1,1,1)$,
and proves the claimed pair, not just its square.

\end{enumerate}
\end{proof}




\begin{theorem}[a complex structure on the stated connected sum]\label{maps:src:F.4.1}
\StatementRecord{maps:src:F.4.1}{proved}{Historical events: Claude Fable 5.2 (Anthropic), effort high, 15 November 2026: carried as background without fault found. Claude Fable 5.2 (Anthropic), effort high, 16 November 2026: targeted proof check. Claude Fable 5.2 (Anthropic), effort high, 17 November 2026: carried as background without fault found. GPT-6 Astra (OpenAI), ChatGPT Pro mode, effort Pro, S6\_09\_PAPERS, 17 November 2026: targeted proof check. Also checked in this writing. Pre-continuation writing events: GPT-6 Astra (OpenAI), ChatGPT Pro mode, effort Pro, S6\_10\_PAPERS, 18 November 2026 02:02:02.394029+00:00: targeted proof check. Continuation mathematical and source checks: GPT-6 Astra (OpenAI), ChatGPT Pro mode, effort Pro, S6\_10\_PAPERS continuation, 19 November 2026: computation or formal certificate; independent re-proof; targeted proof check; used and rechecked in this run. Continuation metadata events: GPT-6 Astra (OpenAI), ChatGPT Pro mode, effort Pro, S6\_10\_PAPERS continuation, 19 November 2026: used and rechecked in this run; status-vocabulary reconciliation; metadata reconciliation only; no new mathematical check. \emph{Continuation recorded proof scope.} The complete current six-step proof was independently derived before its printed target was opened. The actual integral reflection product, quaternionic normal-clutching coordinates and stabilization, Poincare-duality map of extended forms, and orientation of homotopy-sphere absorption are explicit. Fresh exact computations verify the recorded matrix, descent and clutching scopes. The former proof is preserved with its separate checking scope and the exact coordinate transformation relating its convention.}{the hypotheses, notation and definitions stated here; \cref{maps:src:F.3.7}; \cref{maps:src:F.2.6}; \S\ref{sec:V-known-inputs}; the explicit integral reflections below; \cite[Corollary 1.4, p.~4]{ffMandal}; \cite[Definition 3.6 and Theorem 3.7(1), pp.~12--14]{IVVWallForm}; \cite[Theorem A, p.~178]{IVVKramerStolz}; \cref{night:P1.13}}


The compact complex fourfold of Theorem \ref{maps:src:F.3.7} is three-connected,
has Euler number six, and is orientation-preservingly diffeomorphic to
\[
 3\mathbb H\mathbb P^2\ \#\ (-\mathbb H\mathbb P^2).
 \tag{F.25}\label{maps:eq:F.25}
\]
Its integral middle lattice is $I_{3,1}$ and, in a diagonal integral
basis, $p_1/2=(1,1,1,1)$.  In particular the conclusion is an exact
smooth diffeomorphism, with no remaining homotopy-sphere ambiguity.

\end{theorem}


\begin{proof}
\begin{enumerate}[label=\arabic*.]
\item The integral calculation in \cref{maps:src:F.2.6} gives
$\pi_1(M)=0$, $H_2(M;\mathbb Z)=H_3(M;\mathbb Z)=0$, and
torsion-free homology of ranks $1,4,1$ in degrees $0,4,8$.
Hurewicz applied successively gives $\pi_2(M)=\pi_3(M)=0$, and
$e(M)=6$. The complex orientation is the one in
\cref{maps:src:F.3.7}. The characteristic calculations of \cref{night:P1.13} give
\[
Q=I_{3,1},\qquad c_2^2=2,\qquad
p:=p_1/2=-c_2,\qquad p^2=2.
\]
Here $c_1=c_3=0$ because the corresponding integral groups vanish.
The spin Wu identity makes $p$ characteristic for $Q$. Evaluating
its characteristic congruence on each diagonal basis vector shows
that all four coordinates of $p$ are odd.

\item Put $v=(1,1,1,1)^t$. The actual integral reflection is
\[
R=I-v(v^tQ)=
\begin{pmatrix}
0&-1&-1&1\\
-1&0&-1&1\\
-1&-1&0&1\\
-1&-1&-1&2
\end{pmatrix}.
\]
Since $v^tQv=2$, direct multiplication gives $R^2=I$ and
$R^tQR=Q$. Independent coordinate sign changes are also integral
isometries. Use them to write the current vector as $(a,b,c,d)$,
with positive odd entries and
$a^2+b^2+c^2=d^2+2$. For $d>1$, set $s=a+b+c$. Then
\[
s^2-d^2=2(ab+ac+bc)+2>0,\qquad
s^2\leq3(d^2+2)<9d^2.
\]
Thus $d<s<3d$. Reflection subtracts the even integer $s-d$ from
each coordinate; the new last coordinate is $2d-s$, a nonzero odd
integer of absolute value less than $d$. Make the signs positive
again and repeat. The last coordinate strictly decreases until
$d=1$, when the norm equation forces $a=b=c=1$. The resulting
product $U$ of the displayed reflections and sign matrices satisfies
$Up=v$. Its inverse is the reversed product of these involutions.

\item We compute the complete normal clutching datum in fixed
coordinates. Orient $\mathbb H$ by $(1,i,j,k)$ and use
\[
\operatorname{Spin}(4)=S^3_L\times S^3_R,
\qquad (u,w)h=uhw^{-1}.
\]
The covering induces an isomorphism on $\pi_3$. Let $(a,b)$ be
the degrees of the two clutching maps. Evaluation on $1$ gives
$uw^{-1}$, of degree $a-b$, so $e=a-b$, with the sphere orientation
fixed by the left degree-one Euler generator. The left action is
complex linear for right multiplication by $i$, whose real complex
basis $(1,i,j,-k)$ has the opposite quaternion orientation. Its
degree-one generator therefore has $c_2=-1$, and hence $p_1=2$.
The right inverse action is complex linear for left multiplication
by $i$, with real basis $(1,i,j,k)$. Its Euler number and $c_2$ are
both $-1$, so its $p_1$ is again $2$. The determinant-one complex
actions have $c_1=0$. Consequently the exact homomorphism and its
inverse on its image are
\[
(a,b)\longmapsto(e,p_1/2)=(a-b,a+b),
\qquad
(a,b)=\left(\frac{e+p_1/2}{2},\frac{p_1/2-e}{2}\right).
\]
The image consists of pairs of the same parity, and the map is
injective. The tangent clutching class of $S^4$ has $(e,p_1)=(2,0)$,
hence coordinates $(1,-1)$. The homotopy sequence of
$SO(4)\to SO(5)\to S^4$ identifies its span with the kernel of
stabilization and gives surjectivity on $\pi_3$. Subsequent
stabilizations are isomorphisms there, since
$\pi_3(S^n)=\pi_4(S^n)=0$ for $n\geq5$.
The stable coordinate is therefore $a+b=p_1/2$, with the left
generator sent to $1$.

\item For $x\in H_4(M;\mathbb Z)$, use the embedded sphere
representative in the setup of
\cite[\S3.3, pp.~11--12]{IVVWallForm}. Orient its normal bundle
$\nu_x$ by the ambient orientation. Self-intersection and the
tangent splitting give
\[
e(\nu_x)=x\cdot x,\qquad
\langle p_1(\nu_x),[S^4]\rangle=\langle p_1(TM),x\rangle,
\]
because $TS^4$ is stably trivial. Thus its two clutching coordinates
are exactly
\[
\left(\frac{x\cdot x+\langle p,x\rangle}{2},
      \frac{\langle p,x\rangle-x\cdot x}{2}\right),
\qquad S\alpha_M(x)=\langle p,x\rangle.
\]
Their integrality is the characteristic congruence in step~1.
This is the stable normal homomorphism of
\cite[Definition~3.6, p.~12]{IVVWallForm} with its actual generator.
For an oriented middle-cohomology isometry $\psi$ from $M$ to a
comparison manifold $N$, the corresponding homology map is
\[
\Phi=\operatorname{PD}_N\circ\psi\circ\operatorname{PD}_M^{-1}.
\]
When $\psi(p_M)=p_N$, the cup-product evaluation identity gives
$\langle p_N,\Phi(x)\rangle=\langle p_M,x\rangle$.
The map $\Phi$ preserves the intersection pairing and both displayed
normal coordinates, so it intertwines the complete extended forms.

\item Take $N=3\mathbb H\mathbb P^2\#(-\mathbb H\mathbb P^2)$.
On each summand choose the generator $u$ with $p_1=2u$, as in
\cite[Lemma~2.1, p.~5]{ffMandal}. Reversing orientation changes
the evaluation of $u^2$ and leaves the tangent Pontryagin class
unchanged. The collapse maps of the oriented connected sum give
\[
H^4(N;\mathbb Z)=\bigoplus_{i=1}^4\mathbb Z u_i,
\qquad Q_N=I_{3,1},\qquad p_N=u_1+u_2+u_3+u_4;
\]
the tangent formula is
\cite[Theorem~2.4, p.~7]{ffMandal}.
Van Kampen, the connected-sum homology sequence and Hurewicz show
that $N$ is three-connected. The isometry $U$ of step~2 is an
isometry $\psi:H^4(M;\mathbb Z)\to H^4(N;\mathbb Z)$ carrying
$p_M$ to $p_N$. Step~4 supplies the isomorphism of their smooth
extended forms. Wall's classification in
\cite[Theorem~3.7(1), pp.~13--14]{IVVWallForm} now gives an
orientation-preserving diffeomorphism
$M\cong N\#\Sigma$ for an oriented homotopy eight-sphere $\Sigma$.

\item Kramer--Stolz's
\cite[Theorem~A, p.~178]{IVVKramerStolz} applies to
$\mathbb H\mathbb P^2$: it is smooth, simply connected, has dimension
eight, and has total integral homology $\mathbb Z\oplus\mathbb Z
\oplus\mathbb Z$. It gives
$\mathbb H\mathbb P^2\#\Sigma\cong\mathbb H\mathbb P^2$.
This diffeomorphism preserves orientation: its map on $H^4$ sends
a generator to that generator or its negative, whose square still
evaluates to $1$, forcing its degree to be $+1$.
The oriented associativity and permutation maps for connected sum
place $\Sigma$ beside one positive summand. Choose the punctured
balls compatibly with that diffeomorphism and use the identity on
the remaining punctured summands. These maps give
\[
\begin{aligned}
N\#\Sigma
&\cong(\mathbb H\mathbb P^2\#\Sigma)\#\mathbb H\mathbb P^2
       \#\mathbb H\mathbb P^2\#(-\mathbb H\mathbb P^2)\\
&\cong N.
\end{aligned}
\]
Composing with step~5 supplies the required oriented smooth
diffeomorphism. Transporting the integrable structure of
\cref{maps:src:F.3.7} along it gives the asserted complex structure.
\end{enumerate}
\end{proof}


\section{The normalized connection and the exact parameter domain}



\begin{theorem}\label{night:P1.15}
\StatementRecord{night:P1.15}{proved}{Historical events: GPT-6 Astra (OpenAI), ChatGPT Pro mode, effort Pro, S6\_06\_BLIND, 16--18 November 2026: targeted proof check. Also checked in this writing. Pre-continuation writing events: GPT-6 Astra (OpenAI), ChatGPT Pro mode, effort Pro, S6\_10\_PAPERS, 18 November 2026 02:02:02.394029+00:00: computation or formal certificate; targeted proof check. Continuation mathematical and source checks: GPT-6 Astra (OpenAI), ChatGPT Pro mode, effort Pro, S6\_10\_PAPERS continuation, 19 November 2026: computation or formal certificate. Continuation metadata events: GPT-6 Astra (OpenAI), ChatGPT Pro mode, effort Pro, S6\_10\_PAPERS continuation, 19 November 2026: used and rechecked in this run; status-vocabulary reconciliation; metadata reconciliation only; no new mathematical check. }{\cref{night:P1.1}, \cref{night:P1.8}}
The rank-six flat extension, as a complex flat bundle with its invariant
line, has the following explicit algebraic connection model.  Set
$\omega_k=x^kdx/y$, $0\leq k\leq3$, and use the quotient frame
\[
(w_2,w_1,w_0,w_{-1},w_{-2})
 =(\omega_{01},\omega_{02},\omega_{03},\omega_{13},\omega_{23}),
\qquad \omega_{12}=-3\omega_{03}.
\]
These are meromorphic frames of
$(L^2,L,\mathcal O,L^{-1},L^{-2})$, where $L=T^2$ on the root
orbifold of orders $5,4$.  If a frame is written as a row of vectors,
its derivative is right multiplication by
\[
\Gamma_6=\begin{pmatrix}0&b\\0&\Gamma_W\end{pmatrix},\qquad
b=\frac1{\lambda-\lambda_*}
 \left(-\frac{25}{6\lambda^2},-\frac{125}{48\lambda},0,1,0\right),
\tag{1.15.1}
\]
where
\[
\resizebox{\textwidth}{!}{$\displaystyle
\Gamma_W=\frac1{3125\lambda+256}
\begin{pmatrix}
-8(625\lambda+48)/\lambda&-8&10\lambda&75\lambda^2/2&0\\
-125/(2\lambda)&-4(625\lambda+48)/\lambda&8&0&75\lambda^2/2\\
300/\lambda^2&-375/\lambda&0&-48&60\lambda\\
120/\lambda^3&0&125/(2\lambda)&4(625\lambda+48)/\lambda&-8\\
0&120/\lambda^3&50/\lambda^2&-125/(2\lambda)&8(625\lambda+48)/\lambda
\end{pmatrix}.$}
\tag{1.15.2}
\]
This model specifies a canonical algebraic splitting, denoted
$\sigma_b$.  The compatible contraction used for the real parameter is
$\alpha=(625\pi i/128)\alpha_{\rm alg}$, where
$\alpha_{\rm alg}(w_0)=4/3$ and it vanishes on $w_1,w_2$.
\end{theorem}
\begin{proof}
\begin{enumerate}[label=\arabic*.]
\item Reduction modulo exact differentials is completely explicit.  To
differentiate $\omega_k$, solve the polynomial identity
\[
-\tfrac12x^k\partial_\lambda f
 =f\sum_{j=0}^3G_{kj}x^j+R_k'f-\tfrac12R_kf',
 \qquad\deg R_k\leq4,
\]
where $f=x^5+\lambda^3x+\lambda^4$.  The last two terms are the
numerator of $d(R_k/y)$.  Matching the nine coefficients gives the
four-by-four Gauss--Manin matrix exactly.
At the branch point at infinity, a local coordinate $z=x^{-1/2}$
gives $\omega_k=-2z^{2-2k}(1+O(z^8))\,dz$.
Residues of a primitive times a differential give the alternating
pairings $J_{03}=4/3$, $J_{12}=4$, with their negatives in reverse
order and all other entries zero.  This form is nondegenerate.
The invariant bivector is therefore $3\omega_{03}+\omega_{12}$.
Taking the exterior-square connection and this quotient gives
(1.15.2), as checked symbolically in \texttt{de\_rham\_connection.py}.

\item The root-chart powers computed from $x^kdx/y$ are
$u^{6k-9}$ at zero and $u^{6-4k}$ at infinity.
Their nondegenerate pairing shows that the four resulting regular
frames form the full de Rham bundle on those charts.  At the node
the computed connection has a simple pole and square-zero residue;
hence these are the nilpotent logarithmic extension frames there.
Thus the four line bundles are
$T^3,T,T^{-1},T^{-3}$, and their quotient exterior-square lines are
the stated powers of $L$.

\item All $H^1(L^j)$ for $j=0,\pm1,\pm2$ vanish by the coarse-degree
calculation already used in P1.8.  Consequently every holomorphic
extension of this $W$ by $\mathcal O$ splits as a bundle.
In such a splitting its flat connection is an upper block extension
by a row $b$.  Regularity at the finite root charts and a logarithmic
pole at the cusp require, entry by entry,
\[
b=\frac1{\lambda-\lambda_*}
 (a\lambda^{-2},b'\lambda^{-1},0,c,d\lambda).
\]
For example the $L^j$ entry, after $\lambda=u^4$, has exponent
$3j+3+4p$ for a term $\lambda^p/ (\lambda-\lambda_*)$;
the infinity chart requires $-4j-1-5p\geq0$.
These two inequalities leave exactly the four displayed terms.
The equation that the full cusp residue have square zero is
$b_{\rm res}(\Gamma_W)_{\rm res}=0$.
Its solution space has basis
\[
(-25/6,-125/48,1,0),\qquad(5/3,-25/3,0,1).
\]
The sole nonconstant holomorphic splitting change is the row selecting
$w_0$, and its connection coboundary is the third row of $\Gamma_W$.
In these two basis vectors that row is $(-48b_1+60b_2)/3125$.
Thus there is one nontrivial flat extension direction modulo splitting
changes and rescaling of the invariant line.  The first vector has
rank-three square-zero cusp residue, so it represents that direction.
The supplied extension is nontrivial, since its dual has no common
invariant vector by P1.1.  It is therefore isomorphic to (1.15.1).
This proves the connection identification without inferring it from
the subsequent numerical matching.

\item For the last normalization, the kernel of the quotient cusp residue
on $(w_2,w_1,w_0)$ is generated by
$v=(-4096/1171875,64/1875,1)^t$.
It satisfies $b_{\rm res}(v)=625/192$ and
$\alpha_{\rm alg}(v)=4/3$.
The logarithmic monodromy is $2\pi i$ times this residue.
The exact cusp identity of P1.8 characterizes the contraction on this
line by $N_V(\widetilde v)=\alpha(v)e$.
Since the space of global contraction forms is one-dimensional, the
ratio $(625/192)/(4/3)=625/256$ fixes its scalar everywhere and gives
the displayed $625\pi i/128$.

\end{enumerate}
\end{proof}


\paragraph{An explicit formula for $h$ and the normalization comparison.}
Let $Y'=Y\Gamma_6$.  Normalize $Y(1)=X$ by requiring that continuation
around the specified meridians intertwine exactly with $A,B$, and that
$Xe_0=e_0$.  Existence follows from P1.15; uniqueness follows from the
scalar centralizer, whose exact linear equations are in
\texttt{centralizer\_torsion.py}.
At each nonspecial point there is a unique pair $v_2,v_1\in\mathbb C$
such that the lower five coordinates of
$Y(0,v_2,v_1,-i,0,0)^t$ are real.
The intersection-line proof in P1.8 and the normalization of $\alpha$
prove existence and uniqueness.  Therefore the following is a fully
specified period formula for the real function associated with
$\sigma_b$:
\[
h_b(\lambda)=\frac{96}{625\pi}
 \operatorname{Im}\bigl(Y_{01}v_2+Y_{02}v_1-iY_{03}\bigr).
\tag{1.15.3}
\]
Its periods for the real parameter $c$ are the three columns
$Y_{*,1},Y_{*,2},Y_{*,3}-(625\pi/96)c\,e_0$.

This is a different splitting normalization from the one chosen in
P1.8.  To compare them exactly, parallel-transport the fixed marking
from $\lambda=1$ to $\lambda=\lambda_*/2$ along a specified path in the
upper half-plane, and put
$a=(\sigma_b(g))_e$ at that point, where
$g=(96/(625\pi i))w_0$ and $\alpha(g)=1$.
Then the splitting of P1.8 is
$\sigma=\sigma_b-a\alpha e$, so $h=h_b-\operatorname{Im}a$.
For the full complex splitting parameter $\xi$, the comparison is
$\xi_b=\xi-a$.  This gives the actual translation between the two
families, including its possible real part; it does not silently
identify their real parameter slices.



\begin{theorem}\label{night:P1.16}
\StatementRecord{night:P1.16}{proved}{Historical events: GPT-6 Astra (OpenAI), ChatGPT Pro mode, effort Pro, S6\_06\_BLIND, 16--18 November 2026: targeted proof check. Also checked in this writing. Pre-continuation writing events: GPT-6 Astra (OpenAI), ChatGPT Pro mode, effort Pro, S6\_10\_PAPERS, 18 November 2026 02:02:02.394029+00:00: computation or formal certificate; targeted proof check. Continuation mathematical and source checks: GPT-6 Astra (OpenAI), ChatGPT Pro mode, effort Pro, S6\_10\_PAPERS continuation, 19 November 2026: computation or formal certificate. Continuation metadata events: GPT-6 Astra (OpenAI), ChatGPT Pro mode, effort Pro, S6\_10\_PAPERS continuation, 19 November 2026: used and rechecked in this run; status-vocabulary reconciliation; metadata reconciliation only; no new mathematical check. }{\cref{night:P1.15}, \cref{night:P1.18}, \cref{night:P1.8}}
In the canonical splitting $\sigma_b$ of P1.15,
\[
\begin{aligned}
-1.0342766224163217994942488453783461&<h_b(1),\\
h_b(1)&<-1.0342766224163217994942488453783460.
\end{aligned}
\]
In particular $h_b(1)<-1$.  The parameter $c=5/(3\pi)$ is outside the
global transversality domain.  The lifted period matrices at
$\lambda=1,2$ have been evaluated by two different continuation methods
agreeing to more than thirty decimal digits.
\end{theorem}
\begin{proof}
\begin{enumerate}[label=\arabic*.]
\item The file \texttt{certified\_lift\_bound.py} propagates complex balls
without dropping their radii.  We describe the explicit truncation
bound so the numerical inequality has a mathematical certificate.
On a Taylor step from $z$ to $z+h$, choose a disk of radius $R$ avoiding
the two finite poles.  Use a fixed positive diagonal matrix $S$ on
this disk.  Each entry of the connection is a polynomial divided by
$\lambda^p(\lambda-\lambda_*)$; the triangle inequality gives a
rigorous bound $M$ for the maximum row-sum norm of
$S^{-1}\Gamma_6S$ on the disk.  Along a radius, the integral equation
and Gronwall give a bound $\exp(MR)$ for the fundamental solution
normalized to the identity at its center.  Cauchy's coefficient
estimate then bounds the omitted Taylor tail after degree $N$ by
\[
\exp(MR)\frac{(|h|/R)^{N+1}}{1-|h|/R}.
\tag{1.16.1}
\]
The program adds this error to each entry with factor $S_i/S_j$.
All coefficient evaluations, matrix products, and linear solves retain
their outward-rounded complex-ball enclosures.

\item For both finite meridians we use rational polygonal paths, and choose
the order adaptively until (1.16.1) is below $10^{-65}$.
There are $152$ steps in total and the largest Taylor order is $239$.
The thirty independent equations determining the normalized
intertwiner are solved by interval matrix inversion; exact existence
and uniqueness are P1.15 and the scalar-centralizer calculation.
The real intersection vector in (1.15.3) is likewise obtained by an
invertible four-equation subsystem, whose exact consistency follows
from P1.8.  Substitution into (1.15.3) gives the certified enclosure
\[
\begin{aligned}
h_b(1)\in\bigl[&
-1.03427662241632179949424884537834605259496477790218821132\\
&\qquad\mathbin{\pm}8.11\cdot10^{-57}\bigr].
\end{aligned}
\]
which implies the stated slightly wider decimal bounds.
Every step radius, norm bound, order and tail bound is exported in
\texttt{certified\_lift\_bound.json}.  Since $5/(3\pi)<1$, the value
$h_b(1)+5/(3\pi)$ is negative, whereas $h_b+c$ tends to positive
infinity at the cusp.  A path joining these points must contain a
zero.  This proves inadmissibility, not merely failure of a numerical
test at one point.

\item For the independent period evaluation,
\texttt{frobenius\_period\_check.py} uses regular finite-root frames
and their deck eigencharacters, without integrating a closed loop.
At zero the frame powers are $(0,6,3,0,-3,-6)$ in $u$ with
$\lambda=u^4$, and at infinity they are
$(0,-8,-4,0,4,8)$ with $\lambda=u^{-5}$.
The transformed connections are regular at $u=0$ and have denominators
$3125u^4+256$ and $3125+256u^5$, respectively.
Their Taylor recursions therefore give the connection matrices from
the center to $u=1/4$ and $u=1/2$.  Conjugating the exact diagonal
deck characters by these matrices gives the finite monodromies.
Their normalized intertwiner agrees with the closed-loop calculation
to $4.5\cdot10^{-94}$ in maximum entry norm.
The resulting period planes at $\lambda=1,2$ and their real
transversality determinants are all recorded in the accompanying
JSON files; for example at $c=10$ the determinants are purely imaginary
with positive imaginary parts
\[
\begin{gathered}
88279007523.2167451846066607088299709479845\ldots,\\
3178499608.3637608517583236010752861132441\ldots.
\end{gathered}
\]
These last two local checks do not by themselves prove that $c=10$
is globally admissible.  Global admissibility is exactly P1.8's
inequality involving the minimum.  The independent root-chart
calculation also gives
\[
\begin{split}
h_b(\infty)&=-1.094540250311402780116268263054945703725533\ldots,\\
h_b(0)&=-0.931793376356523505430795105796011114990828\ldots.
\end{split}
\]
P1.18 proves, by a complete interval covering argument, that the
first value is the unique global minimum.

\end{enumerate}
\end{proof}




\begin{theorem}\label{night:P1.18}
\StatementRecord{night:P1.18}{proved}{Historical events: GPT-6 Astra (OpenAI), ChatGPT Pro mode, effort Pro, S6\_06\_BLIND, 16--18 November 2026: targeted proof check. Also checked in this writing. Pre-continuation writing events: GPT-6 Astra (OpenAI), ChatGPT Pro mode, effort Pro, S6\_10\_PAPERS, 18 November 2026 02:02:02.394029+00:00: computation or formal certificate; targeted proof check. Continuation mathematical and source checks: GPT-6 Astra (OpenAI), ChatGPT Pro mode, effort Pro, S6\_10\_PAPERS continuation, 19 November 2026: computation or formal certificate; targeted proof check; used and rechecked in this run. Continuation metadata events: GPT-6 Astra (OpenAI), ChatGPT Pro mode, effort Pro, S6\_10\_PAPERS continuation, 19 November 2026: used and rechecked in this run; status-vocabulary reconciliation; metadata reconciliation only; no new mathematical check. \emph{Continuation recorded proof scope.} Every step of the current seven-step proof has renewed prereading coverage through the root geometric derivation, the Paper VI analytic derivation, and two isolated ledger derivations for the actual polynomial-bound and anchor maps. Each installed scope was compared after its derivation was frozen. Independently authored programs verify the complete 12911-leaf domain certificate, central inverse and Hessian, exact symmetry and cusp maps; a separate period method verifies the normalized constants. The renewed main-result and primary-input credit comparison is recorded. The full former proof is preserved without retrospective prereading credit.}{\cref{night:P1.15}; the rational connection and root-gauge definitions displayed in \cref{night:P1.16}; \cref{night:P1.8}; \cite[Lemma~2.7, pp.~42--43]{TeschlODE2012}; \cite[\S1.9(iii), (1.9.30)--(1.9.31)]{DLMFCauchy2026}}
In the algebraic splitting $\sigma_b$, the function $h_b$ has its
unique global minimum at the finite-root point $\lambda=\infty$.
The minimum is the exact Frobenius period constant
\[
 h_{b,\min}=-\frac{96}{625\pi}\operatorname{Re}(XH_\infty)_{03},
 \qquad
 h_{b,\min}=-1.0945402503114027801162682630549457037\ldots,
\tag{1.18.1}
\]
where $X$ is the normalized integral intertwiner of P1.15 and
$H_\infty$ is the regular infinity-root limit of its fundamental
solution, normalized to be the identity at $\lambda=1$ before
multiplication by $X$.  Formula (1.18.1), the rational connection,
and those two normalizations specify the constant exactly.
A certified enclosure has radius $4.04\cdot10^{-56}$ about the
56-digit midpoint retained in
\path{code/part1/minimum_certificate_summary.json}.
Consequently the exact domain of completed transverse periods is
\[
 c>-h_{b,\min},\qquad
 \operatorname{Im}\xi>-h_{b,\min}
\]
for the real and complex parameters, respectively.
The member $c=10$ used in the numerical period checks is therefore
admissible on the whole completed base.

For the point-normalized splitting of P1.8, continued from
$1$ to $\lambda_*/2$ by the polygon
$1,1+i,\lambda_*/2+i,\lambda_*/2$, the exact comparison constant
$a_{\rm split}$ of (1.15.4) has certified value
\[
\begin{split}
 \operatorname{Re}a_{\rm split}
   &=-0.7099885679126427771723473455012308149\ldots,\\
 \operatorname{Im}a_{\rm split}
   &=-0.9134017059364153362143229223828373269\ldots.
\end{split}
\]
Thus the minimum in that splitting is
$-0.1811385443749874439019453406721083768\ldots$,
attained at the same base point.  Its real parameter interval is
\[
 c>0.1811385443749874439019453406721083768\ldots.
\]
The full parameter comparison remains
$\xi_b=\xi_{\rm original}-a_{\rm split}$; both its real and imaginary
parts are necessary to identify the two normalized families.
\end{theorem}
\begin{proof}
\begin{enumerate}[label=\arabic*.]
\item \textbf{The actual height, parameter, and real coordinate maps.}
Put $k=625\pi/96$, $c_0=k^{-1}$, and use the specified right
connection $H'=H\Gamma_6$, $H(1)=I$, with $Y=XH$.
For the uniquely normalized real quotient-intersection vector write
\[
 q=(0,v_2,v_1,-i,0,0)^t,\qquad z=Yq,
 \qquad h_b=c_0\operatorname{Im}z_0.
\]
The lower five entries of $z$ are real. Every marked monodromy
matrix $T$ is real and fixes $e_0$. Its top entry on $z$ is
$z_0$ plus a real linear combination of those lower entries, so
$\operatorname{Im}(Tz)_0=\operatorname{Im}z_0$. This proves
branch independence in the given marking.
The normalization gives $\alpha(q)=k$, since the coefficient of
$w_0$ is $-i$ and $\alpha(w_0)=ki$. The lifted vector for parameter
$\xi$ therefore has top coordinate $z_0+k\xi$ and the same lower
coordinates. The exact obstruction in the real quotient-intersection
line is
\[
 \operatorname{Im}(z_0+k\xi)=k(h_b+\operatorname{Im}\xi).
\]
The transversality criterion of \cref{night:P1.8} consequently
uses this same function for complex and real parameters.

The numerical preconditioner has an explicit real coordinate map.
For a constant limiting root frame $H_c$, let $B$ have lower-five
columns
\[
 \operatorname{Im}(H_c)_1,\quad\operatorname{Re}(H_c)_1,
 \quad\operatorname{Im}(H_c)_2,\quad\operatorname{Re}(H_c)_2,
 \quad\operatorname{Im}(H_c)_3.
\]
The certified inverse of $B$ defines
\[
 r=(\operatorname{Im}(H_c)_{01},\operatorname{Re}(H_c)_{01},
       \operatorname{Im}(H_c)_{02},\operatorname{Re}(H_c)_{02},0)B^{-1},
 \qquad
 T_c=\begin{pmatrix}1&-r\\0&B^{-1}\end{pmatrix}.
\]
This is a real invertible matrix fixing $e_0$.
Realness of $B^{-1}z_{\rm low}$ is equivalent to realness of
$z_{\rm low}$, and the change of top entry is real on that vector.
Thus it preserves both the normalized pair and $h_b$ exactly.
Writing $v_2=a+ib$, $v_1=c+id$ gives the actual real equations
with row
$(\operatorname{Im}Y_{i1},\operatorname{Re}Y_{i1},
\operatorname{Im}Y_{i2},\operatorname{Re}Y_{i2})$
and right-hand entry $\operatorname{Re}Y_{i3}$.
Approximate row selection only chooses four rows; the inverse
and every ensuing enclosure are certified in ball arithmetic.

\item \textbf{The marked conjugation map.}
The connection has real rational coefficients. Conjugating the
specified ordered polygonal meridians gives their inverse paths.
The actual matrix $K_0=X\overline X^{-1}$ therefore satisfies
$K_0A=A^{-1}K_0$, $K_0B=B^{-1}K_0$, and $K_0e_0=e_0$.
The normalized intertwiner uniqueness in \cref{night:P1.15}
applies: the quotient of two invertible solutions is in the
common centralizer fixing $e_0$. Direct exact solution gives
\[
 K=\begin{pmatrix}
 1&2&-1&-1&2&-1\\
 0&0&0&0&-1&1\\
 0&0&-1&0&0&1\\
 0&0&0&-1&0&1\\
 0&-1&0&0&0&1\\
 0&0&0&0&0&1
 \end{pmatrix}.
\]
Multiplication verifies the two intertwiner equations,
$Ke_0=e_0$, $K^2=I$, and $\det K=-1$, hence $K=K_0$.
On conjugate continuation branches,
$Y(\overline\lambda)=K\overline{Y(\lambda)}$.
The vector $q'=-\overline q$ has the required $-i$ normalization,
and its image is $-K\overline z$, with real lower coordinates.
Uniqueness therefore identifies the normalized intersection
vector and gives
$h_b(\overline\lambda)=h_b(\lambda)$.
This is the actual coordinate map used in reflecting the covered
regions.

\item \textbf{Rigorous analytic enclosures.}
The following formulas specify the complete-domain computation.
All matrix coefficients come from the displayed rational
connection and the specified gauges of
\cref{night:P1.15,night:P1.16}. A regular root frame is the
pullback followed by its stated diagonal gauge $D$; its right
connection is exactly
$D^{-1}\lambda'\Gamma_6D+D^{-1}D'$.
At a positive anchor the constant limiting frame is obtained by
the actual map $H_c=Y_{\rm root}(z_c)U(z_c)^{-1}$, using the
normalized continuation and a certified matrix inverse.

\paragraph{The analytic enclosures and their actual matrix maps.}
For positive weights $w_i$ let $W=\operatorname{diag}(w_i)$. The
specified norm is the maximum absolute row-sum norm of the actual
conjugate $W^{-1}MW$:
\[
 \|M\|_w=\max_i\sum_j|M_{ij}|\frac{w_j}{w_i}.
\]
The equality $W^{-1}MNW=(W^{-1}MW)(W^{-1}NW)$ and the rowwise
triangle inequality prove
$\|MN\|_w\le\|M\|_w\|N\|_w$ and
$|M_{ij}|\le\|M\|_w w_i/w_j$. For matrices with centres $C,D$
and entrywise error bounds $E,F$, multiplication has entrywise
error at most
\[
 |C|F+E|D|+EF,
\]
with the computed rounding error also retained. In particular, a
left factor $H_c$ acts on an entrywise error matrix by the actual
sum $|H_c|_+E$, where $|H_c|_+$ bounds its entries including their
ball radii.

\paragraph{Root series.}
The listed regular-root connection matrices already include the
specified diagonal gauges. Their frame map is
$E_{\rm root}=E_{\rm original}D_{\rm root}$, and the numerical
coordinate change is exactly
$G_z(z)=\alpha G_{\rm root}(\alpha z)$, where
$\alpha^4=256/3125$ at zero and $\alpha^5=3125/256$ at infinity.
Multiplying $G_{\rm root}(t)$ by $1+t^m/\alpha^m$ gives the exact
rational polynomial $\sum_{j=0}^{m-1}C_jt^j$, so
$B_j=\alpha^{j+1}C_j$. Thus
\[
 U'=U\frac{\sum_{j=0}^{m-1}B_jz^j}{1+z^m},\qquad U(0)=I.\tag{1.18.2}
\]
Coefficient comparison in $(1+z^m)U'=UB$ gives
\[
 (n+1)U_{n+1}
 =\sum_{j=0}^{\min(m-1,n)}U_{n-j}B_j
   -(n-m+1)U_{n-m+1},
\]
where the last term is absent for $n<m$. The order of every
matrix product is as displayed. Picard iteration on any smaller
disc converges by the bound $(Lr)^k/k!$ on its $k$-fold integral,
so these are the coefficients of the actual normalized solution.

Use $w=(1,1/128,1/128,1/32,1/64,1/32)$ and
\[
 \mathfrak b=\max_i\sum_{j,k}|(B_k)_{ij}|\frac{w_j}{w_i}.
\]
On a ray of radius $t<1$ the connection norm is at most
$\mathfrak b/(1-t^m)$. The scalar Gronwall inequality
\cite[Lemma~2.7, pp.~42--43]{TeschlODE2012} applies to the actual
matrix integral equation. The elementary estimate
\[
 \frac1{1-t^m}\le\frac1{m(1-t)}+\frac{m-1}{m}
\]
holds because its numerator difference is
$\bigl(\sum_{j=1}^{m-1}t^j-(m-1)t^m\bigr)/m\ge0$.
Consequently
\[
 \|U(z)\|_w\le K(R):=
 \exp\!\left\{\mathfrak b\left[-\frac{\log(1-R)}m
                 +\frac{m-1}{m}R\right]\right\}
 \quad (|z|\le R<1).
\]
Cauchy's coefficient integral, with the positively oriented circle
of radius $R$ \cite[\S1.9(iii), (1.9.30)--(1.9.31)]{DLMFCauchy2026},
gives $\|U_n\|_w\le K(R)R^{-n}$. Hence the error after degree $N$
on the whole disc $|z|\le r<R$ is
\[
 \tau_N(r;R)=K(R)\frac{(r/R)^{N+1}}{1-r/R}.\tag{1.18.3}
\]
Its $ij$ entry is bounded by $\tau_Nw_i/w_j$; after left
multiplication by $H_c$ the bound is
$\tau_N\sum_l|(H_c)_{il}|_+w_l/w_j$.
The current root computations use $N=3000$, $R=0.999$ and each
rectangle's actual maximum radius. Every coefficient and point
rounding radius remains in the corresponding ball calculation.

\paragraph{Enclosing an entire rectangle.}
Let $z_c$ be its centre, $\rho$ its Euclidean half-diagonal, and
$\mathcal A$ a nonnegative entrywise bound for $|G(z)|$ on the
whole rectangle. For any point $z$ the segment
$z_c+t(z-z_c)$ remains in it. Iteration of the integral equation
on that segment bounds the $k$-th ordered integral by
$|Y(z_c)|_+|z-z_c|^k\mathcal A^k/k!$. Subtracting its zeroth term
and summing proves the entrywise simultaneous enclosure
\[
 |Y(z)-Y(z_c)|
 \le |Y(z_c)|_+\bigl(\exp(\rho\mathcal A)-I\bigr).
\]
The centre enclosure and matrix-exponential rounding radii are
retained, so this bounds every point of the rectangle.

The midpoint-polynomial bound has the same explicit domain
control. For $p(z)=\sum_{j=0}^Nc_jz^j$, coefficient discs with
midpoints $\widehat c_j$ and radii $\epsilon_j$, and midpoint
evaluation rounding radius $\eta$, the centre error is
$\eta+\sum_j\epsilon_j|z_c|^j$. Integrating $p'$ on each segment
bounds the remaining variation by
\[
 \rho\sum_{j=1}^N j\bigl(|\widehat c_j|+\epsilon_j\bigr)r^{j-1},
\]
where $r$ is the rectangle's maximum modulus. Add the analytic
tail at that same $r$. This formula includes coefficient
uncertainty in the derivative bound as well as at the centre.

\paragraph{The cusp holomorphic factor.}
The exact connection here is
$R/q+A_0/(1+q)$, with $R^2=0$. In
$Y=H(I+R\log q)P(q)$, $P(0)=I$, the inverse of the logarithmic
factor is $I-R\log q$. Multiplication by that inverse gives
\[
 P'=\frac{\operatorname{ad}(P)}q+\frac{PA_0}{1+q},
 \qquad\operatorname{ad}(Q)=QR-RQ.
\]
Thus
\[
 (nI-\operatorname{ad})P_n
  =\sum_{j=0}^{n-1}(-1)^{n-1-j}P_jA_0.
\]
Equivalently its right side is
$\operatorname{ad}(P_{n-1})+P_{n-1}A_0-(n-1)P_{n-1}$.
Direct expansion gives
$\operatorname{ad}^2(Q)=-2RQR$ and
$\operatorname{ad}^3(Q)=0$, so the actual inverse is
\[
 (nI-\operatorname{ad})^{-1}
  =\frac1n\left(I+\frac{\operatorname{ad}}n
                   +\frac{\operatorname{ad}^2}{n^2}\right).
\]
This fixes both the commutator sign and the side of the $A_0$
multiplication.

Use the actual cusp weights $(100,1,1,1,1,1)$ and put
$r_0=\|R\|_w$, $a_0=\|A_0\|_w$,
$C_0=1+2r_0+4r_0^2$, $\gamma=C_0a_0$. For the supplied exact
matrices,
\[
 r_0=\frac{5103}{8000},\quad a_0=\frac{367}{512},\quad
 C_0=\frac{62452609}{16000000},\quad
 \gamma=\frac{22920107503}{8192000000}.
\]
The norm of the inverse operator is at most $C_0/n$. Therefore
$\|P_n\|_w\le(\gamma/n)\sum_{j<n}\|P_j\|_w$.
The nonnegative scalar majorant $c_0=1$,
$c_n=(\gamma/n)\sum_{j<n}c_j$ satisfies
$c_n=(\gamma+n-1)c_{n-1}/n$ and has generating function
$(1-s)^{-\gamma}$. Its radius is at least one by the coefficient
ratio. Thus the actual matrix series converges normally on
$|q|<1$, solves the equation term by term, and obeys
\[
 \|P(q)\|_w\le K_0(s):=(1-s)^{-\gamma}\quad(|q|\le s<1).
\]
The coefficient integral and geometric sum give its degree-$N$
tail on $|q|\le r<s$ as
\[
 K_0(s)\frac{(r/s)^{N+1}}{1-r/s},
\]
with the exact entry factor $w_i/w_j$. The current cusp series
uses $N=600$ and $s=0.4$ for $r\le0.25$. On the inner disc
$r=0.01$ the enclosure retains all actual coefficient absolute
sums through degree600 and this tail. It bounds $P$ itself;
the subsequent logarithmic factor retains its specified branch
and its own actual matrix entries.

\paragraph{Collar continuation.}
At a centre $z_c$ put $t=z-z_c$ and expand
\[
 1+(z_c+t)^5=\sum_{j=0}^5q_jt^j,
 \quad q_j=\binom5jz_c^{5-j}+\delta_{j0},
 \qquad
 \widetilde B_j=\sum_{k=j}^4\binom kjB_kz_c^{k-j}.
\]
For $T(0)=I$ the exact transition coefficients satisfy
\[
 T_{n+1}=
 \frac{\displaystyle
   \sum_{j=0}^{\min(4,n)}T_{n-j}\widetilde B_j
  -\sum_{j=1}^{\min(5,n)}q_j(n-j+1)T_{n-j+1}}
 {(n+1)q_0}.
\]
Let $\zeta_j=e^{(2j+1)\pi i/5}$ and
$R_{\rm loc}=\tfrac12\min_j|z_c-\zeta_j|$. On the closed disc
of this radius,
\[
 |1+(z_c+t)^5|\ge D_{\rm lower}
   :=\prod_j\bigl(|z_c-\zeta_j|-R_{\rm loc}\bigr)>0.
\]
The numerator norm is at most
\[
 B_{\max}=\max_i\sum_{j,k}|(\widetilde B_k)_{ij}|
                  \frac{w_j}{w_i}R_{\rm loc}^k.
\]
The actual complete-domain code uses the further upper bound
\[
 \widehat B_{\max}=\sum_{k=0}^4\|B_k\|_w
                           (|z_c|+R_{\rm loc})^k.
\]
To see the exact comparison, substitute the displayed formula for
$\widetilde B_l$ in each row sum and use the triangle inequality.
Reindex the finite set $0\le l\le k\le4$ by $(l,k)\mapsto(k,l)$.
The sum in $l$ is
$\sum_{l=0}^k\binom kl|z_c|^{k-l}R_{\rm loc}^l
=(|z_c|+R_{\rm loc})^k$.
Each remaining row sum is at most $\|B_k\|_w$; taking the row
maximum afterward proves $B_{\max}\le\widehat B_{\max}$.
Thus the bound actually used below is
$M=\widehat B_{\max}/D_{\rm lower}$. Radial Gronwall and Cauchy
give the exact degree-$N$ transition tail
\[
 e^{MR_{\rm loc}}
 \frac{(|\delta|/R_{\rm loc})^{N+1}}
      {1-|\delta|/R_{\rm loc}}.
\]
The current steps use $N=160$ and
$|\delta|/R_{\rm loc}<0.3$. The actual continuation order is
$Y(z_c+\delta)=Y(z_c)T_{z_c}(\delta)$, as follows by checking
the right ODE and its initial value. Every transition tail is
added entrywise before multiplying on the right, and the
entrywise product rule propagates all earlier errors.

\paragraph{The real implicit system and Hessian.}
Use the specified fixed four rows on the central square and
write the actual system as $Mv=b$, with
\[
 M_{i,*}=(\operatorname{Im}Y_{i1},\operatorname{Re}Y_{i1},
          \operatorname{Im}Y_{i2},\operatorname{Re}Y_{i2}),
 \qquad b_i=\operatorname{Re}Y_{i3}.
\]
For the top row define $c$ by the same four coordinate projections,
$d=\operatorname{Re}Y_{03}$ and $h=c_0(cv-d)$,
$c_0=96/(625\pi)$. The certified determinant enclosure excludes
zero on the entire central square and on its sixteen subsquares;
the matrix inverse is therefore the actual one there.
For $a,b\in\{x,y\}$, direct differentiation gives
\[
 \begin{split}
 v_a&=M^{-1}(b_a-M_av),\\
 v_{ab}&=M^{-1}(b_{ab}-M_{ab}v-M_av_b-M_bv_a),\\
 h_a&=c_0(c_av+cv_a-d_a),\\
 h_{ab}&=c_0(c_{ab}v+c_av_b+c_bv_a+cv_{ab}-d_{ab}).
 \end{split}
\]
The repeated-index cases retain both cross terms. The actual
complex derivatives supplying these real projections are
\[
 Y_x=YG,\quad Y_y=iYG,\quad
 Y_{xx}=Y(G^2+G'),\quad Y_{xy}=iY(G^2+G'),\quad
 Y_{yy}=-Y(G^2+G').
\]
Taking the stated real and imaginary parts and carrying out these
products with rigorous full-square balls consequently encloses
all Hessian entries throughout the domain. Approximate row
selection merely chooses a minor; it supplies neither its inverse
nor its error bound.

\item \textbf{The infinity value and its strict local minimum.}
Put $a=256/3125$, $\lambda=u^{-5}$ and
\[
 D_\infty(u)=\operatorname{diag}(1,u^{-8},u^{-4},1,u^4,u^8),
 \qquad \widetilde Y(u)=XH(u^{-5})D_\infty(u).
\]
The specified root extension makes $\widetilde Y$ holomorphic
and invertible at zero. The $w_0$ column normalization is unchanged.
In the preconditioned root frame the chosen four-by-four real
system has determinant in
\[
 [\,0.997929591387219260,\ 1.002070408612780740\,]
\]
on the entire central square, with its inverse certified there.
Its inverse defines a real analytic normalized coefficient pair
at zero. The fifth reality equation extends by continuity from
the punctured chart, and any solution must satisfy the four
invertible equations. This proves existence and uniqueness there.

For $\zeta=e^{2\pi i/5}$ set
$T_\zeta=\operatorname{diag}(1,\zeta^{-8},\zeta^{-4},1,
\zeta^4,\zeta^8)$. The actual deck formula is
\[
 \widetilde Y(\zeta u)=A\widetilde Y(u)T_\zeta.
\]
At zero it gives
$A\widetilde Y(0)=\widetilde Y(0)T_\zeta^{-1}$.
Applying $A$ to the normalized real-lower vector changes its
coefficient pair to $(\zeta^8v_2,\zeta^4v_1)$ and preserves
the $-i$ coefficient. Uniqueness forces $v_2=v_1=0$.
Consequently
\[
 h_b(\infty)=-c_0\operatorname{Re}\widetilde Y(0)_{03}
            =-\frac{96}{625\pi}\operatorname{Re}(XH_\infty)_{03}.
\]
For general $u$, the same deck formula and the real monodromy
map give $h_b(\zeta u)=h_b(u)$. Its real differential at zero
is fixed by a rotation through $2\pi/5$, so it vanishes.

The numerical infinity coordinate is exactly
$w=a^{1/5}u$, hence $\lambda=aw^{-5}$.
The sixteen closed subsquares cover
$Q=[-1/10,1/10]^2$ in $(\operatorname{Re}w,\operatorname{Im}w)$.
The implicit derivative formulas above give on their union
\[
 h_{xx}>0.3215,\qquad h_{yy}>0.3220,
 \qquad |h_{xy}|<0.0218.
 \tag{1.18.5}
\]
For every real vector $t=(t_1,t_2)$, the inequality
$2|t_1t_2|\le t_1^2+t_2^2$ yields
$t^tD^2h_b\,t>(299/1000)|t|^2$.
The gradient is still zero after the positive scalar coordinate
change. Integrating twice along the segment in the convex square
therefore gives
\[
 h_b(aw^{-5})-h_b(\infty)
 =\int_0^1(1-s)w^tD^2h_b(sw)w\,ds
 \ge\frac{299}{2000}|w|^2.
\]
Thus the center is the unique minimum on $Q$.

The independent normalized continuations give the following
outward enclosure, whose endpoints are exact terminating decimals:
\[
 \begin{gathered}
 -1.094540250311402780116268263054945703726
       <h_b(\infty),\\
 h_b(\infty)<
 -1.094540250311402780116268263054945703725.
 \end{gathered}
 \tag{1.18.4}
\]
The complete stored enclosure has radius less than
$4.38\cdot10^{-100}$. It lies in the statement's earlier
$4.04\cdot10^{-56}$ enclosure. The independently authored
root-connection computation agrees with the independently
propagated normalized period calculation to more than89 digits;
both enclose the exact expression just proved.

\item \textbf{The closed covers and their strict lower bounds.}
Let $R=49/50$. The infinity sector
$0\le\arg w\le2\pi/5$, $|w|\le R$ maps under
$\lambda=aw^{-5}$ onto $|\lambda|\ge aR^{-5}$.
The square $[0,1/10]^2$ removed from its leaf calculation lies
in $Q$ and has already been treated.
At zero put $t=a^{1/4}v$, so $\lambda=t^4=av^4$.
The first-quadrant sector $|v|\le R$ covers
$|\lambda|\le aR^4$, including the completed zero root.
The marked real matrices $A$ and $B$ identify the respective
sector boundaries and preserve the height.

The intervening radial region is precisely the image of
\[
 R\le|w|\le R^{-4/5},\qquad 0\le\arg w\le2\pi/5.
\]
Indeed its endpoint radii map to $aR^{-5}$ and $aR^4$.
Here the cusp coordinate is exactly
$q=\lambda/\lambda_*-1=-w^{-5}-1$.
The collar certificate treats its closed part $|q|\ge1/4$.
The remaining part is in the cusp disc; the annulus certificate
treats $1/100\le|q|\le1/4$ for all arguments. The punctured
inner disc is treated in the next item. These closed regions
include their shared boundaries and cover the whole punctured
base, with the two finite-root points included as specified.

For each required rectangle, a whole-domain enclosure is either
strictly accepted or the rectangle is subdivided. A discarded
rectangle has empty intersection with the required region or is
contained in the already treated central square. The complete
finite subdivision and domain checks have no unresolved leaves.
The current computation uses110-digit ball arithmetic. Its
unrounded lower endpoints have been exported and replayed against
the same exact rational threshold
\[
 U=-1.0945402503114027801162682630549457037
       >h_b(\infty).
\]
The following terminating decimals are exact uniform lower
bounds, strictly smaller than every recorded lower endpoint
in the indicated region and strictly greater than $U$:
\[
 \begin{array}{lrr}
 \text{region}&\text{accepted leaves}&\text{strict lower bound for }h_b\\
 \hline
 \text{infinity sector outside the square}&1203&-1.094400523035\\
 \text{zero sector}&592&-1.083329522544\\
 \text{cusp annulus}&10780&-1.094402439774\\
 \text{collar outside the cusp disc}&336&-1.092202452969
 \end{array}
\]
All12911 precise endpoint comparisons pass. The endpoint replay
uses the same newly authored complete-domain computation and
retains its full interval endpoints. It adds a directly
replayable inequality record to that computation.

\item \textbf{The actual cusp graph and the inner disc.}
Use $\lambda=\lambda_*(1+q)$ and the actual diagonal gauge
\[
 G=Y\operatorname{diag}(1,\lambda^2,\lambda,1,
                         \lambda^{-1},\lambda^{-2})
   =H_c(I+R_c\log q)P(q).
\]
The cusp series and continuation above give the normalized
$H_c$ and $P$. The matrix $N=2\pi iH_cR_cH_c^{-1}$ is the
logarithm of the actual marked cusp monodromy. Its certified
enclosure contains the following unique integer matrix, and
exact multiplication verifies the stated marking:
\[
 N=\begin{pmatrix}
 0&0&0&-1&1&-1\\0&0&0&0&0&0\\0&0&0&0&0&0\\
 0&1&0&0&0&0\\0&0&1&0&0&0\\0&-1&1&0&0&0
 \end{pmatrix},\qquad
 J=\begin{pmatrix}
 -1&0&0&0&0&0\\0&0&0&0&1&0\\0&0&0&0&0&1\\
 0&1&0&1&0&0\\0&0&1&0&0&0\\0&-1&1&0&0&0
 \end{pmatrix}.
\]
In particular $N^2=0$, $J$ is unimodular,
$J^{-1}NJ=\left(\begin{smallmatrix}0&I_3\\0&0\end{smallmatrix}\right)$,
and $e_J:=J^{-1}e_0=(-1,0,0,0,0,0)^t$.
Put $c_*=-h_b(\infty)$, $s=kc_*$, and $V=J^{-1}H_cP(q)$.
The three actual period columns are
$V_1,V_2,V_3-se_J$. Their upper and lower blocks will be
denoted by $U_0,L$. The right scaling from the algebraic period
columns is exactly $\operatorname{diag}(\lambda^2,\lambda,1)$;
it preserves their span and the third-column shift.
The inverse of $L$ is certified on the entire closed disc
$|q|\le1/100$.
Since $Ne_0=0$, the logarithmic factor acts on these shifted
columns by
$\left(\begin{smallmatrix}I&\ell I\\0&I\end{smallmatrix}\right)$,
where $\ell=\log q/(2\pi i)$. The actual graph is therefore
\[
 Z(q)=\frac{\log q}{2\pi i}I_3+C(q),\qquad C(q)=U_0L^{-1}.
\]
For $B(q)=\operatorname{Im}C(q)$ and
$y=-\log|q|/(2\pi)$, expansion of the determinant gives
\[
 \det\operatorname{Im}Z
 =y^3+\operatorname{tr}(B)y^2
  +\sum_{i<j}(B_{ii}B_{jj}-B_{ij}B_{ji})y+\det B.
 \tag{1.18.6}
\]
Put $y_0=\log100/(2\pi)$ and $t=y-y_0\ge0$.
The four coefficients of
$\det(tI_3+y_0I_3+B(q))$, in increasing degree, have the
uniform lower bounds
\[
 1.9918,\qquad4.9602,\qquad3.9580,\qquad1,
\]
with the last coefficient exactly one. Thus the determinant is
positive on the punctured inner disc.
In the actual real $J$ coordinates, a graph vector $(Zv,v)$
is real precisely when $v$ is real and $(\operatorname{Im}Z)v=0$.
The determinant therefore proves transversality of the
$c_*$ plane there. By the first item's exact obstruction map,
$h_b-h_b(\infty)$ never vanishes. The punctured disc is connected,
and at $q=1/100$ the certified value
$h_b=-0.4504422009074\ldots$ is greater than $h_b(\infty)$.
Consequently $h_b>h_b(\infty)$ throughout this disc.

For completeness these same graph maps prove the cusp limit.
The lower block of $e_J$ vanishes, so $L$ is independent of a
real parameter $c$, and $C_c$ is holomorphic at zero and affine
in $c$. For any fixed real bound $M$, it is uniformly bounded
on a smaller closed disc and the compact interval between
$c_*$ and $-M$. As $y\to\infty$, the Neumann inverse makes
$yI+\operatorname{Im}C_c$ invertible uniformly on that interval.
The real quantity $h_b+c$ cannot cross zero along it. Its sign
at $c_*$ is positive, so $h_b-M>0$ at the other endpoint.
This proves $h_b\to+\infty$ at the cusp. The covers and the
strict local estimate now prove the global minimum and its
uniqueness.

\item \textbf{The parameter domain and the full splitting translation.}
The first item's obstruction map and \cref{night:P1.8} give
$c>-h_b(\infty)$ and $\operatorname{Im}\xi>-h_b(\infty)$
on the completed base. The enclosure above gives
$-h_b(\infty)<1.095<10$.
For the stated upper-half-plane polygon, let
$g=(96/(625\pi i))w_0$, so $\alpha(g)=1$.
The exact top coordinate of its algebraic lift is
\[
 a_{\rm split}=\frac{96}{625\pi i}Y_{03}(\lambda_*/2).
\]
The defining equality
$\sigma_{\rm original}=\sigma_b-a_{\rm split}\alpha e$
gives the actual plane equality
\[
 (\sigma_{\rm original}+\xi_{\rm original}\alpha e)(P)
  =(\sigma_b+(\xi_{\rm original}-a_{\rm split})\alpha e)(P).
\]
Hence $\xi_b=\xi_{\rm original}-a_{\rm split}$ and, on the
normalized vector with $\alpha(q)=k$,
$h_{\rm original}=h_b-\operatorname{Im}a_{\rm split}$ pointwise.
The continuation gives the stated real and imaginary parts of
$a_{\rm split}$ with radii less than $4.41\cdot10^{-100}$ and
$1.41\cdot10^{-100}$, respectively. The resulting minimum is
\[
 -0.181138544374987443901945340672108376803615204418668779889\ldots,
\]
with stored radius less than $5.77\cdot10^{-100}$, attained
at the same base point. The original real line
$\xi_{\rm original}=ic$ maps exactly to
\[
 \xi_b=-\operatorname{Re}a_{\rm split}
          +i(c-\operatorname{Im}a_{\rm split}).
\]
Both coordinates are retained, and its criterion is precisely
$c>-h_{{\rm original},\min}$.

The exact coefficient arrays, finite subdivision trees, actual
matrix inverses and unrounded endpoint certificates are retained
in \path{code/continuation_V/}; in particular the complete
region summary is
\path{minimum_precise_endpoint_summary.json} there.
The separate analytic certificate uses256-bit balls and713
checks of the displayed recurrences and matrix maps; its tested
representative domains are recorded separately from the complete
110-digit leaf partition used above.
\end{enumerate}
\end{proof}


\section{Cohomology, deformations and isomorphisms}

We retain the complex orientation and the integral marking of P1.1.
Let $\widehat B$ be the root orbifold of orders $5,4$, let
$h:M\to\widehat B$ be the map supplied by the two multiple divisors,
and let $\pi:\widehat B\to\mathbb P^1$ be its coarse map.
Write $L=T^2$; in the pencil coordinate it is
\[
 L=\mathcal O_{\widehat B}\bigl(\tfrac45[\infty]
                  -\tfrac34[0]\bigr),\qquad \deg L=\tfrac1{20}.
\]
The cusp point has line bundle $L^{20}$, and
$T_{\widehat B}=L^9$. These identities include the stabilizer
characters: their quotients have integral divisors of degree zero
on $\mathbb P^1$. Denote the cusp by $W$, its normalization by $Y$,
and its reduced singular locus by $D$.

\begin{theorem}\label{night:P1A.1}
\StatementRecord{night:P1A.1}{proved}{Historical events: GPT-6 Astra (OpenAI), ChatGPT Pro mode, effort Pro, S6\_06\_BLIND, 16--18 November 2026: targeted proof check. Also checked in this writing. Pre-continuation writing events: GPT-6 Astra (OpenAI), ChatGPT Pro mode, effort Pro, S6\_10\_PAPERS, 18 November 2026 02:02:02.394029+00:00: computation or formal certificate; targeted proof check. Continuation mathematical and source checks: GPT-6 Astra (OpenAI), ChatGPT Pro mode, effort Pro, S6\_10\_PAPERS continuation, 19 November 2026: computation or formal certificate; independent re-proof; used and rechecked in this run. Continuation metadata events: GPT-6 Astra (OpenAI), ChatGPT Pro mode, effort Pro, S6\_10\_PAPERS continuation, 19 November 2026: used and rechecked in this run; status-vocabulary reconciliation; metadata reconciliation only; no new mathematical check. }{\cref{night:P1.10}, \cref{night:P1.8}, \cref{night:P1.9}}
For every admissible member, its vertical translation bundle on
$\widehat B$ and its first coherent direct image are
\[
 E_0=F^\vee=\mathcal O\oplus L^{-1}\oplus L^{-2},\qquad
 G=R^1h_*\mathcal O_M=\mathcal O\oplus L^{-1}\oplus L^{-2}.
\]
Moreover $R^qh_*\mathcal O_M=\bigwedge^qG$ for $0\le q\le3$,
$K_M=h^*L^{-6}$, and
\[
 (h^0(\mathcal O_M),\ldots,h^4(\mathcal O_M))=(1,1,0,0,0).
\]
\end{theorem}
\begin{proof}
\begin{enumerate}[label=\arabic*.]
\item The evaluation of the marked lattice on the lifted plane $F$
identifies the Lie bundle of the smooth tori with $F^\vee$.
P1.8 gives $F=\mathcal O\oplus L\oplus L^2$; its lift as a bundle
has the same splitting. The logarithmic toric vector fields extend
this Lie bundle over the cusp, so the identification holds on the
whole orbifold, including both finite root charts.

\item On the smooth locus the constant-cocycle description of torus
cohomology gives $G=V/F$. In the rank-five orthogonal quotient,
$K=\ker(\alpha:P\to\mathcal O)=L\oplus L^2$ is the radical of
the restricted pairing on $P$. Indeed every wedge containing the
fixed line pairs to zero with another such wedge except through
their contraction components; removing that scalar component
leaves exactly $K$. The nondegenerate pairing consequently
identifies $W/P$ with $K^\vee$. Thus
\[
 0\longrightarrow\mathcal O e\longrightarrow V/F
   \longrightarrow L^{-1}\oplus L^{-2}\longrightarrow0.
\]
For every integer $j$, descent by invariant local Laurent series gives
\begin{equation}\label{eq:P1Adegree}
 \pi_*L^j=\mathcal O_{\mathbb P^1}(d(j)),\qquad
 d(j)=\lfloor4j/5\rfloor+\lfloor-3j/4\rfloor.
\end{equation}
The extension group $H^1(L\oplus L^2)$ therefore vanishes, so this
sequence splits as a holomorphic vector-bundle sequence.

\item We verify that this quotient is the coherent extension at the cusp.
The normalization resolution of $\mathcal O_W$ has one toric
threefold, seven toric surfaces, twelve toric curves and six points.
All positive coherent cohomology of those smooth complete toric
strata vanishes \cite[Theorem 9.2.3, p.~410]{CoxLittleSchenck2011}. Their constants form precisely the cellular
cochain complex of the dual three-torus, with the incidences in
\path{code/part1/cusp_complex.json}. Its cohomology is
$\bigwedge^q\mathbb C^3$, and its cup product is the exterior product.
The invariant flat cohomology $\ker N$ maps to these classes by
constant Cech cocycles on the toroidal covering. A component
translation in that cover evaluates the corresponding lattice
coordinate. In the nilpotent logarithmic frame, $\ker N$ maps
isomorphically to $V/F$, since $N:F\to\ker N$ is an isomorphism.
Thus these are the actual extending quotient classes, rather than
an unrelated basis of the same dimension.

\item Flatness over each smooth root disc follows because its base is a
DVR and the smooth total space is torsion free over it. Properness
was proved in P1.9--P1.10. The coherent fibre dimensions are the
constant binomial dimensions on smooth tori, finite root fibres,
and the cusp. Grauert's \S7 Satz~5, with these checked hypotheses,
gives local freeness and base change. The cup-product maps are
isomorphisms on every fibre, including the preceding central
cellular model, so they identify the remaining direct images with
$\bigwedge^qG$.

\item There is an exact sequence of locally free logarithmic bundles
\begin{equation}\label{eq:P1Alogcotangent}
0\longrightarrow h^*L^{11}\longrightarrow\Omega_M^1(\log W)
 \longrightarrow h^*F\longrightarrow0.
\end{equation}
In a chart $q=x_0\cdots x_r$, its first arrow sends
$dq/q$ to $\sum dx_i/x_i$ and its quotient is the height-zero
logarithmic cotangent bundle of the torus. These formulas also
prove exactness on every stratum. Taking determinants and
cancelling $\mathcal O_M(W)=h^*L^{20}$ yields
$K_M=h^*L^{11+3-20}=h^*L^{-6}$.

\item For any bundle which is a sum of powers of $L$, projection and
the exact finite averaging projector compute its direct images
by tensoring with $\bigwedge^qG$ and applying
\eqref{eq:P1Adegree}. All exponents $-1,\ldots,-5$ have coarse
degree $-1$, while exponent zero has degree zero. The Leray
sequence on the curve then gives the asserted structure-sheaf
cohomology. This includes the nonreduced coarse fibres, since
the computation was made before exact finite descent.

\end{enumerate}
\end{proof}




\begin{theorem}\label{night:P1A.2}
\StatementRecord{night:P1A.2}{proved}{Historical events: GPT-6 Astra (OpenAI), ChatGPT Pro mode, effort Pro, S6\_06\_BLIND, 16--18 November 2026: targeted proof check. Also checked in this writing. Pre-continuation writing events: GPT-6 Astra (OpenAI), ChatGPT Pro mode, effort Pro, S6\_10\_PAPERS, 18 November 2026 02:02:02.394029+00:00: computation or formal certificate; targeted proof check. Continuation mathematical and source checks: GPT-6 Astra (OpenAI), ChatGPT Pro mode, effort Pro, S6\_10\_PAPERS continuation, 19 November 2026: computation or formal certificate; independent re-proof; targeted proof check; used and rechecked in this run. Continuation metadata events: GPT-6 Astra (OpenAI), ChatGPT Pro mode, effort Pro, S6\_10\_PAPERS continuation, 19 November 2026: used and rechecked in this run; status-vocabulary reconciliation; metadata reconciliation only; no new mathematical check. \emph{Continuation recorded proof scope.} Complete written prereading derivation of the Hodge table by the exact logarithmic connecting maps and a direct residue map; the six-dimensional symmetrization and seven-edge evaluation are computed in \path{code/continuation_root/P1A2_hodge_fresh.py}. The exact change of basis identifies its cusp matrix with the displayed matrix. The additional printed Riemann--Roch step has its separate reviewer receipt; its chronology is not inferred from the alternative proof. The renewed main-result credit comparison positively identifies the earlier conductor and KS method.}{\cref{night:P1.15}, \cref{night:P1A.1}}
The complete Hodge table of every admissible fourfold is the
following, with rows indexed by $p$ and columns by $q$:
\[
\begin{array}{c|ccccc}
h^{p,q}&0&1&2&3&4\\\hline
0&1&1&0&0&0\\
1&0&2&1&0&0\\
2&0&1&6&1&0\\
3&0&0&1&2&0\\
4&0&0&0&1&1
\end{array}
\]
\end{theorem}
\begin{proof}
\begin{enumerate}[label=\arabic*.]
\item All maps needed for the two nontrivial rows can be evaluated at
the cusp. We give their matrices and the reason evaluation
determines their global ranks. In the algebraic splitting of
P1.15 write the three generators of $F$ as $(f_0,f_1,f_2)$ of
weights $(0,1,2)$, and those of $G$ as $(g_0,g_1,g_2)$ of
weights $(0,-1,-2)$. The residue of the quotient connection is
\[
R:F_0\longrightarrow G_0,\qquad
R=\begin{pmatrix}0&A&B\\C&0&D_0\\E&D_0&0\end{pmatrix},
\]
where
\begin{align*}
A&=-125(192z-3125)/12288,&
B&=-78125(96z+3125)/393216,\\
C&=-125/512,&E&=78125/32768,\\
D_0&=-146484375/2097152.
\end{align*}
Here the lifted first vector is $w_0+(4/3)ze$.
These entries result from the exact Gauss--Manin reduction of
P1.15, with the quotient projection taken after lifting $F$.
They are independently recomputed in
\path{code/part1/analytic_cohomology.py}.

\item The connecting map of \eqref{eq:P1Alogcotangent}, on relative
cohomology, sends $f\otimes\eta$ to $R(f)\wedge\eta$ at the
cusp. To see the actual extension class, trivialize logarithmic
forms on the toric covering. A deck translation changes a
relative logarithmic form by its period derivative times
$dq/q$. The constant Cech class of this derivative is exactly
the residue of the quotient connection. This proves the
formula with the indicated frames. It also proves its exterior
power version, with the usual alternating signs.

\item The floor formula gives the following absolute cohomology:
\[
\begin{array}{c|ccccc}
&h^0&h^1&h^2&h^3&h^4\\\hline
h^*F&1&3&2&0&0\\
h^*L^{11}&0&2&3&1&0\\
h^*(\bigwedge^2F)&0&2&3&1&0\\
h^*(L^{11}\otimes F)&2&6&7&3&0.
\end{array}
\]
Every group in this display comes from global sections of the
indicated direct image; none comes from $H^1$ of a line on the
coarse base. Every contributing summand has coarse degree zero,
so evaluation at the cusp is injective on its space of sections.
For the equal-dimensional source and target spaces used below
it is an isomorphism onto the displayed weight subspaces.
Thus fibrewise evaluation computes the actual ranks of the
absolute connecting maps.

\item For one-forms, the three middle source classes are $f_i\otimes g_i$.
In the ordered target $(g_0g_1,g_0g_2,g_1g_2)$, the matrix is
\[
C_1=\begin{pmatrix}-C&A&0\\-E&0&B\\0&-D_0&D_0\end{pmatrix},
\qquad
\det C_1=
\frac{476837158203125}{281474976710656}(384z+3125).
\]
The degree-zero connecting map is nonzero because $R(f_0)$
has nonzero $g_1,g_2$ coefficients. The degree-two connecting
map, from the classes $f_1g_0g_1,f_2g_0g_2$ to $g_0g_1g_2$,
has rank one because $D_0\ne0$.

\item We must exclude the sole possible zero of $\det C_1$ on the
whole admissible parameter domain. P1.15 identifies
$z=-(625\pi/128)c$ on the real slice in its algebraic splitting.
The exact exceptional parameter is therefore $c=5/(3\pi)<1$.
The certified continuation in
\path{code/part1/certified_lift_bound.py} proves
\[
\begin{aligned}
h_b(1)\in\bigl[&
-1.03427662241632179949424884537834605259496477790218821132\\
&\qquad\mathbin{\pm}8.11\cdot10^{-57}\bigr].
\end{aligned}
\]
so $h_b(1)+5/(3\pi)<0$. Since $h_b\to+\infty$ at the cusp,
the intermediate value theorem makes this parameter
nontransverse somewhere on the base. For the full complex
parameter the determinant has the same single zero, and the
same argument excludes it. The certification propagates complex
balls and adds an explicit Cauchy bound for every omitted
Taylor tail; no midpoint replacement discards an error radius.
Its exact connection and all step bounds are retained with the
program. Consequently the three connecting ranks are $1,3,1$.
The long exact sequence of \eqref{eq:P1Alogcotangent} gives
\[
h^q(\Omega_M^1(\log W))=(0,1,1,0,0).
\]

\item The residue sequence
\[
0\to\Omega_M^1\to\Omega_M^1(\log W)\to\nu_*\mathcal O_Y\to0
\]
is exact even at self-intersections: in a local branch chart it
records independently the coefficient of each $dx_i/x_i$ on
$x_i=0$. The normalization $Y$ is smooth complete toric, so its
only coherent structure cohomology is $H^0=\mathbb C$ \cite[Theorem 9.2.3, p.~410]{CoxLittleSchenck2011}.
The one-form row follows immediately.

\item For two-forms take the second exterior power of
\eqref{eq:P1Alogcotangent}. Its short exact sequence has left
term $h^*(L^{11}\otimes F)$ and right term
$h^*(\bigwedge^2F)$. In degree two the three source classes
are $f_if_j\otimes g_ig_j$, $i<j$. The target is
$F\otimes g_0g_1g_2$. Its connecting matrix is
\[
C_2=\begin{pmatrix}-D_0&D_0&0\\E&0&-B\\0&-C&A\end{pmatrix},
\qquad \det C_2=-\det C_1.
\]
Thus it is an isomorphism, and
$h^3(\Omega_M^2(\log W))=1$.

\item For completeness the ordinary/logarithmic comparison in this
degree is the exact residue resolution
\[
0\to\Omega_M^2\to\Omega_M^2(\log W)
 \to\nu_*\Omega_Y^1(\log\partial Y)
 \to\nu^{[2]}_*\mathcal O_{W^{[2]}}\to0.
\]
The last space is the disjoint normalization of the seven
double surfaces. Local exactness is verified by writing a
logarithmic two-form as a sum of
$a_{ij}\,dx_i/x_i\wedge dx_j/x_j$ and terms with ordinary
differentials. The two residues of an $ij$ term have opposite
sign; their cancellation is the sole compatibility. Every local
double-surface coefficient lifts to one branch, which proves
surjectivity at the last term.
Both final sheaves have zero positive cohomology: the first
is the pushforward of the trivial rank-three toric logarithmic
bundle, and the second has complete toric surface components.
Writing $Q=\Omega_M^2(\log W)/\Omega_M^2$, we obtain
$H^q(Q)=0$ for $q\ge2$. Hence
$h^3(\Omega_M^2)=1$.
In degree zero the last residue map is zero: the fourteen toric
boundary divisors of $Y$ occur in seven opposite pairs; their
primitive normal vectors have opposite signs and their gluing
is a torus translation. This gives the zero $7$-by-$3$ residue
matrix on global logarithmic one-forms. Thus $H^0(Q)=\mathbb C^3$.
The two global logarithmic two-forms come from the
$L^{12},L^{13}$ summands of $L^{11}\otimes F$; their residues
are the independent $f_1,f_2$ directions at the cusp.
Therefore $h^0(\Omega_M^2)=0$.

\item Serre duality \cite[Theorem 4, p.~22]{Serre1955Duality} on this compact smooth complex fourfold gives
$h^{2,1}=h^{2,3}=1$. Finally the analytic Riemann--Roch theorem
\cite[Theorem 25.3.1, p.~187; \S25.4(a), pp.~188--190]{HirzebruchTopologicalMethods} applies to the Dolbeault elliptic
operator of $\Omega_M^p$. No Kähler hypothesis occurs there.
Using $c_1=c_3=0$, the Chern-root product gives
\[
\chi_y(M)=
\frac{(1+y)^4}{240}\,c_2^2+
\left(\frac{(-1+10y-y^2)^2}{144}
       -\frac{(1+y)^4}{120}\right)c_4
=-y+4y^2-y^3.
\]
The program verifies this polynomial exactly. It gives
$\chi(\Omega_M^2)=4$, hence $h^{2,2}=6$.
P1A.1 and Serre duality supply all other entries of the table.

\end{enumerate}
\end{proof}




\begin{theorem}\label{night:P1A.3}
\StatementRecord{night:P1A.3}{proved}{Historical events: GPT-6 Astra (OpenAI), ChatGPT Pro mode, effort Pro, S6\_06\_BLIND, 16--18 November 2026: targeted proof check. Also checked in this writing. Pre-continuation writing events: GPT-6 Astra (OpenAI), ChatGPT Pro mode, effort Pro, S6\_10\_PAPERS, 18 November 2026 02:02:02.394029+00:00: computation or formal certificate; targeted proof check. Continuation mathematical and source checks: GPT-6 Astra (OpenAI), ChatGPT Pro mode, effort Pro, S6\_10\_PAPERS continuation, 19 November 2026: computation or formal certificate; independent re-proof; targeted proof check; used and rechecked in this run. Continuation metadata events: GPT-6 Astra (OpenAI), ChatGPT Pro mode, effort Pro, S6\_10\_PAPERS continuation, 19 November 2026: used and rechecked in this run; status-vocabulary reconciliation; metadata reconciliation only; no new mathematical check. \emph{Continuation recorded proof scope.} Complete tangent-duality and Kodaira--Spencer derivation written before reading the prior target proof. The six-step proof printed here is that derivation. Its exact bundle pairing, residue, normal tangent sequence and analytic-ring maps are explicit; the fresh program \path{code/continuation_root/P1A3_tangent_fresh.py} checks all weights, the base determinant and the parameter rank. The superseded conductor argument is preserved with its original checking scope.}{\cref{night:P1.9}, \cref{night:P1.10}, \cref{night:P1.15}, \cref{night:P1A.1}, \cref{night:P1A.2}}
For every admissible fourfold,
\[
 (h^0(T_M),h^1(T_M),h^2(T_M),h^3(T_M),h^4(T_M))
 =(1,1,1,2,0).
\]
The complex parameter family has a nonzero Kodaira--Spencer map
at every member and supplies a smooth one-dimensional miniversal
Kuranishi germ.
\end{theorem}
\begin{proof}
\begin{enumerate}[label=\arabic*.]
\item Put $a=-256/3125$, and denote the whole reduced cusp fibre by
$W$.  The logarithmic sequence established in the proof of
\cref{night:P1A.2}, tensored with $K_M=h^*L^{-6}$ from
\cref{night:P1A.1}, is
\[
0\longrightarrow h^*L^5
\longrightarrow\Omega^1_M(\log W)\otimes K_M
\longrightarrow h^*(F\otimes L^{-6})\longrightarrow0.
\]
Here $K_{\widehat B}(a)=L^{11}$.  More precisely, in the meromorphic
$L$-frames of \cref{night:P1.15},
$\kappa=d\lambda/(\lambda-a)$ is $\lambda^9$ times the
$L^{11}$-frame: both sides have divisor $3D_0-D_\infty$.
The logarithmic sequence is exact in each chart
$q=z_0\cdots z_r$, where its first arrow sends $dq/q$ to
$\sum_{j=0}^r dz_j/z_j$; the quotient is the relative translation
cotangent bundle $h^*F$.  These arrows agree on the specified
monomial and deck transitions and descend through the free finite
root-chart actions.

\item For a summand $L^n$, the coarse degree is
$d(n)=\lfloor4n/5\rfloor+\lfloor-3n/4\rfloor$.
Projection formula and
$R^jh_*\mathcal O_M=\bigwedge^jG$ compute the two end terms.
For $h^*L^5$ the weights in relative degrees $j=0,1,2,3$ are
\[
\{5\},\quad\{5,4,3\},\quad\{4,3,2\},\quad\{2\}.
\]
Weights $4,5$ have degree zero and weights $2,3$ have degree $-1$.
Thus $h^q(h^*L^5)=(1,2,1,0,0)$.
The three weights of $F\otimes L^{-6}$ are $-4,-5,-6$.
After tensoring with $\bigwedge^jG$, every weight is between
$-9$ and $-4$.  We have $d(-9)=-2$ and $d(n)=-1$ for
$-8\le n\le-4$.  Weight $-9$ occurs once in each of relative
degrees $2$ and $3$, so
\[
h^q\bigl(h^*(F\otimes L^{-6})\bigr)=(0,0,0,1,1).
\]
The base has cohomological dimension one, so the two-column Leray
spectral sequences have no further differential.  Every connecting
map for the displayed logarithmic sequence has zero target.
Consequently
\[
h^q\bigl(\Omega^1_M(\log W)\otimes K_M\bigr)=(1,2,1,1,1).
\]

\item Let $\nu:Y\to W$ be the branch normalization.  The residue
sequence, with its actual branchwise maps, is
\[
0\longrightarrow\Omega^1_M\otimes K_M
\longrightarrow\Omega^1_M(\log W)\otimes K_M
\longrightarrow\nu_*\mathcal O_Y\otimes K_M|_W
\longrightarrow0.
\]
Since $K_M$ is pulled back from the base, its restriction to $W$ is
a constant line.  The complete smooth toric variety $Y$ has
$H^0(\mathcal O_Y)=\mathbb C$ and $H^{>0}(\mathcal O_Y)=0$
\cite[Theorem~9.2.3, p.~410, with divisor zero]{CoxLittleSchenck2011}.
The sole global section of the logarithmic middle term comes from
$h^*L^5$.  Its meromorphic expression is
$\lambda^{-5}\kappa$ tensored with the $L^{-6}$-frame: the
unique $L^5$ section is $\lambda^4$ times its frame, and
$\kappa$ contributes $\lambda^9$.  Its residue at $a$ is
$a^{-5}\ne0$.  Hence the map on global sections is an isomorphism
onto the residue line, and
\[
h^q(\Omega^1_M\otimes K_M)=(0,2,1,1,1).
\]
The pointwise perfect evaluation pairing
$T_M\otimes(\Omega^1_M\otimes K_M)\to K_M$ induces analytic
Serre duality
\cite[\S13, Theorem~4, p.~22; \S15, p.~23]{Serre1955Duality}.
It identifies $H^q(T_M)^\vee$ with
$H^{4-q}(\Omega^1_M\otimes K_M)$ and gives the stated tangent
cohomology dimensions.

\item We check the class of the actual parameter family on the
smooth fibre $A=h^{-1}(1)$.  In the fixed frame its period plane
has basis
\[
w_2,\quad w_1,\quad w_0+\frac{625\pi i}{96}\xi e.
\]
Differentiation in $\xi$, followed by projection to $G=H/F$,
gives the rank-one map $P:F\to G$ with sole nonzero matrix entry
$P_{13}=625\pi i/96$ in the ordered bases above and
$(e,w_{-1},w_{-2})$ of the quotient.
Differentiation in the base coordinate gives, by projecting the
connection in \cref{night:P1.15}, the matrix
\[
Q=\frac1{3125(1-a)}
\begin{pmatrix}
-78125/6+300r&-390625/48-375r&0\\
120&0&125/2\\
0&120&50
\end{pmatrix},
\qquad r=\frac{625\pi}{96}(-i\xi).
\]
Its determinant is $3/[625(1-a)^3]\ne0$.
The translation-valued \v Cech cocycles obtained by differentiating
the lattice gluing identify
$H^1(T_A)=\operatorname{Hom}(F,G)$; contraction with a relative
one-form sends the two cocycles to the displayed projected period
derivatives, in the same convention.

\item The normal tangent sequence on this fibre is
\[
0\longrightarrow T_A\longrightarrow T_M|_A
\longrightarrow\mathcal O_A\otimes T_{\widehat B,1}
\longrightarrow0.
\]
Its connecting map sends a base tangent vector to its
Kodaira--Spencer class $Q$.  As $H^0(\mathcal O_A)=\mathbb C$,
the kernel of $H^1(T_A)\to H^1(T_M|_A)$ is exactly
$\mathbb C Q$.  The nonzero matrix $P$ has determinant zero;
every nonzero scalar multiple of $Q$ has nonzero determinant.
Therefore the image of $P$ in $H^1(T_M|_A)$ is nonzero.
The global parameter Kodaira--Spencer cocycle restricts to this
image, so its class in $H^1(T_M)$ is nonzero.
The finite affine actions are fixed, and the cusp multipliers
$\exp(2\pi iC(q,\xi)\nu)$ vary holomorphically in $\xi$.
The same finite fan charts, with $C$ bounded on a small parameter
neighbourhood, give the proper local family of
\cref{night:P1.9,night:P1.10}.  Thus the cocycle just computed is
the derivative of this compact family.

\item Douady's versality and semiuniversality theorems
\cite[\S VII.8, pp.~598--599; \S VIII.1, pp.~599--600]{Douady1974}
give a miniversal germ with tangent space $H^1(T_M)=\mathbb C$.
Its local ring has a minimal presentation $\mathbb C\{u\}/I$,
where $I\subset(u)^2$.  The actual parameter disc classifies by
$u\mapsto g(s)$ with $g'(0)\ne0$, by the preceding step.
Substitution $\mathbb C\{u\}\to\mathbb C\{s\}$ is therefore an
isomorphism of convergent power-series rings.  Since this map factors
through $\mathbb C\{u\}/I$, it annihilates $I$; its injectivity
gives $I=0$.  The miniversal germ is reduced and smooth of dimension
one, and the classifying map is a local isomorphism.  This proves
miniversality of the given family at each member.
\end{enumerate}
\end{proof}




\begin{theorem}\label{night:P1A.4}
\StatementRecord{night:P1A.4}{proved}{Historical events: GPT-6 Astra (OpenAI), ChatGPT Pro mode, effort Pro, S6\_06\_BLIND, 16--18 November 2026: targeted proof check. Also checked in this writing. Pre-continuation writing events: GPT-6 Astra (OpenAI), ChatGPT Pro mode, effort Pro, S6\_10\_PAPERS, 18 November 2026 02:02:02.394029+00:00: computation or formal certificate; targeted proof check. Continuation mathematical and source checks: GPT-6 Astra (OpenAI), ChatGPT Pro mode, effort Pro, S6\_10\_PAPERS continuation, 19 November 2026: computation or formal certificate; independent re-proof; targeted proof check; used and rechecked in this run. Continuation metadata events: GPT-6 Astra (OpenAI), ChatGPT Pro mode, effort Pro, S6\_10\_PAPERS continuation, 19 November 2026: used and rechecked in this run; status-vocabulary reconciliation; metadata reconciliation only; no new mathematical check. \emph{Continuation recorded proof scope.} Complete alternative automorphism and parameter proof written before reading the target proof. The current divisor restriction, special translation stalks and integral quotient are explicit; \path{code/continuation_root/P1A4_integral_fresh.py} checks the matrices. A separate post-comparison program constructs inverse maps and an integral homotopy to the prior cochain presentation. The printed half-period representative has its own fresh reviewer receipt; it is not inferred from the alternative abstract splitting.}{\cref{night:P1.1}, \cref{night:P1.9}, \cref{night:P1.15}, \cref{night:P1A.1}}
The fourfold has algebraic dimension one and
\[
\Aut^0(M)=\mathbb C,\qquad
\Aut(M)=\mathbb C\times\mathbb Z/2.
\]
In a fixed holomorphic splitting, two members of the full
admissible complex parameter half-plane are biholomorphic
exactly when their parameters are equal.
\end{theorem}
\begin{proof}
\begin{enumerate}[label=\arabic*.]
\item An effective divisor dominating the base would restrict to
a nonempty effective divisor on a general complex torus.
Pairing its Chern class with the square of a Kähler form
on that torus is strictly positive. The restriction class is an alternating integral form $J$ on $\Lambda$
fixed by both monodromies:
\[
M_1^tJM_1=J,\qquad M_2^tJM_2=J.
\]
The exact linear system in its fifteen alternating entries has rank
fifteen, so $J=0$.  This contradicts the positive pairing, and every
divisor is vertical.
The pole divisor of a meromorphic function is vertical;
its restriction to a general compact fibre is holomorphic
and therefore constant. The compact reduced graph of the function maps properly to
$\mathbb P^1\times\mathbb P^1$ by $h$ on its first factor.
Remmert's proper-image theorem, in
\cite[\S7, Satz~1, p.~60]{Grauert1960D}, makes its image analytic.
It is generically a single graph and therefore defines the descended
meromorphic function on the normal base curve, including the three
special values.
Conversely pullbacks from the base occur. Hence the
meromorphic function field is $\mathbb C(t)$ and $a(M)=1$.
Every automorphism consequently induces a Möbius map on
the base. The two distinct multiplicities and the cusp
identify all three values separately, so this map is the
identity.

\item On smooth fibres an automorphism is affine. Its integral
linear part is locally constant and commutes with $M_1,M_2$.
Solving the two exact matrix equations in
\path{code/continuation_root/P1A4_integral_fresh.py} gives the rational
centralizer $\mathbb Q I$; its integral invertible elements
are $\pm I$. At the order-five filling an affine lift with
linear part $\epsilon I$ would have to conjugate the same
finite affine generator. For its translation $b$ this
requires
\[
(I-M_1)b=(1-\epsilon)v_1/5\pmod\Lambda.
\]
Applying the invariant integral covector $\ell$ makes the
left side zero. For $\epsilon=-1$ the right side is
$2/5\pmod\mathbb Z$, which is nonzero. Thus every
automorphism is a translation.

\item Let $\mathcal A$ be the sheaf of translations extending to
the specified cusp model and finite root charts. There is
an exact analytic sheaf sequence
\[
0\to j_*\Lambda\to E_0\to\mathcal A\to0
\]
on the root orbifold. Here $j$ omits the cusp and the stalk
of $j_*\Lambda$ there is $\ker N$.
On a finite root cover the torus evaluation sequence is exact and
equivariant for the prescribed affine action; its translation part
cancels under conjugation.  We descend this sequence on the orbifold,
retaining its finite stabilizer cohomology.

To check the last surjectivity at the cusp, lift a translation
to the toroidal cover. It preserves the fan precisely when
its valuation is an integral translation of the periodic
fan. Subtract the corresponding period; the remaining
characters are nonvanishing holomorphic functions on a
disc and possess logarithms. These give a local section
of $E_0$. The kernel consists exactly of the invariant
integral periods. This also proves that the sequence
accounts for all extendable translations, including
possible component permutations.

\item Write $A=M_1$, $B=M_2$ and
$N_1=I+A+\cdots+A^4$, $N_2=I+B+B^2+B^3$.
The orbifold fundamental group away from the cusp has generators
$g_1,g_2$ with relations $g_1^5=g_2^4=1$.  A lattice cocycle is
therefore a pair $(x_1,x_2)$ with $N_1x_1=N_2x_2=0$.
The degree-one exact sequence for $j_*$ says that it extends across
the cusp precisely when its value on $g_1g_2$ is a coboundary:
\[
x_1+Ax_2=(AB-I)u,\qquad u\in\Lambda.
\]
Subtract the coboundary of $u$.  The resulting pair is $(-Ay,y)$,
where $N_1y=N_2y=0$.  The ambiguity in $u$ is
$t\in\ker(AB-I)$, which changes $y$ by $(B-I)t$.
Thus the exact cohomology group, retaining the finite-root torsion, is
\[
H^1(\widehat B,j_*\Lambda)=
\frac{\ker N_1\cap\ker N_2}{(B-I)\ker(AB-I)}.
\]
Primitive integral column bases for its numerator and for
$\ker(AB-I)$ are respectively
\[
K=\begin{pmatrix}
0&0&0\\0&1&0\\1&0&0\\0&0&0\\0&-1&0\\0&0&1
\end{pmatrix},\qquad
T=\begin{pmatrix}
0&0&0\\0&1&0\\1&0&0\\0&0&-1\\0&0&1\\0&0&-1
\end{pmatrix}.
\]
Direct multiplication gives
\[
(B-I)T=KZ,\qquad
Z=\begin{pmatrix}-2&-1&1\\0&-1&1\\1&0&1\end{pmatrix}.
\]
The determinant of $Z$ has absolute value two; an entry and a
$2$-by-$2$ minor have absolute value one.  Its Smith factors are
$1,1,2$, proving $H^1(\widehat B,j_*\Lambda)=\mathbb Z/2$.
All displayed bases, the integral elimination and the coordinate
identity are checked independently in
\path{code/continuation_root/P1A4_integral_fresh.py}.

\item Now $H^0(j_*\Lambda)=0$ by P1.1,
$H^0(E_0)=\mathbb C$, and $H^1(E_0)=0$ by the floor
calculation. The exact sequence gives
\[
0\to\mathbb C\to H^0(\mathcal A)\to\mathbb Z/2\to0.
\]
The half-period
$\tau=(0,1,1,1,1,1)^t/2$ is fixed modulo $\Lambda$
by every monodromy matrix. It extends over the finite
quotients because the translations commute there.
At the cusp $N\tau$ is integral; primitiveness of $N$
provides $\lambda\in\Lambda$ with $N\lambda=N\tau$.
The representative $\tau-\lambda\in\ker N$ is a bounded
angular half-period and extends on every toric chart.
It is a nonzero order-two global section. Since the
additive group $\mathbb C$ has no nonzero torsion, it
splits the displayed exact sequence. Translations
commute, giving the asserted full group.

\item For an isomorphism between two members, the same intrinsic
base argument applies. Its integral linear part still
lies in the common centralizer. Either scalar $\pm I$
preserves a complex period plane exactly by equality
of the two planes. The graphs
$(\sigma+\xi\alpha e)(P)$ are equal precisely when their
parameters $\xi$ agree, since $e\notin F$ and $\alpha$
is nonzero. Translation does not change this condition.
This proves the parameter assertion.

\end{enumerate}
\end{proof}




\begin{theorem}\label{night:P1.17}
\StatementRecord{night:P1.17}{proved}{Historical events: GPT-6 Astra (OpenAI), ChatGPT Pro mode, effort Pro, S6\_06\_BLIND, 16--18 November 2026: targeted proof check. Also checked in this writing. Pre-continuation writing events: GPT-6 Astra (OpenAI), ChatGPT Pro mode, effort Pro, S6\_10\_PAPERS, 18 November 2026 02:02:02.394029+00:00: targeted proof check. Continuation mathematical and source checks: GPT-6 Astra (OpenAI), ChatGPT Pro mode, effort Pro, S6\_10\_PAPERS continuation, 19 November 2026: computation or formal certificate; targeted proof check; used and rechecked in this run. Continuation metadata events: GPT-6 Astra (OpenAI), ChatGPT Pro mode, effort Pro, S6\_10\_PAPERS continuation, 19 November 2026: used and rechecked in this run; status-vocabulary reconciliation; metadata reconciliation only; no new mathematical check. }{\cref{night:P1.1}}
For the cusp marking $S$ fixed in P1.1 and the corresponding Kuhn fan,
none of the completed fourfolds in this family admits an
anti-holomorphic automorphism.  In particular none admits a real
structure.  More strongly, two members with this same cusp marking
are never anti-biholomorphic.
\end{theorem}
\begin{proof}
\begin{enumerate}[label=\arabic*.]
\item We first recover the base intrinsically.  An effective divisor on
$M_c$ cannot be horizontal: its first Chern class vanishes by
$H^2(M_c;\mathbb Z)=0$, whereas its nonzero effective restriction to a
general torus fibre would have positive pairing with the square of a
flat K\"ahler class.  The pole divisor of any meromorphic function is
therefore vertical.  On every regular fibre outside that divisor the
function is holomorphic and hence constant.  It factors through a
meromorphic function on the base, which is rational on
$\mathbb P^1$.  Thus the function field is exactly $\mathbb C(\lambda)$.
An anti-biholomorphism consequently induces an anti-M\"obius map of
the base.  The two distinct multiplicities $5,4$ and the singular cusp
support distinguish the three special values, so this base map fixes
$\infty,0,\lambda_*$.  As those points are real, the map is precisely
$\lambda\mapsto\bar\lambda$.

\item A holomorphic map between compact complex tori is affine: after
lifting to their vector spaces, its derivative is periodic and
therefore constant, since it descends to holomorphic functions on a
compact torus.  Applying this to an anti-holomorphic map after
conjugating its source gives the same assertion for its anti-linear
lift.  Its induced integral map on $H_1$ is constant along the
regular base.  Complex conjugation reverses the two finite meridians,
so this matrix must simultaneously conjugate $M_i$ to $M_i^{-1}$.
Solving these exact integral linear equations gives only $R$ and
$-R$ among unimodular matrices, where
\[
R=\begin{pmatrix}
1&0&0&0&0&0\\
2&0&0&0&-1&0\\
-1&0&-1&0&0&0\\
-1&0&0&-1&0&0\\
2&-1&0&0&0&0\\
-1&1&1&1&1&1
\end{pmatrix},\qquad R^2=I.
\]
Indeed the simultaneous intertwiner space over $\mathbb Q$ has
dimension one and this primitive matrix is unimodular; determinant
one in absolute value permits only the scalar factors $\pm1$.
The equations and this certificate are in
\texttt{centralizer\_torsion.py} and
\texttt{real\_structure\_check.py}.

\item The finite affine cap excludes $G=R$.  Normalize the pullback by
$t=u^{m_i}$; any base-preserving isomorphism lifts to this canonical
root cover and conjugates its deck generator to the inverse.
At its central torus the translation part $t_i$ would satisfy
\[
(I-M_i^{-1})t_i=-(v_i+Gv_i)/m_i\pmod\Lambda.
\]
Applying $\ell$, and using $\ell R=\ell$, gives
$0=-2\ell(v_i)/m_i$ in $\mathbb R/\mathbb Z$, impossible already
for $m_i=5$.  Hence $G=-R$ is necessary.

\item We now use the cusp marking, retaining the conjugation of based
meridians.  Conjugation sends
$g_0=(g_1g_2)^{-1}$ to $g_2g_1=g_2g_0^{-1}g_2^{-1}$.
Thus the matrix in a fixed local cusp frame is
$H=M_2^{-1}G$, and it anti-conjugates $M_0$.
Exact multiplication in the marking $S$ gives
\[
S^{-1}HS=
\begin{pmatrix}
1&0&0&-1&-1&0\\
0&1&0&-2&0&0\\
1&-1&-1&1&-1&-1\\
0&0&0&-1&0&0\\
0&0&0&0&-1&0\\
0&0&0&-1&1&1
\end{pmatrix}
=\begin{pmatrix}P&Q\\0&-P\end{pmatrix},\quad
P=\begin{pmatrix}1&0&0\\0&1&0\\1&-1&-1\end{pmatrix}.
\tag{1.17.1}
\]
Changing the based cusp path by a cusp loop multiplies by a power of
$M_0$ and does not change $P$.

\item We explain why an arbitrary translation cannot remove the obstruction
in (1.17.1).  In the additive torus coordinates whose first periods
are the $a_j$, the anti-linear map must have the form
$\zeta' =P\bar\zeta+t(q)$, because its values on these three real
periods already determine its linear part.  In multiplicative
coordinates it is
$z'=u(q)\bar z^{-P}$.
The source zero section limits to a smooth point of the central fibre.
An isomorphism takes it to another smooth point.  The central fibre
is irreducible, and its smooth locus is exactly the dense torus of
its toric normalization: every toric boundary divisor is a
normal-crossings self-identification.  In a chart about this image,
the translation therefore has the form $q^v$ times nonvanishing
holomorphic units, with $v\in\mathbb Z^3$.  Translating by the period
$v$ changes the selected component chart; the units act by torus
translation.  Neither changes the linear action $-P$ on the fan.

\item The primitive rays of the normal fan of that component are
\[
\{\pm v:0\ne v\in\{0,1\}^3\}.
\]
An extending torus monomial isomorphism must permute these rays.
This can be seen directly on a one-parameter arc $z=t^v$ meeting the
generic point of its boundary divisor: if the image vector lies in
the interior of a cone of dimension at least two, that divisor is
contracted, which is impossible for an isomorphism.  Complex
conjugation itself preserves the toric charts, so the same criterion
applies to the anti-monomial map above.
But $-Pe_2=(0,-1,1)^t$ is not one of the displayed rays.
This is a contradiction.  We have exhausted both possible integral
linear parts, so no anti-biholomorphism exists.
The assertion is explicitly about the chosen $S$ and its fan; it does
not identify different toroidal compactifications without a map
between them.

\end{enumerate}
\end{proof}


\section{The boundary of the family}

Fix one splitting and write the complex parameter as $\xi$, so that
$F_\xi=(\sigma+\xi\alpha e)(P)$. Put
$c_*=-\min h$ and $\mathcal H=\{\operatorname{Im}\xi>c_*\}$.
The real family in the prompt is $\xi=ic$.
By an analytic extension across a boundary point we mean a proper
Hausdorff complex-analytic family on a full complex disc whose
restriction to the nonempty open part lying in $\mathcal H$ is
the prescribed family. At infinity we use the coordinate $u=1/\xi$.
This definition permits singular central spaces, multiplicities,
and additional central components.

The stability of complex tori under deformation in the large is
established in \cite[definition, p.~261; Theorem 2.1, p.~264]{Catanese2004}.
The relative-form calculation below writes the actual torus quotient
on the regular analytic open set used in the boundary argument.

\begin{theorem}\label{night:P1A.5}
\StatementRecord{night:P1A.5}{proved}{Historical events: GPT-6 Astra (OpenAI), ChatGPT Pro mode, effort Pro, S6\_06\_BLIND, 16--18 November 2026: targeted proof check. Also checked in this writing. Pre-continuation writing events: GPT-6 Astra (OpenAI), ChatGPT Pro mode, effort Pro, S6\_10\_PAPERS, 18 November 2026 02:02:02.394029+00:00: targeted proof check. Continuation mathematical and source checks: GPT-6 Astra (OpenAI), ChatGPT Pro mode, effort Pro, S6\_10\_PAPERS continuation, 19 November 2026: computation or formal certificate; independent re-proof; targeted proof check; used and rechecked in this run. Continuation metadata events: GPT-6 Astra (OpenAI), ChatGPT Pro mode, effort Pro, S6\_10\_PAPERS continuation, 19 November 2026: used and rechecked in this run; status-vocabulary reconciliation; metadata reconciliation only; no new mathematical check. }{\cref{night:P2.4,night:P1A.1,night:P1.8}; \cite[Theorem 1.6, preprint p.~5]{BM97}; \cite[\S6, Hauptsatz I; \S7, Satz 1 and Satz 5]{Grauert60}}
There is no such analytic extension across any finite boundary
point of $\mathcal H$, or across infinity. The same conclusion
holds after any finite ramified change of the local boundary
coordinate. In particular neither end of the real interval
$c_*<c<+\infty$ has a proper Hausdorff analytic extension in
this sense.
\end{theorem}
\begin{proof}
\begin{enumerate}[label=\arabic*.]
\item We first identify the actual moving pencil on each known member.
By \cref{night:P1A.1}, $K_M^{-4}=h^*L^{24}$ and the fibres of $h$
are connected.  For the coarse map
$\rho:\widehat B\to\mathbb P^1$, the root-divisor formula gives
\[
 \rho_*L^{24}
 =\mathcal O_{\mathbb P^1}
   \bigl(\lfloor96/5\rfloor\,\infty+
                \lfloor-72/4\rfloor\,0\bigr)
 =\mathcal O_{\mathbb P^1}(19\infty-18\,0).
\]
Its sections are spanned by $\lambda^{18},\lambda^{19}$.
On a root chart the invariant sections have a common local
monomial factor, and their ratio cancels that factor.  Thus
$h^0(M,K_M^{-4})=2$, and the moving part of this complete
system is exactly the marked base map, up to a projective
coordinate change.  In particular its base ideal is invertible
on every member of the prescribed family.

\item Suppose that a proper Hausdorff analytic extension over a full
parameter disc exists.  Take the reduced component containing
the given nonempty open family and resolve it.  Analytic
resolution applies on a relatively compact neighbourhood of
the compact central fibre and can be kept unchanged over the
given smooth locus
\cite[Theorem 1.6 and its following analytic-neighbourhood
paragraph, preprint p.~5]{BM97}.
Write $p:Z\to\Delta$ for the resulting proper map from a smooth
irreducible space.  The proper images of the exceptional and
critical loci are analytic
\cite[\S7, Satz 1, p.~60]{Grauert60}; they omit the given open
parameter set.  After shrinking the disc they are supported
at its centre.  Hence $p$ is a proper holomorphic submersion
over $\Delta^*$, with compact fourfold fibres.

Put $\mathcal J=\omega_{Z/\Delta}^{-4}$.
The sheaves $R^jp_*\mathcal J$ are coherent
\cite[\S6, Hauptsatz I, p.~59]{Grauert60}.
Multiplication by a nonzero parameter germ is injective on
$\mathcal J$ and on $p_*\mathcal J$, so $p_*\mathcal J$ is locally free
over the disc and free after shrinking.
The actual exact sequence for a fibre over $t$ gives
\[
 0\longrightarrow(p_*\mathcal J)_t/\mathfrak m_t(p_*\mathcal J)_t
 \longrightarrow H^0(Z_t,\mathcal J_t)
 \longrightarrow
 \ker\bigl((t'-t): (R^1p_*\mathcal J)_t
                       \longrightarrow(R^1p_*\mathcal J)_t\bigr)
 \longrightarrow0.
\]
The last term is supported in the torsion support of the
coherent sheaf $R^1p_*\mathcal J$, a closed discrete analytic
set.  At a point of the given open family outside this support,
the displayed comparison shows that $p_*\mathcal J$ has rank
two.  Shrinking makes the support empty away from the centre.
Thus a frame of $p_*\mathcal J$ gives the complete two-dimensional
linear system on each punctured fibre.  This uses the displayed
coherent maps, without assuming closedness of a general
base-change failure set.

\item The two evaluation sections define a coherent base ideal and
a meromorphic map $Z\dashrightarrow\mathbb P^1$ over the full
disc.  The ideal is locally invertible on the given open family
by the first step.  Its non-invertibility locus is analytic;
its proper image in $\Delta$ is therefore discrete and can be
removed away from the centre.  Blowing up that ideal and
resolving the graph consequently changes only the central
fibre.  It gives a proper holomorphic map to
$\mathbb P^1\times\Delta$.  The construction uses the graph
of the coherent evaluation map, so its coefficients have
finite meromorphic order in the parameter.

Its critical image is proper analytic.  Every component
dominating the parameter disc meets the given open family,
where there are precisely three critical values.  Their
degrees over the parameter are one.  The multiplicities five
and four distinguish two of them, and the third is the cusp.
After shrinking, their normalizations give three distinct
labelled meromorphic sections, and all vertical exceptional
components lie over the centre.  The projective transformation
taking these sections to $\infty,0,\lambda_*$ has meromorphic
coefficients and is invertible on $\Delta^*$.  Resolving its
graph gives the actual marked map $H$ over the disc, still
unchanged on punctured fibres.  Its regular restriction is a
proper holomorphic submersion
\[
 q:W\longrightarrow
 B=\bigl(\mathbb P^1\setminus\{0,\infty,\lambda_*\}\bigr)
          \times\Delta^*.
\]
Ehresmann transport makes the number of components of a fibre
locally constant.  The base is connected and the known fibres
are connected, so every fibre of $q$ is connected.

\item We need the torus structure on an analytically open dense
part of this regular family.  The coherent sheaf
$\mathcal E=q_*\Omega^1_{W/B}$ has generic rank three.
Generic cohomology and base change for this flat bundle gives
a proper analytic exceptional set outside which $\mathcal E$
is locally free of rank three and its fibre is the space of
holomorphic one-forms.  The constant-dimension comparison is
\cite[\S7, Satz 5, pp.~63--64]{Grauert60}.
Use the global coherent evaluation map
\[
 q^*\mathcal E\longrightarrow\Omega^1_{W/B}.
\]
Its cokernel has analytic support.  The proper image of that
support is analytic in $B$ and omits the nonempty part of the
known torus open set outside the generic exceptional set.
Enlarge the exceptional set by this image.  On its complement
evaluation is an isomorphism of rank-three bundles.

Relative exterior differentiation defines an
$\mathcal O_B$-linear map
$\mathcal E\to q_*\Omega^2_{W/B}$.  It vanishes on the known
torus open set.  The target is torsion-free, since it is the
proper direct image of a vector bundle under a submersion.
The identity theorem therefore makes this map zero.  Each
remaining compact connected fibre has three everywhere
independent closed holomorphic one-forms.  Their dual vector
fields commute and are complete by compactness.  Their flows
give a transitive holomorphic $\mathbb C^3$-action: every orbit
is open and the fibre is connected.  Its stabilizer $\Gamma$
is discrete, and compactness of $\mathbb C^3/\Gamma$ makes
$\Gamma$ a full real lattice.  This is the actual complex-torus
structure on each such fibre.

\item Let $E\subset B$ be the proper analytic exceptional set just
removed.  On the universal cover of $B\setminus E$, use
Gauss--Manin transport from the given open family to mark
integral $H^1$.  Integration of the closed relative forms
gives a holomorphic map to $\operatorname{Gr}(3,6)$.  On the
known open set it is the explicit graph
$F_\xi=(\sigma+\xi\alpha e)(P)$.  The identity theorem on the
connected cover identifies these two maps everywhere.
The explicit graph is a rank-three holomorphic plane for
every finite $\xi$, including parameters for which it is not
real-transverse.  At infinity the same assertion is used with
$\xi=1/u$ on the punctured parameter disc.  All changes of
continuation marking are real integral and preserve the
condition $F\cap\bar F=0$.

\item Let $\xi_0$ be a finite boundary point, so
$\operatorname{Im}\xi_0=c_*=-\min h$.
The minimum may occur at a finite root point.  Continuity
there nevertheless gives regular $\lambda$ with $h(\lambda)$
arbitrarily close to $\min h$.  Every sufficiently small full
parameter disc about $\xi_0$ therefore contains noncentral
points of
\[
 \operatorname{Im}\xi+h(\lambda)=0.
\]
This equation defines nonempty smooth real hypersurface
pieces in the complex surface $B$, since its derivative in
$\operatorname{Im}\xi$ is one.  A proper complex analytic
subset of $B$ has real dimension at most two; it cannot
contain those three-dimensional pieces.  Choose such a
point outside $E$.  Its fibre is a complex torus by the fourth
step, so its period plane is real-transverse.  The fifth step
identifies that plane with $F_\xi$, whereas the exact criterion
of \cref{night:P1.8} gives
$F_\xi\cap\overline{F_\xi}\ne0$ at the chosen point.
This is a contradiction.

At infinity fix a regular $\lambda$ and set
\[
 \xi=R-i h(\lambda),\qquad
 u=\frac1{R-i h(\lambda)},\qquad R\in\mathbb R.
\]
These are noncentral points of the same failure locus and
tend to $u=0$ as $|R|\to\infty$.  Varying $\lambda$ and $R$
again gives real hypersurface pieces in every punctured
$u$-disc, so one can choose a point outside $E$ and obtain the
same contradiction.  Under a finite ramified change of the
boundary coordinate, the parameter map has nonzero derivative
off the centre after shrinking.  The inverse images of these
hypersurface pieces remain nonempty real hypersurfaces there,
and the argument is unchanged.

All contradictions occur on a noncentral compact smooth
torus fibre of the recovered marked map.  The reductions and
modifications are proper and preserve the prescribed open
family; they impose no smoothness, multiplicity or component
condition on the original central fibre.  The permitted
central spaces therefore do not avoid the contradiction.
Both nonextension assertions follow.
\end{enumerate}
\end{proof}




\begin{theorem}\label{night:P1A.6}
\StatementRecord{night:P1A.6}{proved}{Historical events: GPT-6 Astra (OpenAI), ChatGPT Pro mode, effort Pro, S6\_06\_BLIND, 16--18 November 2026: targeted proof check. Also checked in this writing. Pre-continuation writing events: GPT-6 Astra (OpenAI), ChatGPT Pro mode, effort Pro, S6\_10\_PAPERS, 18 November 2026 02:02:02.394029+00:00: targeted proof check. Continuation mathematical and source checks: GPT-6 Astra (OpenAI), ChatGPT Pro mode, effort Pro, S6\_10\_PAPERS continuation, 19 November 2026: independent re-proof; used and rechecked in this run. Continuation metadata events: GPT-6 Astra (OpenAI), ChatGPT Pro mode, effort Pro, S6\_10\_PAPERS continuation, 19 November 2026: used and rechecked in this run; status-vocabulary reconciliation; metadata reconciliation only; no new mathematical check. }{\cref{night:P1A.3}, \cref{night:P1A.4}}
A holomorphic punctured-disc family whose fibres are all
members of this fourfold family has an interior parameter
limit. After a base change of degree at most two, it is
isomorphic to the pullback of the constructed family by
a holomorphic map of the full disc to $\mathcal H$, and
therefore has a proper smooth filling. Such a family
cannot approach either parameter end.
\end{theorem}
\begin{proof}
\begin{enumerate}[label=\arabic*.]
\item At every punctured parameter, P1A.3 supplies local
miniversality and a local holomorphic classifying map to
$\mathcal H$. P1A.4 makes this map unique on each geometric
point. Thus these local maps agree on overlaps and define
one holomorphic function $\xi:\Delta^*\to\mathcal H$.
There is no unmentioned quotient by parameter translations:
P1A.4 proved their absence by the integral centralizer.

\item The Cayley coordinate
\[
 z=\frac{\xi-i(c_*+1)}{\xi-i(c_*-1)}
\]
has $|z|<1$. It extends holomorphically across the puncture
by the removable-singularity theorem. Its value there has
modulus strictly less than one, since a boundary value
would force $z$ to be constant by the maximum-modulus
principle. The inverse Cayley map therefore extends
$\xi$ with value in $\mathcal H$.

\item Local identifications with the pulled-back constructed
family differ by its automorphisms. These form the constant
holomorphic group $\mathbb C\times\mathbb Z/2$ over the
parameter half-plane: the translation field and the
half-period in P1A.4 extend over the entire family.
The finite component torsor on $\Delta^*$ is killed by
the double cover. The remaining additive holomorphic
torsor is zero because $H^1(\Delta^*,\mathcal O)=0$;
the punctured disc is Stein and its additive Cousin
problem has a solution. Thus after that cover the local
identifications glue to a global isomorphism. Pulling
back the already constructed proper smooth family along
the extended classifying map supplies the stated filling.

\end{enumerate}
\end{proof}


\section{The effective integral input classification}



\begin{theorem}[the orbifold quotient]\label{maps:src:F.2.7}
\StatementRecord{maps:src:F.2.7}{proved}{Historical events: Claude Fable 5.2 (Anthropic), effort high, 16 November 2026: carried as background without fault found. Claude Fable 5.2 (Anthropic), effort high, 17 November 2026: carried as background without fault found. GPT-6 Astra (OpenAI), ChatGPT Pro mode, effort Pro, S6\_09\_PAPERS, 17 November 2026: targeted proof check. Also checked in this writing. Pre-continuation writing events: GPT-6 Astra (OpenAI), ChatGPT Pro mode, effort Pro, S6\_10\_PAPERS, 18 November 2026 02:02:02.394029+00:00: targeted proof check. Continuation mathematical and source checks: GPT-6 Astra (OpenAI), ChatGPT Pro mode, effort Pro, S6\_10\_PAPERS continuation, 19 November 2026: independent re-proof; used and rechecked in this run. Continuation metadata events: GPT-6 Astra (OpenAI), ChatGPT Pro mode, effort Pro, S6\_10\_PAPERS continuation, 19 November 2026: used and rechecked in this run; status-vocabulary reconciliation; metadata reconciliation only; no new mathematical check. \emph{Pre-continuation recorded proof scope.} Claude Fable 5.2 (Anthropic), effort high, 2026-11-17; Claude Fable 5.2 (Anthropic), effort high, 16 November 2026; also checked in this writing by GPT-6 Astra (OpenAI), ChatGPT Pro mode, effort Pro, 18 November 2026}{the hypotheses, notation and definitions stated here; the explicit constructions and calculations in the proof}


For the specified torus fillings in F.2, with affine
multiplicities $m_1,\ldots,m_k$, the fundamental group of the
total space has a canonical surjection onto
\[
 \langle x_1,\ldots,x_k\mid
       x_i^{m_i}=1,\ x_1\cdots x_k=1\rangle.
 \tag{F.27}\label{maps:eq:F.27}
\]
Consequently a simply connected example with $k=2$ requires
$\gcd(m_1,m_2)=1$.  There is no simply connected example with
$k=3$ and $m_1,m_2,m_3\geq2$.  An affine multiplicity-one
point is an ordinary smooth point and reduces to the problem
with one fewer multiple fibre.

\end{theorem}


\begin{proof}
\begin{enumerate}[label=\arabic*.]
\item Over the complement of small discs at the special values,
the torus bundle projects surjectively on fundamental groups
to the punctured base.  Quotient its fundamental group by the
normal closure of all regular-fibre loops.  At a finite
filling its root cover is a smooth torus bundle over a disc,
and its free deck group is cyclic of order $m_i$.
The covering exact sequence shows that the quotient of the
filling group by its fibre subgroup is precisely this cyclic
group.  The boundary base meridian maps to its generator.
Thus adjoining that filling imposes exactly
$x_i^{m_i}=1$ after the fibre loops have been killed.
At the index-one Kuhn cusp the fundamental group is generated
by the three deck loops of its dual torus; these are regular
fibre loops.  Its base meridian contracts.  Hence its filling
imposes the ordinary meridian relation and leaves no extra
generator in this quotient.  The base relation gives the
product relation in \eqref{maps:eq:F.27}.  Van Kampen proves that \eqref{maps:eq:F.27} is
the resulting quotient, in particular the asserted surjection.

\item For two points, substitute $x_2=x_1^{-1}$; the quotient is
$\mathbb Z/\gcd(m_1,m_2)$.
For three points of orders at least two it is a nontrivial
triangle group.  One can see nontriviality without assuming
a classification of the total space.  Take a geodesic triangle
with angles $\pi/m_1,\pi/m_2,\pi/m_3$ in the sphere, Euclidean
plane or hyperbolic plane according as their sum is greater
than, equal to or less than $\pi$.
In the spherical case the possible orders are, after sorting,
$(2,2,n),(2,3,3),(2,3,4),(2,3,5)$, and the corresponding
spherical triangles exist.  The Euclidean and hyperbolic
angle conditions give the other triangles directly.
If $r_1,r_2,r_3$ are reflections in the opposite sides,
the rotations $r_2r_3,r_3r_1,r_1r_2$ have the stated
orders and product one.  They therefore give a nontrivial
representation of \eqref{maps:eq:F.27} into the orientation-preserving
isometry group.  Its quotient cannot be trivial.
Finally a free cyclic action of order one is the identity
construction on a root disc of order one, hence an ordinary
torus bundle filling and not a multiple fibre.

\end{enumerate}
\end{proof}




\begin{proposition}[exclusion of an involutory finite part]\label{maps:src:F.2.8}
\StatementRecord{maps:src:F.2.8}{proved}{Historical events: Claude Fable 5.2 (Anthropic), effort high, 16 November 2026: carried as background without fault found. GPT-6 Astra (OpenAI), ChatGPT Pro mode, effort Pro, S6\_09\_PAPERS, 17 November 2026: targeted proof check. Also checked in this writing. Pre-continuation writing events: GPT-6 Astra (OpenAI), ChatGPT Pro mode, effort Pro, S6\_10\_PAPERS, 18 November 2026 02:02:02.394029+00:00: targeted proof check. Continuation mathematical and source checks: GPT-6 Astra (OpenAI), ChatGPT Pro mode, effort Pro, S6\_10\_PAPERS continuation, 19 November 2026: computation or formal certificate; independent re-proof; used and rechecked in this run. Continuation metadata events: GPT-6 Astra (OpenAI), ChatGPT Pro mode, effort Pro, S6\_10\_PAPERS continuation, 19 November 2026: used and rechecked in this run; status-vocabulary reconciliation; metadata reconciliation only; no new mathematical check. \emph{Pre-continuation recorded proof scope.} Claude Fable 5.2 (Anthropic), effort high, 16 November 2026; also checked in this writing by GPT-6 Astra (OpenAI), ChatGPT Pro mode, effort Pro, 18 November 2026}{the hypotheses, notation and definitions stated here; \cref{maps:src:F.2.7}; \cref{maps:src:F.2.15}; \cref{maps:src:F.3.1}; \cref{maps:eq:F.3}; \cref{maps:src:F.2.9}}


A simply connected datum with $b_2=0$, $k=2$ and the stated
finite affine fillings has neither finite linear
monodromy of order one or two.  Its unordered pair of finite
linear orders must therefore belong to
\[
 \{(3,4),(3,5),(3,8),(3,10),(4,5),
       (5,6),(5,8),(5,12)\}.
 \tag{F.28}\label{maps:eq:F.28}
\]
The affine multiplicities may still be larger than these
linear orders.

\end{proposition}


\begin{proof}
\begin{enumerate}[label=\arabic*.]
\item The linear orders divide the affine multiplicities and are
therefore coprime by \ref{maps:src:F.2.7}.  The fixed-rank calculation
of \ref{maps:src:F.2.9}, which uses only \ref{maps:src:F.3.1} and the
evenness in \ref{maps:src:F.2.15}, gives fixed rank two at each finite
generator.  An identity generator would make the product finite and
unipotent, hence the identity, contrary to the rank-three cusp.

\item Suppose the cohomology generators are $A,B$ with $A^2=I$.
The order of $B$ is odd.  The $(-1)$-eigenspace of $A$ has dimension
four, whereas $\ker(AB-I)$ has dimension three.  Their intersection
in the six-dimensional space contains a nonzero vector $v$.
Then $ABv=v$ and $Av=-v$ imply
\[
 Bv=Av=-v,
\]
which contradicts the odd order of $B$.  Interchanging the two
generators proves the same assertion for either finite site.

\item Delete one and two from the complete finite-order list
\eqref{maps:eq:F.3} and enumerate its coprime unordered pairs.
They are exactly the eight pairs in \eqref{maps:eq:F.28}.
Only the divisibility of the linear orders into the affine
multiplicities has been used.
\end{enumerate}
\end{proof}




\begin{proposition}[the only first-cohomology pattern for two points]\label{maps:src:F.2.9}
\StatementRecord{maps:src:F.2.9}{proved}{Historical events: Claude Fable 5.2 (Anthropic), effort high, 16 November 2026: carried as background without fault found. GPT-6 Astra (OpenAI), ChatGPT Pro mode, effort Pro, S6\_09\_PAPERS, 17 November 2026: targeted proof check. Also checked in this writing. Pre-continuation writing events: GPT-6 Astra (OpenAI), ChatGPT Pro mode, effort Pro, S6\_10\_PAPERS, 18 November 2026 02:02:02.394029+00:00: targeted proof check. Continuation mathematical and source checks: GPT-6 Astra (OpenAI), ChatGPT Pro mode, effort Pro, S6\_10\_PAPERS continuation, 19 November 2026: computation or formal certificate; independent re-proof; used and rechecked in this run. Continuation metadata events: GPT-6 Astra (OpenAI), ChatGPT Pro mode, effort Pro, S6\_10\_PAPERS continuation, 19 November 2026: used and rechecked in this run; status-vocabulary reconciliation; metadata reconciliation only; no new mathematical check. \emph{Pre-continuation recorded proof scope.} Claude Fable 5.2 (Anthropic), effort high, 16 November 2026; also checked in this writing by GPT-6 Astra (OpenAI), ChatGPT Pro mode, effort Pro, 18 November 2026}{the hypotheses, notation and definitions stated here; \cref{maps:src:F.3.1}; \cref{maps:src:F.2.15}}


For a simply connected two-multiple-fibre datum with $b_2=0$,
the rational
cohomology system $V$ has one invariant, no coinvariants,
and fixed rank two at each finite generator.  Its quotient
$W=V/\mathbb Qe$ by the invariant line has neither invariants
nor coinvariants.  If one finite characteristic polynomial
on $W$ is $(T-1)\Phi_d(T)$ with $\deg\Phi_d=4$, then $W$
is irreducible over $\mathbb Q$.

\end{proposition}


\begin{proof}
\begin{enumerate}[label=\arabic*.]
\item Write $a=\dim V^\Gamma$, $b=\dim(V^*)^\Gamma$ and $f_i$
for the finite fixed ranks.  Theorem \ref{maps:src:F.3.1} gives $a=1$,
vanishing first parabolic cohomology, and
$f_1+f_2=4+b$.  By \ref{maps:src:F.2.15} each $f_i$ is positive and even, so $b$ is
even.  The maximal square-zero cusp has
$\ker N=\operatorname{im}N$.  Therefore every globally
fixed vector belongs to $\operatorname{im}N$, whereas every
globally fixed covector kills $\operatorname{im}N$.
The global invariant and coinvariant spaces pair to zero.

\item Restriction of a globally invariant covector to the fixed
space of a finite generator is injective: averaging that
generator is the projection onto its fixed space, and the
covector is unchanged by averaging.  Its restriction kills
the $a$ global invariants in that space.  Hence
$a+b\leq f_i$ for each $i$.
If $b\geq2$ is even, the even ranks satisfy
$f_i\geq b+2$, so $4+b=f_1+f_2\geq2b+4$, a contradiction.
Thus $b=0$ and $f_1=f_2=2$.

\item The quotient $W$ has no invariant covectors, since its dual
embeds into $V^*$.  If it had a fixed line, its preimage in
$V$ would be an extension of the trivial line by the trivial
line.  Every finite generator on this two-dimensional
space would be unipotent of finite order, hence identity.
This would give two invariants in $V$, a contradiction.
Thus $W$ also has no invariant vectors.

\item If one finite characteristic polynomial is
$(T-1)\Phi_d(T)$ with the second factor irreducible of degree
four, every proper nonzero rational submodule has dimension
one or four.  In dimension one that generator acts as the
identity.  Since the product of the two finite generators
is unipotent, their induced one-dimensional characters are
inverse; hence the other generator also acts trivially.
This would be an invariant line.  In dimension four the
one-dimensional quotient is trivial by the same argument,
giving a coinvariant.  Both have just been excluded.

\end{enumerate}
\end{proof}


\paragraph{The classical primitive exterior-square comparison.}
Here is an elementary construction of the classical five-dimensional
orthogonal representation, with the precise unipotent comparison.
Let $E=\mathbb C^4$ have a symplectic basis $e_1,e_2,f_1,f_2$.
The invariant bivector $\omega=e_1\wedge f_1+e_2\wedge f_2$ has
nonzero square. The wedge pairing on $\bigwedge^2E$ is symmetric
and nondegenerate. Its orthogonal complement
$E_0=\omega^\perp$ therefore has dimension five and is nondegenerate.
Preservation of $\omega^2$ gives determinant one on $E$;
$\det(\bigwedge^2g)=(\det g)^3$ and the fixed line
$\mathbb C\omega$ therefore give determinant one on $E_0$.
The action gives
\[
 \rho:\operatorname{Sp}(E)\longrightarrow\operatorname{SO}(E_0).
\]
Its differential is injective. Indeed, a symplectic infinitesimal
transformation acting trivially on $E_0$ also acts trivially on
$\mathbb C\omega$, hence on all of $\bigwedge^2E$. In a basis,
the equations $Xe_i\wedge e_j+e_i\wedge Xe_j=0$ for every $i<j$
first make every off-diagonal coefficient of $X$ zero, and then give
$X_{ii}+X_{jj}=0$ for every distinct pair. Since there are four
indices, these equations make every diagonal coefficient zero.
Both Lie algebras have dimension ten, so the image contains a
neighbourhood of the identity in $\operatorname{SO}(E_0)$.
That complex special orthogonal group is connected. For completeness,
orthogonal elimination expresses an orthogonal matrix as a product
of reflections in vectors of nonzero norm: to fix a nonisotropic
vector, reflection in the difference of its image and itself works
when that difference is nonisotropic; otherwise reflection in their
sum followed by reflection in the vector works. At least one of
sum and difference is nonisotropic. Continue on its nondegenerate
orthogonal complement. Determinant one makes the number of
reflections even. The nonisotropic vectors of a complex vector
space of dimension at least two form a path-connected set, since
two generic complex line segments can avoid the zero set of the
norm polynomial except at finitely many parameter values, around
which the paths can be detoured. Varying one vector in a pair of
reflections to the other connects that pair to the identity.
The one-dimensional final complement contributes only the already
accounted-for sign. This proves connectedness. Thus the open image
of $\rho$ is the whole special orthogonal group.

The kernel is $\{I,-I\}$. If an element acts trivially on
$E_0$, it also fixes $\omega$ and hence every bivector.
It preserves every two-plane, as follows by applying it to a
nonzero decomposable bivector of that plane. Intersections of
two-planes show that it preserves every line. A linear map
preserving every line is scalar: apply it to the coordinate lines
and their pairwise sums. Its scalar has square one. In particular
every finite orthogonal element of order $m$ has lifts whose
$m$th powers belong to $\{I,-I\}$; the lifts are finite and
semisimple.

The symplectic transvection fixing $e_1,e_2,f_2$ and sending
$f_1$ to $f_1+e_1$ acts on the following primitive basis:
\[
 e_1\wedge e_2,\quad e_1\wedge f_2,\quad
 e_2\wedge f_1,\quad f_1\wedge f_2,\quad
 e_1\wedge f_1-e_2\wedge f_2.
\]
Its difference from the identity sends the third vector to
the negative of the first and the fourth to the second, and
annihilates the other three. Its orthogonal Jordan type is
therefore $(2,2,1)$. Conversely, for an orthogonal unipotent
$I+N$ of that type, $N^2=0$ implies $N^*=-N$, because
$(I+N)^*=(I+N)^{-1}=I-N$. Its image is an isotropic plane.
After choosing a basis of that plane and a dual isotropic
plane, use the isometry on their residual nondegenerate line
to make the identification of two such decompositions have
determinant one. Changing its sign changes that determinant
and commutes with $N$, which kills the residual line.
In these decompositions, $N$ is determined by a nonzero alternating form on
the dual two-dimensional space. All such forms are equivalent
under its general linear group. Extending that change by the
inverse transpose on the first plane gives an orthogonal
change of basis; its determinant is one. Thus these unipotents
form one special-orthogonal conjugacy class. Surjectivity of
$\rho$ lifts the conjugacy. The two lifts of $I+N$ are a
symplectic transvection and its negative. 

Finally a symplectic diagonal element with eigenvalues
$x,x^{-1},y,y^{-1}$ acts on $E_0$ with eigenvalues
\[
 1,\quad xy,\quad x/y,\quad y/x,\quad (xy)^{-1}.
\]
This gives the stated eigenvalue comparison through the actual
primitive exterior-square representation.

\begin{theorem}[exclusion of three further order pairs]\label{maps:src:F.2.10}
\StatementRecord{maps:src:F.2.10}{proved}{Historical events: Claude Fable 5.2 (Anthropic), effort high, 16 November 2026: carried as background without fault found. GPT-6 Astra (OpenAI), ChatGPT Pro mode, effort Pro, S6\_09\_PAPERS, 17 November 2026: targeted proof check. Also checked in this writing. Pre-continuation writing events: GPT-6 Astra (OpenAI), ChatGPT Pro mode, effort Pro, S6\_10\_PAPERS, 18 November 2026 02:02:02.394029+00:00: targeted proof check. Continuation mathematical and source checks: GPT-6 Astra (OpenAI), ChatGPT Pro mode, effort Pro, S6\_10\_PAPERS continuation, 19 November 2026: computation or formal certificate; independent re-proof; targeted proof check; used and rechecked in this run. Continuation metadata events: GPT-6 Astra (OpenAI), ChatGPT Pro mode, effort Pro, S6\_10\_PAPERS continuation, 19 November 2026: used and rechecked in this run; status-vocabulary reconciliation; metadata reconciliation only; no new mathematical check. \emph{Pre-continuation recorded proof scope.} Claude Fable 5.2 (Anthropic), effort high, 16 November 2026; also checked in this writing by GPT-6 Astra (OpenAI), ChatGPT Pro mode, effort Pro, 18 November 2026}{the hypotheses, notation and definitions stated here; \cref{maps:src:F.2.9}; \cref{maps:eq:F.31}; \cref{maps:src:F.2.8}; \cref{maps:src:F.3.1}}


No simply connected datum with $k=2$, $b_2=0$, and finite linear orders
$(3,5)$, $(3,8)$ or $(3,10)$ has the required maximal
square-zero cusp.  Among such data, the
remaining necessary unordered linear order pairs are
\[
 \{(3,4),(4,5),(5,6),(5,8),(5,12)\}.
 \tag{F.30}\label{maps:eq:F.30}
\]
When one order is five, the other finite characteristic
polynomial cannot be $(T-1)^2\Phi_4(T)^2$ or
$(T-1)^2\Phi_6(T)^2$.

\end{theorem}


\begin{proof}
\begin{enumerate}[label=\arabic*.]
\item For each proposed pair, \ref{maps:src:F.2.9} gives an
irreducible rational rank-five quotient $W$.  Its cusp has Jordan
type $(2,2,1)$.  One finite polynomial is
$(T-1)\Phi_3(T)^2$ and the other is $(T-1)\Phi_d(T)$ for
$d=5,8,10$.  For the two additional assertions, replace
$\Phi_3^2$ by $\Phi_4^2$ or $\Phi_6^2$, with $d=5$.
In all these cases a commuting endomorphism of $W$ restricts to a
rational scalar on the one-dimensional fixed space of the second
generator.  This restriction is injective: a nonzero endomorphism
of an irreducible representation is invertible and cannot vanish
on that line.  Consequently
\[
 \operatorname{End}_{\Gamma}(W)=\mathbb Q.
\]

\item Consider $\mathcal H=\operatorname{Hom}(W,W^*)$.
The local fixed dimensions are $9,5,13$.  At the repeated finite
site, pairing reciprocal eigenspaces of dimensions $1,2,2$ gives
$1+4+4=9$.  At the distinct finite site it gives five.
For the cusp, the ordered pairs of blocks of sizes $2,2,1$
contribute the sum of their minimum sizes, namely thirteen.
The invariant-sheaf Euler formula therefore gives
\[
 \chi(\mathbb P^1,j_*\mathcal H)=-25+9+5+13=2.
 \tag{F.31}\label{maps:eq:F.31}
\]
Here degree zero is $\operatorname{Hom}_{\Gamma}(W,W^*)$,
and degree two is dual to
$\operatorname{Hom}_{\Gamma}(W^*,W)$.  At least one is nonzero,
so irreducibility gives a nondegenerate invariant bilinear form
$J:W\longrightarrow W^*$, inverting the map if necessary.
Its alternating part is equivariant; if nonzero it would be
invertible, which is impossible in odd dimension.  Thus $J$ is
symmetric.

\item An invariant bivector is an alternating equivariant map
$K:W^*\longrightarrow W$.  The actual comparison map
\[
 K\longmapsto KJ\in\operatorname{End}_{\Gamma}(W)
\]
makes it a scalar, so $K$ is a scalar multiple of the symmetric
map $J^{-1}$.  Hence $K=0$.  Applying the same argument to an
alternating map $W\longrightarrow W^*$ gives
\[
 (\Lambda^2W)^\Gamma=(\Lambda^2W^*)^\Gamma=0.
\]
Both degree zero and degree two of $j_*\Lambda^2W$ vanish.
Its Euler characteristic must therefore be nonpositive.

\item The repeated finite spectrum has four fixed bivectors and
the distinct spectrum has two.  At the cusp, write the unipotent
module as two blocks of size two and one block of size one.
Their exterior square has two trivial blocks from the individual
wedges, blocks of sizes three and one from the product of the
two size-two blocks, and two further size-two blocks.  Its fixed
rank is six.  Consequently
\[
 \chi(\mathbb P^1,j_*\Lambda^2W)=-10+4+2+6=2,
\]
contradicting the preceding step.  This proves all five stated
exclusions.  Removing the three excluded order pairs from
\eqref{maps:eq:F.28} leaves exactly \eqref{maps:eq:F.30}.
\end{enumerate}
\end{proof}




\begin{proposition}[rational uniqueness of the extension]\label{maps:src:F.2.11}
\StatementRecord{maps:src:F.2.11}{proved}{Historical events: Claude Fable 5.2 (Anthropic), effort high, 16 November 2026: carried as background without fault found. GPT-6 Astra (OpenAI), ChatGPT Pro mode, effort Pro, S6\_09\_PAPERS, 17 November 2026: targeted proof check. Also checked in this writing. Pre-continuation writing events: GPT-6 Astra (OpenAI), ChatGPT Pro mode, effort Pro, S6\_10\_PAPERS, 18 November 2026 02:02:02.394029+00:00: targeted proof check. Continuation mathematical and source checks: GPT-6 Astra (OpenAI), ChatGPT Pro mode, effort Pro, S6\_10\_PAPERS continuation, 19 November 2026: computation or formal certificate; independent re-proof; used and rechecked in this run. Continuation metadata events: GPT-6 Astra (OpenAI), ChatGPT Pro mode, effort Pro, S6\_10\_PAPERS continuation, 19 November 2026: used and rechecked in this run; status-vocabulary reconciliation; metadata reconciliation only; no new mathematical check. \emph{Pre-continuation recorded proof scope.} Claude Fable 5.2 (Anthropic), effort high, 16 November 2026; also checked in this writing by GPT-6 Astra (OpenAI), ChatGPT Pro mode, effort Pro, 18 November 2026}{the hypotheses, notation and definitions stated here; \cref{maps:src:F.2.9}; \cref{maps:eq:F.32}}


Let $W$ be a rational rank-five representation of a
twice-punctured affine line, with two finite generators
having fixed rank one and product with unipotent Jordan type
$(2,2,1)$.  Suppose $W$ has neither invariants nor
coinvariants.  Up to rescaling the invariant line, there is
exactly one rational extension
$0\to\mathbb Qe\to V\to W\to0$ in which the two finite
generators retain their finite orders and the cusp has
square-zero logarithm of rank three.  Thus, for the
two-point candidates of \ref{maps:src:F.2.9}, the remaining rational
representation problem is exactly the problem of the
rank-five quotients.  This statement does not identify
integral lattices in those representations.

\end{proposition}


\begin{proof}
\begin{enumerate}[label=\arabic*.]
\item Write the finite generators of $W$ as $S_1,S_2$ and the
extension generators as
\[
 A_i=\begin{pmatrix}1&a_i\\0&S_i\end{pmatrix}.
\]
If $S_i$ has order $m_i$, retaining that finite order is equivalent
to $a_i(I+S_i+\cdots+S_i^{m_i-1})=0$.  Averaging identifies this
condition with $a_i\in\operatorname{row}(S_i-I)$, a space of
dimension four.  A change of splitting by $b\in W^*$ sends the
two rows to $a_i+b(S_i-I)$.  This coboundary map is injective
because $W$ has no invariant covectors.  Thus the extension
space subject to the two finite conditions has dimension $8-5=3$.

\item Put $U=S_1S_2=I+N$, so $N^2=0$ and $\operatorname{rank}N=2$.
The product $A_1A_2$ has upper row $c=a_1S_2+a_2$.
Changing the splitting sends $c$ to $c+bN$.  Hence cusp restriction
is the explicit map
\[
 [(a_1,a_2)]\longmapsto[c]\in W^*/\operatorname{row}N.
\]
It is an isomorphism.  Indeed, a class in its kernel can be
represented with $c=0$, so $a_2=-a_1S_2$.  The finite conditions
then say
\[
 a_1\in\operatorname{row}(S_1-I)
              \cap\operatorname{row}(S_2-I).
\]
This intersection has dimension three: the two four-dimensional
row spaces span $W^*$ because $W$ has no common fixed vector.
The rows $b$ satisfying $bN=0$ also form a three-dimensional
space.  Their map $b\mapsto b(S_1-I)$ lands in the displayed
intersection, since $bU=b$ gives
$b(S_1-I)S_2=b-bS_2$.  It is injective, because
$bS_1=b$ and $bU=b$ together give $bS_2=b$, forcing $b=0$.
It is therefore onto.  Every pair with $c=0$ is a coboundary,
proving injectivity of cusp restriction.  Both its source and
target have dimension three, so it is an isomorphism.

\item The square of the lifted product minus identity vanishes
exactly when $cN=0$.  Under this restriction isomorphism the
square-zero extensions form
\[
 \frac{\operatorname{ann}(\operatorname{im}N)}
      {\operatorname{row}N},
 \tag{F.32}\label{maps:eq:F.32}
\]
of dimension $3-2=1$.  The lifted nilpotent has the rows of $N$
together with the row $c$, so its rank is three exactly when
$c\notin\operatorname{row}N$.  These are precisely the nonzero
classes in \eqref{maps:eq:F.32}.  Rescaling the invariant line
rescales this one-dimensional rational space and is transitive
on its nonzero classes.  This proves the asserted existence and
uniqueness with the stated finite orders and maximal cusp.
\end{enumerate}
\end{proof}




\begin{theorem}[an integral obstruction for a rational $(5,12)$ family]\label{maps:src:F.2.12}
\StatementRecord{maps:src:F.2.12}{proved}{Historical events: Claude Fable 5.2 (Anthropic), effort high, 16 November 2026: carried as background without fault found. GPT-6 Astra (OpenAI), ChatGPT Pro mode, effort Pro, S6\_09\_PAPERS, 17 November 2026: targeted proof check. Also checked in this writing. Pre-continuation writing events: GPT-6 Astra (OpenAI), ChatGPT Pro mode, effort Pro, S6\_10\_PAPERS, 18 November 2026 02:02:02.394029+00:00: targeted proof check. Continuation mathematical and source checks: GPT-6 Astra (OpenAI), ChatGPT Pro mode, effort Pro, S6\_10\_PAPERS continuation, 19 November 2026: computation or formal certificate; independent re-proof; used and rechecked in this run. Continuation metadata events: GPT-6 Astra (OpenAI), ChatGPT Pro mode, effort Pro, S6\_10\_PAPERS continuation, 19 November 2026: used and rechecked in this run; status-vocabulary reconciliation; metadata reconciliation only; no new mathematical check. \emph{Pre-continuation recorded proof scope.} Claude Fable 5.2 (Anthropic), effort high, 16 November 2026; also checked in this writing by GPT-6 Astra (OpenAI), ChatGPT Pro mode, effort Pro, 18 November 2026}{the displayed matrices of \eqref{maps:eq:F.33}; \cref{maps:src:F.2.1,maps:src:F.2.7,night:P1.11}; the actual displayed finite-cover, Wang and cellular maps; Hatcher2002 and StacksProject at the cited locators}


Let a rational rank-six cohomology representation have generators
\[
 A=\begin{pmatrix}1&0\\0&A_5\end{pmatrix},\qquad
 B=\begin{pmatrix}1&c\\0&B_5\end{pmatrix},\qquad
 c=(0,0,0,1,0),
\]
where
\[
 A_5=\begin{pmatrix}
 1&0&1&0&0\\1&0&0&0&0\\-1&0&0&1&1\\
 1&-1&0&-1&0\\0&1&0&0&0
 \end{pmatrix},\qquad
 B_5=\begin{pmatrix}
 0&1&0&0&0\\0&0&0&0&1\\1&-2&0&0&0\\
 0&1&0&-1&-1\\0&0&1&1&2
 \end{pmatrix}.
 \tag{F.33}\label{maps:eq:F.33}
\]
For every invariant integral lattice with an index-one maximal
square-zero cusp, and every free affine filling of multiplicities
$m_1,m_2$ divisible respectively by five and twelve, a simply
connected total space has
\[
 \operatorname{Tor}_3 H_2
       \simeq\mathbb Z/3^{v_3(m_2)}.
 \tag{F.34}\label{maps:eq:F.34}
\]
In particular none of these data is three-connected.
The standard lattice nevertheless has a primitive cusp,
free affine fillings of orders five and twelve, and a simply
connected total space with rational Betti vector
$(1,0,0,0,4,0,0,0,1)$.

\end{theorem}


\begin{proof}
\begin{enumerate}[label=\arabic*.]
\item Put $R=\mathbb Z_{(3)}$ and use the rational coordinates of
\eqref{maps:eq:F.33}.  Let $\mathcal P_R$ be the algebra of blocks
$\left(\begin{smallmatrix}a&b\\0&D\end{smallmatrix}\right)$,
where $a\in R$, $b\in R^{1\times5}$ and $D\in\operatorname{Mat}_5(R)$.
We first determine the actual integral algebra generated by $A,B$.
The following ordered words have $31$ coordinate columns in the order
$(0,j)$, $0\leq j\leq5$, followed by $(i,j)$,
$1\leq i,j\leq5$:
\[
\begin{gathered}
I,A,B,A^2,AB,BA,B^2,A^3,A^2B,ABA,AB^2,BA^2,\\
BAB,B^2A,B^3,A^3B,A^2BA,A^2B^2,ABA^2,AB^2A,\\
AB^3,BA^3,BA^2B,BAB^2,B^2A^2,B^3A,A^3BA,\\
A^3B^2,A^2BA^2,BA^3B,B^2A^3.
\end{gathered}
\]
Their coordinate determinant is $-32$.  Its adjugate divided by
$-32$ therefore expresses every elementary matrix of
$\mathcal P_R$ as an $R$-linear combination of these words.
Since $A,B\in\mathcal P_R$, this proves
$R[A,B]=\mathcal P_R$ with the specified coordinate map.

For a stable lattice $L_R$, the elementary idempotents split its
six coordinate fractional ideals; the matrices $E_{ij}$ for
$i,j\geq1$ make the last five equal.  Thus
\[
 L_R=3^aRe_0\oplus3^bR^5,\qquad b\geq a,
\]
where the inequality follows from the matrices $E_{0j}$.
The matrix $N=AB-I$ has square zero and nonzero Smith factors
$1,1,1$, while $A_5B_5-I$ has factors $1,1$.
In this basis of $L_R$, its upper-right row is multiplied by
$3^{b-a}$.  For $b>a$ its reduction modulo three has rank two,
contradicting the primitive rank-three cusp.  Hence $b=a$.
The actual scalar map $3^{-a}I$ identifies $L_R$ with $R^6$
and intertwines every exterior power of the two generators.
The given cusp operator has the same elementary divisors:
$BA-I=A^{-1}(AB-I)A$, and
$(BA)^{-1}-I=-(BA)^{-1}(BA-I)$.
This proves the asserted local marking for every stable lattice.

All arithmetic in this item is independently reproduced by
\path{F2_12_three_primary_fresh.py} in \path{code/continuation_V/};
its output includes the full word-coordinate matrix, its determinant,
and the exact Smith and invariant-basis certificates used below.
The argument from that determinant treats all lattices, without a
bound on their indices.

\item Write $v_i=g_i^{m_i}$ for the actual finite power vectors.
They are fixed by the homology matrices $M_i$.
The base-orbifold quotient in \cref{maps:src:F.2.7} shows that
simple connectivity forces $\gcd(m_1,m_2)=1$.
After the cusp meridian relation the two base generators are
inverse up to a fibre translation.  Their actions on the abelian
fibre lattice are therefore inverse.  Their linear orders five
and twelve force both actions to be the identity in this quotient.
The exact Smith factors of $[A^t-I\ B^t-I]$ are five units,
with first row zero, so its fibre coinvariant module over $R$
is $R$ through the primitive character $\ell=x_0$.

Set $a_i=\ell(v_i)$.  An arbitrary cusp clutch contributes a scalar
$\delta$ to $g_1+g_2=\delta z$.  The remaining relation matrix is
\[
 \begin{pmatrix}m_1&0&-a_1\\0&m_2&-a_2\\1&1&-\delta\end{pmatrix}.
\]
Its determinant is a unit for a simply connected total space.
Since $3\mid m_2$ and $3\nmid m_1$, reduction modulo three makes
$a_2$ a unit in $R$.  This conclusion includes arbitrary clutch
translations and twists.

\item We determine the actual second finite-cap restriction over $R$.
Choose a primitive integral invariant character $\alpha$ whose
$R$-span is $\ell$.  The stable sublattice
\[
 \Lambda'=\mathbb Z v_2\oplus\ker\alpha
\]
has index $|\alpha(v_2)|$, prime to three.  It gives an actual
finite torus cover; the affine action lifts because it preserves
$v_2$, preserves $\ker\alpha$, and translates by $v_2/m_2$.
The corresponding reduced cap $F'_2$ covers $F_2$ with the same
degree.  In its basis beginning with $v_2$, the action is
\[
 (u,y)\longmapsto(u+1/m_2,Dy).
\]
Thus $F'_2$ is the mapping torus of $D:T^5\to T^5$, and its
regular six-torus is the $m_2$-fold cover in the circle direction.
The coefficient change has determinant a unit in $R$; under the
scalar marking of item~1, the vertical cohomology action $D^t$
is the displayed $B_5$.

The matrices $D^t-I$ and $\bigwedge^2D^t-I$ have no nonunit
three-primary Smith factors.  Their Wang sequences therefore make
$H^q(F'_2;R)$ free for $q\leq3$.  Pullback to the generic torus
is injective in these degrees, because transfer kills its kernel
by $m_2$ and the source is free.  Deck transformations of
$F'_2\to F_2$ act by translations on the generic torus cover.
Their action there is trivial on cohomology, so the injectivity
just proved makes their action on $H^q(F'_2;R)$ trivial as well.
Transfer for this prime-to-three cover now identifies
$H^q(F_2;R)$ with $H^q(F'_2;R)$, compatibly with their actual
restriction maps.  In particular the degree-two and degree-three
restrictions have zero kernel.

Put $T=D^t$.  Its polynomial is $(x-1)\Phi_{12}(x)$, and
\[
 E=\Phi_{12}(T)=T^4-T^2+I,\qquad E^2=E,\qquad
 \sum_{j=0}^{11}T^j=12E.
\]
The image of this integral rank-one idempotent is the primitive
invariant line.  Hence the norm for every stated $m_2$ is $m_2E$.
The Wang sequence in degree two has quotient
$\operatorname{Inv}(\bigwedge^2T)$ and kernel
$\operatorname{coker}(T-I)$.  Restriction to the $m_2$-fold cover
projects identically onto the former group.  On the latter it is
the horizontal wedge with the norm: a horizontal cell lifts to
$m_2$ successive transported prisms.
Consequently the restriction image $J_2$ projects onto the vertical
invariants, and its intersection with the horizontal invariant
summand is $e_0\wedge m_2\operatorname{Inv}(T)$.
The actual horizontal inclusion induces the isomorphism
\[
 \frac{\operatorname{Inv}(T)}{m_2\operatorname{Inv}(T)}
 \ \xrightarrow{\ \sim\ }\\
 \frac{\operatorname{Inv}(\bigwedge^2V_R)}{J_2}
 \ \cong\ R/3^rR.
\]
No choice of lifts of the vertical Wang classes changes this
quotient or its kernel.  At the first finite cap, $m_1$ is a unit
in $R$, so transfer and averaging give the full invariant image
and zero kernel.  At the cusp, \cref{night:P1.11} gives injective
restrictions onto the invariant lattices in degrees one, two
and three, by dualizing its actual saturated-kernel collapse.

\item The integral specialization-sheaf approach is used in
\cite[Lemma 7.25 and Propositions 7.26--7.27, pp.~61--63]{AC26}
and \cite[section ``The integral homology'']{E26}.
Here are the coefficient maps and calculations for the present lattice.  Put $L_q=\bigwedge^qV_R$, $A_q=\bigwedge^qA$,
$B_q=\bigwedge^qB$, and let $j$ include the thrice-punctured base.
For the full invariant-extension sheaf $j_*L_q$ the oriented
cellular complex is
\[
\begin{gathered}
 J_{1,q}\oplus J_{2,q}\oplus J_{0,q}
 \xrightarrow{d_0}L_q^2\xrightarrow{\Delta_q}L_q,\\
 J_{1,q}=\ker(A_q-I),\quad J_{2,q}=\ker(B_q-I),\quad
 J_{0,q}=\ker(B_qA_q-I),\\
 d_0(j_1,j_2,j_0)=(j_1-j_0,j_2-A_qj_0),\\
 \Delta_q(a,b)=(A_q-I)a+(B_q-I)b.
\end{gathered}
\]
These are the actual base attachment maps: the two cusp paths
carry $j_0$ and $A_qj_0$, whose combined boundary is
$(B_qA_q-I)j_0=0$.  Primitive invariant bases and a unimodular
kernel basis of $\Delta_q$ give the following certificates:
\[
\begin{array}{c|c|c|c|c}
q&\operatorname{rank}H^0&\operatorname{rank}H^1&
 \operatorname{SNF}(d_0\text{ in }\ker\Delta_q)&
 \operatorname{SNF}(\Delta_q)\\\hline
1&1&0&1^6&1^5,2\\
2&0&0&1^{15}&1^{15}\\
3&0&3&1^{16},2&1^{19},2.
\end{array}
\]
Here $1^s$ denotes $s$ entries equal to one.  The row-coordinate
matrices and their unimodular HNF transformations are specified
by the preceding formulas and retained in the independent output.
In particular $H^0(j_*L_2)=H^1(j_*L_2)=0$,
$H^2(j_*L_1)=0$ over $R$, and $H^0(j_*L_3)=0$.

Write $\mathcal R_q=R^qf_*R$.  The actual stalk/fibre
identification is proper base change
\cite[Tag 09V5, Lemma 20.18.1]{StacksProject}.
The local maps of item~3 give
\[
 0\longrightarrow\mathcal R_2\longrightarrow j_*L_2
 \longrightarrow Q\longrightarrow0,
 \qquad Q_{p_2}=R/3^rR,
\]
with no other support.  Hence its connecting homomorphism gives
$H^1(\mathcal R_2)=R/3^rR$.
The degree-one quotient is also point supported, so
$H^2(\mathcal R_1)=H^2(j_*L_1)=0$.
The degree-three sheaf map is injective, so
$H^0(\mathcal R_3)=0$.

\item The only total-degree-three term in the Leray sequence
\cite[Tag 01F2, Lemma 20.13.4]{StacksProject} is therefore
$E_2^{1,2}=R/3^rR$.  It has no possible differential because
the base has cohomological dimension two.  Thus
$H^3(M;R)=R/3^rR$.
Universal coefficients
\cite[Theorem 3.2, p.~195]{Hatcher2002} identify this torsion
module with $\operatorname{Ext}_R(H_2(M;R),R)$: the quotient
$\operatorname{Hom}_R(H_3(M;R),R)$ is free and must vanish.
The elementary divisor decomposition over the PID $R$ proves
\eqref{maps:eq:F.34}.

\item For the standard lattice take
\[
 v_1=(3,0,0,0,0,0)^t,\qquad
 v_2=(-7,7,-7,7,0,7)^t,\qquad (m_1,m_2)=(5,12).
\]
They are fixed by $A^t,B^t$ respectively.
Their $\ell$-values $3,-7$ are units modulo the respective
multiplicities; applying $\ell$ to a fixed-point equation proves
that every nonidentity affine power is free.
The primitive cusp in item~1 has the integral marking of
\cref{maps:src:F.2.1}.  Glue the two finite quotient caps and
the index-one Kuhn cap to the marked torus bundle over a pair
of pants, with zero additional clutch translation.  The finite
boundary identification is the prescribed logarithmic fibre
translation; the cusp is attached in that integral marking.
These actual local models give a closed smooth oriented
manifold.  Its remaining van Kampen matrix is
\[
 \begin{pmatrix}5&0&-3\\0&12&7\\1&1&0\end{pmatrix},
\]
with determinant of absolute value one.  It is simply connected.

Over $\mathbb Q$ the finite images are full invariants.
Item~4 gives vanishing $H^1(j_*V_{\mathbb Q})$,
$H^0(j_*\bigwedge^2V_{\mathbb Q})$,
$H^1(j_*\bigwedge^2V_{\mathbb Q})$ and
$H^2(j_*V_{\mathbb Q})$, and the degree-three invariant group
is zero.  The five-term Leray sequence and $H^1(M;\mathbb Q)=0$
make the map from the one-dimensional $H^0(\mathcal R_1)$ to
$H^2(\mathbb P^1;\mathbb Q)$ an isomorphism.
Thus $b_2=b_3=0$.
Euler additivity gives $\chi(M)=6$ from the Kuhn cusp;
all other fibres have Euler number zero.
Poincar\'e duality therefore gives
$(1,0,0,0,4,0,0,0,1)$, as claimed.
\end{enumerate}
\end{proof}




\begin{theorem}[two further rational families excluded integrally]\label{maps:src:F.2.13}
\StatementRecord{maps:src:F.2.13}{proved}{Historical events: Claude Fable 5.2 (Anthropic), effort high, 16 November 2026: carried as background without fault found. GPT-6 Astra (OpenAI), ChatGPT Pro mode, effort Pro, S6\_09\_PAPERS, 17 November 2026: targeted proof check. Also checked in this writing. Pre-continuation writing events: GPT-6 Astra (OpenAI), ChatGPT Pro mode, effort Pro, S6\_10\_PAPERS, 18 November 2026 02:02:02.394029+00:00: targeted proof check. Continuation mathematical and source checks: GPT-6 Astra (OpenAI), ChatGPT Pro mode, effort Pro, S6\_10\_PAPERS continuation, 19 November 2026: computation or formal certificate; targeted proof check; used and rechecked in this run. Continuation metadata events: GPT-6 Astra (OpenAI), ChatGPT Pro mode, effort Pro, S6\_10\_PAPERS continuation, 19 November 2026: used and rechecked in this run; status-vocabulary reconciliation; metadata reconciliation only; no new mathematical check. \emph{Continuation recorded proof scope.} Every step of the current nine-step proof has renewed prereading coverage: the complete universal norm and sheaf argument was frozen by the Paper V reviewer before its target, and the explicit smooth nonvacuity construction by the Paper VI reviewer before its isolated target. The actual quotient, norm, affine collar and cohomology maps are retained. Newly authored exact programs verify the current computations; separate post-target programs reproduce the former periods, Gram forms and lattice data with explicit coordinate changes. Former proof routes are preserved without retrospective prereading credit. The supporting source scope was renewed.}{the two stated companion constructions; \cref{maps:src:F.2.11}; \cref{maps:src:F.2.7}; \cref{night:P1.9}; \cref{night:P1.11}; the displayed actual integral norm and invariant-extension complexes; \cite[Tag~09V5, Lemma~20.18.1; Tag~01F2, Lemma~20.13.4]{StacksProject}; \cite[Theorem~3.2, p.~195]{Hatcher2002}}


Let $C_f,C_g$ denote the degree-four companion matrices.
Take $f=\Phi_5$ and either
\[
 \text{case A: }g=\Phi_{12},\qquad
 \text{case B: }g=\Phi_3\Phi_6.
 \tag{F.36}\label{maps:eq:F.36}
\]
Let $W$ be the quotient of
$\bigwedge^2(C_f^{-1},C_g)$ by its invariant line, and
let $V$ be its unique rational extension of \ref{maps:src:F.2.11}.
For every invariant integral lattice in $V$ with primitive
maximal square-zero cusp, every simply connected free affine
filling, with multiplicities $m_1,m_2$ divisible by five and
six respectively, fails three-connectivity.
More precisely, put $r=v_3(m_2)\geq1$.
In case A, $\operatorname{Tor}_3 H_3$ has order
$3^{2r-1}$.  In case B,
$\operatorname{Tor}_3 H_2\simeq\mathbb Z/3^r$.
For $m_2=6$ the respective obstructing group is $\mathbb Z/3$.

\end{theorem}


\begin{proof}
\begin{enumerate}[label=\arabic*.]
\item We first give actual rational models, including the quotient
maps.  Use the exterior basis $(01,02,03,12,13,23)$ for the companion
pair.  Its common primitive invariant vectors are
\[
 \omega_{\mathrm A}=(1,2,1,2,2,1)^t,\qquad
 \omega_{\mathrm B}=(1,0,1,0,0,1)^t.
\]
Let $T$ have columns $\omega,e_1,\ldots,e_5$; its determinant is one.
The last five rows of $T^{-1}$ give the quotient map.  In that
quotient the two matrices $a,b$, and an extension row $c$, are
\[
\begin{array}{c|c|c}
 &a&b\\\hline
\mathrm A&
\begin{pmatrix}2&-1&-2&1&0\\1&0&-1&0&0\\3&-1&-2&0&1\\
3&0&-2&0&0\\1&1&-1&0&0\end{pmatrix}&
\begin{pmatrix}0&-2&0&1&0\\0&-1&0&0&1\\0&-1&0&0&0\\
1&-2&0&0&0\\0&-1&1&0&-1\end{pmatrix}\\[3pt]
\mathrm B&
\begin{pmatrix}0&-1&0&1&0\\1&0&-1&0&0\\1&-1&0&0&1\\
1&0&0&0&0\\1&1&-1&0&0\end{pmatrix}&
\begin{pmatrix}0&0&0&1&0\\0&-1&0&0&1\\0&-1&0&0&0\\
1&0&0&0&0\\0&-1&1&0&1\end{pmatrix}
\end{array}
\]
\[
 c_{\mathrm A}=(-2,-3,1,2,0),\qquad
 c_{\mathrm B}=(0,-1,1,0,0),\qquad
 A=\operatorname{diag}(1,a),\quad
 B=\begin{pmatrix}1&c\\0&b\end{pmatrix}.
\]
Exterior eigenvalue products give characteristic polynomial
$(x-1)\Phi _5$ for $a$, and $(x-1)(x+1)^2\Phi _3$ or
$(x-1)(x+1)^2\Phi _6$ for $b$.  Each has fixed rank one.
Direct multiplication gives $A^5=B^6=I$, $(AB-I)^2=0$, and
nonzero Smith factors $1,1,1$ for $AB-I$ and $1,1$ for $ab-I$.
The row conditions are $c\sum_{j=0}^5b^j=0$ and $c(ab-I)=0$.
The full quotient matrix algebra verified in the next item implies
that this rank-five quotient is irreducible, with no invariant
vectors or covectors.  Thus \cref{maps:src:F.2.11} identifies these
models with the specified rational extensions, up to rescaling
their invariant line.

\item Put $R=\mathbb Z_{(3)}$ and let $\mathcal P_R$ be the algebra
of blocks $\left(\begin{smallmatrix}u&v\\0&D\end{smallmatrix}\right)$,
with $u\in R$, $v\in R^{1\times5}$ and $D\in\operatorname{Mat}_5(R)$.
In the coordinate order $(0,j)$, $0\leq j\leq5$, followed by
$(i,j)$, $1\leq i,j\leq5$, the following $31$ words have coordinate
determinants $512$ and $-512$ in cases A and B:
\[
\begin{aligned}
&I,A,B,A^2,AB,BA,B^2,A^3,A^2B,ABA,AB^2,BA^2,BAB,B^2A,\\
&A^3B,A^2BA,A^2B^2,ABA^2,AB^2A,BA^3,BA^2B,BAB^2,B^2A^2,B^2AB,\\
&A^3BA,A^3B^2,A^2BA^2,ABA^3,ABA^2B,BA^3B,A^3BA^2.
\end{aligned}
\]
Their adjugates divided by the displayed determinants express
every elementary matrix of $\mathcal P_R$ as an $R$-linear
combination of these actual words.  Consequently $R[A,B]=\mathcal P_R$.
Its lower block also proves the irreducibility used in item~1.

For an invariant lattice, the elementary projectors and lower-block
matrix units give
\[
 L_R=3^\alpha Re_0\oplus3^\beta R^5,\qquad \beta\geq\alpha;
\]
the inequality is forced by $E_{0j}$.  In this basis the upper
nilpotent row is multiplied by $3^{\beta-\alpha}$.  If
$\beta>\alpha$, reduction modulo three has rank two, contradicting
the primitive rank-three cusp.  Thus $\beta=\alpha$ and the actual
map $3^{-\alpha}I$ identifies $L_R$ with $R^6$, intertwining every
exterior power.  The other product conventions are related by
\[
 BA-I=A^{-1}(AB-I)A,\qquad
 (BA)^{-1}-I=-(BA)^{-1}(BA-I).
\]
This proves the assertion for all invariant lattices, with no
bound on their indices.

\item We record the consequence of simple connectivity for the
actual affine translations.  Write a finite affine generator,
after a translation of the torus, as
$x\mapsto M_ix+v_i/m_i$, where $M_iv_i=v_i$ and $v_i$ is its
integral power vector.  This change has an explicit map.  If $t$
is a lift of the original translation and $d$ is the linear order,
put $t_0=d^{-1}\sum_{j=0}^{d-1}M_i^jt$ and $b_0=t-t_0$.
Then
\[
 h=\frac1d\sum_{j=0}^{d-1}(d-1-j)M_i^jb_0,
 \qquad (I-M_i)h=b_0.
\]
For $T_h(x)=x+h$, the actual conjugate $T_h^{-1}gT_h$ has translation $t_0$, and
$v_i=m_it_0$ is integral because the given affine action has
period $m_i$.

The base-orbifold quotient in \cref{maps:src:F.2.7} shows that
$\gcd(m_1,m_2)=1$.  In particular $m_1$ is a unit in $R$.
The homology augmentation matrix $[A^t-I\mid B^t-I]$ has first
row zero and rank-five minors $-4$ and $-1$ in its last five
rows, respectively.  Its coinvariant quotient over $R$ is
therefore $R$ through $\ell(x)=x_0$.  With $a_i=\ell(v_i)$,
the finite power relations and cusp clutch give the quotient
relation matrix
\[
 \begin{pmatrix}m_1&0&-a_1\\0&m_2&-a_2\\1&1&-\delta\end{pmatrix}.
\]
Here $\delta$ is the character value of the actual cusp clutch.
The character annihilates the collapsed angular cycles, so these
maps on the finite caps and cusp agree on the boundary and define
a quotient of the abelianized fundamental group.  Simple
connectivity makes its determinant a unit.  Reduction modulo
three forces $a_2$ to be a unit.  Arbitrary clutch translations
are included in $\delta$.

\item Let $\alpha$ be the primitive integral invariant character
whose $R$-span is $\ell$, and let $K=\ker\alpha$.  The invariant
sublattice $\Lambda'=\mathbb Zv_2\oplus K$ has index
$|\alpha(v_2)|$, prime to three.  It gives an actual torus cover
and an actual cover $F'_2\to F_2$ of the second reduced finite
fibre.  In coordinates beginning with $v_2$, its affine map is
\[
 (\theta,y)\longmapsto(\theta+1/m_2,Dy).
\]
Thus $F'_2$ is a five-torus mapping torus, and its generic
six-torus is the $m_2$-fold cover in the circle direction.
Restriction to $K$ is the actual dual quotient
$(x_0,x')\mapsto x'$ from $V_R$ to $W_R$, intertwining $B$ with $b$.
Reversing the mapping-torus circle replaces $b$ by $b^{-1}$;
the identities $b^{-1}-I=-b^{-1}(b-I)$ and
$\sum_{j=0}^5b^{-j}=\sum_{j=0}^5b^j$ give the exact comparison.

Put $b_k=\bigwedge^kb$.  The nonzero Smith factors of $b_k-I$
for $0\leq k\leq4$ are units in $R$.  The Wang sequences make
$H^q(F'_2;R)$ free for $q\leq4$.  Pullback to its generic torus
is injective, since transfer kills its kernel by $m_2$ and the
source is free.  Deck transformations of $F'_2\to F_2$ act by
translations on that generic cover.  Their action on its
cohomology is trivial, and injectivity makes it trivial on
$H^q(F'_2;R)$ as well.  Transfer for this prime-to-three cover
therefore identifies the restriction maps with those for $F_2$.

The actual covering map on the horizontal term of the Wang
sequence is the norm
\[
 \mathcal N_{m_2,k}=\sum_{j=0}^{m_2-1}b_k^j:
 \operatorname{coker}(b_k-I)\longrightarrow\operatorname{Inv}(b_k).
\]
Indeed, the lift is the union of the $m_2$ successive prisms
with their transported coefficients.  On the vertical invariant
quotient pullback is the identity.  Consequently horizontal
inclusion induces the actual isomorphism
\[
 \frac{\operatorname{Inv}(b_{q-1})}
      {\mathcal N_{m_2,q-1}\bigwedge^{q-1}R^5}
 \ \xrightarrow{\ \sim\ }\
 \frac{\operatorname{Inv}(\bigwedge^q V_R)}{J_{2,q}}.
\]
Surjectivity follows by subtracting a restriction with the same
vertical component; its kernel is exactly the horizontal norm
image.  Thus no choice of a Wang splitting changes this quotient.

\item Write $\lambda=m_2/6$.  Since $b^6=I$,
$\mathcal N_{m_2,k}=\lambda\mathcal N_{6,k}$.
In case A a primitive basis of $\operatorname{Inv}(b_2)$,
in lexicographic exterior order, is the column matrix
\[
 Q=\begin{pmatrix}
 1&0\\1&4\\2&4\\0&2\\0&3\\1&4\\0&3\\1&0\\0&0\\0&2
 \end{pmatrix}.
\]
The norm coordinates in this basis are
\[
 \begin{pmatrix}
 0&3&0&-3&0&0&-4&3&4&3\\
 0&0&0&0&0&0&2&0&-2&0
 \end{pmatrix}.
\]
Their image has basis $\left(\begin{smallmatrix}1&0\\4&6\end{smallmatrix}\right)$,
and multiplication by
$\left(\begin{smallmatrix}1&0\\-4&1\end{smallmatrix}\right)$
gives $\operatorname{diag}(1,6)$.  This gives the explicit quotient
map $(x,y)\mapsto(x,y-4x)$ modulo $(\lambda,6\lambda)$.
In case B the primitive invariant vector of $b$ is
$(1,0,0,1,0)^t$, and the norm coordinates are $(3,0,0,3,0)$.
After multiplication by the remaining units of $R$, these give
\[
 Q_{2,3}^{\mathrm A}\cong R/3^{r-1}R\oplus R/3^rR,
 \qquad Q_{2,2}^{\mathrm B}\cong R/3^rR.
\]
The HNF transformations giving the stated image bases and their
inverse coordinate maps are retained in the fresh certificate.

\item We use the integral specialization and base-sheaf method
of \cite[Lemma~7.25 and Propositions~7.26--7.27, pp.~61--63]{AC26}
and \cite[section ``The integral homology'']{E26}, with the
following explicitly computed coefficients.  Put
$L_q=\bigwedge^qR^6$, $A_q=\bigwedge^qA$, $B_q=\bigwedge^qB$.
The actual based-sphere complex for $j_*L_q$ is
\[
\begin{gathered}
 J_{1,q}\oplus J_{2,q}\oplus J_{0,q}
 \xrightarrow{d_0} L_q^2\xrightarrow{\Delta_q}L_q,\\
 J_{1,q}=\ker(A_q-I),\quad J_{2,q}=\ker(B_q-I),\quad
 J_{0,q}=\ker(B_qA_q-I),\\
 d_0(j_1,j_2,j_0)=(j_1-j_0,j_2-A_qj_0),\\
 \Delta_q(x,y)=(A_q-I)x+(B_q-I)y.
\end{gathered}
\]
The cusp paths carry $j_0$ and $A_qj_0$, so their combined
boundary is zero.  Inverse generator conventions give a chain
isomorphism: use cusp basis $A_qJ_{0,q}$ for the inverse pair,
and maps $-I$, $\operatorname{diag}(-A_q,-B_q)$, $I$ in degrees
zero, one and two.  Both chain identities follow by multiplication.

Unimodular HNF transformations give primitive invariant bases
and kernel coordinates for $d_0$ in $\ker\Delta_q$.  In both cases
the exact Smith records are
\[
\begin{array}{c|ccc|c|c}
q&\operatorname{rank}H^0&\operatorname{rank}H^1&\operatorname{rank}H^2
 &\operatorname{SNF}(d_0\text{ in }\ker\Delta_q)
 &\operatorname{SNF}(\Delta_q)\\\hline
1&1&0&0&1^6&1^5,2\\
2&0&0&0&1^{15}&1^{14},2\\
3&0&3&0&1^{17}&1^{19},2\\
4&0&0&0&1^{15}&1^{14},2
\end{array}
\]
All displayed groups are therefore free over $R$.  The independently
authored programs
\path{F2_13_companion_explore_fresh.py},
\path{F2_13_algebra_fresh.py}, and \path{F2_13_cellular_fresh.py}
in \path{code/continuation_V/} specify every exterior matrix,
kernel-coordinate map, unit rank minor, and norm-image HNF
transformation.  Their computations use the printed models.

\item Write $\mathcal R_q=R^qf_*R$.  Proper base change
\cite[Tag~09V5, Lemma~20.18.1]{StacksProject} identifies the stalks
with the actual fibres.  At the first finite cap, $m_1$ is a
unit in $R$, so transfer gives the full invariant restriction
image and zero kernel.  At the second finite cap, item~4 gives
injectivity and the norm quotient.  At the prescribed cusp,
dualizing the actual primitive-image collapse of
\cref{night:P1.11} gives the invariant restriction image, with
zero kernel in degrees one through three and a free kernel $R$
in degree four.  The latter is dual to the free degree-four
cokernel of the collapse and must be retained.

For $q=2,3$ there is consequently an exact sequence
\[
 0\longrightarrow\mathcal R_q\longrightarrow j_*L_q
 \longrightarrow Q_q\longrightarrow0,
\]
where $Q_q$ is supported at the second finite point.  Since
$H^0(j_*L_q)=0$, its connecting map embeds $Q_q$ in
$H^1(\mathcal R_q)$, with the free quotient $H^1(j_*L_q)$
computed above.  Its torsion is exactly $Q_q$.
Also $H^2(\mathcal R_1)=H^2(\mathcal R_2)=H^2(\mathcal R_3)=0$.
Injectivity gives $H^0(\mathcal R_3)=0$.  The degree-four
free cusp kernel, together with $H^0(j_*L_4)=0$, gives
$H^0(\mathcal R_4)=R$.

\item In the Leray sequence
\cite[Tag~01F2, Lemma~20.13.4]{StacksProject}, total degree three
has only $H^1(\mathcal R_2)$.  Its middle base column has no
incoming or outgoing differential.  Thus case B gives
$H^3(M;R)=R/3^rR$.
In total degree four, the surviving terms are
$H^1(\mathcal R_3)$ and $H^0(\mathcal R_4)=R$;
the degree-two base term is zero and the possible differential
from the latter term lands in $H^2(\mathcal R_3)=0$.
The extension by the free module $R$ splits.  Hence in case A
\[
 \operatorname{Tor}H^4(M;R)
   \cong R/3^{r-1}R\oplus R/3^rR.
\]
Universal coefficients
\cite[Theorem~3.2, p.~195]{Hatcher2002} identify this torsion
with the torsion of $H_3(M;R)$, and similarly identify the
case-B torsion with that of $H_2(M;R)$.  Localization retains
the finite three-primary subgroups by the identity on each
$\mathbb Z/3^s$ summand.  This proves the stated orders and
the case-B isomorphism.  At $m_2=6$ the first case-A summand
vanishes and the second is $\mathbb Z/3$; case B is also
$\mathbb Z/3$.  These groups obstruct three-connectivity.
\item The two displayed rational models have explicit smooth marked
realisations with the required free caps and simply connected filling.
Use their displayed standard cohomology lattice $\mathbb Z^6$ and its
actual dual homology basis, so $M_1=A^t$, $M_2=B^t$.
The independent exact calculation gives $A^5=B^6=I$, exact orders five
and six, and determinant one. Put $K=BA-I$. Then $K^2=0$, and the
third minor on rows $(0,3,4)$ and columns $(1,2,3)$ is one in both
models. Hence $K$ has rank three and all its nonzero Smith factors
are one. The actual marked cusp is
\[
 M_0-I=((M_1M_2)^{-1}-I)=-K^t.
\]
Set $b_1=e_0,b_2=e_3,b_3=e_4$, $a_j=(M_0-I)b_j$ and
$S=(a_1,a_2,a_3,b_1,b_2,b_3)$. Direct multiplication gives
$S\in\operatorname{GL}_6(\mathbb Z)$ and
\[
 (M_0-I)S=S\begin{pmatrix}0&I_3\\0&0\end{pmatrix}.
\]
This is the exact integral marking used to transport the index-one
Kuhn cap and its collapse maps.

The affine power vectors in these same homology coordinates are
\[
\begin{array}{c|cc}
 &v_1&v_2\\\hline
 \mathrm A&(1,0,0,0,0,0)&(-1,2,0,-1,0,0)\\
 \mathrm B&(1,0,0,0,0,0)&(-1,0,0,0,0,1).
\end{array}
\]
They satisfy $M_iv_i=v_i$. The actual first-coordinate covector
$\ell$ is fixed by both $M_i$ and has values $1,-1$ on $v_1,v_2$.
Consequently every nonidentity power of
$x\mapsto M_1x+v_1/5$ or $x\mapsto M_2x+v_2/6$ changes the
$\ell$-circle by $j/5$ or $-j/6$, respectively. These are free
cyclic actions of the exact stated orders. Their diagonal actions
on the product with the rotating disc therefore give the actual
smooth finite quotient caps.

Take the affine torus bundle over the three-holed sphere with these
two boundary holonomies and their inverse product on the third
boundary. Glue the two finite caps and the proper index-one Kuhn
cap of \cref{night:P1.9,night:P1.11} by the
specified marked mapping-torus collars and zero cusp clutch.
For an affine cusp monodromy $x\mapsto M_0x+c$, a path of translations
$d(s)$ with $d(0)=0$, $d(1)=c$ gives the actual collar map
$(x,s)\mapsto(x+d(s),s)$ from the linear mapping torus. The prescribed
zero clutch fixes its integral lift. The lattice map on this collar
is the displayed $S$.

The cap fundamental-group maps retain these choices. On the finite
cap the universal covering $\mathbb R^6\times\Delta$ has the deck
transformation
$(x,t)\mapsto(M_ix+v_i/m_i,e^{2\pi i/m_i}t)$; its $m_i$-th power
is the lattice translation $v_i$. Thus the boundary generators
$a,b$ satisfy $a^5=v_1$, $b^6=v_2$. On the Kuhn cap the polar
projection onto its base three-torus and angular-identity section
induce inverse maps on fundamental groups. In its dual-cell model
the one-skeleton consists of the base edges, and each angular loop
contracts after transport to a dual vertex where that angular
fibre is collapsed. The fibre map is therefore exactly
$\Lambda\to\Lambda/\operatorname{im}(M_0-I)$, expressed by the last
three rows of $S^{-1}$. The meridian bounds the toric coordinate
disc at a smooth point of a central component; the zero clutch
identifies it with $(ab)^{-1}$. Hence the cusp imposes $ab=1$.

For $J=[M_1-I\mid M_2-I]$, the first row is zero. Its lower five
rows have a unit minor on columns $(1,2,3,6,10)$ in case A and
$(1,2,3,4,6)$ in case B, with determinants $-1$ and $1$.
The inverse of that minor supplies actual integral preimages of
$e_1,\ldots,e_5$. Thus
$\operatorname{im}J=\ker\ell$ integrally. Van Kampen now gives
simple connectivity at the group level: from $b=a^{-1}$ and the
two power relations,
\[
 a^{-1}=a^5a^{-6}=v_1v_2\in\Lambda.
\]
So $a,b$ lie in the abelian image of the fibre lattice, their
conjugation actions become trivial, and $J$ is killed. The exact
remaining fibre quotient is $\mathbb Z$ via $\ell$, with generator
$h$ and images $v_1=h$, $v_2=-h$. The displayed identity forces
$a=1$, and $a^5=v_1$ then kills $h$. Every generator is trivial.
Equivalently, after this nonabelian reduction the residual relation
matrix on $(h,a,b)$ is
\[
 \begin{pmatrix}-1&5&0\\1&0&6\\0&1&1\end{pmatrix},
 \qquad\det=1.
\]

Finally use the checked rational invariant-extension table for these
exact models. The actual finite covering restrictions identify
rational cap cohomology with invariants. The primitive cusp collapse
of~\ref{night:P1.11} gives the same identification through degree
three and a one-dimensional kernel in degree four. Hence
$R^q p_*\mathbb Q=j_*\Lambda^qV$ for $1\le q\le3$, and
\[
 0\longrightarrow\mathbb Q_{\mathrm{cusp}}
 \longrightarrow R^4p_*\mathbb Q
 \longrightarrow j_*\Lambda^4V\longrightarrow0.
\]
The table therefore leaves, in total degree at most four, only
$E_2^{0,0}=E_2^{2,0}=E_2^{0,1}=\mathbb Q$,
$E_2^{1,3}=\mathbb Q^3$ and $E_2^{0,4}=\mathbb Q$.
Simple connectivity forces the actual
$d_2:E_2^{0,1}\to E_2^{2,0}$ to be an isomorphism. There are no
other possible nonzero differentials on these terms. Thus
$(b_0,\ldots,b_4)=(1,0,0,0,4)$; oriented duality on the closed
smooth eight-manifold gives the full tuple
$(1,0,0,0,4,0,0,0,1)$.
\end{enumerate}
\end{proof}




\begin{theorem}[four order pairs and eight local spectral pairs]\label{maps:src:F.2.14}
\StatementRecord{maps:src:F.2.14}{proved}{Historical events: Claude Fable 5.2 (Anthropic), effort high, 16 November 2026: carried as background without fault found. GPT-6 Astra (OpenAI), ChatGPT Pro mode, effort Pro, S6\_09\_PAPERS, 17 November 2026: targeted proof check. Also checked in this writing. Pre-continuation writing events: GPT-6 Astra (OpenAI), ChatGPT Pro mode, effort Pro, S6\_10\_PAPERS, 18 November 2026 02:02:02.394029+00:00: targeted proof check. Continuation mathematical and source checks: GPT-6 Astra (OpenAI), ChatGPT Pro mode, effort Pro, S6\_10\_PAPERS continuation, 19 November 2026: computation or formal certificate; independent re-proof; used and rechecked in this run. Continuation metadata events: GPT-6 Astra (OpenAI), ChatGPT Pro mode, effort Pro, S6\_10\_PAPERS continuation, 19 November 2026: used and rechecked in this run; status-vocabulary reconciliation; metadata reconciliation only; no new mathematical check. \emph{Pre-continuation recorded proof scope.} Claude Fable 5.2 (Anthropic), effort high, 16 November 2026; also checked in this writing by GPT-6 Astra (OpenAI), ChatGPT Pro mode, effort Pro, 18 November 2026}{the hypotheses, notation and definitions stated here; \cref{maps:src:F.2.15}; \cref{maps:src:F.2.9}; \cref{maps:src:F.2.7}; \cref{maps:eq:F.40}; \cref{maps:src:F.2.2}; \cref{night:P1.11}}


Every simply connected two-finite-point datum with $b_2=0$ has
one finite linear order equal to five and the other equal
to four, six, eight or twelve.
After suppressing the common factor $(T-1)^2$, its unordered
pair of finite characteristic polynomials is exactly one
of the following necessary possibilities:
\[
\begin{array}{c|c|l}
 \text{orders}&\text{order-five factor}&\text{other factor}\\\hline
 (5,4)&\Phi_5&\Phi_2^2\Phi_4\\
 (5,6)&\Phi_5&\Phi_2^2\Phi_3,\ \Phi_2^2\Phi_6,\ \Phi_3\Phi_6\\
 (5,8)&\Phi_5&\Phi_8\\
 (5,12)&\Phi_5&\Phi_{12},\ \Phi_3\Phi_4,\ \Phi_4\Phi_6.
\end{array}
 \tag{F.39}\label{maps:eq:F.39}
\]
Its rank-five quotient is rationally irreducible.
In particular all order-three pairs, including $(3,4)$,
are impossible.  The table is a necessary spectral filter;
it does not identify the relative rational or integral
positions of the matrices having those polynomials.

\end{theorem}


\begin{proof}
\begin{enumerate}[label=\arabic*.]
\item By \cref{maps:src:F.2.9,maps:src:F.2.15}, the cohomology
system $V$ has a one-dimensional invariant space $I$, no invariant
covectors, and two-dimensional finite fixed spaces $F_1,F_2$.
The finite orders are among the coprime pairs of
\cref{maps:src:F.2.8}.  At a finite quotient, averaging identifies
the degree-two specialization with the fixed part of $\Lambda^2V$.
At the cusp, \cref{night:P1.11} gives the actual surjective homology
collapse with kernel the saturation of
$(\Lambda^2M_0-I)\Lambda^2\Lambda$.  Its dual has image
$\ker(\Lambda^2M_0^t-I)$ under the evaluation pairing.  Consequently
the natural specialization morphism identifies
\[
 R^2g_*\mathbb Q=j_*\Lambda^2V.
\]
The term $H^0(j_*\Lambda^2V)$ has no incoming Leray differential,
and its only possible outgoing differential targets
$H^2(j_*V)=H^0(V^*)^*=0$.  It survives to $H^2(M;\mathbb Q)=0$.
Thus $a_2':=\dim(\Lambda^2V)^\Gamma=0$.

\item Let $\beta:V\to V^*$ be an invariant alternating form.
Since $\beta(I)\subset(V^*)^\Gamma=0$, its radical contains $I$.
For the finite averaging projection $P_i:V\to F_i$, invariance gives
$\beta(u,x)=\beta(u,P_ix)$ when $u\in F_i$.
The alternating restriction to the two-dimensional space $F_i$
has $I$ in its radical and is therefore zero.  Hence
$F_1+F_2\subset\operatorname{rad}\beta$.  This sum has dimension
three because $F_1\cap F_2=I$.  The alternating rank is even,
so every such form has rank at most two.

Suppose $\beta,\gamma$ were independent invariant forms.
They and their sum have rank two, so
$\beta\wedge\gamma=0$.  Their actual image planes in $V^*$
are distinct invariant two-planes meeting in a rational line.
Each finite generator acts on that line by a sign.  At least one
finite order is odd, and the unipotent product acts trivially on the
line; both signs are therefore $+1$.  This contradicts
$(V^*)^\Gamma=0$.  It follows that
$b_2':=\dim(\Lambda^2V^*)^\Gamma\leq1$.

\item At the cusp, the exterior square of three size-two unipotent
blocks has fixed rank nine: three fixed vectors come from their
individual exterior squares and two from each of the three pairwise
tensor products.  If $d_i$ denotes a finite fixed rank on
$\Lambda^2V$, the invariant-sheaf Euler identity gives
\[
 -15+d_1+d_2+9
 =a_2'-h^1(j_*\Lambda^2V)+b_2'\leq1.
 \tag{F.40}\label{maps:eq:F.40}
\]
Here the degree-two group is dual to
$(\Lambda^2V^*)^\Gamma$, by compact-support duality and the
extension-by-zero sequence.

\item The two-dimensional unit eigenspace leaves four dimensions
for the nonunit cyclotomic factors.  Local complex compatibility
requires even multiplicity of $-1$.  Partitioning four into the
cyclotomic degrees gives the following complete list:
\[
\begin{array}{c|l}
3&\Phi_3^2\\
4&\Phi_4^2,\ \Phi_2^2\Phi_4\\
5&\Phi_5\\
6&\Phi_6^2,\ \Phi_2^2\Phi_3,\ \Phi_2^2\Phi_6,\ \Phi_3\Phi_6\\
8&\Phi_8\\
10&\Phi_{10}\\
12&\Phi_{12},\ \Phi_3\Phi_4,\ \Phi_4\Phi_6.
\end{array}
\]
The unit eigenspace contributes one fixed bivector.  A repeated
nonreal quadratic pair contributes four further fixed bivectors,
so $d_i=5$ for $\Phi_3^2,\Phi_4^2,\Phi_6^2$.
Every other listed factor contributes two, so $d_i=3$.
Inequality \eqref{maps:eq:F.40} excludes every occurrence of
$d_i=5$.  In particular all order-three pairs disappear.
The coprime order list of \cref{maps:src:F.2.8} therefore reduces
to $(5,4),(5,6),(5,8),(5,12)$; their remaining factors are exactly
the eight entries of \eqref{maps:eq:F.39}.
Each pair contains the order-five factor, so
\cref{maps:src:F.2.9} proves irreducibility of $V/I$.
\end{enumerate}
\end{proof}




\begin{proposition}[automatic local complex compatibility]\label{maps:src:F.2.15}
\StatementRecord{maps:src:F.2.15}{proved}{Historical events: Claude Fable 5.2 (Anthropic), effort high, 16 November 2026: carried as background without fault found. GPT-6 Astra (OpenAI), ChatGPT Pro mode, effort Pro, S6\_09\_PAPERS, 17 November 2026: targeted proof check. Also checked in this writing. Pre-continuation writing events: GPT-6 Astra (OpenAI), ChatGPT Pro mode, effort Pro, S6\_10\_PAPERS, 18 November 2026 02:27:56.235122+00:00: targeted proof check. Continuation mathematical and source checks: GPT-6 Astra (OpenAI), ChatGPT Pro mode, effort Pro, S6\_10\_PAPERS continuation, 19 November 2026: independent re-proof; used and rechecked in this run. Continuation metadata events: GPT-6 Astra (OpenAI), ChatGPT Pro mode, effort Pro, S6\_10\_PAPERS continuation, 19 November 2026: used and rechecked in this run; status-vocabulary reconciliation; metadata reconciliation only; no new mathematical check. \emph{Pre-continuation recorded proof scope.} Claude Fable 5.2 (Anthropic), effort high, 16 November 2026; also checked in this writing by GPT-6 Astra (OpenAI), ChatGPT Pro mode, effort Pro, 18 November 2026}{the hypotheses, notation and definitions stated here; \cref{maps:src:F.2.7}; \cref{maps:src:F.1.4}}


Every simply connected datum of F.2 with two finite affine
fillings automatically has finite linear actions commuting
with real complex structures.  This assertion concerns the
individual finite actions and does not produce simultaneous
global periods.

\end{proposition}


\begin{proof}
\begin{enumerate}[label=\arabic*.]
\item Let $m_i$ be the order of the free cyclic affine action
defining the $i$th filling.  Its linear order divides $m_i$.
By \ref{maps:src:F.2.7} the integers $m_1,m_2$ are coprime, so at least one
of them is odd.  The corresponding linear determinant is
one, because a finite real matrix of odd order cannot have
determinant minus one.  The relation
$M_1M_2M_0=I$ and the unipotence of $M_0$ imply
$\det M_1\det M_2=1$, so both determinants are one.

\item A finite-order real matrix is semisimple.  Its nonreal
eigenvalues occur in conjugate pairs, and its only possible
real eigenvalues are $1$ and $-1$.  Determinant one says
that the $-1$ eigenspace has even real dimension.  Since the
total dimension is six, the $1$ eigenspace also has even
real dimension.  Choose a complex structure on each of
these real eigenspaces by pairing a real basis.  On each
nonreal isotypical real summand, multiplication by $i$ in
its complex eigenspace supplies a commuting complex
structure.  Their direct sum commutes with $M_i$.
The fixed eigenspace is nonzero because the affine action
is free, as in \ref{maps:src:F.1.4}, so its even dimension is at least two.

\item Each chosen structure makes $\mathbb R^6/\mathbb Z^6$ a
complex torus and the given affine map holomorphic; thus it
supplies precisely the local complex-compatibility
hypothesis used in the subsequent reductions.  For three genuine
finite fillings \ref{maps:src:F.2.7} has already excluded simple
connectivity without this argument.

\end{enumerate}
\end{proof}


Theorem \ref{maps:src:F.2.10} supplies an independent exterior-square
proof for three of the excluded order-three pairs and for
the repeated fourth- and sixth-root cases.
Within the third row of the order-six spectral possibilities,
$\Phi_3\Phi_6$, no exclusion follows from F.2.13: that theorem
treats particular rational representations with the first
two order-six spectra.  Likewise \ref{maps:src:F.2.12} excludes one
rational representation of order pair $(5,12)$, not its
entire spectral locus.

\begin{theorem}[an unbounded integral family]\label{maps:src:F.2.16}
\StatementRecord{maps:src:F.2.16}{proved}{Historical events: Claude Fable 5.2 (Anthropic), effort high, 16 November 2026: carried as background without fault found. GPT-6 Astra (OpenAI), ChatGPT Pro mode, effort Pro, S6\_09\_PAPERS, 17 November 2026: targeted proof check. Also checked in this writing. Pre-continuation writing events: GPT-6 Astra (OpenAI), ChatGPT Pro mode, effort Pro, S6\_10\_PAPERS, 18 November 2026 02:27:56.235122+00:00: computation or formal certificate; targeted proof check. Continuation mathematical and source checks: GPT-6 Astra (OpenAI), ChatGPT Pro mode, effort Pro, S6\_10\_PAPERS continuation, 19 November 2026: computation or formal certificate; targeted proof check; used and rechecked in this run. Continuation metadata events: GPT-6 Astra (OpenAI), ChatGPT Pro mode, effort Pro, S6\_10\_PAPERS continuation, 19 November 2026: used and rechecked in this run; status-vocabulary reconciliation; metadata reconciliation only; no new mathematical check. \emph{Continuation recorded proof scope.} Every step of the current seven-step proof was written before its target was read: the full family derivation by the Paper VI reviewer, and the additional hatted integral-lattice and cusp-unit derivation by the root reviewer. The actual receiving proof combines exactly those frozen arguments. Newly authored programs passed 551 family checks, 74 hatted-unit checks, and 94 further comparisons of the preserved former computations. The supporting source and dependency scope was renewed. The preserved former routes receive no retrospective prereading credit.}{the hypotheses, notation and definitions stated here; \cref{maps:eq:F.U1}; \cref{maps:eq:F.U3}; \cref{maps:eq:F.U4}; \cref{maps:eq:F.U5}; \cref{maps:eq:F.U6}; \cref{maps:eq:F.U7}; \cref{maps:src:F.2.5}}


There are infinitely many simultaneous rational conjugacy classes
of integral rank-six tuples of the linear type $(5,4)$ in F.2,
with primitive maximal square-zero cusp.  An explicit infinite
subfamily has free affine lifts of orders five and four.
Apart from the parameter zero, the rank-five irreducible quotient
in the family is not self-dual.  Thus no finite bound on rational
conjugacy classes is valid for the unrestricted input tuples.
This does not assert infinitely many three-connected survivors.

\end{theorem}


\begin{proof}
\begin{enumerate}[label=\arabic*.]
\item Let $v\in\mathbb Z$, put $q=v^2-v+1$, and let
$Ae_j=e_{j+1}$ on $\mathbb Q^5$, with indices modulo five. Define
\[
 X=(e_0,\ e_2+ve_4),\qquad
 Y=\begin{pmatrix}
 0&-q&-v^2&-1&v\\
 0&1-vq&-vq&v^2-v&q
 \end{pmatrix},\qquad N=XY.
 \tag{F.U1}\label{maps:eq:F.U1}
\]
Multiplying $Y$ by the two columns of $X$ gives $YX=0$, so
$N^2=X(YX)Y=0$. Consequently
\[
 C=(I+N)A,\qquad B=A^4(I-N)=C^{-1}.
\]
The cycle gives $A^5=I$. Exact multiplication and determinant
expansion give
\[
 N^2=0,\qquad A^5=C^4=I_5,\qquad
 \chi_C(x)=(x-1)(x+1)^2(x^2+1).
 \tag{F.U2}\label{maps:eq:F.U2}
\]
For the determinant calculation one may substitute
\[
 (xI-A)^{-1}=
 \frac{x^4I+x^3A+x^2A^2+xA^3+A^4}{x^5-1}
\]
in the rank-two determinant identity
\[
 \det(xI-C)=\det(xI-A)
       \det\bigl(I_2-YA(xI-A)^{-1}X\bigr).
\]
The identity $C^4=I$ is checked separately; it is not inferred
from the repeated factor $x+1$ of the characteristic polynomial.
The factor $x^2+1$ shows that $B$ and $C$ have exact order four.

\item Let $P$ have columns $\mathbf1/q,e_1,e_2,e_3,e_4$, and put
$A'=P^{-1}AP$, $B'=P^{-1}BP$. Since $q>0$ and $\det P=1/q$,
these are actual rational changes of coordinates. Define
\[
 s=v^2+v-1,\qquad r=-\frac{(v+1)(v^2-2)}2,
 \qquad
 \beta=(s,\ s-vr,\ v^2s+vr,\ r,\ -vs-r)
 \tag{F.U3}\label{maps:eq:F.U3}
\]
and
\[
 \widehat A=\begin{pmatrix}1&0\\0&A'\end{pmatrix},\qquad
 \widehat B=\begin{pmatrix}1&\beta\\0&B'\end{pmatrix}.
 \tag{F.U4}\label{maps:eq:F.U4}
\]

The actual lattice $P\mathbb Z^5$ has basis
$f_0=\mathbf1/q,f_1=e_1,\ldots,f_4=e_4$.
Its images under $A$ are $f_0,f_2,f_3,f_4$ and
$qf_0-f_1-f_2-f_3-f_4$, respectively; $A^{-1}=A^4$
therefore also preserves this lattice. The columns of $X$ are
$qf_0-f_1-f_2-f_3-f_4$ and $f_2+vf_4$, while
\[
 Y':=YP=\begin{pmatrix}
 -2&-q&-v^2&-1&v\\
 2(1-v)&1-vq&-vq&v^2-v&q
 \end{pmatrix}.
\]
Thus $X$ and $Y$ are integral maps between the specified lattices.
The already proved $YX=0$ shows that $I+N$ and its inverse $I-N$
preserve $P\mathbb Z^5$. Since $B=A^{-1}(I-N)$, both
$A',B'$ belong to $\operatorname{GL}_5(\mathbb Z)$.
For $r$, odd $v$ makes $v+1$ even, and even $v$ makes
$v^2-2$ even; hence $r$ and all entries of $\beta$ are integers.
The hatted matrices are therefore integral automorphisms.

The exact unit needed in the hatted lattice is
\[
 z=(-v^3+v-1,0,0,2,2(1-v^2))^t.
\]
Writing $a=-v^3+v-1$, the first coordinate of $Y'z$ is
$-2a-2+2v(1-v^2)=0$. Half its second coordinate is
$(1-v)a+(v^2-v)+q(1-v^2)=0$, since the first two terms
sum to $v^4-v^3+v-1$ and the last is its negative. Furthermore
\[
 \beta z=s(v^3-v-1)+2v^2r=1.
 \tag{F.U5}\label{maps:eq:F.U5}
\]
Indeed the two summands expand to
$v^5+v^4-2v^3-2v^2+1$ and
$-v^5-v^4+2v^3+2v^2$.
Since $AB=I-N$ and $NPz=X(Y'z)=0$, the actual six-dimensional
map satisfies
\[
 (\widehat A\widehat B-I)(0,z)^t=e_0.
\]
This establishes the unit in the displayed hatted coordinates.

The row identity
$\beta(I+B'+B'^2+B'^3)=0$ proves
$\widehat B^4=I$; also $\widehat A^5=I$. Their quotient actions
show that these orders are exact.

Use the specified lattice change
\[
 L=\frac14\begin{pmatrix}
 -2&3&-5&-1&5&-3\\
 0&2&-2&2&-2&2\\
 0&-4&0&0&4&-4\\
 0&0&0&-4&4&-4\\
 0&0&4&-4&4&-4\\
 0&0&0&-4&4&0
 \end{pmatrix},\qquad
 A_v=L^{-1}\widehat A L,\quad B_v=L^{-1}\widehat B L.
 \tag{F.U6}\label{maps:eq:F.U6}
\]
Here $\det L=-1/4$. Expanding these products gives entries in
$\mathbb Z[v]$ for both matrices. The complete expanded entries
and the exact identities
$LA_v=\widehat A L$, $LB_v=\widehat B L$ are supplied by the
polynomial certificate described in the last item. Thus the
defined matrices act on the stated $\mathbb Z^6$ for every integer
parameter. Their determinants are one and their inverses are
the integer powers $A_v^4,B_v^3$. Their characteristic polynomials are
\[
 (x-1)^2\Phi_5(x),\qquad (x-1)^2(x+1)^2(x^2+1),
\]
respectively.

\item Put $K=B_vA_v-I$. Multiplication gives $K^2=0$.
The following are actual third minors of $K$, in zero-based indices:
\[
 \begin{array}{c|c|c}
 \text{rows}&\text{columns}&\text{determinant}\\\hline
 (0,1,4)&(1,2,3)&v(v+1)\\
 (0,1,4)&(1,3,4)&v^2\\
 (0,3,4)&(1,2,3)&(v+1)^2.
 \end{array}
\]
Their integer combination $(v+1)^2-2v(v+1)+v^2$ equals one.
Hence $K$ has rank at least three, while $K^2=0$ in dimension six
gives rank at most three. The gcd of its third minors is one,
so its three nonzero Smith factors are all one. This proves
primitivity directly in the specified lattice, for every $v$.
For $M_1=A_v^t$ and $M_2=B_v^t$ the actual cusp map is
\[
 M_0=(M_1M_2)^{-1}=(I+K^t)^{-1}=I-K^t.
\]
Thus $M_0-I=-K^t$, which has the same integral Smith factors
and the required maximal square-zero type.

\item Both $A_v,B_v$ fix $e_0$. The actual quotient morphism to
the original rank-five pair is
\[
 R=P\,\operatorname{pr}_5 L:\mathbb Q^6\longrightarrow\mathbb Q^5,
 \qquad RA_v=AR,\quad RB_v=BR.
\]
It is surjective and has kernel $\mathbb Qe_0$, since
$Le_0=-e_0/2$. We prove that this quotient is irreducible over
$\mathbb Q$. The coprime factors $x-1$ and $\Phi_5(x)$ of
$\chi_A$ give the projector
$(I+A+A^2+A^3+A^4)/5$. Its image is $\mathbb Q\mathbf1$ and
its kernel is $S=\{\sum x_i=0\}$. Irreducibility of $\Phi_5$
shows that these are the only nonzero proper rational
$A$-invariant subspaces.

Neither is $C$-invariant. Indeed,
\[
 N\mathbf1=-2qX\binom{1}{v-1}
\]
has nonzero first entry and zero entries in positions one and three,
so $C\mathbf1=\mathbf1+N\mathbf1$ is not a multiple of
$\mathbf1$. Write $\epsilon=(1,1,1,1,1)$. If $S$ were
$C$-invariant, then $\epsilon N$ would be a multiple of
$\epsilon$; its first entry is zero, so that multiple would
vanish. But its entries in columns three and four are
$v^3-v-1$ and $v^3+v+1$. Their simultaneous vanishing would
give $v=-1$, at which the first is $-1$. This is impossible.
Since $B=C^{-1}$, the pair $(A,B)$ is irreducible, as claimed.

\item The rank-two update gives exact traces
\[
 \operatorname{tr}(AC)=\operatorname{tr}(YA^2X)=1-v^2q,
 \qquad
 \operatorname{tr}((AC)^{-1})
       =-\operatorname{tr}(YA^3X)=1+v^2q.
\]
Here $\operatorname{tr}(A^2)=\operatorname{tr}(A^3)=0$ and
$B=A^4(I-N)$ give the displayed equalities. An equivariant
isomorphism $F$ to the contragredient representation would
satisfy $FA=A^{-t}F$ and $FB=B^{-t}F$, and therefore identify
the character of the word $AC=AB^{-1}$ with the character
of its contragredient matrix. The two displayed traces would
be equal. Their difference $-2v^2q$ is nonzero when $v\ne0$,
proving the required nonselfduality.

For completeness, at $v=0$ the symmetric matrix
$Q_{ij}=1$ for $j-i\equiv1,4\pmod5$ and $Q_{ij}=0$ otherwise
has determinant two and satisfies $A^tQA=Q$, $B_0^tQB_0=Q$.
It gives the actual equivariant isomorphism in that case.

For the rank-six matrices, the same word has trace
\[
 h(v):=\operatorname{tr}(A_vB_v^{-1})=2-v^2q,
 \qquad h(v+1)-h(v)=-4v^3-3v^2-3v-1.
\]
Thus $h$ is strictly decreasing on the nonnegative integers,
in particular on the nonnegative multiples of twenty.
A simultaneous rational conjugacy of the ordered matrices
would preserve this trace. There are consequently infinitely
many such rational conjugacy classes already in that subfamily.

\item For $v\in20\mathbb Z$, put
$a=(1,3,-2,-1,2,0)^t$, $e=(1,0,0,0,0,0)^t$, and define
\[
 w_1=\frac15\sum_{j=0}^4(A_v^t)^j a,
 \qquad
 w_2=-\frac14\sum_{j=0}^3(B_v^t)^j e.
 \tag{F.U7}\label{maps:eq:F.U7}
\]
The polynomial matrices reduce to their parameter-zero
values modulo five and four, respectively. At zero,
$M_1a=a$ and $M_2e=e$. The two sums are therefore divisible
by their indicated denominators, so $w_1,w_2$ are integral.
Their defining averages give $M_iw_i=w_i$. The common
first-coordinate covector $\ell$ is fixed by both $M_i$ and
satisfies $\ell(w_1)=1$, $\ell(w_2)=-1$.

For $(m_1,m_2)=(5,4)$ let
$g_i(x)=M_ix+w_i/m_i$ on $\mathbb R^6/\mathbb Z^6$.
Then $g_i^j(x)=M_i^jx+jw_i/m_i$, so $g_i^{m_i}$ is the
identity. Every power with $0<j<m_i$ changes the actual
$\ell$-circle coordinate by $j/5$ or $-j/4$, respectively.
It has no fixed point. These are free actions of the exact
orders five and four, with the specified affine lifts.

\item Substitution of $v=0$ in the displayed construction gives
exactly the marked pair in Theorem~\ref{maps:src:F.2.5}, and
the averages give $w_1(0)=a$, $w_2(0)=-e$. This identifies
both the matrices and the affine numerators without a further
change of basis.

The independent certificate
\texttt{F2\_16\_family\_fresh.py} records the complete expanded
polynomial matrices and checks every matrix, quotient, trace
and affine identity above over $\mathbb Q[v]$, with coefficient
integrality in $\mathbb Z[v]$. It checks all four hundred
third minors as well as the displayed integer combination.
Its separate Smith computations at ten stated integer
parameters are additional finite checks; the universal
primitivity assertion follows from the displayed identity
equal to one. The conclusions concern the unrestricted
input tuples in the statement.
\end{enumerate}
\end{proof}




\begin{proposition}[an odd torsion obstruction in this lattice]\label{maps:src:F.2.17}
\StatementRecord{maps:src:F.2.17}{proved}{Historical events: Claude Fable 5.2 (Anthropic), effort high, 16 November 2026: carried as background without fault found. GPT-6 Astra (OpenAI), ChatGPT Pro mode, effort Pro, S6\_09\_PAPERS, 17 November 2026: targeted proof check. Also checked in this writing. Pre-continuation writing events: GPT-6 Astra (OpenAI), ChatGPT Pro mode, effort Pro, S6\_10\_PAPERS, 18 November 2026 02:27:56.235122+00:00: computation or formal certificate; targeted proof check. Continuation metadata events: GPT-6 Astra (OpenAI), ChatGPT Pro mode, effort Pro, S6\_10\_PAPERS continuation, 19 November 2026: used and rechecked in this run; status-vocabulary reconciliation; metadata reconciliation only; no new mathematical check. \emph{Pre-continuation recorded proof scope.} Claude Fable 5.2 (Anthropic), effort high, 16 November 2026; also checked in this writing by GPT-6 Astra (OpenAI), ChatGPT Pro mode, effort Pro, 18 November 2026}{the hypotheses, notation and definitions stated here; \cref{maps:eq:F.U6}; \cref{maps:src:F.2.19}; \cref{maps:eq:F.U5}; \cref{maps:src:F.2.12}}


For the integral tuple $(A_v,B_v)$ in \eqref{maps:eq:F.U6}, any simply
connected filling with free finite affine caps and the fixed
primitive Kuhn cusp has a subgroup
\[
 \mathbb Z/(v^2-v+1)\ \subset\operatorname{Tor}H_2(M;\mathbb Z).
 \tag{F.U9}\label{maps:eq:F.U9}
\]
In particular, in this specified family of lattices only
$v=0$ and $v=1$ can be three-connected, regardless of the
affine multiplicities and twists.  The stronger exclusion for every stable lattice in each
rational class is proved in Theorem \ref{maps:src:F.2.19} below.

\end{proposition}


\begin{proof}
\begin{enumerate}[label=\arabic*.]
\item First compute the global coinvariants in the hatted lattice.
In the rank-five quotient, the $A'$ relations identify
$f_1=f_2=f_3=f_4=t$ and $qf_0=5t$.
The number $q=v^2-v+1$ is relatively prime to five, as its
five residue values show.  Thus these coinvariants are
infinite cyclic, with $f_0=5h$ and $t=qh$.
The $B'$ relations become $q\kappa_j h$, where
\[
 (\kappa_0,\ldots,\kappa_4)=
 (2v^2,\ v(v^3+v-1),\ v(v^3+v+1),
       -v^3+v+1,\ -v^3-v-1).
 \tag{F.U10}\label{maps:eq:F.U10}
\]
Their gcd is one: it divides $2v$ by subtracting the second
and third entries, and it divides $v^3-v-1$ by the fourth;
the latter is odd and is $-1$ modulo $v$.
Hence the rank-five global coinvariants are $\mathbb Z/q$.
The extension line is killed in rank-six coinvariants by
\eqref{maps:eq:F.U5}, since
$(\widehat A\widehat B-I)(0,z)^t=e$.
Consequently the hatted rank-six coinvariants are also
$\mathbb Z/q$.  The lattice change \eqref{maps:eq:F.U6} is an isomorphism
over every $\mathbb Z_{(p)}$ with $p$ odd, so the odd part of
the actual rank-six global coinvariants is precisely
$\mathbb Z/q$.

\item Let $V$ be the degree-one fibre cohomology local system.
The only difference between $R^1g_*\mathbb Z$ and $j_*V$
is a finite cokernel at finite caps; the cusp degree-one
specialization is onto.  Thus
\[
 H^2(\mathbb P^1,R^1g_*\mathbb Z)
 =H^2(\mathbb P^1,j_*V)=V_\Gamma,
\]
including its $q$-primary part just computed.  The space of
global degree-two invariants is zero: the irreducible odd
rank-five quotient has no invariant alternating tensor,
and the extension filtration
$0\to W\to\bigwedge^2V\to\bigwedge^2W\to0$
gives $(\bigwedge^2V)^\Gamma=0$.

\item It remains to check every possible source of the differential
$d_2:H^0(\mathbb P^1,R^2g_*\mathbb Z)
\to H^2(\mathbb P^1,R^1g_*\mathbb Z)$.
The cusp has no kernel in degree two.  At a finite cap the
kernel is torsion in $H^2$ of that cap, namely torsion in its
first homology.  Simple connectivity of the total filling
forces the distinguished affine numerator to be prime to the
actual affine multiplicity, by the determinant argument of
\ref{maps:src:F.2.12}.  The cap is therefore a mapping torus of a torus
automorphism whose order divides the linear order, five or
four.  Its first-homology torsion is supported only at the
primes five or two, respectively: averaging over the finite
order kills the torsion of the vertical coinvariants.
Hence the source of this $d_2$ has no torsion at any prime
dividing $q$, since $q$ is odd and prime to five.

\item The subgroup $\mathbb Z/q$ in bidegree $(2,1)$ survives.
There are no longer differentials on a two-dimensional base,
and this terminal filtration term is a subgroup of
$H^3(M;\mathbb Z)$.  Universal coefficients identify its
torsion with a subgroup of $\operatorname{Tor}H_2(M;\mathbb Z)$.
This proves \eqref{maps:eq:F.U9}.  Finally, $v^2-v+1=1$ for an integer $v$
exactly when $v=0$ or $v=1$.

\end{enumerate}
\end{proof}




\begin{theorem}[a survivor with a non-self-dual quotient]\label{maps:src:F.2.18}
\StatementRecord{maps:src:F.2.18}{proved}{Historical events: Claude Fable 5.2 (Anthropic), effort high, 16 November 2026: carried as background without fault found. GPT-6 Astra (OpenAI), ChatGPT Pro mode, effort Pro, S6\_09\_PAPERS, 17 November 2026: targeted proof check. Also checked in this writing. Pre-continuation writing events: GPT-6 Astra (OpenAI), ChatGPT Pro mode, effort Pro, S6\_10\_PAPERS, 18 November 2026 02:27:56.235122+00:00: computation or formal certificate; targeted proof check. Continuation metadata events: GPT-6 Astra (OpenAI), ChatGPT Pro mode, effort Pro, S6\_10\_PAPERS continuation, 19 November 2026: used and rechecked in this run; status-vocabulary reconciliation; metadata reconciliation only; no new mathematical check. \emph{Pre-continuation recorded proof scope.} Claude Fable 5.2 (Anthropic), effort high, 16 November 2026; also checked in this writing by GPT-6 Astra (OpenAI), ChatGPT Pro mode, effort Pro, 18 November 2026}{the hypotheses, notation and definitions stated here; \cref{maps:eq:F.U6}; \cref{maps:src:F.3.11}; \cref{maps:src:F.2.5}; \cref{maps:src:F.2.16}; \cref{maps:src:F.2.6}; \cref{maps:eq:F.6}; \cref{maps:src:F.4.1}}


The parameter $v=1$ in \eqref{maps:eq:F.U6}, with affine orders five and four
and twists
\[
 w_1=(1,3,-2,-1,2,0)^t,\qquad
 w_2=(-1,0,0,0,0,0)^t,
\]
gives a three-connected smooth eight-manifold with
$H_4=\mathbb Z^4$ and all other reduced homology zero except
$H_8=\mathbb Z$.  It is almost complex and is oriented
diffeomorphic to $3\mathbb{HP}^2\#(-\mathbb{HP}^2)$.
Its irreducible rank-five quotient has no invariant bilinear
form, so integral three-connectivity does not force that
quotient to be orthogonal.  It has bounded-cusp holomorphic
periods by Theorem \ref{maps:src:F.3.11}; simultaneous finite holomorphic
extensions are not asserted here.

\end{theorem}


\begin{proof}
\begin{enumerate}[label=\arabic*.]
\item The first cohomological matrix is the same $A$ as in
Theorem \ref{maps:src:F.2.5}.  The second is
\[
 B_1=\begin{pmatrix}
 1&0&0&0&0&0\\
 0&0&-1&0&1&-1\\
 0&0&-1&0&0&0\\
 0&0&-2&0&1&-2\\
 0&1&-2&0&1&-3\\
 0&0&-1&1&0&-1
 \end{pmatrix}.
 \tag{F.U11}\label{maps:eq:F.U11}
\]
Take $M_1=A^t$, $M_2=B_1^t$ and
$M_0=(M_1M_2)^{-1}$.  The displayed twists are fixed integral
vectors with $\ell(w_1)=1$, $\ell(w_2)=-1$.
They therefore give free affine actions.  The cusp is
primitive by \ref{maps:src:F.2.16}.  Exact Hermite reduction gives
\[
 \operatorname{col}_{\mathbb Z}
 [\,M_0-I\mid M_1(M_0-I)\mid M_2(M_0-I)\,]
       =\langle e_1,\ldots,e_5\rangle.
\]
Van Kampen leaves only $5a=h$, $4b=-h$, $a+b=0$,
whose relation matrix is unimodular.  This proves
$\pi_1=1$ for the specified gluing, including zero cusp
clutching.

\item For completeness we give the full new integral certificates
needed to reuse the already proved differential calculation
in \ref{maps:src:F.2.6}.  The rational exterior invariant and parabolic ranks
are exactly the table there.  The finite cap first homology
groups are $\mathbb Z^2$ and $\mathbb Z^2\oplus\mathbb Z/2$.
On integral invariant two-form bases their suspension periods
divided by the covering orders are, respectively,
$(-1,0,-1)$ and $(0,0,1)$.  The determinants of the invariant
degree-three skew pairings are $5^4$ and $4^4$, while the
degree-two/degree-four pairings have absolute determinants
$5^3$ and $4^3$.  Thus all the finite covering indices used
in degrees two, three and four are one.

\item For exterior coefficient degrees two, three and four, the
global coboundary Smith form has one factor two and all other
nonzero factors one.  Its torsion class has peripheral
triviality flags $(\mathrm{true},\mathrm{false},\mathrm{true})$,
so the parabolic groups are torsion free, of ranks $0,3,0$.
The degree-one coboundary is primitive and the degree-one
finite covering cokernels are killed by the common invariant
covector and the Chinese remainder theorem, exactly as in
\ref{maps:src:F.2.6}.  The global coinvariants in degrees one through four
are each $\mathbb Z/2$.

\item Write $M_2=\operatorname{diag}(1,D)$.  This time use the
unimodular marking
\[
 P_1=\begin{pmatrix}
 -1&1&-1&0&-1\\
 -1&2&-1&0&-1\\
 -1&0&0&2&-1\\
 1&-1&0&-1&1\\
 0&0&2&-1&1
 \end{pmatrix},\qquad \det P_1=-1.
\]
It satisfies $DP_1=P_1\operatorname{diag}(\text{four-cycle},-1)$.
The cap therefore has the same cellular mapping-torus model
used to prove \eqref{maps:eq:F.6}.  With the unchanged twist $w_2=-e_0$,
the boundary coordinates remain $\theta=a+4k$, $h=k$;
the suspension representative restricts to $w$ in the base
direction and to $-\operatorname{Norm}_4(w)=0$ in the
auxiliary fibre direction.  Consequently the coefficient in
the torsion differential is again one.

\item The new vertical representatives in lexicographic exterior
bases are
\[
\begin{split}
 q=2:\ &(0,-1,-1,-1,0),\\
 q=3:\ &(0,0,1,0,0,0,0,0,0,1),\\
 q=4:\ &(0,0,0,0,0,0,0,-1,-1,0).
\end{split}
\]
Each has residue one in its global coinvariant
$\mathbb Z/2$.  Thus all three torsion differentials in \eqref{maps:eq:F.6}
are isomorphisms.  The integral Leray table and its free
degree-one differential are therefore exactly those of
\ref{maps:src:F.2.6}.  They give $H^2=H^3=0$ and $H^4=\mathbb Z^4$.
Universal coefficients, Hurewicz and Poincare duality give
the stated integral homology and three-connectivity.

\item The relative almost-complex extension over the pair of pants
in \ref{maps:src:F.2.5} applies: the finite matrices have the same even
real eigenspace dimensions and the cusp is the same marked
Kuhn type.  The purely smooth characteristic-vector and
normal-invariant argument in \ref{maps:src:F.4.1} then gives the stated
oriented diffeomorphism type.  No integrability at the finite
punctures is used for this smooth classification.

\item Finally, Theorem~\ref{maps:src:F.2.16} gives nonselfduality of the irreducible rank-five quotient at $v=1$.
The exact program
\texttt{code/fourfolds\_v1\_homology.py} runs every matrix,
Smith, Hermite, suspension-period, pairing and differential
certificate just displayed.  Its monomial marking was found
inside the stated finite cube and then verified directly;
the proof uses that explicit marking, not search completeness.

\end{enumerate}
\end{proof}




\begin{theorem}[a residual obstruction for every lattice]\label{maps:src:F.2.19}
\StatementRecord{maps:src:F.2.19}{proved}{Historical events: Claude Fable 5.2 (Anthropic), effort high, 16 November 2026: carried as background without fault found. GPT-6 Astra (OpenAI), ChatGPT Pro mode, effort Pro, S6\_09\_PAPERS, 17 November 2026: targeted proof check. Also checked in this writing. Pre-continuation writing events: GPT-6 Astra (OpenAI), ChatGPT Pro mode, effort Pro, S6\_10\_PAPERS, 18 November 2026 02:27:56.235122+00:00: computation or formal certificate; targeted proof check. Continuation metadata events: GPT-6 Astra (OpenAI), ChatGPT Pro mode, effort Pro, S6\_10\_PAPERS continuation, 19 November 2026: used and rechecked in this run; status-vocabulary reconciliation; metadata reconciliation only; no new mathematical check. \emph{Pre-continuation recorded proof scope.} Claude Fable 5.2 (Anthropic), effort high, 16 November 2026; also checked in this writing by GPT-6 Astra (OpenAI), ChatGPT Pro mode, effort Pro, 18 November 2026}{the hypotheses, notation and definitions stated here; \cref{maps:src:F.2.16}; \cref{maps:eq:F.U6}; \cref{maps:src:F.2.6}; \cref{maps:src:F.2.18}; \cref{maps:eq:F.U12}; \cref{maps:src:F.2.17}; \cref{maps:eq:F.U1}}


Let a rational rank-six datum have finite linear orders five
and four, a primitive maximal cusp on a stable integral lattice,
and invariant line $L$, with irreducible rank-five quotient $W$
and no rational global coinvariants.  Suppose that at a prime
$p\ne2,5$ the semisimplified reduction of a stable lattice in
$W$ is the sum of the trivial module and an absolutely
irreducible module of dimension four.  Then no stable integral
lattice in that rational datum, with primitive maximal cusp,
and no free finite affine twists can give
$\pi_1=1$ and $H_2=0$ simultaneously.

For every $v\notin\{0,1\}$, the rational family in \ref{maps:src:F.2.16}
satisfies this hypothesis at each prime dividing
$v^2-v+1$.  Consequently the exclusion of those parameters
holds for \emph{all} stable lattices and all twists, not just
the lattice in \eqref{maps:eq:F.U6}.  The two remaining parameters are the
verified survivors of \ref{maps:src:F.2.6} and \ref{maps:src:F.2.18}.

\end{theorem}


\begin{proof}
\begin{enumerate}[label=\arabic*.]
\item We first give the lattice-independent mechanism.  For any
stable lattice $V_\mathbb Z$, use its primitive invariant
line $L_\mathbb Z$ and quotient lattice $W_\mathbb Z$.
The semisimplified reduction of a rational representation is
independent of its stable lattice.  One elementary way to see
this is to compare two commensurable lattices by their finite
quotient $T$.  The exact sequence obtained by reduction modulo
$p$ has end terms $T[p]$ and $T/pT$, which have the same
composition factors by the filtration of $T$ by powers of $p$.
Hence the two lattice reductions have equal classes in the
Grothendieck group of finite-length modules.  Chaining such
comparisons proves the assertion for any two lattices.

\item Thus $W_\mathbb Z/pW_\mathbb Z$ has length two, with one
trivial constituent and one irreducible four-dimensional
constituent.  It has either a trivial submodule or a trivial
quotient (and possibly both).

\item Suppose first that it has a trivial submodule.  Its inverse
image in $V_\mathbb Z/pV_\mathbb Z$ is an extension of a
trivial module by a trivial module.  Each finite generator
acts on this two-dimensional space by a unipotent matrix.
Since its order is five or four, prime to $p$, it acts as
the identity.  There are therefore at least two global
invariant vectors in degree-one fibre cohomology modulo $p$.
Write their dimension as $a_p\ge2$; this is also the
dimension of the global coinvariants of fibre homology over
$\mathbb F_p$, by the perfect pairing between fibre homology
and cohomology.

\item The abelianized van Kampen presentation now shows
\[
 b_1(M;\mathbb F_p)\ge a_p-1.
 \tag{F.U12}\label{maps:eq:F.U12}
\]
Indeed, after the monodromy relations the fibre contributes
$a_p$ dimensions, and the two finite meridians contribute
two more.  The two finite power relations and the single
base relation impose at most three linear relations.
The cusp introduces no further fibre relation in this
abelianization.  On the dual homology lattice its nilpotent
operator is still primitive and maximal, so its angular lattice is
$\ker(M_0-I)=\operatorname{im}(M_0-I)$ by primitivity, and
the latter image already vanishes among the global
monodromy relations.  Its base meridian is the base relation
just counted.  Changing affine twists changes coefficients
of these three relations, not their number.  Thus \eqref{maps:eq:F.U12}
contradicts simple connectivity.

\item Suppose instead that the reduction of $W$ has a trivial
quotient.  Its integral global coinvariants have nonzero
$p$-primary torsion, since they have rational rank zero.
The natural map from the global coinvariants of $V$ to
those of $W$ is an isomorphism.  To check its kernel, a
primitive generator $e$ of $L_\mathbb Z$ lies in
$\ker(M_0-I)=\operatorname{im}(M_0-I)$, and hence is already
zero among global coinvariants.  Therefore $H^2(R^1)$ has
nonzero $p$-primary torsion.  If the total space were simply
connected, the finite-cap argument of \ref{maps:src:F.2.17} would make
every possible torsion source in $H^0(R^2)$ supported only
at two and five.  The invariant part of $\bigwedge^2V$ is
zero by irreducibility of the odd-dimensional quotient.
Consequently the nonzero $p$-primary subgroup survives to
$H^3(M;\mathbb Z)$ and then to
$\operatorname{Tor}H_2(M;\mathbb Z)$, again a contradiction.
This proves the general criterion.

\item For its application to \ref{maps:src:F.2.16}, let $p\mid q=v^2-v+1$.
We have $p\ne2,5$.  In the original cycle lattice of \eqref{maps:eq:F.U1},
$Y\mathbf1=2q(-1,1-v)^t$, so $\mathbf1$ gives a trivial
submodule after reduction.  On the four-dimensional quotient
the first finite generator has characteristic polynomial
$\Phi_5$, and $C=B^{-1}$ has polynomial
$(x+1)^2(x^2+1)$.  This quotient is absolutely irreducible.
Here is a characteristic-independent proof over
$\overline{\mathbb F}_p$.

\item A one-dimensional composition factor would make the
eigenvalues of $A$ and $C$ equal, because $CA^{-1}$ is
unipotent.  An eigenvalue of order dividing both five and
four must be one, which is absent from $\Phi_5$.
If the quotient were reducible, its only remaining possible
composition-factor dimensions would be two and two.
On each factor, the determinants of $A$ and $C$ must agree
and have orders dividing five and four, so both determinants
are one.  The two $C$ eigenvalue pairs must therefore be
$(-1,-1)$ and $(i,-i)$ in some order.  On the former
factor $C=-I$, since its order is prime to $p$.
Then the unipotence of $CA^{-1}=-A^{-1}$ would force the
eigenvalues of $A$ to be $-1$, impossible for $\Phi_5$.
This proves absolute irreducibility and establishes the
required residual composition factors at every $p\mid q$.
When $v\notin\{0,1\}$, the integer $q$ is greater than one
and has such a prime divisor.  The criterion applies.

\end{enumerate}
\end{proof}




\begin{theorem}[the complete rational $(5,4)$ normal form]\label{maps:src:F.2.20}
\StatementRecord{maps:src:F.2.20}{proved}{Historical events: Claude Fable 5.2 (Anthropic), effort high, 16 November 2026: carried as background without fault found. GPT-6 Astra (OpenAI), ChatGPT Pro mode, effort Pro, S6\_09\_PAPERS, 17 November 2026: targeted proof check. Also checked in this writing. Pre-continuation writing events: GPT-6 Astra (OpenAI), ChatGPT Pro mode, effort Pro, S6\_10\_PAPERS, 18 November 2026 02:27:56.235122+00:00: computation or formal certificate; targeted proof check. Continuation metadata events: GPT-6 Astra (OpenAI), ChatGPT Pro mode, effort Pro, S6\_10\_PAPERS continuation, 19 November 2026: used and rechecked in this run; status-vocabulary reconciliation; metadata reconciliation only; no new mathematical check. \emph{Pre-continuation recorded proof scope.} Claude Fable 5.2 (Anthropic), effort high, 16 November 2026; also checked in this writing by GPT-6 Astra (OpenAI), ChatGPT Pro mode, effort Pro, 18 November 2026}{the hypotheses, notation and definitions stated here; \cref{maps:src:F.2.11}; \cref{maps:eq:F.R4}; \cref{maps:eq:F.R6}; \cref{maps:eq:F.R7}; \cref{maps:eq:F.R8}}


Consider ordered rational rank-five pairs $(A,C)$ with
\[
 \chi_A=(x-1)\Phi_5,
 \qquad \chi_C=(x-1)(x+1)^2(x^2+1).
\]
Put
\[
 N=CA^{-1}-I,
 \qquad N^2=0,\quad\operatorname{rank}N=2,
\]
where $A^5=C^4=I$ and the pair is irreducible over $\mathbb Q$.
Their simultaneous rational conjugacy classes are parametrized
bijectively by
\[
 \bigl\{[r:s:t]\in\mathbb P^2(\mathbb Q):D(r,s,t)\ne0\bigr\}
 \ \sqcup\ \mathbb A^1(\mathbb Q),
 \tag{F.R1}\label{maps:eq:F.R1}
\]
where
\[
\begin{split}
 D={}&r^3-4r^2s-r^2t+7rs^2+rst-3rt^2\\
     &{}-7s^3+11s^2t-11st^2+7t^3.
\end{split}
 \tag{F.R2}\label{maps:eq:F.R2}
\]
The second summand lies over the single exceptional point
$[4:1:3]$ of $D=0$.  The construction below is explicit and
uses only a ten-variable linear system.  By \ref{maps:src:F.2.11} it also
parametrizes the rational rank-six maximal-cusp extensions
of all these pairs, with $B=C^{-1}$.
This is a rational classification; the stable integral
lattices, affine twists and surviving homology still require
the arithmetic conditions of F.2.

\end{theorem}


\begin{proof}
\begin{enumerate}[label=\arabic*.]
\item Write $K=\mathbb Q(\zeta)$, $\zeta^4+\zeta^3+\zeta^2+
\zeta+1=0$, and identify the rational space with
$\mathbb Q\oplus K$.  The matrix of $A$ is multiplication by
$(1,\zeta)$:
\[
 A=\operatorname{diag}(1,J),\qquad
 J=\begin{pmatrix}
 0&0&0&-1\\1&0&0&-1\\0&1&0&-1\\0&0&1&-1
 \end{pmatrix},\qquad R=(A+I)^{-1}.
 \tag{F.R3}\label{maps:eq:F.R3}
\]
Its rational centralizer is $\mathbb Q^*\times K^*$.
Let $I_N=\operatorname{im}N$.  If $X$ is any basis matrix
for this plane and $N=XY$, then $YX=0$.
The identity
\[
 C+I=(A+I)+XYA
\]
shows that its two-dimensional kernel has the same dimension
as the kernel of
$I_2+YARX=I_2-YRX$.  Hence
\[
 YRX=I_2,
 \qquad I_N\cap RI_N=0.
 \tag{F.R4}\label{maps:eq:F.R4}
\]
Conversely the two matrix equations in \eqref{maps:eq:F.R4}, together with
$YX=0$, give precisely that two-dimensional $-1$ eigenspace.

\item The plane $I_N$ projects nontrivially to $\mathbb Q$:
otherwise the first-coordinate covector is invariant under
both $A$ and $C$, contrary to irreducibility.
It contains no nonzero vector of $\mathbb Q\oplus0$,
because $R$ acts as $1/2$ on that line and \eqref{maps:eq:F.R4} would fail.
Write a basis of the plane as $(0,u),(1,v)$, where
$u\in K^*$.  Multiplication on $K$ by $u^{-1}$ puts it into
\[
 I_N=\langle(0,1),(1,\beta)\rangle_{\mathbb Q}.
\]
Changing the second basis vector by a rational multiple of
the first changes $\beta$ by a rational constant.
Two such planes are in the same centralizer orbit exactly
when $\beta'=a\beta+b$, with $a\in\mathbb Q^*$ and
$b\in\mathbb Q$.  Indeed a centralizer element carrying
one to the other must preserve the kernel of the projection
to $\mathbb Q$, so its $K^*$ component is itself rational;
the displayed affine change follows, and every such change
is realized.  We may therefore take
\[
 \beta=r\zeta+s\zeta^2+t\zeta^3,
 \qquad [r:s:t]\in\mathbb P^2(\mathbb Q),
 \qquad
 X=\begin{pmatrix}0&1\\1&0\\0&r\\0&s\\0&t\end{pmatrix}.
 \tag{F.R5}\label{maps:eq:F.R5}
\]
The triple is nonzero; a rational $\beta$ would make the
plane contain the forbidden first-coordinate line.
The stabilizer of such a plane in the centralizer consists
only of common scalar matrices.  Thus it introduces no
further identifications of the resulting nilpotent operators.

\item Write $Y$ as an unknown two-by-five matrix.  It is determined
by the following ten linear equations:
\[
 YX=0,\qquad YRX=I_2,\qquad
 \operatorname{tr}(YAX)=-1,\qquad
 \operatorname{tr}(YA^{-1}X)=1.
 \tag{F.R6}\label{maps:eq:F.R6}
\]
The first eight are \eqref{maps:eq:F.R4}.  The last two are exactly
$\operatorname{tr}C=\operatorname{tr}C^{-1}=-1$, since
$\operatorname{tr}A=\operatorname{tr}A^{-1}=0$ and
$C^{-1}=A^{-1}(I-XY)$.
They suffice to obtain the required characteristic
polynomial: the two $-1$ roots and determinant one leave a
monic cubic of constant term $-1$, and the two traces fix
that cubic to $(x-1)(x^2+1)$.
These remaining three roots are distinct, and the $-1$
eigenspace already has dimension two, so $C$ is semisimple
and $C^4=I$.  Also $YRX=I_2$ ensures that $N$ has rank two.
Thus every solution of \eqref{maps:eq:F.R6} gives all the asserted local
linear conditions.

\item In row-major order on the ten entries of $Y$, the determinant
of the coefficient matrix in \eqref{maps:eq:F.R6} is $D/4$.
For $D\ne0$ this proves uniqueness and gives the explicit
formula $Y=\operatorname{reshape}(E^{-1}b)$, where $E$ is
the coefficient matrix of the displayed linear equations
and $b$ their displayed right side.
The exact expanded matrix and all ten rational functions
are printed by \texttt{fourfolds\_relative\_normalform.py}.
There is no search or degree cutoff in this calculation.

\item We check the entire exceptional locus, rather than discard it.
Set
\[
 Q_1=r^2-rs-rt+s^2+t^2,
 \qquad Q_2=r^2-2rs-3rt+4s^2-4st+4t^2.
\]
The first column of the unique generic solution is
\[
 y_{10}=\frac{2Q_1Q_2}{D},\qquad
 y_{20}=\frac{2(r-s-t)Q_1}{D}.
 \tag{F.R7}\label{maps:eq:F.R7}
\]
Here $y_{10},y_{20}$ are the first column of $Y$.
The two equalities after multiplication by $D$ are polynomial
linear combinations of \eqref{maps:eq:F.R6}, by its adjugate matrix.
The program prints the corresponding two polynomial rows
and verifies those identities, so they remain necessary
also when $D=0$.
Moreover
\[
 2Q_1=(r-s)^2+(r-t)^2+s^2+t^2>0
\]
for every nonzero rational triple.  Consequently consistency
on $D=0$ forces $r=s+t$ and $Q_2=0$.
Substitution gives
\[
 D(s+t,s,t)=(t-3s)(2t-s)^2,
 \qquad Q_2(s+t,s,t)=(3s-t)(s-2t).
 \tag{F.R8}\label{maps:eq:F.R8}
\]
There are precisely two resulting projective points,
$[3:2:1]$ and $[4:1:3]$.
At the former, $[X\mid RX]$ has rank less than four, so
\eqref{maps:eq:F.R4} is impossible.  At the latter the coefficient and
augmented matrices of \eqref{maps:eq:F.R6} both have rank nine, and every
solution is
\[
 Y_z=\begin{pmatrix}
 4z+16/5&0&-2z-1&4/5-2z&2z\\
 2z-2/5&0&-z&2/5-z&z
 \end{pmatrix},\qquad z\in\mathbb Q.
 \tag{F.R9}\label{maps:eq:F.R9}
\]
Direct substitution proves both necessity and sufficiency.
This establishes the exceptional affine line in \eqref{maps:eq:F.R1}.

\item Every pair just constructed is irreducible over $\mathbb Q$.
An $A$-stable proper rational subspace has dimension one or
four.  A one-dimensional subrepresentation or quotient must
have trivial character, since its finite eigenvalues of
orders dividing five and four must agree under the
unipotent product.  The only $A$-invariant covector is the
first-coordinate covector.  Its product with $N$ is the
second row of $Y$, which is nonzero by $YRX=I_2$.
The only $A$-invariant vector is the first-coordinate vector;
its image under $N$ vanishes exactly if both entries in
\eqref{maps:eq:F.R7} vanish.  For $D\ne0$, positivity of $Q_1$ and
\eqref{maps:eq:F.R8} rule this out.  On the exceptional line the two
entries $4z+16/5$ and $2z-2/5$ never vanish simultaneously.
There are therefore neither invariant vectors nor covectors,
which excludes both possible dimensions.

\item Finally, rescaling $(r,s,t)$ is exactly the centralizer
change already described; it gives a conjugate pair.
Conversely any conjugacy of two normalized pairs centralizes
$A$ and hence identifies their projective parameters as proved
above.  At the exceptional point its stabilizer is scalar,
so distinct $z$ give distinct pairs.  This proves the stated
bijection, including completeness and uniqueness.
The saved program checks all determinant, adjugate,
characteristic-polynomial and finite-order identities over
the rational function field, and checks the exceptional
line symbolically for every $z$.

\end{enumerate}
\end{proof}




\begin{corollary}[all good-prime lattices are forced]\label{maps:src:F.2.21}
\StatementRecord{maps:src:F.2.21}{proved}{Historical events: Claude Fable 5.2 (Anthropic), effort high, 16 November 2026: carried as background without fault found. GPT-6 Astra (OpenAI), ChatGPT Pro mode, effort Pro, S6\_09\_PAPERS, 17 November 2026: targeted proof check. Also checked in this writing. Pre-continuation writing events: GPT-6 Astra (OpenAI), ChatGPT Pro mode, effort Pro, S6\_10\_PAPERS, 18 November 2026 02:43:58.020152+00:00: targeted proof check. Continuation metadata events: GPT-6 Astra (OpenAI), ChatGPT Pro mode, effort Pro, S6\_10\_PAPERS continuation, 19 November 2026: used and rechecked in this run; status-vocabulary reconciliation; metadata reconciliation only; no new mathematical check. \emph{Pre-continuation recorded proof scope.} Claude Fable 5.2 (Anthropic), effort high, 16 November 2026; also checked in this writing by GPT-6 Astra (OpenAI), ChatGPT Pro mode, effort Pro, 18 November 2026}{the hypotheses, notation and definitions stated here; \cref{maps:src:F.2.20}; \cref{maps:src:F.2.19}}


For a $(5,4)$ datum with $\pi_1=1$ and $H_2=0$, the
rank-five quotient has absolutely irreducible reduction
at every prime $p\ne2,5$.  At each such prime its stable
$\mathbb Z_p$-lattice is unique up to homothety.
The same is true of the rank-six extension among lattices
with primitive maximal cusp.  Thus, once a rational pair
in \ref{maps:src:F.2.20} admits one survivor, all remaining integral-lattice
ambiguity is supported at the primes two and five.
This assertion supplies no finite bound on those bad-prime
lattices or on the rational parameters.

\end{corollary}


\begin{proof}
\begin{enumerate}[label=\arabic*.]
\item Work first over an algebraic closure of $\mathbb F_p$.
Both finite generators are semisimple because $p\ne2,5$.
A one-dimensional composition factor must be trivial:
its $A$ and $C$ eigenvalues agree by the unipotent product,
and their orders divide five and four.
There can be at most one such factor, since the $1$
eigenvalue of $A$ occurs once.
If it occurs, the remaining four-dimensional module is
absolutely irreducible by the spectral argument in \ref{maps:src:F.2.19}.
The only other possible reducible composition pattern is
$2+3$.
On each factor, determinants of $A$ and $C$ agree and
have orders dividing five and four, hence are one.
The two-dimensional $C$ eigenvalue pair is therefore
$(-1,-1)$ or $(i,-i)$.
The first would give $C=-I$ on that factor, contradicting
the unipotence of $CA^{-1}$ and the absence of $-1$
among the eigenvalues of $A$.
In the second case the three-dimensional factor has
$C$ eigenvalues $1,-1,-1$, so $C^2=I$ there; it has
an $A$-fixed vector because the other factor has determinant
one and cannot contain the simple $1$ eigenvalue of $A$.

\item Here is the required elementary involutory identity in the
present characteristic.  Put $B=A^{-1}$ on this factor.
From $(CB-I)^2=0$ and $C^2=I$ one obtains
\[
 BC+CB^{-1}=2I.
\]
For $Bv=v$ this says $(B+I)Cv=2v$.
The operator $B+I$ is invertible, since $B^5=I$ and
$p\ne2$, so $Cv=v$.  The resulting common fixed line
contradicts simplicity of the three-dimensional factor.
We have proved that any reducible residual representation
has exactly the constituents $1$ and an absolutely
irreducible four-dimensional module.  Theorem \ref{maps:src:F.2.19}
excludes that possibility for a survivor, proving absolute
irreducibility at every indicated prime.

\item Let $R=\mathbb Z_p$ and $K_p=\mathbb Q_p$.
For two stable quotient lattices, multiply one by a power
of $p$ until it is contained in the other but not in its
$p$-multiple.  Its nonzero residual image must then be the
whole quotient by absolute irreducibility.  Nakayama's
lemma, or the Smith form of the inclusion, makes the
lattices equal.  This proves quotient-lattice uniqueness.

\item For completeness the extension-lattice comparison is also
integral.  Fix a primitive-cusp lattice $V_R$, with invariant
line $Re$ and quotient $W_R$.
The coboundary map
\[
 W_R^*\longrightarrow W_R^*\oplus W_R^*
\]
is primitive, because its reduction has trivial kernel.
Thus its cokernel, the degree-one cohomology of the free
group on the two generators, is torsion free, and an
integral coboundary over $K_p$ is already integral over $R$.
The extension class $\alpha$ is primitive in this cokernel.
Indeed, divisibility by $p$ would make the extra cusp map
\[
 \ker N_W/\operatorname{im}N_W\longrightarrow Re
\]
divisible by $p$, contrary to the three unit Smith factors
of the full cusp.  Here $N_W$ has primitive image:
it is the quotient of the primitive image of $N_V$ by
$Re\subset\operatorname{im}N_V$.

\item Normalize any other stable lattice by the already proved
quotient homothety so that its quotient is $W_R$.
Its invariant line is $p^rRe$ for some integer $r$.
The extension class in that basis is $p^{-r}\alpha$;
primitivity makes it integral only when $r\leq0$.
Primitivity of the coboundary map eliminates any additional
rational change of splitting.  The extra cusp map is then
multiplied by $p^{-r}$, so its Smith factor is a unit exactly
when $r=0$.  The lattices coincide after the common homothety.
This proves the full assertion and retains the bad-prime
and rational-parameter problems explicitly.

\end{enumerate}
\end{proof}




\begin{theorem}[a terminating lattice algorithm for each rational pair]\label{maps:src:F.2.22}
\StatementRecord{maps:src:F.2.22}{proved}{Historical events: Claude Fable 5.2 (Anthropic), effort high, 16 November 2026: carried as background without fault found. GPT-6 Astra (OpenAI), ChatGPT Pro mode, effort Pro, S6\_09\_PAPERS, 17 November 2026: targeted proof check. Also checked in this writing. Pre-continuation writing events: GPT-6 Astra (OpenAI), ChatGPT Pro mode, effort Pro, S6\_10\_PAPERS, 18 November 2026 02:43:58.020152+00:00: computation or formal certificate; targeted proof check. Continuation metadata events: GPT-6 Astra (OpenAI), ChatGPT Pro mode, effort Pro, S6\_10\_PAPERS continuation, 19 November 2026: used and rechecked in this run; status-vocabulary reconciliation; metadata reconciliation only; no new mathematical check. \emph{Pre-continuation recorded proof scope.} Claude Fable 5.2 (Anthropic), effort high, 16 November 2026; also checked in this writing by GPT-6 Astra (OpenAI), ChatGPT Pro mode, effort Pro, 18 November 2026}{the hypotheses, notation and definitions stated here; \cref{maps:src:F.2.14}; \cref{maps:eq:F.L3}; \cref{maps:eq:F.L5}; \cref{maps:src:F.2.11}; \cref{maps:eq:F.L6}; \cref{maps:eq:F.L7}}


Let $(A,C)$ be any irreducible rational rank-five pair in
one of the eight spectral cases of \ref{maps:src:F.2.14}, with
$A^5=I$, $B=C^{-1}$, and $(AB-I)^2=0$ of rank two.
There is an explicit terminating algorithm which decides
whether the pair preserves an integral lattice, and, if so,
lists all such lattices up to simultaneous integral conjugacy.
It then lists all rank-six extensions with the same finite
orders and primitive maximal square-zero cusp, up to
simultaneous integral conjugacy.  The bounds below depend
on the input rational pair and are proved to lose no lattice.
They do not bound the rational parameter space or classify
all affine twists and resulting homology.

\end{theorem}


\begin{proof}
\begin{enumerate}[label=\arabic*.]
\item We give the algorithm and its termination proof.
The fixed line of $A$ has dimension one.  Choose a nonzero
vector $e$ on it and a nonzero fixed covector $\varphi$.
Irreducibility implies that the vectors $ge$, with $g$ a
word in $A,C$, span the whole rational space; the analogous
statement holds for the covectors $\varphi h$.
In fact words of length at most four suffice.  The spans of
words of lengths at most $j$ increase from dimension one;
if one step does not increase the span, it is invariant
under both generators and hence already the whole space.
There are only four possible increases.

\item Choose words $g_1,\ldots,g_5$ and $h_1,\ldots,h_5$ giving
bases in this way, and put
\[
 P=I+A+\cdots+A^4,\qquad T_{ij}=g_iPh_j.
 \tag{F.L1}\label{maps:eq:F.L1}
\]
The rank-one operator $P$ is a nonzero scalar multiple of
$e\varphi$.  Hence the 25 matrices $T_{ij}$ are a rational
basis of $\operatorname{Mat}_5(\mathbb Q)$.
Each belongs to $\mathcal R=\mathbb Z[A,C]$.
Let $\mathcal R_0$ be their integral span and set
\[
 \mathcal R_0^\#=
 \{X:\operatorname{tr}(XT)\in\mathbb Z
                     \text{ for all }T\in\mathcal R_0\}.
 \tag{F.L2}\label{maps:eq:F.L2}
\]
The trace pairing on the full matrix algebra is
nondegenerate, as is seen from the matrix units, so this
is another full rational lattice, computed by inverting
the 25-by-25 Gram matrix.

\item If any stable lattice exists, every element of $\mathcal R$
acts integrally on it.  Consequently
$\operatorname{tr}(XY)\in\mathbb Z$ for all
$X,Y\in\mathcal R$, and necessarily
\[
 \mathcal R_0\subset\mathcal R\subset\mathcal R_0^\#.
 \tag{F.L3}\label{maps:eq:F.L3}
\]
A nonintegral Gram entry therefore proves nonexistence.
Otherwise let
$d=[\mathcal R_0^\#:\mathcal R_0]$, the absolute Gram
determinant.  Start with the lattice generated by
$\mathcal R_0$ and $I$, and repeatedly replace a current
lattice $L$ by $L+AL+CL$, using exact Hermite reduction.
If it leaves $\mathcal R_0^\#$, \eqref{maps:eq:F.L3} proves that there
is no stable lattice.
While it remains inside, every strict enlargement reduces
its index by a factor at least two.
There can be at most $\lfloor\log_2d\rfloor$ such
enlargements starting from $\mathcal R_0$, so a bound of
$1+\lfloor\log_2d\rfloor$ tests after adjoining $I$ is
sufficient.  Stabilization gives exactly $\mathcal R$:
the stable lattice contains $I$ and all words in $A,C$,
and every element introduced belongs to their ring.
It is then an order, and
$\Lambda_0=\mathcal R\mathbb Z^5$ is an explicitly
computed stable lattice.  This proves the decision step,
including termination in the nonexistence case.

\item Use a basis of $\Lambda_0$ from now on, so
$\mathcal R\subset\operatorname{Mat}_5(\mathbb Z)$.
Smith reduction gives an integer $c\ge1$ with
\[
 c\operatorname{Mat}_5(\mathbb Z)\subset\mathcal R;
 \tag{F.L4}\label{maps:eq:F.L4}
\]
one may take the largest Smith factor of the inclusion
of the order in the full matrix lattice.
Normalize any stable lattice by a rational homothety so
that it is contained in $\mathbb Z^5$ and the gcd of all
coordinates of a lattice basis is one.
The operators $cE_{ij}$ preserve it.  The coordinate values
of its vectors generate $\mathbb Z$, so Bezout's identity,
applied to $cE_{ij}$, shows that it contains every $ce_i$.
Therefore every normalized lattice lies in the finite interval
\[
 c\mathbb Z^5\subset\Lambda\subset\mathbb Z^5.
 \tag{F.L5}\label{maps:eq:F.L5}
\]
Conversely every lattice in this interval stable under
$A,C$ is retained by the algorithm.
The only rational operators commuting with the pair are
scalars: they act by a scalar on the one-dimensional fixed
line of $A$, and then by the same scalar on the spanning
vectors $g_ie$.  Thus rational equivalence of two stable
lattices is precisely homothety.  The coordinate-gcd
normalization is unique, so \eqref{maps:eq:F.L5} gives the integral
conjugacy classification without further identifications.

\item Here is a finite enumeration bound which is also implemented.
For each $p^a\Vert c$, enumerate all subspaces of
$\mathbb F_p^5$ in reduced echelon form.
Beginning with $\mathbb Z^5$, pass through $a$ levels of
stable lattices between $pH$ and $H$, retaining only those
containing $p^a\mathbb Z^5$ and identifying equal Hermite
forms.  Every desired lattice occurs: for any target
$\Lambda$, the chain
$\Lambda+p^j\mathbb Z^5$, $0\le j\le a$, is such a
chain and ends at $\Lambda$.
If $s_p$ denotes the number of all subspaces of
$\mathbb F_p^5$, there are at most $s_p^a$ branches before
deduplication.  Intersect the resulting lattices for the
different primes and discard those of coordinate gcd
greater than one.  The Chinese remainder theorem proves
that this enumerates \eqref{maps:eq:F.L5} completely.

\item For each quotient lattice $W_{\mathbb Z}$ so obtained,
put $N_W=AB-I$.
Its image must be primitive: for a primitive full cusp,
$\operatorname{im}N_V=\ker N_V$ contains the primitive
invariant line, and quotienting that image by the line
gives a primitive image in $W_{\mathbb Z}$.
Reject a quotient whose two nonzero Smith factors are
not one.  For the remaining quotients use extension rows
$c_1,c_2\in\mathbb Z^{1\times5}$ and impose
\[
\begin{gathered}
 c_1(I+A+\cdots+A^4)=0,
 \qquad c_2(I+B+\cdots+B^{m-1})=0,\\
 (c_2+c_1B)N_W=0,
 \qquad m=\operatorname{ord}(B).
\end{gathered}
 \tag{F.L6}\label{maps:eq:F.L6}
\]
These are integral linear equations.  Their integral kernel
modulo the coboundaries
\[
 z\longmapsto\bigl(z(A-I),z(B-I)\bigr),
 \qquad z\in\mathbb Z^{1\times5},
 \tag{F.L7}\label{maps:eq:F.L7}
\]
is a finitely generated group $G$ of rank one by \ref{maps:src:F.2.11}.
Compute its free and torsion factors by Smith reduction.
Choose an integral generator $w$ of the free rank-one
quotient $\ker N_W/\operatorname{im}N_W$.
The homomorphism
\[
 \kappa:G\longrightarrow\mathbb Z,
 \qquad [(c_1,c_2)]\longmapsto(c_2+c_1B)w
 \tag{F.L8}\label{maps:eq:F.L8}
\]
is well defined and nonzero: \eqref{maps:eq:F.L6} kills the image, and
\eqref{maps:eq:F.L7} kills $w$; rational nonvanishing is \ref{maps:src:F.2.11}.
The full cusp has unit Smith factors exactly when
$\kappa=\pm1$.
Indeed the quotient rank-two image is already primitive,
and the remaining invariant-line image is generated by
that one integer.
If $\operatorname{im}\kappa\ne\mathbb Z$, reject the
lattice.  Otherwise the inverse image of $1$ is a finite
coset of $\operatorname{Tor}G$, and lists all extensions.
Changing the sign of the invariant-line generator identifies
$\kappa=-1$ with $\kappa=1$.
There are no further identifications: an integral isomorphism
on the quotient is $\pm I$, on the invariant line is
$\pm1$, and, after $\kappa=1$ is fixed, their ratio is one.
The remaining row is exactly an integral change of splitting
\eqref{maps:eq:F.L7}.
This proves completeness and termination of the extension
step as well.

\item The order, trace-dual bound, conductor and finite-interval
steps are implemented in
\texttt{fourfolds\_lattice\_algorithm.py};
the extension-kernel and Smith steps are implemented in
\texttt{fourfolds\_two\_classes\_lattices.py} with the
finite order as an argument.  Their applications and exact
exhaustive bounds are recorded next.

\end{enumerate}
\end{proof}




\begin{theorem}[complete lattice classes for the two family survivors]\label{maps:src:F.2.23}
\StatementRecord{maps:src:F.2.23}{proved}{Historical events: Claude Fable 5.2 (Anthropic), effort high, 16 November 2026: carried as background without fault found. GPT-6 Astra (OpenAI), ChatGPT Pro mode, effort Pro, S6\_09\_PAPERS, 17 November 2026: targeted proof check. Also checked in this writing. Pre-continuation writing events: GPT-6 Astra (OpenAI), ChatGPT Pro mode, effort Pro, S6\_10\_PAPERS, 18 November 2026 02:43:58.020152+00:00: computation or formal certificate; targeted proof check. Continuation metadata events: GPT-6 Astra (OpenAI), ChatGPT Pro mode, effort Pro, S6\_10\_PAPERS continuation, 19 November 2026: used and rechecked in this run; status-vocabulary reconciliation; metadata reconciliation only; no new mathematical check. \emph{Pre-continuation recorded proof scope.} Claude Fable 5.2 (Anthropic), effort high, 16 November 2026; also checked in this writing by GPT-6 Astra (OpenAI), ChatGPT Pro mode, effort Pro, 18 November 2026}{the hypotheses, notation and definitions stated here; \cref{maps:src:F.2.16}; \cref{maps:src:F.2.6}; \cref{maps:src:F.2.18}; \cref{maps:eq:F.U1}; \cref{maps:eq:F.L4}; \cref{maps:eq:F.L5}; \cref{maps:eq:F.L6}; \cref{maps:eq:F.L7}; \cref{maps:src:F.2.12}; \cref{maps:eq:F.L10}; \cref{maps:eq:F.U6}; \cref{maps:src:F.2.19}}


For each of the two rational representations $v=0$ and $v=1$
of \ref{maps:src:F.2.16}, there are exactly two homothety classes of stable
rank-five quotient lattices, and exactly three integral
conjugacy classes of rank-six extensions with primitive
maximal cusp.  Exactly one of those three can have simply
connected finite affine fillings.  It is the integral class
already realized in \ref{maps:src:F.2.6}, respectively \ref{maps:src:F.2.18}.
Thus arbitrary stable-lattice changes produce no further
linear survivor classes anywhere in the family \ref{maps:src:F.2.16}.
This assertion retains the classification of all affine
twists and multiplicities as a separate problem.

\end{theorem}


\begin{proof}
\begin{enumerate}[label=\arabic*.]
\item Apply the preceding algorithm in the original five-cycle
lattice \eqref{maps:eq:F.U1}.  For both parameters the order is obtained
after four strict enlargements.  Its Smith factors in the
full matrix lattice are 21 factors one and four factors two;
thus the conductor $c=2$ in \eqref{maps:eq:F.L4} is exact.
All 374 subspaces of $\mathbb F_2^5$, including zero and
the whole space, are enumerated, with no height restriction
or neighbor-depth assumption.
After imposing stability and coordinate gcd one, the two
lattice representatives are
\[
 \mathbb Z^5,\qquad
 \Lambda_* =2\mathbb Z^5+\mathbb Z(1,1,1,1,1)^t.
 \tag{F.L9}\label{maps:eq:F.L9}
\]
The latter has index $16$.
The generic prime-by-prime implementation of \eqref{maps:eq:F.L5} and
an independent bit-vector enumeration give exactly the
same two Hermite forms.

\item For $\mathbb Z^5$, the extension group in \eqref{maps:eq:F.L6}--\eqref{maps:eq:F.L7}
is $\mathbb Z$ and $\kappa$ is a unit on its generator.
For $\Lambda_*$, that group is
$\mathbb Z\oplus\mathbb Z/2$, again with primitive
$\kappa$.  Hence there are exactly one and two primitive
extensions, respectively.
For each extension the program computes the image of the
primitive global invariant homology covector $\ell$ on
the two finite fixed lattices.  The resulting positive
divisibilities are
\[
\begin{array}{c|c|c}
 \text{quotient lattice}&\text{extension class}&
       (\operatorname{div}_1,\operatorname{div}_2)\\\hline
 \mathbb Z^5&\text{unique}&(1,2)\\
 \Lambda_*&\text{first torsion coset}&(1,1)\\
 \Lambda_*&\text{second torsion coset}&(1,2).
\end{array}
 \tag{F.L10}\label{maps:eq:F.L10}
\]
The table is identical for $v=0$ and $v=1$.
All six full matrices, the integral kernel bases, the
Smith factors and the finite divisibilities are computed
exactly by \texttt{fourfolds\_two\_classes\_lattices.py}
and saved in its output and JSON certificate.

\item For a simply connected filling the determinant quotient
of \ref{maps:src:F.2.12} requires $\gcd(\ell(v_i),m_i)=1$ at each
finite cap.  The second affine multiplicity is divisible
by four.  A divisibility of two in \eqref{maps:eq:F.L10} therefore
excludes that class for every choice of twists and larger
multiplicities.
The sole remaining class has an explicit integral
intertwiner of determinant $\pm1$ with the corresponding
pair in \eqref{maps:eq:F.U6}, as printed and checked by the program.
Those pairs already have the verified simply connected,
three-connected fillings.  \ref{maps:src:F.2.19} excludes every other
rational parameter for all lattices.
Together these facts prove all the claims, including the
absence of any additional linear survivor class in the
unbounded family.

\end{enumerate}
\end{proof}




\begin{theorem}[every surviving twist in the unbounded family]\label{maps:src:F.2.24}
\StatementRecord{maps:src:F.2.24}{proved}{Historical events: Claude Fable 5.2 (Anthropic), effort high, 16 November 2026: carried as background without fault found. GPT-6 Astra (OpenAI), ChatGPT Pro mode, effort Pro, S6\_09\_PAPERS, 17 November 2026: targeted proof check. Also checked in this writing. Pre-continuation writing events: GPT-6 Astra (OpenAI), ChatGPT Pro mode, effort Pro, S6\_10\_PAPERS, 18 November 2026 02:43:58.020152+00:00: computation or formal certificate; targeted proof check. Continuation metadata events: GPT-6 Astra (OpenAI), ChatGPT Pro mode, effort Pro, S6\_10\_PAPERS continuation, 19 November 2026: used and rechecked in this run; status-vocabulary reconciliation; metadata reconciliation only; no new mathematical check. \emph{Pre-continuation recorded proof scope.} Claude Fable 5.2 (Anthropic), effort high, 16 November 2026; also checked in this writing by GPT-6 Astra (OpenAI), ChatGPT Pro mode, effort Pro, 18 November 2026}{the hypotheses, notation and definitions stated here; \cref{maps:src:F.2.12}; \cref{maps:src:F.2.16}; \cref{maps:src:F.2.23}; \cref{maps:src:F.2.20}; \cref{maps:src:F.2.19}; \cref{maps:eq:F.U6}; \cref{maps:src:F.2.5}; \cref{maps:src:F.2.18}; \cref{maps:eq:F.T4}; \cref{maps:src:F.2.6}; \cref{maps:eq:F.T6}; \cref{maps:eq:F.T7}; \cref{maps:eq:F.T5}; \cref{maps:eq:F.T8}; \cref{maps:eq:F.T3}; \cref{maps:eq:F.T2}; \cref{maps:eq:F.T9}}


Use the fixed Kuhn cusp and the affine translation clutching
convention of \ref{maps:src:F.2.12}.  In the entire rational family \ref{maps:src:F.2.16},
allow every stable lattice, every free affine finite lift,
and every pair of actual finite multiplicities.
The data with $\pi_1=1$ and $H_2=H_3=0$ are precisely the
following, up to the integral linear classification already
proved and the translation identifications described below:
\[
\begin{gathered}
 v\in\{0,1\},\qquad
 \text{the unique eligible integral class of \ref{maps:src:F.2.23}},\\
 (m_1,m_2)=(5,4),\qquad
 20\kappa+4r_1+5r_2=\pm1.
\end{gathered}
 \tag{F.T1}\label{maps:eq:F.T1}
\]
Here $r_i=\ell(v_i)$ is the invariant affine numerator, and
$\kappa\in\mathbb Z$ is the invariant coordinate of the
cusp translation clutching, with signs fixed by the displayed
van Kampen relation below.  The remaining fixed-direction
components of the twists impose no condition.
Every manifold in \eqref{maps:eq:F.T1} is oriented diffeomorphic to its
corresponding known survivor and hence to
$3\mathbb{HP}^2\#(-\mathbb{HP}^2)$.
This completes the survivor classification for this entire
unbounded rational family; it does not classify the other
rational points of \ref{maps:src:F.2.20} or the other spectral loci.

\end{theorem}


\begin{proof}
\begin{enumerate}[label=\arabic*.]
\item \ref{maps:src:F.2.19} excludes all other rational parameters, and \ref{maps:src:F.2.23}
excludes all other integral linear classes.
We therefore work with either known integral pair, in the
marking of \eqref{maps:eq:F.U6}.  Its primitive invariant homology
covector is $\ell(x)=x_0$.
Each finite fixed lattice has an integral basis
$u_i,\lambda_i$, where $\ell(u_i)=1$,
$\ell(\lambda_i)=0$, and $\lambda_i$ is primitive in
$\ker\ell$.
We may write the affine numerator as
$v_i=r_i u_i+s_i\lambda_i$.
This covers arbitrary finite affine lifts: translating the
origin makes their translation vector invariant by averaging,
and its $m_i$th iterate is then the integral vector $v_i$.
The resulting boundary translation change is included in the
specified affine clutching.

\item For both linear pairs, the normal closure of the cusp angular
lattice in the fibre lattice is exactly $\ker\ell$, by
\ref{maps:src:F.2.5} and \ref{maps:src:F.2.18}.  Van Kampen therefore leaves the abelian
presentation
\[
 m_1a=r_1h,\qquad m_2b=r_2h,\qquad
 a+b+\kappa h=0.
 \tag{F.T2}\label{maps:eq:F.T2}
\]
This also defines the sign of $\kappa$.
Its determinant is, up to sign,
\[
 p=m_1m_2\kappa+m_2r_1+m_1r_2.
 \tag{F.T3}\label{maps:eq:F.T3}
\]
Thus simple connectivity is exactly $p=\pm1$.
It forces $\gcd(m_1,m_2)=1$ and
$\gcd(r_i,m_i)=1$ by reduction modulo each multiplicity.

\item We first prove that a surviving datum has the minimal
multiplicities.  Set $n_1=5$, $n_2=4$ and $m_i=n_i k_i$.
In the splitting $\Lambda=\mathbb Zu_i\oplus\ker\ell$,
write the finite linear map as $1\oplus D_i$.
The reduced finite fibre is a mapping torus of $T^5$ over
the circle $m_i\ell$; its fibre is connected because
$r_i$ is prime to $m_i$.
Its vertical monodromy is a power of $D_i$ prime to $n_i$,
possibly followed by a translation, which changes no
integral Wang group.

\item The primitive fixed vertical loop $\lambda_i$ gives the
only additional free two-cycle beyond vertical cycles.
Consequently an invariant integral two-form $\omega$ on
the covering torus descends integrally exactly when
\[
 \frac{\omega(v_i,\lambda_i)}{m_i}
       =\frac{r_i\omega(u_i,\lambda_i)}{m_i}
       \quad\text{is an integer}.
 \tag{F.T4}\label{maps:eq:F.T4}
\]
For precision, the torus made from the affine generator
and $\lambda_i$ projects with degree $r_i$ to the mapping-
torus base, whereas the ordinary $u_i$ suspension projects
with degree $m_i$.  A Bezout combination is a primitive
suspension.  Its integrality condition is equivalent to
\eqref{maps:eq:F.T4}, because the ordinary covering periods are already
integral and $\gcd(r_i,m_i)=1$.
The Wang sequence shows that these vertical cycles and one
primitive suspension generate the free second homology;
there is no omitted period condition.

\item The integral invariant two-form computations in \ref{maps:src:F.2.6} and
\ref{maps:src:F.2.18} give
\[
 \gcd_{\omega}\omega(u_i,\lambda_i)=n_i.
 \tag{F.T5}\label{maps:eq:F.T5}
\]
They are also independently printed for both pairs by
\texttt{fourfolds\_family\_twists.py}.
Hence the covering-image cokernel in degree two at the
$i$th finite point is exactly $\mathbb Z/k_i$.
The global invariant group of $\bigwedge^2\mathcal V$
and its integral first parabolic group both vanish.
The skyscraper cokernel sequence therefore gives
\[
 H^1(B,R^2g_*\mathbb Z)
       =\mathbb Z/k_1\oplus\mathbb Z/k_2.
 \tag{F.T6}\label{maps:eq:F.T6}
\]
Point-supported torsion kernels do not change this degree-one
base cohomology, and the cusp specialization is onto.
The term \eqref{maps:eq:F.T6}, in bidegree $(1,2)$, survives every Leray
differential on the two-dimensional base.
All rational terms of total degree three vanish, independently
of the multiplicities and twists, by the same exterior
local-system table as \ref{maps:src:F.2.6}.  Thus $H^3$ is finite.
If either $k_i>1$, its surviving nonzero filtration quotient
makes $H^3$ nonzero; universal coefficients then give
nonzero torsion in $H_2$.
This proves $m_1=5$, $m_2=4$ for every survivor.

\item We now identify all twists with these multiplicities.
The following integral fact permits removal of each $s_i$:
\[
 \operatorname{Norm}_{n_i}:
 \ker\ell\longrightarrow\mathbb Z\lambda_i
 \quad\text{is onto}.
 \tag{F.T7}\label{maps:eq:F.T7}
\]
Here the norm is $I+D_i+\cdots+D_i^{n_i-1}$.
One can read \eqref{maps:eq:F.T7} directly from \eqref{maps:eq:F.T5}.
Let $\varphi_i$ generate the primitive invariant covector
lattice on $\ker\ell$.
The splitting by $u_i$ shows that the gcd in \eqref{maps:eq:F.T5} is
$|\varphi_i(\lambda_i)|$: its forms are
$\ell\wedge\varphi_i$ plus invariant vertical forms.
Thus that value is $n_i$.
The norm has the form
$x\mapsto d_i\varphi_i(x)\lambda_i$ for an integer
$d_i$; evaluating $\varphi_i$ gives
$|d_i|n_i=n_i$, hence $|d_i|=1$.
This proves surjectivity.  The saved program additionally
gives integral preimages $q_i\in\ker\ell$ with
$\operatorname{Norm}_{n_i}q_i=\lambda_i$.

\item Put
\[
 w_i=-\frac{s_i}{n_i}\sum_{j=0}^{n_i-1}jM_i^j q_i.
\]
Then
\[
 (I-M_i)w_i=-s_i\lambda_i/n_i+s_iq_i.
 \tag{F.T8}\label{maps:eq:F.T8}
\]
Translation by $w_i$ conjugates the finite affine action
with numerator $r_i u_i+s_i\lambda_i$ to that with
numerator $r_i u_i$: the remaining $s_iq_i$ is integral.
It gives a diffeomorphism of the entire finite cap.
In the standard linear boundary marking, its translation
section is
$w_i+\theta s_i\lambda_i/n_i$.
Its endpoint mismatch with the linear deck map is exactly
$s_iq_i$ by \eqref{maps:eq:F.T8}, so its integral boundary residue has
$\ell$-coordinate zero.

\item We explain why these local changes extend simultaneously,
including the prescribed cusp marking.
Let $P$ be the complementary pair of pants and
$\mathcal T\to P$ its flat torus group bundle.
A section on the boundary extends over the relative
one-skeleton because the torus is connected.
The only remaining obstruction is the two-cell obstruction
with values in the fibre fundamental-group local system:
\[
 H^2(P,\partial P;\Lambda)
        \simeq H_0(P;\Lambda)=\Lambda_\Gamma.
 \tag{F.T9}\label{maps:eq:F.T9}
\]
The equality follows from the oriented relative cellular
complex, equivalently integral Poincare--Lefschetz duality
with local coefficients.  Its boundary residues map to
the indicated global coinvariants.
For both linear pairs the horizontal monodromy relation
matrix has five unit Smith factors, so
$\Lambda_\Gamma=\mathbb Z$, detected exactly by $\ell$.
This is independently checked in the saved program.
Thus the residues $s_iq_i$, with the identity section at the
cusp, have zero obstruction and extend over $P$.
A section may be smoothed relative to its boundary collars;
translation by it is a fibre-preserving diffeomorphism.
Together with the finite-cap maps it removes all $s_i$
without changing the cusp.
The same argument removes all components of translation
clutching whose $\ell$-coordinate is zero.

\item It remains to handle the integer numerators and the remaining
cusp coordinate.  If $p=-1$, the congruences from \eqref{maps:eq:F.T3}
give
\[
 r_1=1+5a,\qquad r_2=-1+4b,
 \qquad a+b+\kappa=0.
 \tag{F.T10}\label{maps:eq:F.T10}
\]
Adding $5a u_1$ or $4b u_2$ changes no finite affine action;
in its linear boundary marking it changes the translation
residue by $a u_1$ or $b u_2$.
The difference of the cusp marking has residue of
$\ell$-coordinate $\kappa$, by \eqref{maps:eq:F.T2}.
Their total obstruction in \eqref{maps:eq:F.T9} is therefore
$a+b+\kappa=0$.
The same extension argument identifies the glued manifold
with the original twists $(u_1,-u_2)$ and zero cusp
clutching.
If $p=1$, the identical calculation starts from
$(-u_1,u_2)$.

\item Finally simultaneous fibre inversion identifies these two
sign choices by an orientation-preserving diffeomorphism.
It conjugates every affine translation to its negative.
It extends over the Kuhn cusp: the lattice involution
$(x,h)\mapsto(-x,h)$ preserves its toric fan, since a
simplex given by a chain in a unit cube is taken to the
chain with reversed coordinate order in the opposite cube.
It normalizes the cusp lattice action and therefore descends
to its quotient, while fixing the base parameter.  The resulting toric map is holomorphic,
and in the regular real six-torus inversion has determinant
$(-1)^6=1$.  All translation maps also preserve orientation.
Thus every datum in \eqref{maps:eq:F.T1} gives the stated oriented smooth
type and has $H_2=H_3=0$, proving sufficiency as well as
necessity.

\end{enumerate}
\end{proof}




\begin{theorem}[the universal degree-two suspension obstruction]\label{maps:src:F.2.25}
\StatementRecord{maps:src:F.2.25}{proved}{Historical events: Claude Fable 5.2 (Anthropic), effort high, 16 November 2026: carried as background without fault found. GPT-6 Astra (OpenAI), ChatGPT Pro mode, effort Pro, S6\_09\_PAPERS, 17 November 2026: targeted proof check. Also checked in this writing. Pre-continuation writing events: GPT-6 Astra (OpenAI), ChatGPT Pro mode, effort Pro, S6\_10\_PAPERS, 18 November 2026 02:43:58.020152+00:00: computation or formal certificate; targeted proof check. Continuation metadata events: GPT-6 Astra (OpenAI), ChatGPT Pro mode, effort Pro, S6\_10\_PAPERS continuation, 19 November 2026: used and rechecked in this run; status-vocabulary reconciliation; metadata reconciliation only; no new mathematical check. \emph{Pre-continuation recorded proof scope.} Claude Fable 5.2 (Anthropic), effort high, 16 November 2026; also checked in this writing by GPT-6 Astra (OpenAI), ChatGPT Pro mode, effort Pro, 18 November 2026}{the hypotheses, notation and definitions stated here; \cref{maps:src:F.2.14}; \cref{maps:src:F.2.7}; \cref{maps:src:F.2.9}; \cref{maps:src:F.2.17}; \cref{maps:src:F.2.12}; \cref{maps:eq:F.T4}; \cref{maps:eq:F.S3}; \cref{maps:eq:F.S6}; \cref{maps:src:E.1.3}; \cref{maps:eq:F.S2}; \cref{maps:eq:F.S7}; \cref{maps:eq:F.S5}; \cref{maps:eq:F.S8}}


Every datum of F.2 with $\pi_1=1$ and $H_2=H_3=0$ has
exactly two genuine finite fillings and has minimal affine
multiplicities.  After ordering the order-five point first,
its finite characteristic polynomials are one of precisely
the following two remaining necessary possibilities:
\[
\begin{array}{c|c|c}
 (m_1,m_2)&\chi_{M_1}&\chi_{M_2}\\\hline
 (5,4)&(x-1)^2\Phi_5&(x-1)^2\Phi_2^2\Phi_4\\
 (5,6)&(x-1)^2\Phi_5&(x-1)^2\Phi_2^2\Phi_3.
\end{array}
 \tag{F.S1}\label{maps:eq:F.S1}
\]
In particular all order-eight and order-twelve spectral
loci, and the other two order-six spectral loci of \ref{maps:src:F.2.14},
are excluded for every rational relative position, every
stable lattice and every free affine twist.
There is no multiplicity search bound in this statement:
the only possible actual affine multiplicities are the
two displayed pairs.

\end{theorem}


\begin{proof}
\begin{enumerate}[label=\arabic*.]
\item Theorem \ref{maps:src:F.2.7} excludes three genuine finite fillings.
For two fillings, \ref{maps:src:F.2.9} and \ref{maps:src:F.2.14} give a primitive global
invariant homology covector $\ell$, finite fixed rank two,
and an irreducible rank-five cohomology quotient.
In particular
\[
 H^0(B,j_*\bigwedge^2\mathcal V)=0.
 \tag{F.S2}\label{maps:eq:F.S2}
\]
This is also immediate from the exterior filtration and
irreducibility in odd dimension, as in \ref{maps:src:F.2.17}.
Simple connectivity gives the determinant quotient of
\ref{maps:src:F.2.12}, so each invariant affine numerator $r_i=\ell(v_i)$
is prime to the actual multiplicity $m_i$, and the two
multiplicities are coprime.
As before one may conjugate any finite affine lift by a
translation so that its numerator $v_i$ is integral and
fixed by its linear part.

\item Consider one finite point.  Let $M$ have linear order $n$,
so $n\mid m$, and let its complementary degree-four factor
be $f(x)$, with
$\chi_M=(x-1)^2f(x)$.
The image of $\ell$ on the finite fixed lattice is $g\mathbb Z$
for some $g\mid n$: averaging $I+M+\cdots+M^{n-1}$ on
any vector of $\ell$-value one proves this divisibility.
Since $g\mid r$ and $\gcd(r,m)=1$, we have $g=1$.
Choose a fixed integral vector $u$ with $\ell(u)=1$.
It gives an integral invariant splitting
\[
 \Lambda=\mathbb Zu\oplus H,
 \qquad H=\ker\ell,\qquad M=1\oplus D,
 \qquad\chi_D=(x-1)f(x).
 \tag{F.S3}\label{maps:eq:F.S3}
\]
The order of $D$ is $n$.
Let $\lambda$ generate its primitive fixed lattice and
let $\varphi$ generate its primitive invariant covector
lattice.  Choose their signs so that
\[
 h=\varphi(\lambda)>0.
 \tag{F.S4}\label{maps:eq:F.S4}
\]
Their rational pairing is nonzero, since the finite map
is semisimple and its averaging projection has rank one.

\item Both $\operatorname{Norm}_n(D)$ and $f(D)$ are integral
operators, have image on the line $\mathbb Q\lambda$,
and have invariant defining covector proportional to
$\varphi$.  Primitivity of $\lambda$ and $\varphi$
therefore implies
\[
 \operatorname{Norm}_n(D)=\frac n h\lambda\varphi,
 \qquad
 f(D)=\frac{f(1)}h\lambda\varphi,
 \qquad h\mid n,\quad h\mid f(1).
 \tag{F.S5}\label{maps:eq:F.S5}
\]
For example, the integer coefficient in the second formula
is seen by applying the integral operator to a vector on
which $\varphi$ is one; evaluating on $\lambda$ determines
that coefficient.  No lattice splitting over a localization
or semisimplification modulo a prime is used here.

\item Write $v=ru+s\lambda$.  The quotient finite fibre is a
connected $T^5$ mapping torus over the circle $m\ell$,
because $r$ is prime to $m$.
Its vertical cycles already have integral periods for an
invariant integral two-form on the covering torus.
The Wang sequence leaves one additional primitive free
suspension cycle, from $\lambda$.
Exactly as in \eqref{maps:eq:F.T4}, descent of an invariant form $\omega$
is equivalent to
$m\mid r\omega(u,\lambda)$, and hence to
$m\mid\omega(u,\lambda)$.
In the integral splitting \eqref{maps:eq:F.S3}, every invariant two-form
is an integral multiple of $\ell\wedge\varphi$ plus an
invariant vertical two-form.  Its possible values on
$(u,\lambda)$ are therefore exactly $h\mathbb Z$.
It follows that the degree-two specialization image at
this point has cokernel
\[
 C^2_p\simeq\mathbb Z/(m/h).
 \tag{F.S6}\label{maps:eq:F.S6}
\]
The integer $m/h$ is defined because $h\mid n\mid m$.
Torsion in the finite fibre cohomology is a specialization
kernel, not an additional restriction on these free periods.
The primitive suspension and vertical cycles generate the
free second homology, so \eqref{maps:eq:F.S6} is the full cokernel.

\item Let $I^2$ be the image sheaf of
$R^2g_*\mathbb Z\to j_*\bigwedge^2\mathcal V$.
The cusp specialization is onto by \ref{maps:src:E.1.3}, and all kernels
and cokernels at finite points are skyscraper sheaves.
Thus $H^1(R^2g_*\mathbb Z)=H^1(I^2)$.
Equation \eqref{maps:eq:F.S2} and the cokernel exact sequence give an
injection
\[
 \bigoplus_{p\text{ finite}}C^2_p
           \lhook\joinrel\longrightarrow H^1(B,R^2g_*\mathbb Z).
 \tag{F.S7}\label{maps:eq:F.S7}
\]
The Leray term of bidegree $(1,2)$ has no incoming or
outgoing differential on a real two-dimensional base.
If $H_2=H_3=0$, the universal coefficient theorem gives
$H^3=0$, so that term must vanish.
Every cyclic cokernel in \eqref{maps:eq:F.S7} is therefore zero.
Combining \eqref{maps:eq:F.S5}--\eqref{maps:eq:F.S6} yields the universal necessary
condition
\[
 m=h=n,\qquad n\mid f(1).
 \tag{F.S8}\label{maps:eq:F.S8}
\]
This both fixes the multiplicity and tests the spectrum.

\item For completeness, apply \eqref{maps:eq:F.S8} to the entire eight-case
list of \ref{maps:src:F.2.14}.  The order-five factor gives $\Phi_5(1)=5$.
At the other point the exact values are
\[
\begin{array}{c|c|r|r}
 n&f&f(1)&n/\gcd(n,f(1))\\\hline
 4&\Phi_2^2\Phi_4&8&1\\
 6&\Phi_2^2\Phi_3&12&1\\
 6&\Phi_2^2\Phi_6&4&3\\
 6&\Phi_3\Phi_6&3&2\\
 8&\Phi_8&2&4\\
 12&\Phi_{12}&1&12\\
 12&\Phi_3\Phi_4&6&2\\
 12&\Phi_4\Phi_6&2&6.
\end{array}
 \tag{F.S9}\label{maps:eq:F.S9}
\]
The last column is a lower bound for the nonzero local
covering-cokernel order already at the smallest allowed
multiplicity.  Only the first two rows pass.
The complete table is checked by
\texttt{fourfolds\_suspension\_filter.py}; its input is
exactly these eight spectra, with no matrix-entry or
relative-position cutoff.  This proves \eqref{maps:eq:F.S1} and all
claimed universal exclusions.

\end{enumerate}
\end{proof}




\begin{theorem}[the complete rational remaining $(5,6)$ locus]\label{maps:src:F.2.26}
\StatementRecord{maps:src:F.2.26}{proved}{Historical events: Claude Fable 5.2 (Anthropic), effort high, 16 November 2026: carried as background without fault found. GPT-6 Astra (OpenAI), ChatGPT Pro mode, effort Pro, S6\_09\_PAPERS, 17 November 2026: targeted proof check. Also checked in this writing. Pre-continuation writing events: GPT-6 Astra (OpenAI), ChatGPT Pro mode, effort Pro, S6\_10\_PAPERS, 18 November 2026 02:43:58.020152+00:00: computation or formal certificate; targeted proof check. Continuation metadata events: GPT-6 Astra (OpenAI), ChatGPT Pro mode, effort Pro, S6\_10\_PAPERS continuation, 19 November 2026: used and rechecked in this run; status-vocabulary reconciliation; metadata reconciliation only; no new mathematical check. \emph{Pre-continuation recorded proof scope.} Claude Fable 5.2 (Anthropic), effort high, 16 November 2026; also checked in this writing by GPT-6 Astra (OpenAI), ChatGPT Pro mode, effort Pro, 18 November 2026}{the hypotheses, notation and definitions stated here; \cref{maps:eq:F.S1}; \cref{maps:eq:F.R2}; \cref{maps:src:F.2.25}; \cref{maps:eq:F.R3}; \cref{maps:eq:F.R5}; \cref{maps:src:F.2.20}; \cref{maps:eq:F.R6}; \cref{maps:eq:F.R7}; \cref{maps:eq:F.R8}; \cref{maps:eq:F.R12}; \cref{maps:src:F.2.22}}


The irreducible rational rank-five pairs in the second row
of \eqref{maps:eq:F.S1}, with square-zero rank-two cusp, are classified
up to simultaneous rational conjugacy by
\[
 \{[r:s:t]\in\mathbb P^2(\mathbb Q):
             D(r,s,t)\ne0,\ Q_-(r,s,t)\ne0\}
       \ \sqcup\ \mathbb A^1(\mathbb Q),
 \qquad Q_-=r^2-rs+st,
 \tag{F.R10}\label{maps:eq:F.R10}
\]
where $D$ is the same cubic as in \eqref{maps:eq:F.R2}.
The exceptional line again lies over $[4:1:3]$.
Thus the two rational spectral loci left by \ref{maps:src:F.2.25} both
have complete explicit normal forms, including their
entire rational exceptional loci.

\end{theorem}


\begin{proof}
\begin{enumerate}[label=\arabic*.]
\item Use precisely the space $\mathbb Q\oplus\mathbb Q(\zeta_5)$,
matrices $A,R,X$ and projective plane normalization in
\eqref{maps:eq:F.R3}--\eqref{maps:eq:F.R5}.  The centralizer and image-plane arguments
in \ref{maps:src:F.2.20} require only irreducibility and the two-dimensional
$-1$ eigenspace of $C$; both hold here.
Consequently every pair has a unique projective image-plane
parameter, with only scalar stabilizer.
The unknown $Y$ now satisfies
\[
 YX=0,\qquad YRX=I_2,\qquad
 \operatorname{tr}(YAX)=-2,\qquad
 \operatorname{tr}(YA^{-1}X)=2.
 \tag{F.R11}\label{maps:eq:F.R11}
\]
Its coefficient matrix is unchanged from \eqref{maps:eq:F.R6}, so has
determinant $D/4$.
The two traces and the $-1$ eigenspace imply
\[
 \chi_C=(x-1)(x+1)^2(x^2+x+1).
\]
Its other three roots are distinct, and the $-1$ eigenspace
has its full multiplicity, so $C^6=I$.
The cusp $XY$ still has square zero and rank two.
For $D\ne0$, the ten linear equations give the unique
explicit solution.

\item The two first-column adjugate identities become
\[
 D y_{10}=2Q_-Q_2,
 \qquad D y_{20}=2(r-s-t)Q_-,
 \tag{F.R12}\label{maps:eq:F.R12}
\]
with the same $Q_2$ as in \eqref{maps:eq:F.R7}.
For $D\ne0$, these entries vanish simultaneously exactly
when $Q_-=0$.
Indeed, if $Q_-\ne0$ they would instead force
$r=s+t$ and $Q_2=0$, and \eqref{maps:eq:F.R8} would give $D=0$.
The invariant-covector argument of \ref{maps:src:F.2.20} is unchanged.
Thus the pair is rationally irreducible exactly off the
conic $Q_-=0$ in the generic locus.

\item No nonzero rational triple satisfies $D=Q_-=0$.
If $s=0$, the conic forces $r=0$, and then $D=7t^3$.
If $s\ne0$, rescale to $s=1$ and put
$r=u$, $t=u-u^2$.  Direct substitution gives
\[
 D(u,1,u-u^2)=
 -(7u^2-11u+7)(u^4-u^3+u^2-u+1).
 \tag{F.R13}\label{maps:eq:F.R13}
\]
The first factor has discriminant $-75$ and positive
leading coefficient.  The second is strictly positive,
as the identity
\[
 u^4-u^3+u^2-u+1
 =\left(u^2-\frac u2\right)^2
  +\frac34\left(u-\frac23\right)^2+\frac23
\]
shows.  In particular neither has a rational zero.
On the exceptional cubic the necessary identities \eqref{maps:eq:F.R12}
therefore force the same two points $[3:2:1]$ and $[4:1:3]$
as before.  The first is again impossible because
$[X\mid RX]$ has rank less than four.
At the second, the complete solution is
\[
 Y_z=\begin{pmatrix}
 4z+26/5&0&-2z-2&-2z-1/5&2z+1\\
 2z-2/5&0&-z&2/5-z&z
 \end{pmatrix},\qquad z\in\mathbb Q.
 \tag{F.R14}\label{maps:eq:F.R14}
\]
Its first-column entries cannot vanish simultaneously,
so each pair on this line is irreducible.
There are no further identifications, by the scalar
stabilizer argument already proved in \ref{maps:src:F.2.20}.
This establishes the claimed bijection.

\item The exact program is
\texttt{fourfolds\_relative\_six\_probe.py}.
It checks the modified characteristic polynomials, finite-order
identities, first-column adjugate expressions, conic
substitution and entire exceptional line symbolically.
Its auxiliary trace-parameter computations for the other
repeated-negative spectrum are recorded as well, although
that spectrum is now excluded integrally by \ref{maps:src:F.2.25}.
No finite sample is used to infer a rational classification.
The finite lattice and extension algorithm of \ref{maps:src:F.2.22} applies
to every point of \eqref{maps:eq:F.R10}; it does not by itself decide
which rational parameters survive the remaining integral
homology tests.

\end{enumerate}
\end{proof}




\begin{theorem}[the entire remaining order-six locus is excluded]\label{maps:src:F.2.27}
\StatementRecord{maps:src:F.2.27}{proved}{Historical events: Claude Fable 5.2 (Anthropic), effort high, 16 November 2026: carried as background without fault found. GPT-6 Astra (OpenAI), ChatGPT Pro mode, effort Pro, S6\_09\_PAPERS, 17 November 2026: targeted proof check. Also checked in this writing. Pre-continuation writing events: GPT-6 Astra (OpenAI), ChatGPT Pro mode, effort Pro, S6\_10\_PAPERS, 18 November 2026 02:43:58.020152+00:00: computation or formal certificate; targeted proof check. Continuation metadata events: GPT-6 Astra (OpenAI), ChatGPT Pro mode, effort Pro, S6\_10\_PAPERS continuation, 19 November 2026: used and rechecked in this run; status-vocabulary reconciliation; metadata reconciliation only; no new mathematical check. \emph{Pre-continuation recorded proof scope.} Claude Fable 5.2 (Anthropic), effort high, 16 November 2026; also checked in this writing by GPT-6 Astra (OpenAI), ChatGPT Pro mode, effort Pro, 18 November 2026}{the hypotheses, notation and definitions stated here; \cref{maps:eq:F.S1}; \cref{maps:src:F.2.25}; \cref{maps:eq:F.S3}; \cref{maps:eq:F.S5}; \cref{maps:eq:F.X2}; \cref{maps:eq:F.X3}; \cref{maps:eq:F.S7}; \cref{maps:eq:F.X5}; \cref{maps:eq:E.9}}


Every simply connected datum in the second row of \eqref{maps:eq:F.S1},
at every rational relative position and every stable integral
lattice with primitive maximal cusp, has
\[
                 \operatorname{Tor}_3 H_3(M;\mathbb Z)\ne0
\]
if its degree-two suspension cokernels vanish.
Consequently no such datum is three-connected.  Together with
\ref{maps:src:F.2.25} this leaves only the minimal order pair $(5,4)$ as a
possible three-connected datum in the full problem F.2.
\end{theorem}


\begin{proof}
\begin{enumerate}[label=\arabic*.]
\item We use cohomology with coefficients in the discrete valuation
ring $R=\mathbb Z_3$.  Suppose the degree-two suspension
cokernels vanish; otherwise \ref{maps:src:F.2.25} already gives nonzero
$H_2$ or $H_3$.  That theorem forces $m=n=h=6$ at the
order-six point, in the notation \eqref{maps:eq:F.S3}--\eqref{maps:eq:F.S5}.
We first determine its local integral representation over $R$.

\item The splitting by the fixed vector $u$ is integral:
$\Lambda_R=Ru\oplus H_R$.  On $H_R$, the finite map $D$
has characteristic polynomial
$(x-1)(x+1)^2\Phi_3(x)$.  Since two is a unit in $R$,
the idempotents $(I\pm D^3)/2$ give an integral direct sum
\[
 H_R=L_+\oplus L_-,\qquad
 \operatorname{rank}L_+=3,\quad
 \operatorname{rank}L_-=2.
\]
On $L_-$ the map is exactly $-I$, by its rational spectrum.
On $L_+$ it has order three.  The primitive fixed vector
$\lambda$ lies in $L_+$, and the primitive invariant
covector $\varphi$ vanishes on $L_-$, so their restrictions
remain primitive.  Equation \eqref{maps:eq:F.S5} with $h=6$ gives
\[
 2(I+D+D^2)|_{L_+}=\lambda\varphi|_{L_+}.
 \tag{F.X1}\label{maps:eq:F.X1}
\]
This operator is nonzero modulo three.  But modulo three
$I+D+D^2=(D-I)^2$ and $(D-I)^3=0$.  The reduction of $D$
on the rank-three space therefore has one Jordan block
of length three.  Choose $q\in L_+$ whose image has
$(D-I)^2q\ne0$ modulo three.  Then $q,Dq,D^2q$ are a
basis modulo three and hence an $R$-basis.  In this basis
$D$ is the three-cycle permutation matrix, since $D^3=I$.
We have proved, without a global lattice classification,
that the full finite representation over $R$ is
\[
                   1\oplus\operatorname{Perm}_3\oplus(-I_2).
 \tag{F.X2}\label{maps:eq:F.X2}
\]
Its dual has the same form, by reversing the three-cycle.

\item Let $u,a,b,c,d,e$ be a cohomology basis for \eqref{maps:eq:F.X2}, with
$a,b,c$ cycled and $d,e$ negated.  A full integral basis
of the invariant degree-three lattice is
\[
 u(ab+bc+ca),\qquad abc,\qquad ude,\qquad(a+b+c)de.
 \tag{F.X3}\label{maps:eq:F.X3}
\]
Here juxtaposition denotes exterior multiplication.
To check completeness, terms with an odd number of
$d,e$ are negated by the cube of the finite map and
have no invariants.  The terms with zero or two such
factors give exactly the four orbit sums in \eqref{maps:eq:F.X3},
each primitive.  With volume $uabcde$, their skew
pairing has one block with coefficient three and one
block with coefficient a unit.  Thus its determinant
has three-adic valuation two.  The exact certificate
\texttt{fourfolds\_universal\_56\_obstruction.py} checks
the full invariant kernel, its primitivity, and the
pairing independently by exterior multiplication and
six-column determinants; its Smith factors are
$(1,1,3,3)$.

\item Let $B_6$ be the actual reduced finite fibre.  Its torus
cover has degree six.  Its free middle cohomology has
rank four and unimodular skew pairing, by integral
Poincare duality.  If $I$ denotes the index of its
pullback lattice in the invariant lattice, multiplication
of pairings by the covering degree gives
\[
 2v_3(I)+2=4v_3(6)=4,
 \qquad v_3(I)=1.
 \tag{F.X4}\label{maps:eq:F.X4}
\]
Consequently the three-primary degree-three covering
cokernel is $R/3R$.  Translations of the affine action
do not enter this calculation: the cover degree,
linear invariant lattice and oriented pairings are
those of the actual free quotient for every allowed
affine numerator.

\item We explain why this local group must appear in total
homology even when the global integral lattice varies.
Let $\mathcal V$ be the rational rank-six cohomology
local system.  It is an extension of the irreducible
rank-five system $W$ by a trivial line.  The determinant
of $W$ is trivial.  Irreducibility in odd dimension
implies that neither $\bigwedge^2 W$ nor its dual
has a nonzero invariant: an invariant alternating
form would have an invariant radical and would have
to be nondegenerate, impossible in odd dimension.
Also $\bigwedge^3W\simeq\bigwedge^2W^*$ by its
trivial determinant.  The exterior filtration therefore
gives
\[
 (\bigwedge^3\mathcal V)^\Gamma=0,\qquad
 ((\bigwedge^2\mathcal V)^*)^\Gamma=0.
 \tag{F.X5}\label{maps:eq:F.X5}
\]
These rational vanishing statements remain true over
$\mathbb Q_3$, since they are kernels of rational
matrices.

\item As in \eqref{maps:eq:F.S7}, take the image of
$R^3g_*R\to j_*\bigwedge^3\mathcal V_R$.
Its kernel and cokernel are point-supported sheaves.
The first vanishing in \eqref{maps:eq:F.X5} makes the local group
$R/3R$ inject into $H^1(B,R^3g_*R)$.  Bidegree $(1,3)$
has no incoming or outgoing Leray differential on
the two-dimensional base, so this injection survives
to $E_\infty^{1,3}$.

\item There is one extension issue, and it cannot remove
this torsion.  The second vanishing in \eqref{maps:eq:F.X5} says
that $H^2(B,R^2g_*R)$ has rational rank zero.
Point-supported changes do not alter this degree-two
base cohomology; it is the coinvariant module, as
in \eqref{maps:eq:E.9}.  It is a finitely generated torsion $R$-module
and hence a finite group.  Therefore
$E_\infty^{2,2}$ is finite as well.  In the filtration
of $H^4(M;R)$, take the inverse image of the injected
$R/3R\subset E_\infty^{1,3}$.  It is an extension
of $R/3R$ by the finite group $E_\infty^{2,2}$.
It is thus a nonzero finite three-primary subgroup
of $H^4(M;R)$; in particular it cannot be absorbed
as a finite quotient of a free filtration kernel.
Universal coefficients identify this torsion with
the three-primary torsion of $H_3(M;\mathbb Z)$.
This proves the claim at every rational position and
every integral lattice, and completes the exclusion.

\end{enumerate}
\end{proof}




\begin{theorem}[universal fundamental group and twist normalization]\label{maps:src:F.2.28}
\StatementRecord{maps:src:F.2.28}{proved}{Historical events: Claude Fable 5.2 (Anthropic), effort high, 16 November 2026: carried as background without fault found. GPT-6 Astra (OpenAI), ChatGPT Pro mode, effort Pro, S6\_09\_PAPERS, 17 November 2026: targeted proof check. Also checked in this writing. Pre-continuation writing events: GPT-6 Astra (OpenAI), ChatGPT Pro mode, effort Pro, S6\_10\_PAPERS, 18 November 2026 02:43:58.020152+00:00: targeted proof check. Continuation metadata events: GPT-6 Astra (OpenAI), ChatGPT Pro mode, effort Pro, S6\_10\_PAPERS continuation, 19 November 2026: used and rechecked in this run; status-vocabulary reconciliation; metadata reconciliation only; no new mathematical check. \emph{Pre-continuation recorded proof scope.} Claude Fable 5.2 (Anthropic), effort high, 16 November 2026; also checked in this writing by GPT-6 Astra (OpenAI), ChatGPT Pro mode, effort Pro, 18 November 2026}{the hypotheses, notation and definitions stated here; \cref{maps:src:F.2.25}; \cref{maps:eq:F.S5}; \cref{maps:eq:F.T8}; \cref{maps:src:F.2.24}}


In the remaining $(5,4)$ locus, fix the marked standard Kuhn
cap and allow affine translation clutching.
Let $\Lambda=\mathbb Z^6$ be fibre homology and let
$\Lambda_\Gamma$ be its global monodromy coinvariants.
Write $\Lambda_\Gamma=\mathbb Zh\oplus T$, with $T$ finite,
and let $\ell$ be the primitive invariant covector detecting
the free factor.  The following are complete reductions for
any datum that could satisfy $\pi_1=1$ and $H_2=H_3=0$.

The actual multiplicities must be $(5,4)$, and the local
vertical fixed pairings must be $h_1=5,h_2=4$, as in \ref{maps:src:F.2.25}.
For any affine twists and clutching the fundamental group
is abelian.  If its free presentation determinant is nonzero,
its order is
\[
 |\pi_1|=|20\kappa+4r_1+5r_2|\,|T|,
 \qquad r_i=\ell(v_i).
 \tag{F.T11}\label{maps:eq:F.T11}
\]
If that determinant is zero the group is infinite.
In particular simple connectivity is equivalent to
\[
 T=0,\qquad 20\kappa+4r_1+5r_2=\pm1.
 \tag{F.T12}\label{maps:eq:F.T12}
\]
Whenever these conditions and $h_1=5,h_2=4$ hold, every
such twist choice is oriented diffeomorphic to the normalized
filling with twists $(u_1,-u_2)$ and zero cusp clutching,
where $u_i$ is any integral fixed vector with $\ell(u_i)=1$.
Thus each eligible integral linear datum has exactly one
normalized smooth filling to test for the remaining
homology conditions; no affine-twist enumeration remains.

\end{theorem}


\begin{proof}
\begin{enumerate}[label=\arabic*.]
\item The image of the fibre fundamental group is a normal
abelian subgroup in the van Kampen presentation.
After quotienting by it, the two finite meridians satisfy
$a^5=b^4=ab=1$.
That quotient is trivial, because five and four are coprime.
The entire fundamental group is therefore equal to the
image of the fibre group and is abelian.
This argument does not assume that the normal closure of
cusp vanishing cycles already exhausts $\ker\ell$.

\item In abelianization the monodromy relations give
$\Lambda_\Gamma$.
The primitive cusp angular lattice is the image of its
nilpotent operator, and this image is already zero among
those monodromy relations.
The remaining presentation is
\[
 5a=v_1,\qquad4b=v_2,\qquad a+b+v_0=0
 \quad\text{in }\Lambda_\Gamma\oplus\mathbb Za\oplus\mathbb Zb,
 \tag{F.T13}\label{maps:eq:F.T13}
\]
where $\ell(v_0)=\kappa$ fixes the sign convention.
Use any Smith presentation of $T$ and split the free part
$\mathbb Zh$.
The full relation matrix is block triangular, with the
presentation of $T$ as one diagonal block and the
three-by-three free matrix of determinant
$\pm(20\kappa+4r_1+5r_2)$ as the other.
Its determinant gives \eqref{maps:eq:F.T11}; when the latter free
matrix has rank less than three the group has positive
rank and is infinite.
For determinant $\pm1$, elimination of the three free
variables leaves precisely $T$, even if the twists have
nonzero torsion components.
This proves the exact simple-connectivity criterion.

\item We give the global twist comparison to make its scope
precise.  Under \eqref{maps:eq:F.T12}, the finite fixed covector image
is one at each point, so choose integral fixed $u_i$ with
$\ell(u_i)=1$.  The primitive vertical fixed vector
$\lambda_i$ and invariant covector $\varphi_i$ satisfy
$\varphi_i(\lambda_i)=m_i$, by the assumed maximal local
pairings.  Formula \eqref{maps:eq:F.S5} makes the finite norm onto
$\mathbb Z\lambda_i$.
Consequently every vertical numerator component
$s_i\lambda_i$ is removed by the explicit cap translation
\eqref{maps:eq:F.T8}, whose boundary residue has $\ell$-coordinate zero.

\item For the complementary pair of pants the full relative
obstruction group for extending boundary translation
sections is
$H^2(P,\partial P;\Lambda)=\Lambda_\Gamma$.
Now $T=0$, so it is exactly $\mathbb Z$ detected by
$\ell$, not just a rational or free quotient.
Prescribing the identity translation at the cusp therefore
allows all the vertical finite changes to extend over $P$.
The same argument removes the zero-$\ell$ component of
any cusp translation clutching.
For determinant $-1$, write
$r_1=1+5a$, $r_2=-1+4b$; then \eqref{maps:eq:F.T12} says
$a+b+\kappa=0$.
The two integer changes of finite boundary marking and
the remaining cusp change thus also have zero total
relative obstruction and extend.
This gives the normalized twists $(u_1,-u_2)$.
For determinant $1$ the same argument gives their negatives.
The orientation-preserving fibre inversion extending over
the Kuhn fan, proved in \ref{maps:src:F.2.24}, identifies the two choices.
All maps used in the comparison are cap translations or
translations by smooth torus-bundle sections over $P$;
the prescribed cusp is retained throughout.
This proves the claimed complete twist normalization.

\end{enumerate}
\end{proof}




\begin{theorem}[a complete finite integral survivor test]\label{maps:src:F.2.29}
\StatementRecord{maps:src:F.2.29}{proved}{Historical events: Claude Fable 5.2 (Anthropic), effort high, 16 November 2026: carried as background without fault found. GPT-6 Astra (OpenAI), ChatGPT Pro mode, effort Pro, S6\_09\_PAPERS, 17 November 2026: targeted proof check. Also checked in this writing. Pre-continuation writing events: GPT-6 Astra (OpenAI), ChatGPT Pro mode, effort Pro, S6\_10\_PAPERS, 18 November 2026 02:43:58.020152+00:00: computation or formal certificate; targeted proof check. Continuation metadata events: GPT-6 Astra (OpenAI), ChatGPT Pro mode, effort Pro, S6\_10\_PAPERS continuation, 19 November 2026: used and rechecked in this run; status-vocabulary reconciliation; metadata reconciliation only; no new mathematical check. \emph{Pre-continuation recorded proof scope.} Claude Fable 5.2 (Anthropic), effort high, 16 November 2026; also checked in this writing by GPT-6 Astra (OpenAI), ChatGPT Pro mode, effort Pro, 18 November 2026}{the hypotheses, notation and definitions stated here; \cref{maps:src:F.2.20}; \cref{maps:src:F.2.22}; \cref{maps:src:F.2.25}; \cref{maps:src:F.2.28}; \cref{maps:src:F.2.6}; \cref{maps:eq:F.HC7}; \cref{maps:src:E.1.3}; \cref{maps:src:F.2.16}}


Take any rational pair in \ref{maps:src:F.2.20} and any of its integral,
primitive-cusp rank-six extensions supplied by \ref{maps:src:F.2.22}.
Write $M_1,M_2$ for geometric transport on fibre homology
$\Lambda=\mathbb Z^6$, in the peripheral convention
$M_0M_1M_2=I$.
First require the preliminary conditions
\[
 \Lambda_\Gamma=\mathbb Z,\qquad
 \ell(\Lambda^{M_i})=\mathbb Z,\qquad h_1=5,\quad h_2=4 .
 \tag{F.HC1}\label{maps:eq:F.HC1}
\]
Here $\ell$ is the primitive global invariant covector, and
$h_i$ is the integral fixed-vector/fixed-covector pairing
on $\ker\ell$ from \ref{maps:src:F.2.25}.
Failure of any condition excludes the datum.
If they hold, use the unique normalized affine filling of
\ref{maps:src:F.2.28}, with twists $(u_1,-u_2)$ and zero cusp clutching.
The following six finite integer-matrix tests are necessary
and sufficient for this filling to have $H_2=H_3=0$.

Let
\[
 T_i=M_i^{-t},\quad T_0=(T_1T_2)^{-1},\qquad
 G_q=\operatorname{coker}
 [\,\bigwedge^qT_1-I\mid\bigwedge^qT_2-I\,]
 \quad(1\le q\le4).
 \tag{F.HC2}\label{maps:eq:F.HC2}
\]
All $G_q$ are finite.  Let $P_q$ be the integral parabolic
cohomology group of the exterior local system
$\bigwedge^q\mathcal V$, for $q=2,4$.
Let $t_q\in G_{q-1}$, $2\le q\le5$, be the natural global
coinvariant image of the unique nonzero torsion suspension
class at the order-four cap, as specified below.
Let $c_4\in G_3$ be the natural image of the unique nonzero
torsion class in
$\operatorname{coker}(\bigwedge^3T_0-I)$.
Then the six tests are
\[
\begin{array}{c|ccc}
 &\text{first test}&\text{second test}&\text{third test}\\ \hline
 H^3=0&G_1=\langle t_2\rangle&P_2=0&t_3\ne0\\
 H^5=0&G_3=\langle t_4,c_4\rangle&P_4=0&t_5\ne0 .
\end{array}
\tag{F.HC3}\label{maps:eq:F.HC3}
\]
Every group, class, and equality here is computed by
finitely many integer kernels, Smith forms, and exterior
minors.  This is a decision criterion for every rational
parameter and every integral lattice in the complete
normal form; it makes no bounded-height assertion about
the rational parameter set.

\end{theorem}


\begin{proof}
\begin{enumerate}[label=\arabic*.]
\item We first establish the uniform local calculations needed
for the criterion, so the test does not assume that a
general finite matrix is globally conjugate to the example
in \ref{maps:src:F.2.6}.
The fixed vectors $u_i$ split the homology lattice as
$\mathbb Zu_i\oplus H_i$, where $H_i=\ker\ell$.
Write $D_i=M_i|_{H_i}$.
The maximal pairing $h_i=m_i$ gives
\[
 \operatorname{Norm}_{m_i}(D_i)=\lambda_i\varphi_i ,
 \tag{F.HC4}\label{maps:eq:F.HC4}
\]
with both factors primitive.

\item At five, work over $\mathbb Z_5$.
Modulo five the norm is $(D_1-I)^4$ and is nonzero.
Choose $q$ with $\operatorname{Norm}_5q$ primitive.
The five vectors $q,D_1q,\ldots,D_1^4q$ are independent
modulo five: otherwise the nilpotent Jordan block generated
by $q$ would have length less than five.
They are therefore a $\mathbb Z_5$ basis.
Thus $D_1$ is the regular five-cycle permutation locally.

\item At four, work over $\mathbb Z_2$.
The norm is $(D_2-I)^3$ modulo two and is nonzero.
The vectors $q,D_2q,D_2^2q,D_2^3q$ similarly generate
a primitive regular four-cycle submodule $Q$ of rank four.
The quotient is rank one with $D_2=-1$, by the
characteristic polynomial.
For a lift $z$ of its generator, put $a=(D_2+I)z\in Q$.
Since $D_2^4=I$,
\[
 (D_2^3-D_2^2+D_2-I)a=0.
\]
On the regular four-cycle module the kernel of the
displayed alternating norm is exactly $(D_2+I)Q$,
integrally.  One can see this by writing its four
coefficients; equivalently $D_2+I$ has three nonzero
Smith factors, all one, and the indicated rational image.
Subtracting a suitable element of $Q$ from $z$ makes
$D_2z=-z$.
Consequently
\[
 M_1\sim 1\oplus\operatorname{perm}_5\quad\text{over }\mathbb Z_5,
 \qquad
 M_2\sim1\oplus\operatorname{perm}_4\oplus(-1)
 \quad\text{over }\mathbb Z_2 .
 \tag{F.HC5}\label{maps:eq:F.HC5}
\]
These local comparisons apply to cap cohomology, not just
to abstract representation lattices.
Indeed the integral lattice of intertwiners is the kernel
of an integer matrix and its tensor with $\mathbb Z_p$
is the lattice of local intertwiners.
Approximate a local invertible intertwiner by an integral
one whose determinant is a $p$-adic unit.
Keeping the $u_i$ summand fixed produces an equivariant
torus isogeny of degree prime to $p$, preserving the
normalized translation $\pm u_i/m_i$.
The resulting cap and boundary coverings induce
isomorphisms on $\mathbb Z_p$ cohomology: the vertical
torus maps are cohomology isomorphisms over $\mathbb Z_p$,
so the mapping-torus Wang sequences give the assertion.
At primes not dividing $m_i$, transfer already makes
finite-cover pullback an isomorphism onto invariants.

\item Here are the full required integral results for these
two monomial models.
The ranks of their invariant exterior lattices in
degrees zero through six are
$(1,2,3,4,3,2,1)$.
The order-five cap has no cohomology torsion in degrees
two through five.
The order-four cap has exactly $\mathbb Z/2$ in each
of those four degrees.
This follows from the Wang sequence for the mapping
torus of the vertical $D_i$:
its torsion in degree $q$ is
\[
 \operatorname{Tor}\operatorname{coker}
       (\bigwedge^{q-1}D_i^t-I).
 \tag{F.HC6}\label{maps:eq:F.HC6}
\]
In particular it introduces no additional torsion primes.

\item Integral pullback from either finite cap is onto the
invariant lattice in degrees two through five.
For degree two this is precisely the $m_i/h_i=1$
suspension calculation of \ref{maps:src:F.2.25}.
For completeness, the determinants of the pairings between
the invariant lattices in complementary degrees
$(1,5),(2,4),(3,3)$ are respectively
$m_i,m_i^3,m_i^4$ in absolute value.
The degree-one cover image has index $m_i$.
The free downstairs complementary pairings are unimodular,
and pullback multiplies each pairing by the covering
degree $m_i$.
The middle pairing first gives degree-three index one.
The degree-two index one then gives degree-four index one;
the degree-one index $m_i$ gives degree-five index one.
These determinant identities and all the Wang Smith
forms are verified by the fixed, exhaustive computation
in \texttt{fourfolds\_local\_caps.py}, for both orders and
all exterior degrees zero through six.
Thus no unexamined finite matrix is being inferred from
an example.

\item We next identify the boundary maps.
Choose any unimodular matrix $P_2=[u_2,H_2]$ realizing
the splitting at four, and put
$J=P_2^{-t}|_{\mathbb Z^5}$ for the inclusion of vertical
cohomology.
Choose the unique nonzero torsion representative
\[
 w_q\in\operatorname{coker}
            (\bigwedge^{q-1}D_2^t-I),\qquad
 t_q=[(\bigwedge^{q-1}J)w_q]\in G_{q-1}.
 \tag{F.HC7}\label{maps:eq:F.HC7}
\]
A Smith transformation supplies an integral $w_q$.
Its norm is exactly zero, since it lies rationally in
the image of $\bigwedge^{q-1}D_2^t-I$.
The finite boundary coordinates of \ref{maps:src:F.2.6} satisfy
$\theta=a+4k$, $h=k$, hence $a=\theta-4h$.
The pullback of the suspension cocycle has $\theta$
component $w_q$ and $h$ component
$-\operatorname{Norm}_4(w_q)=0$.
This holds in every exterior degree, including $q=5$.
It is an integral cellular assertion for the monomial
model, where the torus cellular differentials vanish;
the prime-to-two isogeny comparison transports it to
every eligible cap.
The base connecting map therefore gives exactly the
class $t_q$ in \eqref{maps:eq:F.HC7}.
Different representative choices change it by a global
monodromy relation.  Any sign change is immaterial for
an order-two class.

\item There is one further boundary map, from the free
degree-four cusp kernel.
The primitive cusp is integrally conjugate to
$\left(\begin{smallmatrix}I&I\\0&I\end{smallmatrix}\right)$.
Its degree-three exterior difference has ten Smith
factors one and one Smith factor two; its cokernel is
$\mathbb Z^9\oplus\mathbb Z/2$.
The other exterior differences in degrees one, two,
four and five have torsion-free cokernels.
The exact matrices and a torsion representative are
saved by \texttt{fourfolds\_cusp\_boundary.py}.
Let $C$ be the smooth cusp cap and $Y=\partial C$.
The boundary Wang sequence gives
$\operatorname{Tor}H^4(Y)=\mathbb Z/2$.
In the Poincare--Lefschetz exact sequence its image in
\[
 H^5(C,Y)\cong H_3(C)=\mathbb Z^9
\]
is zero, so this torsion class lifts to $H^4(C)$.
Its restriction to an ordinary fibre is zero.
By \ref{maps:src:E.1.3} the kernel of integral specialization is
a primitive copy $K=\mathbb Z$ in degree four.
The lift is in $K$, and the image of this cyclic group
in $H^4(Y)$ contains the nonzero torsion class.
A cyclic image containing nonzero torsion cannot
also have a nonzero free part.
Its primitive generator therefore restricts to the
unique nonzero torsion class of $H^4(Y)$.
Consequently the cusp contribution to $d_2$ is
exactly $c_4$ as defined above.  In particular this
column must not be discarded merely because it
vanishes in the first examples.

\item It remains to prove that these are all the tests.
For the irreducible odd-dimensional quotient $W$
of $\mathcal V$, invariant alternating forms and
bivectors vanish: a nonzero one would give an
isomorphism $W\cong W^*$ with an odd-dimensional
nondegenerate alternating pairing.
The exact exterior sequences for
$0\to\mathbb Q\to\mathcal V_{\mathbb Q}\to W\to0$,
together with $\det W=1$, give
\[
 H^0(B,j_*\bigwedge^q\mathcal V)=0
       \quad(2\le q\le5),\qquad
 G_q\otimes\mathbb Q=0\quad(1\le q\le4).
 \tag{F.HC8}\label{maps:eq:F.HC8}
\]
For the endpoint assertions, the extension cannot split:
on $\mathbb Q\oplus W$ a square-zero nilpotent has rank at
most two, whereas the stipulated rank is three.
An invariant covector of $\mathcal V_{\mathbb Q}$ would
split this extension, since $W$ has no trivial quotient.
Thus $(\mathcal V_{\mathbb Q}^*)^\Gamma=0$, which gives
both $H^0(\bigwedge^5\mathcal V_{\mathbb Q})=0$ and
$G_1\otimes\mathbb Q=0$.
Here and below the first assertion is integral as well,
since the coefficient lattices are torsion free.
The local invariant ranks in exterior degrees two
and four are $3,3,9$ at the three punctures.
The Euler formula for the rank-fifteen local system
therefore gives $P_2\otimes\mathbb Q=P_4\otimes\mathbb Q=0$.
In degree one it similarly gives rational parabolic
cohomology zero.
The only global degree-one section is the invariant
line.  Simple connectivity and the universal
coefficient theorem make its free transgression
onto $H^2(B,\mathbb Z)$; equivalently the rational
second Betti number is zero, as in \ref{maps:src:F.2.6}.
Thus the normalized filling has rational $b_2=b_3=0$.

\item All specializations in exterior degrees two through
five are onto, by the local calculations and \ref{maps:src:E.1.3}.
Their only point-supported kernels are
\[
\begin{array}{c|cccc}
q&2&3&4&5\\ \hline
\text{kernel}&\mathbb Z/2&\mathbb Z/2&
                 \mathbb Z\oplus\mathbb Z/2&\mathbb Z/2 .
\end{array}
\tag{F.HC9}\label{maps:eq:F.HC9}
\]
The free summand in degree four is at the cusp; the
others are at four.
Consequently $H^1(B,R^qg_*\mathbb Z)=P_q$ in the
degrees used here, while
$H^2(B,R^qg_*\mathbb Z)=G_q$.
The latter equality also holds in degree one:
changing a sheaf at isolated points changes no
degree-two cohomology.  The equality with coinvariants
follows from compactly supported top cohomology on
the punctured sphere and the sequence $j_!\to j_*$.

\item The base has real dimension two.
Thus the only possible nonzero differential in the
integral Leray spectral sequence is
$d_2:E_2^{0,q}\to E_2^{2,q-1}$.
The complete associated graded groups in total
degrees three and five are therefore
\[
\begin{split}
 \operatorname{gr}H^3&:
   G_1/\langle t_2\rangle,\quad P_2,\quad
   \ker(\mathbb Z/2\xrightarrow{\,t_3\,}G_2),\\
 \operatorname{gr}H^5&:
   G_3/\langle t_4,c_4\rangle,\quad P_4,\quad
   \ker(\mathbb Z/2\xrightarrow{\,t_5\,}G_4).
\end{split}
\tag{F.HC10}\label{maps:eq:F.HC10}
\]
These are all finite.  A finite filtered group
vanishes exactly when every associated graded group
vanishes, regardless of extension classes.
Since $\pi_1=1$ and rational $b_2=b_3=0$, the universal
coefficient theorem gives
$H^3=\operatorname{Ext}(H_2,\mathbb Z)$, so
$H_2=0$ exactly when $H^3=0$.
Integral Poincare duality gives $H^5\cong H_3$.
This proves \eqref{maps:eq:F.HC3}.

\item For reproducibility, no sheaf calculation is left
as a black box in the integer test.
Given $A=\bigwedge^qT_1$, $B=\bigwedge^qT_2$ on a
free lattice $L$, form the saturated integer kernel
\[
 K_q=\ker[\,A-I\mid A(B-I)\mid -(AB-I)\,]\subset L^3.
 \tag{F.HC11}\label{maps:eq:F.HC11}
\]
Its image under
$(x,y,z)\mapsto((A-I)x,(B-I)y)$ is precisely the
integral parabolic cocycle lattice.
Quotient by the columns of
$\left(\begin{smallmatrix}A-I\\B-I\end{smallmatrix}\right)$
to obtain $P_q$.
This uses the actual integral images of peripheral
differences, not their saturations.
Smith transformations also supply every $w_q$, the
cusp order-two representative, and the tests for
their images and generated subgroups in $G_q$.

\item The program \texttt{fourfolds\_survivor\_test.py}
implements exactly these operations.
Its executed test set is precisely $v=0,1,2$ from
\ref{maps:src:F.2.16}.  For $v=0,1$ all six tests pass;
$G_1,\ldots,G_4$ are each $\mathbb Z/2$,
$P_2=P_4=0$, all four finite-cap images are nonzero,
and the cusp image is zero.
For $v=2$ they are each $\mathbb Z/6$,
$P_2=P_4=\mathbb Z/3$, and the first, second,
fourth and fifth tests fail.  The computation thus
checks both accepted and rejected data.

\item We finish by making the convention dictionary explicit.
The imported family matrices called $A,B$ in that
program are geometric cohomology pullbacks, so the
geometric homology inputs are $(A^t,B^t)$.
Their genuine left cohomology transport in \eqref{maps:eq:F.HC2}
is $(A^{-1},B^{-1})$, which is what the code uses.
For comparison with computations written using
$(A,B)$, simultaneous reversal of the two peripheral
generators sends
\[
 g_1\mapsto g_1^{-1},\qquad
 g_2\mapsto g_2^{-1},\qquad
 g_0=(g_1g_2)^{-1}\mapsto
     g_2g_1=g_1^{-1}(g_1g_2)g_1 .
 \tag{F.HC12}\label{maps:eq:F.HC12}
\]
This is an orientation-reversing marking of the
pair of pants fixing its three ends: each peripheral
conjugacy class is inverted.
Thus the new cusp is conjugate to the inverse old
cusp, not an independently chosen third matrix.
Inversion preserves invariant and coinvariant
lattices, and the indicated conjugating generator
acts trivially on global coinvariants.
All boundary classes in \eqref{maps:eq:F.HC3} have order two, so
reversal signs do not change the tests.
The implementation uses \eqref{maps:eq:F.HC2} directly and hence
does not depend on this comparison for its validity.

\end{enumerate}
\end{proof}




\begin{corollary}[constructive classification of all survivors]\label{maps:src:F.2.30}
\StatementRecord{maps:src:F.2.30}{proved}{Historical events: Claude Fable 5.2 (Anthropic), effort high, 16 November 2026: carried as background without fault found. GPT-6 Astra (OpenAI), ChatGPT Pro mode, effort Pro, S6\_09\_PAPERS, 17 November 2026: targeted proof check. Also checked in this writing. Pre-continuation writing events: GPT-6 Astra (OpenAI), ChatGPT Pro mode, effort Pro, S6\_10\_PAPERS, 18 November 2026 02:43:58.020152+00:00: targeted proof check. Continuation metadata events: GPT-6 Astra (OpenAI), ChatGPT Pro mode, effort Pro, S6\_10\_PAPERS continuation, 19 November 2026: used and rechecked in this run; status-vocabulary reconciliation; metadata reconciliation only; no new mathematical check. \emph{Pre-continuation recorded proof scope.} Claude Fable 5.2 (Anthropic), effort high, 16 November 2026; also checked in this writing by GPT-6 Astra (OpenAI), ChatGPT Pro mode, effort Pro, 18 November 2026}{the hypotheses, notation and definitions stated here; \cref{maps:src:F.2.20}; \cref{maps:src:F.2.22}; \cref{maps:eq:F.HC1}; \cref{maps:eq:F.HC3}; \cref{maps:eq:F.T12}; \cref{maps:src:F.3.11}; \cref{maps:src:F.2.7}; \cref{maps:src:F.2.14}; \cref{maps:src:F.2.25}; \cref{maps:src:F.2.27}; \cref{maps:src:F.2.28}; \cref{maps:src:F.2.29}; \cref{maps:src:E.1.3}; \cref{maps:src:F.2.15}; \cref{maps:src:F.2.5}; \S\ref{sec:V-known-inputs}; \cref{maps:src:F.4.1}}


For the marked Kuhn cap and affine translation clutchings
used throughout F.2, the following is a complete constructive
classification of the data with
$\pi_1=1$ and $H_2=H_3=0$, up to integral conjugacy of their
linear tuples and oriented diffeomorphism of their fillings.

There are no data with three genuine finite fillings.
With two fillings their actual affine orders are five and
four.  Order the order-five point first.
The complete, unbounded parameter set is the rational
projective normal form of \ref{maps:src:F.2.20} together with its
exceptional rational affine line.
For each parameter, run the terminating order,
conductor, stable-lattice, and primitive-extension algorithm
of \ref{maps:src:F.2.22}.  Discard precisely the integral classes failing
\eqref{maps:eq:F.HC1} or one of \eqref{maps:eq:F.HC3}.
The classes retained by this procedure are exactly all
surviving integral linear tuples.
For each retained class all simply connected affine
twists are exactly the choices satisfying \eqref{maps:eq:F.T12};
their fillings are all oriented diffeomorphic to its
normalized filling.
If ordered labels with the order-four point first are
desired, interchange the finite labels and recompute
$M_0=(M_1M_2)^{-1}$.

Every retained datum has
\[
 \pi_1=1,\qquad
 H_j=\begin{cases}
  \mathbb Z,&j=0,8,\\
  \mathbb Z^4,&j=4,\\
  0,&\text{otherwise},
 \end{cases}
 \tag{F.HC13}\label{maps:eq:F.HC13}
\]
and its oriented smooth manifold is
$3\mathbb{HP}^2\#(-\mathbb{HP}^2)$.
It has bounded-cusp holomorphic periods in the literal
sense of F.3, by \ref{maps:src:F.3.11}.
The classification here concerns the survivor subproblem.
It does not classify the nonsurviving tuples or compute
their full integral homology, as also requested in F.2.

\end{corollary}


\begin{proof}
\begin{enumerate}[label=\arabic*.]
\item The exclusions \ref{maps:src:F.2.7}, \ref{maps:src:F.2.14}, \ref{maps:src:F.2.25}, and \ref{maps:src:F.2.27}
leave exactly the stated minimal orders and rational
quotient problem.
The rational normal form is injective on conjugacy
classes, including the exceptional line.
The finite interval and extension-coset enumeration
in \ref{maps:src:F.2.22} gives every integral class over each rational
point, with all bounds proved for that input.
No distinct rational points can become integrally
conjugate, since an integral conjugacy induces a
rational quotient conjugacy.
Theorem \ref{maps:src:F.2.28} and the necessary-and-sufficient
criterion \ref{maps:src:F.2.29} prove both directions of the
retention rule and account for every affine twist.
This is also an exhaustive enumeration if desired:
enumerate primitive projective integer triples with
first nonzero coordinate positive by increasing height,
and enumerate reduced rational exceptional parameters
by increasing numerator and denominator height.
Every finite stage has finitely many inputs, each
terminates by \ref{maps:src:F.2.22}, and every rational input occurs.
No finite cutoff is imposed or needed.

\item Hurewicz makes the retained simply connected manifolds
three-connected.  The only singular contribution to
Euler number is the Kuhn cusp, whose value is six
by \ref{maps:src:E.1.3}; torus bundles and the free finite caps have
Euler number zero.
Integral Poincare duality and the universal coefficient
theorem then give \eqref{maps:eq:F.HC13}, including torsion freeness.

\item For the smooth classification, \ref{maps:src:F.2.15} supplies a
compatible complex structure at each finite cap, and
the cusp already has its complex model.
Choose their common fibre orientation to be the one
induced by the cusp.
Commuting complex structures at the finite caps can
be chosen in that component: their two-dimensional
fixed real summands allow the orientation to be
reversed when necessary.
The fibre of the bundle of orientation-compatible
complex structures retracts to $\mathbb{CP}^3$.
Its connectedness and simple connectivity extend
the prescribed boundary structures across the
two-dimensional pair of pants, exactly as in \ref{maps:src:F.2.5}.
A compatible horizontal distribution gives an
almost complex structure on the total manifold.
Theorem \S\ref{sec:V-known-inputs} therefore gives signature two,
intersection lattice $I_{3,1}$, and $c_2^2=2$.
Since $H^2=0$, $p_1/2=-c_2$ is primitive,
characteristic, and has square two.
The complete integral reflection, Wall-form and Kramer--Stolz argument in \ref{maps:src:F.4.1}
depends only on these almost complex characteristic
numbers and three-connectivity.
It applies without alteration and gives the claimed
exact oriented diffeomorphism.
Finally \ref{maps:src:F.3.11} applies to the resulting determinant-one
monodromy and primitive cusp.
It asserts the requested punctured-base periods with
bounded correction, not an additional unproved
simultaneous good extension at all finite points.

\end{enumerate}
\end{proof}




\begin{theorem}[complete fundamental-group data for every input]\label{maps:src:F.2.31}
\StatementRecord{maps:src:F.2.31}{proved}{Historical events: Claude Fable 5.2 (Anthropic), effort high, 16 November 2026: carried as background without fault found. GPT-6 Astra (OpenAI), ChatGPT Pro mode, effort Pro, S6\_09\_PAPERS, 17 November 2026: targeted proof check. Also checked in this writing. Pre-continuation writing events: GPT-6 Astra (OpenAI), ChatGPT Pro mode, effort Pro, S6\_10\_PAPERS, 18 November 2026 02:43:58.020152+00:00: computation or formal certificate; targeted proof check. Continuation metadata events: GPT-6 Astra (OpenAI), ChatGPT Pro mode, effort Pro, S6\_10\_PAPERS continuation, 19 November 2026: used and rechecked in this run; status-vocabulary reconciliation; metadata reconciliation only; no new mathematical check. \emph{Pre-continuation recorded proof scope.} Claude Fable 5.2 (Anthropic), effort high, 16 November 2026; also checked in this writing by GPT-6 Astra (OpenAI), ChatGPT Pro mode, effort Pro, 18 November 2026}{the hypotheses, notation and definitions stated here; \cref{maps:src:E.1.3}; \cite[Theorems 2.3--2.4, p.~425; Euler-characteristic comparison, pp.~427--428]{rootScott83}}


Let $k=2$ or $3$, and let an integral rank-six tuple satisfy the
finite-order and primitive maximal-cusp conditions of F.2.  No
irreducibility, determinant-one, or prospective-survivor hypothesis is
imposed.  Specify each actual finite affine action as
\[
 \tau_i(x)=M_ix+v_i/m_i,\qquad v_i\in\Lambda=\mathbb Z^6,
 \quad \tau_i^{m_i}=1\text{ on }\Lambda_{\mathbb R}/\Lambda,
 \tag{F.G1}\label{maps:eq:F.G1}
\]
and require its cyclic action to be free.  The numerator need not be
fixed by $M_i$.  Define its actual integral power vector and the cusp
lattice by
\[
 w_i=\frac1{m_i}\sum_{j=0}^{m_i-1}M_i^jv_i\in\Lambda^{M_i},
 \qquad U=\operatorname{im}\bigl((M_1\cdots M_k)^{-1}-I\bigr).
 \tag{F.G2}\label{maps:eq:F.G2}
\]
Use left conjugation transport for the peripheral generators $s_i$,
and record the integral translation clutching by $\kappa\in\Lambda$
in $s_1\cdots s_kx^\kappa=1$.  Then the full fundamental group is
\[
 \left\langle x^\lambda,s_1,\ldots,s_k\ \left|\
 \begin{array}{l}
 x^\lambda x^\mu=x^{\lambda+\mu},\quad x^u=1\ (u\in U),\\
 s_ix^\lambda s_i^{-1}=x^{M_i\lambda},\quad
 s_i^{m_i}=x^{w_i},\quad s_1\cdots s_kx^\kappa=1
 \end{array}\right.\right\rangle .                         \tag{F.G3}\label{maps:eq:F.G3}
\]
Only a basis of $\Lambda$ and of $U$ is needed in this finite
presentation.  The program below also determines exactly which
quotient of the apparent fibre subgroup injects.

Put
\[
 L=\sum_{g\in\langle M_1,\ldots,M_k\rangle}gU,
 \qquad A=\Lambda/L.                                      \tag{F.G4}\label{maps:eq:F.G4}
\]
This is an integral sum, without saturation.  For two points, put
$d=\gcd(m_1,m_2)$, $M=M_1$, and
$S_j(M)=I+M+\cdots+M^{j-1}$.  Choose integers $a,b$ with
$am_1+bm_2=d$, and set, in $A$,
\[
 \begin{split}
 w'_2&=-w_2-S_{m_2}(M)\kappa,\\
 e&=(m_2/d)w_1-(m_1/d)w'_2,\qquad
 z=aw_1+bw'_2,\qquad \overline A=A/\langle e\rangle .
 \end{split}                                             \tag{F.G5}\label{maps:eq:F.G5}
\]
There is an exact sequence and a complete marked extension description
\[
 1\longrightarrow\overline A\longrightarrow\pi_1
   \longrightarrow C_d\longrightarrow1,
 \qquad txt^{-1}=M(x),\quad t^d=z.                         \tag{F.G6}\label{maps:eq:F.G6}
\]
Thus its order is $d|\overline A|$ if $\overline A$ is finite;
otherwise it is infinite.  In particular $d=1$ gives precisely
the abelian group $\overline A$.

For three genuine points, so $m_i\ge2$, let
\[
 \Gamma=\langle q_1,q_2,q_3\mid q_i^{m_i}=1,
                         \ q_1q_2q_3=1\rangle,
 \qquad \epsilon=\sum_i1/m_i-1.                           \tag{F.G7}\label{maps:eq:F.G7}
\]
The action $q_i\mapsto M_i$ is well defined on $A$.
If $\epsilon\le0$, the exact sequence is
\[
                    1\longrightarrow A\longrightarrow\pi_1
                         \longrightarrow\Gamma\longrightarrow1. 
 \tag{F.G8}\label{maps:eq:F.G8}
\]
If $\epsilon>0$, $\Gamma$ is the finite spherical triangle group.
Using its action on $A$, form the exact transfer vector
\[
 e=\sum_{g\in\Gamma}g\kappa+
       \sum_{i=1}^3\ \sum_{g\in\Gamma/\langle q_i\rangle}g w_i.
 \tag{F.G9}\label{maps:eq:F.G9}
\]
The second sums use left coset representatives and are independent
of their choice.  In this case the exact sequence has kernel
$A/\langle e\rangle$, with the same lifted relations \eqref{maps:eq:F.G3}.
Consequently \eqref{maps:eq:F.G3}, \eqref{maps:eq:F.G6}, \eqref{maps:eq:F.G8}, and \eqref{maps:eq:F.G9} give every requested
fundamental group by finite operations, including the precise
injected abelian kernel and all extension data.  These statements
do not claim a conjugacy classification of the original tuples.

For every input, the full first integral homology is the Smith
cokernel of the integer matrix
\[
 \begin{pmatrix}
 M_1-I&\cdots&M_k-I&-w_1&\cdots&-w_k&\kappa\\
 0&\cdots&0&&\operatorname{diag}(m_1,\ldots,m_k)&&\boldsymbol1
 \end{pmatrix}.                                           \tag{F.G10}\label{maps:eq:F.G10}
\]
Here the bottom left zero block has $k$ rows.  The code records
the full matrix, not only its invariant factors.

\end{theorem}


\begin{proof}
\begin{enumerate}[label=\arabic*.]
\item \emph{Actual finite and cusp caps.}
The affine lift on $\mathbb R^6$ has $m_i$th power translation by
$w_i$, as the displayed sum shows.  It commutes with that power,
so $M_iw_i=w_i$.  The group generated by lattice translations
and this lift is the fundamental group of the reduced flat
six-manifold.  Its presentation consists of the lattice, its
$s_i$ conjugation relations, and $s_i^{m_i}=w_i$.
The finite cap deformation retracts onto that reduced manifold
by the equivariant radial contraction of its root disc.
Thus these are the actual cap relations, even for a numerator
which is not an invariant lattice vector.

\item For the primitive Kuhn cap, the compact toric strata in its
normalization and all their intersection components are simply
connected.  Applying van Kampen successively to their tubular
neighborhoods identifies the fundamental group of the central
fibre with that of its dual three-torus.  The nearby fibre map
kills the angular lattice $U$ and sends the long-period lattice
$\Lambda/U$ isomorphically onto this dual-torus group.  This is
also the explicit toric collapse used in \ref{maps:src:E.1.3}.  The primitive
index-one condition gives $U=\ker(M_0-I)$ and makes $\Lambda/U$
free of rank three.  A section through a smooth central stratum
bounds the base meridian by its parameter disc.  The smooth cap
retracts onto the central fibre, so it adds exactly these
relations, with no extra meridian relation left unspecified.

\item Over the punctured sphere the bundle has a section up to homotopy,
because the base retracts to a graph.  Its group is the appropriate
semidirect product.  Choose the cap sections and retain their
relative translation clutching as $\kappa$.  Van Kampen with the
actual boundary maps just described gives exactly \eqref{maps:eq:F.G3}.
Changing the transport convention causes no group ambiguity:
reverse the ordered list and replace each affine map by its
inverse, and replace $\kappa$ by $-\kappa$.  The new linear
matrices are $M_k^{-1},\ldots,M_1^{-1}$, the power vectors are
$-w_k,\ldots,-w_1$, and the cusp image lattice is unchanged.
The substitutions $s_i\mapsto s_{k+1-i}^{-1}$ give the explicit
isomorphism between these two presentations.  This supplies the
dictionary when geometric transport is recorded with opposite
path-composition order.

\item \emph{Integral closure and two-point elimination.}
In a normal abelian fibre subgroup, the normal closure of $U$
is precisely the integer sum \eqref{maps:eq:F.G4}.  Start with its column
Hermite basis and repeatedly append its images under every
$M_i$ and $M_i^{-1}$, then take Hermite normal form again.
If the rational rank increases, it can do so at most three
times.  At the first nonincrease the rational span is invariant
under every generator and its inverse, hence never increases
again.  If the current lattice has saturation index $D$ in
that span, every subsequent strict enlargement divides this
index by an integer at least two.  There can therefore be at
most $\lfloor\log_2D\rfloor$ such enlargements.  This proves
termination and gives the computable bound used by the code;
it never replaces a finite-index lattice by its saturation.

\item On $A$ one has $M_2=M^{-1}$ in the two-point case.  The product
relation eliminates $s_2=t^{-1}x^{-\kappa}$, where $t=s_1$.
Multiplying this word $m_2$ times gives
\[
        (t^{-1}x^{-\kappa})^{m_2}
            =t^{-m_2}x^{-S_{m_2}(M)\kappa}.
\]
Thus the two power relations become $t^{m_1}=w_1$ and
$t^{m_2}=w'_2$.  Both right sides are fixed by $M$ on $A$.
Also $M^{m_1}=M^{m_2}=1$, hence $M^d=1$ there.
Bezout elimination gives exactly $e=0$ and $t^d=z$ in
\eqref{maps:eq:F.G5}.  Replacing the Bezout pair changes $z$ by a multiple
of $e$ and therefore changes no extension datum.

\item To prove the asserted injection, construct
$\overline A\rtimes_M\mathbb Z$ and quotient by its central
element $t^dz^{-1}$.  This cyclic subgroup has nonzero
projection $d\mathbb Z$ to the infinite cyclic factor, so it
intersects $\overline A$ trivially.  The quotient has precisely
the presentation obtained by elimination.  This proves \eqref{maps:eq:F.G6},
including the absence of any further hidden fibre relation.

\item \emph{The three-point kernel and the finite transfer.}
For completeness this step computes the fibre injection; it
does not infer it from the appearance of a normal subgroup in
a presentation.  Let $\mathcal B$ be the orbifold classifying
space obtained from the three-holed sphere by attaching, at
its three boundary circles, the mapping cylinders of
$S^1\to BC_{m_i}$.  Over the punctured sphere take the bundle
with fibre $BA=K(A,1)$ and the specified action.  For each
boundary, the cyclic extension
\[
 E_i=\langle A,s_i\mid s_i a s_i^{-1}=M_i(a),\ s_i^{m_i}=w_i\rangle
\]
contains $A$ injectively: the same semidirect-product construction
as above applies because $M_i^{m_i}=1$ and $M_iw_i=w_i$.
The associated fibration $BA\to BE_i\to BC_{m_i}$ restricts
on $S^1$ to the required mapping-torus fibration.  Replace the
maps by their mapping-path fibrations before homotopy gluing,
so that the fibre is preserved on every attaching cylinder.
Glue these three fibrations, using the specified relative clutching
$\kappa$ on the punctured sphere.  Van Kampen gives \eqref{maps:eq:F.G3}
for the resulting total space.  Thus the homotopy sequence
of this fibration over $\mathcal B$ calculates precisely
the image of $A$ in the group under discussion.

\item When $\epsilon\le0$, the universal orbifold cover is the plane
or hyperbolic plane.  Its contractibility identifies
$\mathcal B$ with $B\Gamma$: indeed its Borel construction
has the just-described boundary-cylinder model and contractible
universal fibre.  The Euclidean and hyperbolic alternatives
are the oriented two-orbifold classification in
\cite[Theorems 2.3--2.4, p.~425; Euler-characteristic comparison, pp.~427--428]{rootScott83}.
Since $\pi_2(B\Gamma)=0$, the homotopy sequence proves \eqref{maps:eq:F.G8}.

\item When $\epsilon>0$, the universal orbifold cover is $S^2$.
Pulling the constructed $BA$ fibration to this finite cover
trivializes its local system.  Its transgression
$\pi_2(S^2)\to A$ is its integral clutching vector.
There is one lifted ordinary base cell for each $g\in\Gamma$,
contributing $g\kappa$, and one lifted cone cap for each left
coset of $\langle q_i\rangle$, contributing $gw_i$.
To check the sign, the base relation has fibre term $+\kappa$
and each cone relation has term $-w_i$; the oriented sphere
chain subtracts the cone cells from the base cell.  The result
is therefore the plus sign in both sums of \eqref{maps:eq:F.G9}.
Changing a coset representative by $q_i^j$ does not change
its contribution because $w_i$ is fixed.  The result is
$\Gamma$ invariant, as also checked by the code.  The long
homotopy sequence now says that the only additional kernel
is exactly $\langle e\rangle$.  This proves the spherical
assertion, with the orientation fixed by this cell calculation.

\item There is no implicit infinite search for a spherical group.
After sorting its orders the possibilities are
$(2,2,n),(2,3,3),(2,3,4),(2,3,5)$, with groups of orders
$2n,12,24,60$, respectively.  These follow directly from
the positive spherical-triangle angle condition.
The code constructs the dihedral, tetrahedral, octahedral,
or icosahedral permutation group and finds the marked
peripheral generators in that finite group.  Breadth-first
words supply their actions on $A$; every generator edge is
checked modulo the actual lattice $L$.  It then computes
the two finite sums in \eqref{maps:eq:F.G9} and their integer Smith form.

\item An order-one site is not a genuine cone.  Its matrix is the
identity and its relation is $s_i=x^{w_i}$.  Eliminate it
from \eqref{maps:eq:F.G3}; moving this vector to the right of the remaining
peripheral product replaces the clutching by
\[
       \kappa+(M_{i+1}\cdots M_k)^{-1}w_i.                 \tag{F.G11}\label{maps:eq:F.G11}
\]
This reduces the three-point case with such a site to the
two-point computation.  A primitive rank-three cusp cannot
remain with only one finite linear monodromy: a finite-order
unipotent matrix is the identity.  Thus no smaller case has
been omitted within the hypotheses of this theorem.

\item \emph{Freeness, abelianization, and executed certificates.}
Freeness of \eqref{maps:eq:F.G1} is itself decidable exactly.  A nontrivial
stabilizer in a cyclic group contains a subgroup of prime
order, so it is enough to check $\tau_i^{m_i/p}$ for every
prime $p\mid m_i$.  For an affine map $Rx+b$ of a torus,
a fixed point exists precisely when every integral covector
in $\ker(R^t-I)$ takes an integral value on $b$.
Smith form proves this criterion: the image of $R-I$ on
the torus is the connected subtorus annihilated by those
primitive covectors.  The code computes the powers in the
rational affine matrix group and performs these finite tests.

\item Abelianizing \eqref{maps:eq:F.G3} yields \eqref{maps:eq:F.G10}.  No cusp column was lost:
\[
 \operatorname{im}(M_0-I)
       \subset\sum_i\operatorname{im}(M_i-I)
\]
follows by telescoping the product on the right; inversion
preserves the image of its difference from the identity.
The Smith form therefore computes every free and torsion
summand of $H_1$ exactly.

\item The program \texttt{fourfolds\_vankampen.py} saves all input
matrices, power vectors, integral closure traces, group
extension data, and homology presentation matrices.
It checks the nine explicitly listed three-point parameters
of its original test set, four two-point order/clutching
choices, a genuinely nonfixed affine numerator, and a
finite-index closure step.  It also checks a nonabelian
rank-six two-point tuple, a spherical $(2,4,2)$ example
with injected kernel $\mathbb Z^2\oplus\mathbb Z/4$, and
the infinite perfect hyperbolic example with orders
$(5,4,19)$.  The nonabelian test deliberately belongs to
the full $\operatorname{GL}_6$ scope and has orientation-
reversing finite matrices; it is not claimed to admit a
complex compactification.  All bounds above are derived
from each input, rather than imposed as a search cutoff.

\end{enumerate}
\end{proof}




\begin{theorem}[a complete three-point logarithmic family]\label{maps:src:F.2.32}
\StatementRecord{maps:src:F.2.32}{proved}{Historical events: Claude Fable 5.2 (Anthropic), effort high, 16 November 2026: carried as background without fault found. GPT-6 Astra (OpenAI), ChatGPT Pro mode, effort Pro, S6\_09\_PAPERS, 17 November 2026: targeted proof check. Also checked in this writing. Pre-continuation writing events: GPT-6 Astra (OpenAI), ChatGPT Pro mode, effort Pro, S6\_10\_PAPERS, 18 November 2026 02:43:58.020152+00:00: computation or formal certificate; targeted proof check. Continuation metadata events: GPT-6 Astra (OpenAI), ChatGPT Pro mode, effort Pro, S6\_10\_PAPERS continuation, 19 November 2026: used and rechecked in this run; status-vocabulary reconciliation; metadata reconciliation only; no new mathematical check. \emph{Pre-continuation recorded proof scope.} Claude Fable 5.2 (Anthropic), effort high, 16 November 2026; also checked in this writing by GPT-6 Astra (OpenAI), ChatGPT Pro mode, effort Pro, 18 November 2026}{the hypotheses, notation and definitions stated here; \cref{maps:src:F.3.7}; \cref{maps:src:F.2.28}; \cref{maps:eq:F.LT4}; \cref{maps:eq:F.LT5}; \cref{maps:src:F.2.5}; \cref{maps:eq:F.HC8}; \cref{maps:eq:F.HC9}; \cref{maps:eq:F.LT8}; \cref{maps:src:F.2.29}; \cref{maps:eq:F.LT6}; \cref{maps:eq:F.LT9}}


Start with the compact complex fourfold of \ref{maps:src:F.3.7} and its
normalized order-five and order-four affine fillings.
Let $\ell$ be the primitive invariant covector on fibre
homology, with the sign chosen in \ref{maps:src:F.2.28}.
At any additional ordinary fibre choose
$v\in\mathbb Z^6$ and an integer $m\ge2$ such that
\[
 r=\ell(v),\qquad \gcd(r,m)=1.
\tag{F.LT1}\label{maps:eq:F.LT1}
\]
Perform the marked logarithmic transformation with
translation $v/m$ described below.
This gives a smooth compact complex fourfold $X_{m,v}$
with Euler number six and three genuine finite affine
fillings, whose linear orders are five, four and one
and whose actual affine orders are five, four and $m$.
Its fundamental group is
\[
 \pi_1(X_{m,v})=
 \langle a,b,c,h\mid h\text{ central},\
   a^5=h,\ b^4=h^{-1},\ c^m=h^r,\ abc=1\rangle .
\tag{F.LT2}\label{maps:eq:F.LT2}
\]
This is the group of the oriented Seifert manifold
with invariants
$(0;(5,-1),(4,1),(m,-r))$, in the convention
$x_i^{\alpha_i}h^{\beta_i}=1$ and product $x_i=1$.
Its base orbifold is hyperbolic, and its central fibre
$h$ has infinite order.

Put $C_{m,r}=\mathbb Z/(20r-m)\mathbb Z$, with
$\mathbb Z/0\mathbb Z=\mathbb Z$, and let
$\varepsilon=1$ when $20r=m$ and zero otherwise.
All the integral homology groups are
\[
\begin{array}{c|c}
j&H_j(X_{m,v};\mathbb Z)\\ \hline
0,8&\mathbb Z\\
1,6&C_{m,r}\\
2&\mathbb Z^\varepsilon\oplus(\mathbb Z/m)^5\\
3&(\mathbb Z/m)^{10}\\
4&\mathbb Z^4\oplus(\mathbb Z/m)^{10}\\
5&(\mathbb Z/m)^5\\
7&\mathbb Z^\varepsilon .
\end{array}
\tag{F.LT3}\label{maps:eq:F.LT3}
\]
Under \eqref{maps:eq:F.LT1} the exceptional free-rank case is exactly
$(m,r)=(20,1)$.
In particular $(m,r)=(19,1)$ and $(21,1)$ have
$H_1=0$ but nontrivial, infinite fundamental groups.
More generally this occurs whenever $|20r-m|=1$.
The result treats every vertical component of $v$
in $\ker\ell$, not only a chosen fixed representative.

\end{theorem}


\begin{proof}
\begin{enumerate}[label=\arabic*.]
\item We first specify the complex construction and its marking.
Over a small disc at the ordinary value use a holomorphic
period marking
$T_t=\mathbb C^3/\Lambda_t$.
Let $\omega_v(t)$ be the holomorphic period vector
corresponding to $v$.
On the root disc $t=s^m$ quotient the pulled-back family by
\[
 (s,Z)\longmapsto
   (\zeta_m s,\ Z+\omega_v(s^m)/m).
\tag{F.LT4}\label{maps:eq:F.LT4}
\]
Condition \eqref{maps:eq:F.LT1} makes $v$ have exact order $m$
in $\mathbb Z^6/m\mathbb Z^6$.
Every nontrivial power therefore acts freely, including
over $s=0$.
The quotient is smooth, has a proper map to the disc,
and has central fibre of multiplicity $m$.
On the punctured disc its holomorphic identification with
the original family is
\[
 (s,Z)\longmapsto
 \left(s^m,\ Z-\frac{\log s}{2\pi i}\omega_v(s^m)\right).
\tag{F.LT5}\label{maps:eq:F.LT5}
\]
Changing the branch of $\log s$ changes the second
coordinate by a lattice period.
Changing $s$ by $\zeta_m s$ changes its logarithm by
$2\pi i/m$ and cancels the translation in \eqref{maps:eq:F.LT4}.
Thus \eqref{maps:eq:F.LT5} is a well-defined holomorphic isomorphism
of the quotient with the original torus family off zero.
Gluing on a fixed annulus supplies the asserted compact
complex manifold: its charts are the original charts
and the free finite quotient charts \eqref{maps:eq:F.LT4}, and it is
the union of compact pieces over closed base regions.
The gluing is an isomorphism on collars, so the total
space is Hausdorff and smooth.
The Euler number is unchanged, since the removed
ordinary torus cap, the new cap, and their boundary
all have Euler number zero.
This marking fixes the new relation $c^m=v$ and
zero extra cusp clutching.

\item In the original datum the invariant normal closure
of the cusp angular lattice is exactly $\ker\ell$,
as proved in \ref{maps:src:F.2.5}.
Adding the identity linear generator does not change
that lattice.
The fibre image in van Kampen therefore reduces to
the central generator $h$, while the original two
finite relations remain $a^5=h$, $b^4=h^{-1}$.
The new relation is $c^m=h^{\ell(v)}$ and the base
relation is $abc=1$.
This proves \eqref{maps:eq:F.LT2}, including the absence of any
dependence on the vertical part of $v$.

\item One may verify the Seifert identification directly:
take a product circle bundle on a three-holed sphere
and attach its three solid tori with primitive meridian
slopes $(5,-1),(4,1),(m,-r)$.
They are primitive by \eqref{maps:eq:F.LT1}, and van Kampen gives
exactly \eqref{maps:eq:F.LT2}.
Its orbifold base has
$1/5+1/4+1/m<1$ for every $m\ge2$.
Pulling back the circle fibration to its contractible
hyperbolic orbifold cover gives a circle bundle on
the plane, hence the trivial circle bundle.
The subsequent universal cover is the plane times
$\mathbb R$, and the fibre subgroup is infinite cyclic.
This also proves directly that none of these groups
is trivial.
The rational Euler number in these conventions is
$r/m-1/20$.

\item In abelianization the relations $5a=h$, $4b=-h$
give $a=4x$, $b=-5x$, $h=20x$ because five and four
are coprime.
The product relation gives $c=x$.
The remaining relation is $(m-20r)x=0$.
Thus $H_1=C_{m,r}$ exactly, including its integral
generator and the exceptional infinite-cyclic case.

\item We now compute the higher integral groups, including
all extensions.
Write $\mathcal R^q$ for the higher direct-image
sheaf of the original fourfold and
$\mathcal R_m^q$ for that of $X_{m,v}$, after identifying
the bases and their ordinary local systems.
At the new point the reduced finite quotient is a
torus with lattice
$\mathbb Z^6+\mathbb Z(v/m)$.
Its covering by the original torus is cyclic of
degree $m$.
The degree-one pullback has Smith factors
$1,1,1,1,1,m$, so its degree-$q$ cokernel is
\[
 C^q=(\mathbb Z/m)^{\binom5{q-1}},
 \qquad 1\le q\le6 .
\tag{F.LT6}\label{maps:eq:F.LT6}
\]
There is no cohomology torsion in this torus cap.
The sheaf comparison, which does not assert a map
between the total manifolds, is therefore
\[
 0\longrightarrow\mathcal R_m^q
  \longrightarrow\mathcal R^q
  \longrightarrow C^q_{\{p\}}\longrightarrow0 .
\tag{F.LT7}\label{maps:eq:F.LT7}
\]
It is the identity away from the new point and the
covering image at that point.

\item For $2\le q\le5$, all original global sections of
$\mathcal R^q$ are supported at the original finite
point or the cusp, by \eqref{maps:eq:F.HC8}--\eqref{maps:eq:F.HC9}.
Their map to $C^q_{\{p\}}$ is zero.
Consequently their $H^0$ groups are unchanged,
their $H^2$ groups are unchanged, and
\[
 0\longrightarrow C^q
  \longrightarrow H^1(\mathcal R_m^q)
  \longrightarrow H^1(\mathcal R^q)
  \longrightarrow0 .
\tag{F.LT8}\label{maps:eq:F.LT8}
\]
The original degree-one base cohomology of these
sheaves is zero for $q=2,4,5$ and is $\mathbb Z^3$
for $q=3$.
The assertion for $q=5$ also follows immediately
from the original $H^6=0$: its bidegree $(1,5)$
has no incoming or outgoing differential.
Thus \eqref{maps:eq:F.LT8} gives $C^q$ in degrees two, four and
five, and $\mathbb Z^3\oplus C^3$ in degree three.

\item The original four order-two torsion transgressions
remain isomorphisms.
Indeed their boundary suspension representatives and
their images in global coinvariants are the ones in
\ref{maps:src:F.2.29}; adding the identity monodromy changes none
of those coinvariant quotients.
These point-supported boundary maps do not depend
on the free degree-one transgression, which changes
under the new logarithmic gluing.
The free degree-four cusp column likewise remains
zero, since its boundary torsion representative has
zero image in that same unchanged global quotient.
It follows that the full filtrations in degrees
three and five have just one nonzero term:
\[
 H^3=C^2=(\mathbb Z/m)^5,\qquad
 H^5=C^4=(\mathbb Z/m)^{10}.
\tag{F.LT9}\label{maps:eq:F.LT9}
\]
In degree four there is an exact sequence
\[
 0\longrightarrow
       \mathbb Z^3\oplus(\mathbb Z/m)^{10}
  \longrightarrow H^4
  \longrightarrow\mathbb Z\longrightarrow0.
\]
It splits as an abelian-group extension because
its quotient is free, giving the stated middle group.
Thus no torsion extension has been inferred from
rational ranks.

\item For the remaining free ranks,
$H^1(X_{m,v})=\operatorname{Hom}(C_{m,r},\mathbb Z)$
has rank $\varepsilon$.
The degree-one global section group has rank one,
and the only rational differential involving it is
$d_2:H^0(\mathcal R_m^1)\to H^2(\mathcal R_m^0)$.
Its kernel has rank $\varepsilon$, hence its
cokernel also has rank $\varepsilon$.
The degree-one parabolic group is rationally zero,
and there are no rational degree-two global sections.
Therefore $b_2=\varepsilon$.
For an integral check of this degree-one row, its
finite sheaf cokernel is
$\mathbb Z/\gcd(20,m)$: the original invariant
section restricts as $20\ell$, and its value at
the new covering cokernel is $20r$ modulo $m$.
This finite group changes no rational rank.

\item Poincare duality gives
$H_3=H^5$, $H_4=H^4$, and $H_5=H^3$.
The universal coefficient theorem now identifies
the torsion of $H_2$ with $H^3$ and its free rank
with $\varepsilon$.
It also gives $H^2=C_{m,r}$: when $20r\ne m$ this
is $\operatorname{Ext}(H_1,\mathbb Z)$; when
$20r=m$ it is the free rank-one group.
Duality gives $H_6=C_{m,r}$ and
$H_7=\mathbb Z^\varepsilon$.
This proves the complete table \eqref{maps:eq:F.LT3}.

\item The program \texttt{fourfolds\_three\_point\_family.py}
checks the original four torsion images and zero
cusp image once, then computes the actual cyclic
covering Smith forms in all exterior degrees one
through six and the exact van Kampen abelianization
for the thirty specified pairs
$m\in\{2,3,4,5,6,8,10,19,20,21\}$,
$r\in\{-1,1,m+1\}$.
It includes a nonzero vertical component in every
test vector.
The unbounded assertion follows from \eqref{maps:eq:F.LT6}--\eqref{maps:eq:F.LT9},
not from extrapolation of those finite tests.

\end{enumerate}
\end{proof}




\begin{theorem}[arbitrary finite affine caps]\label{maps:src:F.2.33}
\StatementRecord{maps:src:F.2.33}{proved}{Historical events: Claude Fable 5.2 (Anthropic), effort high, 16 November 2026: carried as background without fault found. GPT-6 Astra (OpenAI), ChatGPT Pro mode, effort Pro, S6\_09\_PAPERS, 17 November 2026: targeted proof check. Also checked in this writing. Pre-continuation writing events: GPT-6 Astra (OpenAI), ChatGPT Pro mode, effort Pro, S6\_10\_PAPERS, 18 November 2026 02:43:58.020152+00:00: computation or formal certificate; targeted proof check. Continuation metadata events: GPT-6 Astra (OpenAI), ChatGPT Pro mode, effort Pro, S6\_10\_PAPERS continuation, 19 November 2026: used and rechecked in this run; status-vocabulary reconciliation; metadata reconciliation only; no new mathematical check. \emph{Pre-continuation recorded proof scope.} GPT-6 Astra (OpenAI, ChatGPT Pro mode), effort Pro, 2026-11-17 UTC; GPT-6 Astra (OpenAI, ChatGPT Pro mode, effort Pro),17November2026,/root; Claude Fable 5.2 (Anthropic), effort high, 16 November 2026; also checked in this writing by GPT-6 Astra (OpenAI), ChatGPT Pro mode, effort Pro, 18 November 2026}{\cref{maps:src:F.2.31}; the displayed affine actions and bar-chain maps}


Let $g(x)=Mx+v/m$ be a free cyclic affine action of actual order
$m\geq2$ on $T^r$, with $M\in\mathrm{GL}(r,\mathbb Z)$ and
$v\in\mathbb Z^r$.  For every specified degree bound $d$, there is
an explicit terminating integer algorithm producing chains and
boundary maps through degree $d+1$ for the pair
\[
 B=(D^2\times T^r)/\langle(s,x)\mapsto
 (e^{2\pi i/m}s,g(x))\rangle,
 \qquad Y=\partial B.                                      \tag{F.CA1}\label{maps:eq:F.CA1}
\]
The cap complex has at most $2^r$ generators in each degree and
the boundary complex at most $2^{r+1}$.  These are complexes of
finite rank in each degree; only their indicated finite prefixes
are required to compute $H_j(B)$, $H_j(Y)$ and $H_j(B,Y)$ for
$j\leq d$.  Their maps retain the affine translations and all
integral higher coherence terms.  In particular, the construction
does not replace an equivariant torus chain complex by its
homology without recording the correction terms.

\end{theorem}
\begin{proof}
\begin{enumerate}[label=\arabic*.]
\item We give the complete recursive algorithm.  All chains below are
finite integer dictionaries, so the displayed operations are
effective over $\mathbb Z$.

\item First consider the normalized bar resolution of $\mathbb Z^r$
over $R=\mathbb Z[t_1^{\pm1},\ldots,t_r^{\pm1}]$ and its Koszul
resolution $K=R\otimes\bigwedge^*\mathbb Z^r$.
The latter differential is contraction with $(t_i-1)_i$.
Write $t^a=t_1^{a_1}\cdots t_r^{a_r}$ and put
$S_n(t)=(t^n-1)/(t-1)$, interpreted as its finite integral
Laurent sum for positive, zero or negative $n$.
The following $\mathbb Z$-linear contraction of the augmented
Koszul complex fixes a completely explicit choice:
\[
 s_K(t^a e_I)=
 \sum_{j<\min I}
 t_{j+1}^{a_{j+1}}\cdots t_r^{a_r}S_{a_j}(t_j)
 e_j\wedge e_I,                                           \tag{F.CA2}\label{maps:eq:F.CA2}
\]
where $\min\varnothing=r+1$ and indices in $I$ are increasing.
The bar contraction is
$s_B(t^a[g_1|\cdots|g_q])=[a|g_1|\cdots|g_q]$; terms with a
zero bar entry are zero.  Set $p_0=1$ and define the
$R$-linear projection and homotopy on coefficient-one bar
generators recursively by
\[
 p_q=s_Kp_{q-1}d_B,
 \qquad
 h_q=s_B(1-ip_q-h_{q-1}d_B).                               \tag{F.CA3}\label{maps:eq:F.CA3}
\]
Here $i(e_I)$ is the alternating sum of the bars of the ordered
unit vectors in $I$.  The recursion decreases degree.  The
identities for the two contractions give, by induction,
$pi=1$, $dh+hd=1-ip$, $ph=hi=h^2=0$.
For example $s_K^2=s_B^2=0$, $ps_B=s_Kp$, and the alternating
unit-vector formula gives $pi=1$; substituting these into
\eqref{maps:eq:F.CA3} proves the three side identities in the same induction.
Augmenting $R\to\mathbb Z$ gives a strong chain contraction
of the ordinary torus bar complex to $\bigwedge^*\mathbb Z^r$.

\item Translations must enter before this contraction.  Let
$b_j=\sum_{a=0}^{j-1}M^av/m$ and let $p_j$ be its coordinatewise
fractional part.  The $m$ points $p_j$ are distinct because the
action is free.  Use the full fundamental groupoid of $T^r$
on these points: an arrow from object $a$ to object $b$ is an
integer vector $\lambda$, represented by the affine path from
$p_a$ to $p_b+\lambda$.  Composition adds the integer vectors.
The generator acts strictly by shifting objects and, for its
$s$th power, by
\[
 \lambda\longmapsto M^s\lambda+z_s(b)-z_s(a),
 \qquad z_s(a)=M^sp_a+b_s-p_{a+s}\in\mathbb Z^r.             \tag{F.CA4}\label{maps:eq:F.CA4}
\]
The subscripts on objects are taken modulo $m$.
The nerve of this groupoid is equivariantly homotopy equivalent
to the torus: the simplex with successive arrow labels
$\lambda_i$ maps affinely to the torus with lifted vertices
$p_{a_j}+\sum_{i\leq j}\lambda_i$.
Forgetting the objects contracts its chains to the ordinary
torus bar complex.  An explicit homotopy is the negative
alternating prism which, at each possible position, inserts
the zero-labelled arrow to object zero and replaces all
subsequent objects by zero.  Combining this with \eqref{maps:eq:F.CA3}
gives a contraction $(i,p,h)$ of the groupoid chains to
$\bigwedge^*\mathbb Z^r$ with the same side identities.

\item Use the normalized cyclic bar construction on this strict
groupoid action.  Its differential is $d_v+\delta$, where
$\delta$ is the ordinary group-bar boundary, multiplied by
$(-1)^q$ on vertical degree $q$; its first face applies \eqref{maps:eq:F.CA4}.
The vertical contraction transfers this to cyclic bar chains
with exterior coefficients by the finite formulas
\[
 I=(1+h\delta)^{-1}i,\qquad
 P=p(1+\delta h)^{-1},\qquad D=p\delta I.                  \tag{F.CA5}\label{maps:eq:F.CA5}
\]
Inverses mean their alternating geometric series.  On a chain
of cyclic bar degree $p$ each series has at most $p+1$ terms,
since $\delta$ lowers that degree.  Multiplication and the
contraction identities directly give $(d_v+\delta)I=ID$ and
$P(d_v+\delta)=DP$,
and a chain homotopy inverse.  Equivalently these are the
finite homological perturbation identities, here derived by
expanding the finite series.  The first term of $D$ is the
cyclic bar differential for the exterior action of $M$.
Every remaining term lowers cyclic bar degree by at least two.

\item For completeness, the second contraction used in the algorithm
is also explicit.  The cyclic periodic resolution has one
$\mathbb Z[C_m]$ generator in each degree, with differential
$g-1$ in odd degree and $1+g+\cdots+g^{m-1}$ in even degree.
On $g^a$, $0\leq a<m$, its augmented contraction is
$\sum_{j=0}^{a-1}g^j$ in even degree and is $1$ precisely when
$a=m-1$ in odd degree, zero otherwise.
Apply \eqref{maps:eq:F.CA3} with this contraction to compare it with the
cyclic bar resolution.  Tensor its coefficients with the
exterior action of $M$.  Transferring the higher terms of $D$
once more by \eqref{maps:eq:F.CA5} terminates: each higher term followed by
the horizontal homotopy decreases cyclic degree by at least
one.  The resulting basis in total degree $n$ is
$\{(p,I):p+|I|=n,\ I\subseteq\{1,\ldots,r\}\}$.
This proves the claimed $2^r$ bound and termination in every
specified degree.

\item \par\medskip\noindent\textbf{Explicit comparison with the periodic cyclic resolution.}\quad 
\label{maps:mapsF:repair:cyclic-inclusion}
Write $C_m=\langle g\rangle$, $R_m=\mathbb Z[C_m]$, and let
$e_n$ generate the periodic resolution in degree $n$.
Let $B_*$ be the normalized $R_m$-bar resolution, with the
insertion contraction $s_B$.  Besides the periodic contraction
$s_P$ just specified, fix the inclusion recursively by
\[
 i_0(e_0)=[],\qquad
 i_n(e_n)=s_B\bigl(i_{n-1}(d_Pe_n)\bigr),
 \quad n>0,
\]
and extend it $R_m$-linearly.  This is a finite formula.  More
explicitly, with every summation index between $1$ and $m-1$,
\[
\begin{split}
 i_{2j}(e_{2j})
   &=\sum_{a_1,\ldots,a_j}
       [g^{a_1}|g|\cdots|g^{a_j}|g],\\
 i_{2j+1}(e_{2j+1})
   &=\sum_{a_1,\ldots,a_j}
       [g|g^{a_1}|g|\cdots|g^{a_j}|g].
\end{split}
\]
The empty sum of indices gives $i_0(e_0)=[]$ and
$i_1(e_1)=[g]$.  The contraction identity for $s_B$ proves
$d_Bi=id_P$ by induction: the chain to which $s_B$ is applied
is a cycle, since the preceding comparison commutes with
differentials and $d_P^2=0$.

\item Define $p$ and $h$ on coefficient-one bar generators by the
recursion already used for the torus comparison, substituting
$s_P$ for $s_K$, and extend them $R_m$-linearly.  Their values
can be checked without an unspecified comparison theorem.
For $1\leq a_i<m$, put
$\epsilon(a,b)=\lfloor(a+b)/m\rfloor$.  The projection is
\[
\begin{split}
 p_{2j}[g^{a_1}|\cdots|g^{a_{2j}}]
   &=\prod_{\nu=1}^{j}\epsilon(a_{2\nu-1},a_{2\nu})\,e_{2j},\\
 p_{2j+1}[g^{a_1}|\cdots|g^{a_{2j+1}}]
   &=\left(\sum_{b=0}^{a_1-1}g^b\right)
       \prod_{\nu=1}^{j}\epsilon(a_{2\nu},a_{2\nu+1})\,e_{2j+1}.
\end{split}
\]
These expressions follow by induction from
$p_n=s_Pp_{n-1}d_B$: the terms other than the first bar face
are annihilated after the preceding contraction identity is
used.  In $pi(e_{2j})$ only the summand with every
$a_\nu=m-1$ survives.  The same observation, with initial
factor $1$, proves $pi(e_{2j+1})=e_{2j+1}$.
Thus $pi=1$ in every degree.

\item The recursion
$h=s_B(1-ip-hd_B)$ gives $d_Bh+hd_B=1-ip$ by induction.
Moreover $p s_B=s_Pp$ follows from the displayed projection
formulas, while $s_P^2=s_B^2=0$ and $s_Bi=0$ follow directly
from their definitions.  Substitution in the same recursion
gives $ph=hi=h^2=0$.  We therefore have the strong comparison
required for the second finite perturbation series, including
its inclusion, projection, homotopy and side identities.
After tensoring with the exterior coefficient module, the
same identities hold.  Every higher perturbation followed by
the horizontal homotopy lowers cyclic degree, so this added
description changes neither the termination proof nor the
$2^r$ generator bound.

\item The free action makes the Borel construction equivalent to
$T^r/C_m$, and radial contraction identifies this with $B$.
For $Y$ use exactly the same algorithm on $T^{r+1}$ with
matrix $\operatorname{diag}(1,M)$ and numerator $(1,v)$.
The projection forgetting its first circle is equivariant and
represents $Y\to B$ under the radial contraction.
More generally an equivariant integral matrix $L$ acts on
groupoid arrows as
$\lambda\mapsto L\lambda+k_b-k_a$, where
$k_a=Lp_a-p'_a\in\mathbb Z^{r'}$.
Transferring this strict bar map by $P'LI$ gives the required
integer chain map.  Its mapping cone computes the pair.
Every series is finite with the proved degree bound, and
integer kernels and Smith normal forms finish the calculation.
\end{enumerate}
\end{proof}


\emph{Executed certificates.}
The programs \path{code/fourfolds_bar_koszul.py} and
\path{code/fourfolds_affine_cap_chains.py} implement these
formulas, including the prime-subgroup freeness check of \ref{maps:src:F.2.31}.
The first checks all 281 abelian and 498 cyclic bar simplices
in its explicitly stated alphabets and degree ranges.
The second checks translated circles of orders two and three,
two affine presentations of the Klein bottle (one with a
nonfixed numerator), and the three-dimensional half-turn
manifold.  The separate rank-six program checks every degree
zero through seven for the order-five and order-four caps,
against the independent integral Wang sequence of their
five-torus mapping-torus descriptions.
The boundary-map program checks both translated circles and
the Klein bottle, including the full chain-map identities and
the integral Thom-shift homology of each cone.  These are
finite verification sets, not an asserted exhaustive test over
all input matrices; the terminating formulas prove generality.







\begin{theorem}[the marked Kuhn cusp pair]\label{maps:src:F.2.34}
\StatementRecord{maps:src:F.2.34}{proved}{Historical events: Claude Fable 5.2 (Anthropic), effort high, 16 November 2026: carried as background without fault found. GPT-6 Astra (OpenAI), ChatGPT Pro mode, effort Pro, S6\_09\_PAPERS, 17 November 2026: targeted proof check. Also checked in this writing. Pre-continuation writing events: GPT-6 Astra (OpenAI), ChatGPT Pro mode, effort Pro, S6\_10\_PAPERS, 18 November 2026 02:43:58.020152+00:00: computation or formal certificate; targeted proof check. Continuation metadata events: GPT-6 Astra (OpenAI), ChatGPT Pro mode, effort Pro, S6\_10\_PAPERS continuation, 19 November 2026: used and rechecked in this run; status-vocabulary reconciliation; metadata reconciliation only; no new mathematical check. \emph{Pre-continuation recorded proof scope.} GPT-6 Astra (OpenAI, ChatGPT Pro mode), effort Pro, 2026-11-17 UTC; GPT-6 Astra (OpenAI, ChatGPT Pro mode, effort Pro),17November2026,/root; Claude Fable 5.2 (Anthropic), effort high, 16 November 2026; also checked in this writing by GPT-6 Astra (OpenAI), ChatGPT Pro mode, effort Pro, 18 November 2026}{\cref{maps:src:F.2.33}; \cite[Theorem 4.1.1 and Lemmas 4.1.2--4.1.3, pp.~480--482]{ffAchingerOgus20}; \cite[Theorem 8.7, pp.~57--59, and Corollary 8.8, p.~59, attached source build]{manYang26}}


The fixed primitive rank-three Kuhn cusp admits an explicit
finite integer chain complex for its cap $W$, its boundary
$Y$, and the inclusion $Y\to W$.  It is obtained from a
category with 26 objects, 170 nonidentity morphisms, and
respectively $26,170,288,144$ strict flags of lengths
$0,1,2,3$.  The construction records the full integral
collapse map.  In particular it gives
\[
\begin{array}{c|rrrrrrrrr}
j&0&1&2&3&4&5&6&7&8\\\hline
\operatorname{rank}H_j(W)&1&3&9&9&10&3&1&0&0\\
\operatorname{rank}H_j(Y)&1&4&12&18&18&12&4&1&0\\
\operatorname{rank}H_j(W,Y)&0&0&1&3&10&9&9&3&1
\end{array}                                                \tag{F.CK1}\label{maps:eq:F.CK1}
\]
The only torsion in these three rows is
$\operatorname{Tor}H_3(Y)=\mathbb Z/2$.

\end{theorem}
\begin{proof}
\begin{enumerate}[label=\arabic*.]
\item Let $\mathcal K$ be the periodic Kuhn triangulation of
$\mathbb R^3$.  Normalize every simplex by translating its
lexicographically first vertex to zero.  The objects are its
simplex orbits; a morphism $\tau\to\sigma$ is every strict
translated face inclusion $\tau+w\subset\sigma$.
Keep distinct translated inclusions as distinct morphisms.
Composition adds the translation vectors.  Enumerating all
nonempty faces of the six chains in the unit cube gives
$(1,7,12,6)$ simplex orbits in dimensions zero through three,
and the displayed morphism and flag counts.  The category
nerve is the barycentric cell model of the dual decomposition
of $\mathbb R^3/\mathbb Z^3$; the direction of the face order
is reversed when interpreted as dual cells.

\item For an object $\tau$ let $L_\tau\subset\mathbb Z^3$ be the
primitive span of its vertex differences.  Its primitivity
follows from unimodularity of every Kuhn simplex.  Choose by
Smith reduction a quotient map
$Q_\tau:\mathbb Z^3\to\mathbb Z^{3-\dim\tau}$ with kernel
$L_\tau$, and an integral section $R_\tau$.
Define two strict diagrams of tori on this category:
\[
\begin{array}{c|c|c}
 &\text{object at }\tau&\text{map for }\tau+w\subset\sigma\\\hline
\mathcal W&T^{3-\dim\tau}&Q_\sigma R_\tau\\
\mathcal Y&T^4&S_w=\begin{pmatrix}1&0\\ w&I_3\end{pmatrix}.
\end{array}                                                \tag{F.CK2}\label{maps:eq:F.CK2}
\]
The first maps compose because $Q_\sigma$ vanishes on
$L_\tau$, and the second compose because translations add.
The maps $[\,0\ Q_\tau\,]$ form a natural transformation
$\mathcal Y\to\mathcal W$: the only additional term is
$Q_\sigma w=0$, since $w$ is a vertex of $\sigma$.

\item Here these diagrams describe the actual local topology.
The toric moment quotient on the dual cell to $\tau$ has
angular torus $T^3/(L_\tau\otimes\mathbb R/L_\tau)$;
passing to a face collapses exactly the next angular
subtorus.  Barycentrically subdividing the dual cells gives
the homotopy colimit of $\mathcal W$ for the central fibre.
For its real oriented logarithmic boundary, retain the
argument $\theta_0$ of the base parameter and the three
angular coordinates $\theta$ of the toric chart.  Translating
the chart by $w$ changes them to
$(\theta_0,\theta+w\theta_0)$, which is precisely $S_w$.
Thus the logarithmic boundary is the homotopy colimit of
$\mathcal Y$, and forgetting the logarithmic phases on each
stratum is exactly the natural transformation just given.
This can also be read directly in the unimodular charts
$x_0\cdots x_k=t$: the moment coordinates form the dual
simplex, the four logarithmic phases remain, and the ordinary
stratum retains their quotient by the vertex differences.
The finite face identifications give \eqref{maps:eq:F.CK2}, so there are
no further gluing constants in this fixed marked cap.

\item The Kuhn map with its divisor log structures is proper,
log smooth and vertical after restricting to a small disc.
Achinger--Ogus, Theorem~4.1.1, identifies its topological germ
with the mapping cylinder of this logarithmic collapse
\cite[Theorem 4.1.1 and Lemmas 4.1.2--4.1.3, pp.~480--482]{ffAchingerOgus20}.  Consequently the preceding
homotopy-colimit diagrams and natural transformation model
the pair $(W,Y)$ itself.  Fix the marking by letting $b_i$
be the category loop with translation label $e_i$.  A
homotopy-colimit edge with map $S_{e_i}$ has the relation
$b_i^{-1}\theta_0b_i=\theta_0+a_i$; equivalently
$\theta_0b_i\theta_0^{-1}=b_i+a_i$.  This specifies the
required primitive cusp
$U=\begin{pmatrix}I_3&I_3\\0&I_3\end{pmatrix}$.

\item \par\medskip\noindent\textbf{Hypotheses and marking of the logarithmic comparison.}\quad 
\label{maps:mapsF:repair:log-cylinder-scope}
Give the small disc its standard divisorial log structure and
the Kuhn total space the divisorial structure of the reduced
central fibre.  In a unimodular chart the characteristic
monoid map is
\[
 \mathbb N\longrightarrow\mathbb N^{k+1},\qquad
 1\longmapsto(1,\ldots,1).
\]
Both monoids are fine and saturated.  The group cokernel is
free, and the underlying chart is the monomial semistable
map $x_0\cdots x_k=t$ with additional smooth coordinates;
this is log smooth.  The quotient monoid is a group, since
the negative of each coordinate generator equals the sum of
the other generators after the diagonal generator is killed.
The morphism is therefore vertical.  The periodic toric
construction is proper over the sufficiently small disc.

\item These are precisely the hypotheses of
\cite[Theorem~4.1.1, pp.~480--482]{ffAchingerOgus20}.
Its proof also notes that exactness is automatic over the
standard log disc.  It identifies the germ with the mapping
cylinder of the logarithmic collapse.  This homeomorphism
depends on a trivialization; Lemma~4.1.3 fixes that
trivialization on the logarithmic central fibre.  We keep
the resulting radial boundary identification and its stated
angular and tropical generators.  This suffices to model
the pair and its inclusion up to marked homotopy.  It does
not replace the supplied cusp marking by another integral
basis.

\item We specify finite integer chains, rather than only the
homology spectral sequence of these diagrams.  Replace every
torus by its ordinary normalized bar complex.  The affine
simplex map from the bar construction to the torus is natural
for all the integer homomorphisms in \eqref{maps:eq:F.CK2}, so this replacement
does not change either homotopy colimit or its collapse map.
Take the simplicial replacement over all strict flags of the
finite category.  Its vertical differential is the torus bar
boundary; its horizontal differential composes adjacent
arrows, with the first face applying the diagram map, and
has the sign $(-1)^q$ in vertical degree $q$.
Apply the first contraction and the finite series \eqref{maps:eq:F.CA5}
object by object.  There are no flags longer than three,
so every series has at most four terms.  The resulting
finite basis consists of a strict flag beginning at $\tau$
and an exterior basis vector of its torus.  Transfer the
natural transformation by $P_{\mathcal W}FI_{\mathcal Y}$.
This proves both the chain-map identity and completeness of
the collapse map, including all integral coherence corrections.

\item The raw cap chain ranks in degrees zero through six are
\[
 (26,199,688,1231,1226,648,144),
\]
and the raw boundary ranks in degrees zero through seven are
\[
 (26,274,1124,2420,3010,2186,864,144).
\]
The program enumerates every object, every translated
incidence and every flag, verifies composition, checks
$d^2=0$ on every generator, and checks the collapse-chain
identity on every generator.  Integral cancellation of a
unit pivot is a chain homotopy equivalence: clear its row by
integer column operations, then remove the resulting
contractible two-generator summand.  The adjacent differential
updates are forced by $d^2=0$.  After respectively 2063,
4988 and 7087 such cancellations for the cap, boundary and
collapse cone, Smith reduction gives \eqref{maps:eq:F.CK1}.  All matrices
and the remaining differentials are saved by
\texttt{fourfolds\_kuhn\_cap\_chains.py}; the reduced cap and
relative differentials vanish, and the boundary has one
nonzero Smith factor, equal to two in degree four to three.
This computes the integral groups without a filtration
extension problem.
\end{enumerate}
\end{proof}


The finite-pair and cusp-pair chain models enter the global
gluing cone of \cref{maps:src:F.2.35}. Its maps to the common
torus-bundle complement retain the specified peripheral word
and clutching translation.

The periodic toric quotient used for the cusp is the construction in
Yang's theorem \emph{Higher-dimensional periodic quotient}, and its
Kuhn-grid Euler number is the formula in his corollary
\emph{The standard grid family}
\cite[Theorem 8.7, pp.~57--59, and Corollary 8.8, p.~59, attached source build]{manYang26}.
For the present primitive rank-three grid, this gives Euler number
$3!=6$.  The full cap and boundary complexes, their integral inclusion,
and the torsion groups are calculated in the proof above.







\begin{theorem}[complete integral homology of a marked gluing]\label{maps:src:F.2.35}
\StatementRecord{maps:src:F.2.35}{proved}{Historical events: Claude Fable 5.2 (Anthropic), effort high, 16 November 2026: carried as background without fault found. GPT-6 Astra (OpenAI), ChatGPT Pro mode, effort Pro, S6\_09\_PAPERS, 17 November 2026: targeted proof check. Also checked in this writing. Pre-continuation writing events: GPT-6 Astra (OpenAI), ChatGPT Pro mode, effort Pro, S6\_10\_PAPERS, 18 November 2026 02:43:58.020152+00:00: computation or formal certificate; targeted proof check. Continuation metadata events: GPT-6 Astra (OpenAI), ChatGPT Pro mode, effort Pro, S6\_10\_PAPERS continuation, 19 November 2026: used and rechecked in this run; status-vocabulary reconciliation; metadata reconciliation only; no new mathematical check. \emph{Pre-continuation recorded proof scope.} GPT-6 Astra (OpenAI, ChatGPT Pro mode), effort Pro, 2026-11-17 UTC; GPT-6 Astra (OpenAI, ChatGPT Pro mode, effort Pro),17November2026,/root; Claude Fable 5.2 (Anthropic), effort high, 16 November 2026; also checked in this writing by GPT-6 Astra (OpenAI), ChatGPT Pro mode, effort Pro, 18 November 2026}{\cref{maps:src:F.2.31,maps:src:F.2.33,maps:src:F.2.34}; the displayed complement and peripheral maps}


For every fully marked F.2 datum there is a terminating exact
integer algorithm computing every integral homology group of
the resulting eight-manifold, including every torsion invariant
factor.  The input consists of the finite homology matrices
$M_1,\ldots,M_k$, free affine numerators $v_i/m_i$ of their
actual orders, the integral clutch vector $\kappa$, and the
chosen primitive cusp marking
\[
 M_0=(M_1\cdots M_k)^{-1},\qquad
 M_0P=P\begin{pmatrix}I_3&I_3\\0&I_3\end{pmatrix},
 \qquad P\in\mathrm{GL}(6,\mathbb Z).                       \tag{F.GC1}\label{maps:eq:F.GC1}
\]
Here $k=2$ or $3$ and the genuine orders are at least two;
order-one sites are first eliminated as in \ref{maps:src:F.2.31}.
If a marking is not supplied, the algorithm makes and records
one by Smith reduction.  It does not identify that choice
with an unspecified prescribed marking.
The complete chain prefix through degree nine has at most
4003 generators in any degree.  No spectral-sequence
extension, saturation convention, orientation assumption or
unbounded search enters the homology calculation.

\end{theorem}
\begin{proof}
\begin{enumerate}[label=\arabic*.]
\item We supply both global boundary maps which were left outstanding
after \ref{maps:src:F.2.34}.  Let $F_k$ be the free group with ordered
generators $s_1,\ldots,s_k$, and let $M(w)$ be the product
specified by a reduced word $w$.  The torus-bundle complement
over the punctured sphere retracts to the corresponding bundle
over a graph.  Its fundamental group and multiplication are
\[
 \Gamma=\mathbb Z^6\rtimes F_k,
 \qquad (x,w)(y,u)=(x+M(w)y,wu).                            \tag{F.GC2}\label{maps:eq:F.GC2}
\]
The fibre and base are aspherical and the base has no
two-cells.  Hence a chosen section identifies the complement
up to homotopy with $B\Gamma$.  These section and peripheral
choices are fixed by the input clutch convention below.

\item \emph{A finite complement complex with its full bar projection.}
Use reduced words for $F_k$.  The tree resolution of the free
group has $K_0=\mathbb Z[F_k]$ and
$K_1=\bigoplus_i\mathbb Z[F_k]e_i$, with
$de_i=s_i-1$, and no higher terms.  Its explicit augmented
contraction sends a word to its signed path in the Cayley
tree: for each positive letter $s_i$ add its preceding prefix
times $e_i$, and for a negative letter $s_i^{-1}$ subtract
the prefix including that letter times $e_i$.
This is the Fox derivative, and its boundary is $w-1$.
Together with the bar insertion contraction it gives, by
the recursion \eqref{maps:eq:F.CA3}, a strong contraction of the free-group
bar resolution to this tree resolution.  Every recursion
decreases bar degree and every word is finite.

\item For a bar simplex $[(x_1,w_1)|\cdots|(x_n,w_n)]$ of $\Gamma$,
put $y_j=M(w_1\cdots w_{j-1})x_j$.
The exact Alexander--Whitney chain map to the torus/free-group
double bar complex is
\[
 \sum_{p+q=n}(-1)^{pq}
 [w_1|\cdots|w_p]\otimes[y_{p+1}|\cdots|y_n].              \tag{F.GC3}\label{maps:eq:F.GC3}
\]
Terms with an identity bar entry are zero.  The horizontal
first face acts on torus vectors by
$\rho(w)=M(w)^{-1}$, and the horizontal boundary has sign
$(-1)^q$.  Notice that $\rho$ is a right action:
$\rho(wu)=\rho(u)\rho(w)$.  This order is essential when
the matrices do not commute.  Formula \eqref{maps:eq:F.GC3} follows from
the diagonal bisimplicial bar construction; alternatively,
its chain-map identity follows by pairing adjacent interior
faces, with the first face using
$M(w_1)^{-1}y_j=M(w_2\cdots w_{j-1})x_j$.
The usual shuffle inverse is finite in each degree, so the
map is a chain equivalence.

\item Apply the torus contraction and \eqref{maps:eq:F.CA5}, then the free-group
tree contraction and the same finite perturbation formula.
Higher terms decrease group bar degree at least two, whereas
the tree complex only has degrees zero and one.  Thus the
resulting complement complex is
\[
 C_n(E)=\bigwedge^n\mathbb Z^6\ \oplus\
 \bigoplus_{i=1}^k\bigwedge^{n-1}\mathbb Z^6,
 \qquad
 d(e_i\otimes\xi)=(-1)^{n-1}
   (\bigwedge^{n-1}M_i^{-1}-1)\xi.                         \tag{F.GC4}\label{maps:eq:F.GC4}
\]
Its differential alone has this short form; its projection
from group bars retains the finite higher perturbation terms.
Those projection terms are used for every boundary map.

\item \emph{Finite affine boundary maps.}
For the $i$th finite cap use the orbit groupoid of \ref{maps:src:F.2.33}.
Its one-edge homotopy colimit for the affine generator is
the boundary mapping torus.  Vertically contracting it gives
a bounded complex with
$C_n(Y_i)=\bigwedge^n\mathbb Z^6\oplus
\bigwedge^{n-1}\mathbb Z^6$ and differential
$(-1)^{n-1}(\bigwedge^{n-1}M_i-1)$ on the second summand.
Its inclusion into the cyclic bar construction, followed by
the cap projection in \eqref{maps:eq:F.CA5}, is the full map $Y_i\to B_i$.
In particular this one-edge model avoids an unnecessary
periodic resolution of the seven-dimensional boundary.

\item Here is an explicit group-bar map for its other inclusion.
Send a torus-groupoid arrow labelled $\lambda$ to the
translation $x^\lambda\in\Gamma$.  With the integral offsets
$z_1(a)$ in \eqref{maps:eq:F.CA4}, put
$\eta_a=s_i^{-1}x^{z_1(a)}$.
For an arrow $f:a\to b$, the exact identity
\[
 \eta_a\,\psi(gf)=\psi(f)\,\eta_b                        \tag{F.GC5}\label{maps:eq:F.GC5}
\]
is the natural-transformation identity: it is just
$s_i^{-1}x^{M_i\lambda+z_1(b)}
 =x^\lambda s_i^{-1}x^{z_1(b)}$.
For a $q$-simplex, the edge map is $(-1)^q$ times the
alternating prism which inserts $\eta$ at every position,
uses the original arrows before the insertion and the
transformed arrows afterwards.  Its boundary is precisely
the signed difference of the two end maps, by \eqref{maps:eq:F.GC5}.
Compose this prism with the finite vertical lift of the
one-edge complex, then with \eqref{maps:eq:F.GC3} and the complement
projection.  This proves the full chain-map identity.
The offsets record the actual integer affine lift: replacing
$v_i$ by $v_i+m_i\lambda$ changes the offsets and the map,
even if the affine map on the torus has not changed.
An optional unimodular local basis simply applies its matrix
to every translation in \eqref{maps:eq:F.GC5}.

\item We make the distinction between a bar cell and the geometric
positive root explicit.  Put
$a_i=z_1(0)=\lfloor v_i/m_i\rfloor$, coordinatewise, and
let $\tau_i$ be the positive one-edge bar cell closed using
the zero-labelled return arrow to object zero.  With the
first-face convention above its geometric loop is
\[
 \tau_i=g_i^{-1}x^{a_i},\qquad
 F_E(\tau_i)=-e_i+M_i^{-1}a_i\in C_1(E),\qquad
 g_i=x^{a_i}\tau_i^{-1}.                                  \tag{F.GC5a}\label{maps:eq:F.GC5a}
\]
Thus the bar cell has the opposite peripheral orientation,
with the displayed integer correction when the numerator
has a nonzero integer part.  Formula \eqref{maps:eq:F.GC5} sends its
group element to $s_i^{-1}x^{a_i}$; consequently the
geometric positive root in the last expression of \eqref{maps:eq:F.GC5a}
maps to $s_i$ exactly.  This is the stated conversion between
the boundary and complement chain conventions.  In a local basis
$Q$, replace $a_i$ in the group and complement formulas by
$Q\lfloor Q^{-1}v_i/m_i\rfloor$.
Writing $\mathcal N_i=\sum_{j=0}^{m_i-1}M_i^j$, the cyclic
degree-two cell has boundary
$m_i\tau_i+w_i-\mathcal N_i a_i$.
Indeed $\tau_i^{m_i}=x^{\mathcal N_i a_i-w_i}$ follows
directly from $g_i^{m_i}=x^{w_i}$.
This independently checks the affine power relation and
the degree-one prism sign, including negative numerators.

\item There is an optional exact geometric shortcut.  If a verified
local integral basis gives $M_i=\operatorname{diag}(1,D)$ and
$v_i=\epsilon e_0$, $\epsilon\in\{1,-1\}$, eliminate the
first fibre generator from the cap group using
$g^{m_i}=u^\epsilon$.  The cap is the mapping torus with
group $\mathbb Z^5\rtimes_D\mathbb Z$.
The fibre vector $(u,a)$ maps to $x^a g^{\epsilon m_i u}$,
and the root generator maps to $g$.  Substituting this
homomorphism for $\psi$ in the same prism gives the boundary
map to this bounded geometric cap model.  The shortcut is
used only after checking the displayed integral normal form;
the general periodic construction remains available for
every other input.

\item To check the peripheral convention directly, the universal
cover of $Y_i$ has the root generator $g_i$ with linear part
$M_i$ and affine displacement $v_i/m_i$.  Its integral
translations and $g_i$ map to the same fibre translations
and $s_i$ in \eqref{maps:eq:F.GC2}.  The full base-circle translation obeys
$g_i^{m_i}=x^{w_i}t$, where $w_i$ is the power vector of
\ref{maps:src:F.2.31}.  Filling the disc kills $t$, giving exactly
$s_i^{m_i}=x^{w_i}$.  The straight paths between the chosen
orbit points give \eqref{maps:eq:F.GC5}, so the prism is this marked map,
including numerators with nonzero integer parts.

\item \emph{The marked cusp boundary map.}
Write $P=[A\ B]$ in \eqref{maps:eq:F.GC1}, and set
\[
 w_0=(s_1\cdots s_k)^{-1},\qquad
 c=w_0x^{-\kappa}\in\Gamma.
                                                               \tag{F.GC6}\label{maps:eq:F.GC6}
\]
Send the three angular generators to the columns of $A$,
the three tropical generators to the columns of $B$, and
the base-phase generator to $c$.  Relation \eqref{maps:eq:F.GC1} gives
$c b_i c^{-1}=b_i+a_i$ and all the other nilpotent boundary
relations.  Filling kills the angular generators and $c$,
so \eqref{maps:eq:F.GC6} gives exactly
$s_1\cdots s_kx^\kappa=1$, in agreement with \ref{maps:src:F.2.31}.

\item This homomorphism has a direct chain implementation on the
full Kuhn diagram.  For an incidence with translation label
$w$, assign the group element $x^{Bw}$.  A vertical phase
$(n,a)\in\mathbb Z\oplus\mathbb Z^3$ maps to $c^nx^{Aa}$.
The incidence shear is $S_w(n,a)=(n,a+n w)$, and
\[
 (c^nx^{Aa})x^{Bw}=x^{Bw}(c^nx^{A(a+n w)}).                \tag{F.GC7}\label{maps:eq:F.GC7}
\]
For a flag of $p$ incidence arrows and a vertical bar
simplex of length $q$, sum over all $(p,q)$ shuffles.
Each incidence step contributes its assigned group element;
each vertical step contributes its phase transformed by
all preceding incidence shears.  Give the shuffle the sign
$(-1)^{\#\{\text{incidence steps before vertical steps}\}}$.
Interior face cancellations use exactly \eqref{maps:eq:F.GC7}, and adjacent
incidence faces use addition of their translation labels.
This is therefore a chain map from the double bar complex
to $C_*(B\Gamma)$.  Precompose with the finite Kuhn lift of
\ref{maps:src:F.2.34} and postcompose with \eqref{maps:eq:F.GC3} and the complement
projection.  The other cusp map is the full logarithmic
collapse already computed in \ref{maps:src:F.2.34}.
Both the cusp boundary and the complement are aspherical,
so the displayed marked homomorphism also verifies the
homotopy class of this geometric inclusion.

\item \emph{The final finite integer calculation.}
Take the mapping cone of the just constructed maps
\[
 \bigoplus_{i=0}^k C_*(Y_i)\longrightarrow
 C_*(E)\oplus\bigoplus_{i=0}^k C_*(B_i),
 \qquad y\longmapsto(f_E(y),-f_B(y)),                     \tag{F.GC8}\label{maps:eq:F.GC8}
\]
where $B_0=W$ is the Kuhn cap.  Gluing collars makes this
the homotopy pushout model of the actual closed manifold.
Consequently its cone computes integral homology directly;
there is no remaining filtration extension to choose.

\item For clarity the exact finite size bound can be read off
before running any matrix calculation.  Put
$t_n=\sum_{q=0}^{\min(n,6)}\binom6q$, with binomial
coefficients outside their range interpreted as zero.
Let $w_n$ and $y_n$ be the raw cap and boundary ranks of
\ref{maps:src:F.2.34}, zero outside their listed ranges.  Then the cone rank
in degree $n$ is
\[
 \binom6n+k\binom6{n-1}+k t_n+w_n+
 k\binom7{n-1}+y_{n-1}.                                 \tag{F.GC9}\label{maps:eq:F.GC9}
\]
For $0\leq n\leq9$ its maximum is 3890 when $k=2$ and
4003 when $k=3$.  Every bar comparison and perturbation
series used to produce these matrices terminates by its
degree filtration.  Word reductions and affine orbit
enumerations are finite, with no imposed size cutoff.
Computing integer kernels followed by Smith normal forms
now gives $H_0,\ldots,H_8$; higher groups vanish by dimension.
Unit cancellations are optional exact integral chain
equivalences which substantially reduce the matrices before
Smith reduction.  This proves the algorithm and the bound.
\end{enumerate}
\end{proof}


\emph{Verification and implementation.}
The implementation is in \path{code/fourfolds_global_chains.py}
and \path{code/fourfolds_total_homology.py}.
The generic JSON driver \path{code/fourfolds_marked_gluing.py}
retains the complete supplied marking and affine data, computes
the full group structure of \ref{maps:src:F.2.31}, and independently compares
its abelianization with the chain calculation of $H_1$.
The Fox contraction is checked on all 156 explicitly listed
free-group bar simplices through degree three; the exact
Alexander--Whitney and finite-projection identities are
checked on 156 specified semidirect-product simplices with
two noncommuting rank-two matrices.
Four finite boundary comparisons using the independent
seven-dimensional quotient model have acyclic integral
comparison cones.  The more efficient one-edge model is
checked in six stated low-rank cases, including negative,
nonfixed, and nontrivial integer-part numerators; both its
complement comparison and its full Thom cone pass.
The program \path{code/fourfolds_peripheral_signs.py} also
checks \eqref{maps:eq:F.GC5a} and the full cyclic degree-two relation in
seven cases, including both actual rank-six caps.  It
reconstructs the positive geometric root as a group element
and checks that its complement bar projection is $+e_i$.
For the Kuhn boundary, every one of its 10048 raw generators
passes the chain-map identity in all degrees zero through
seven.  Its comparison to the corresponding one-generator
complement has an integrally acyclic cone, eliminated by
5088 unit cancellations.  Thus the comparison also checks
the nonzero boundary order-two torsion.
The full eight-dimensional integration for the original
order-$(5,4)$ survivor has passed.  Every marked boundary
column and the complete Mayer--Vietoris differential are
checked.  After 7372 integral unit cancellations, its Smith
homology is $\mathbb Z$ in degrees zero and eight,
$\mathbb Z^4$ in degree four, and zero otherwise, exactly
as in the independent calculation \ref{maps:src:F.2.6}.
The corresponding three-point run with logarithmic order
$m=2$ also passes, after 7470 integral unit cancellations.
It gives $H_1=H_6=\mathbb Z/18$,
$H_2=H_5=(\mathbb Z/2)^5$, $H_3=(\mathbb Z/2)^{10}$,
$H_4=\mathbb Z^4\oplus(\mathbb Z/2)^{10}$,
$H_0=H_8=\mathbb Z$ and $H_7=0$, exactly as in the
independent theorem \ref{maps:src:F.2.32}.  It reuses only the identical
first-two finite boundary maps: their group bars involve only
the fibre and their own peripheral generator, so adjoining
a third generator changes neither their columns nor their
target row indices.  The third cap, the new cusp word,
the full cone and every Smith group are recomputed.

A third full integration uses
$A=I_3$, $B=\operatorname{diag}(0,1,0)$ in \eqref{maps:eq:F.IV2},
actual affine orders $4,4$, numerators $a_1,a_2$, and zero
clutching.  These two linear actions reverse orientation.
Its 7345 integral unit cancellations give free ranks
\[
                  (1,0,4,11,10,1,6,3,0).
\]
The nonzero torsion groups are
\[
\begin{aligned}
 \operatorname{Tor}H_1=\operatorname{Tor}H_5
       &=(\mathbb Z/2)^3\oplus\mathbb Z/4,&
 \operatorname{Tor}H_2&=(\mathbb Z/2)^2,\\
 \operatorname{Tor}H_3=\operatorname{Tor}H_7
       &=\mathbb Z/2,&
 \operatorname{Tor}H_4&=(\mathbb Z/2)^4\oplus(\mathbb Z/4)^5.
\end{aligned}
\]
The independent van Kampen $H_1$ agrees.  Its Euler number
is six, its mod-two Betti vector is
$(1,4,10,14,20,14,10,4,1)$, and its integral top groups
have the required nonorientable form.  Thus this integration
also tests the orientation-reversing part of the stated scope.

The \texttt{--fresh} option bypasses all saved map checkpoints.
The free-bar and semidirect tests, all four periodic-boundary
comparisons, all six one-edge comparison/Thom tests, all seven
peripheral-sign checks, the complete 10048-generator cusp
comparison, and the cached generic survivor computation have
also been independently rerun successfully.  The
reduced-chain, van Kampen $H_1$, Euler and mod-two duality
checks have additionally been independently replayed on
all three complete integration outputs.  The general
algorithm is proved by the explicit chain equivalences and
finite operations above; these finite execution sets are
reported with their exact scope.







\begin{theorem}[simultaneous integral conjugacy for every input]\label{maps:src:F.2.36}
\StatementRecord{maps:src:F.2.36}{proved}{Historical events: Claude Fable 5.2 (Anthropic), effort high, 16 November 2026: carried as background without fault found. GPT-6 Astra (OpenAI), ChatGPT Pro mode, effort Pro, S6\_09\_PAPERS, 17 November 2026: targeted proof check. Also checked in this writing. Pre-continuation writing events: GPT-6 Astra (OpenAI), ChatGPT Pro mode, effort Pro, S6\_10\_PAPERS, 18 November 2026 02:02:02.394029+00:00: targeted proof check. Continuation mathematical and source checks: GPT-6 Astra (OpenAI), ChatGPT Pro mode, effort Pro, S6\_10\_PAPERS continuation, 19 November 2026: targeted proof check. Continuation metadata events: GPT-6 Astra (OpenAI), ChatGPT Pro mode, effort Pro, S6\_10\_PAPERS continuation, 19 November 2026: used and rechecked in this run; status-vocabulary reconciliation; metadata reconciliation only; no new mathematical check. \emph{Pre-continuation recorded proof scope.} Claude Fable 5.2 (Anthropic), effort high, 16 November 2026; also checked in this writing by GPT-6 Astra (OpenAI), ChatGPT Pro mode, effort Pro, 18 November 2026}{the hypotheses, notation and definitions stated here; \cite[Algorithm 8.3, p.~21]{rootBHJ22}; \cite[Algorithm 8.3 and its correctness proof, p.~21]{rootBHJ22}; \cite[Theorem 3.4, p.~6; Propositions 8.1--8.2, p.~21]{rootBHJ22}}


There is an unconditional exact algorithm which decides simultaneous
$\mathrm{GL}_6(\mathbb Z)$ conjugacy of any two marked matrix tuples of
F.2 and returns a conjugating matrix when one exists. The algorithm
does not impose a bound on their entries, on the affine multiplicities,
or on the integral conjugator. More generally it applies to any integral
tuple containing a matrix $1+N$ with $N^2=0$ and $\operatorname{rank}N=3$;
the primitive-image assumption is unnecessary for this statement.

\end{theorem}


\begin{proof}
\begin{enumerate}[label=\arabic*.]
\item We use column vectors and seek $P$ satisfying $PX_i=Y_iP$ for every
marked index, including the cusp index. Put
\[
\begin{aligned}
 H&=\{P\in\operatorname{Mat}_6(\mathbb Z):PX_i=Y_iP\ (\forall i)\},\\
 \mathcal O&=\{C\in\operatorname{Mat}_6(\mathbb Z):CX_i=X_iC\ (\forall i)\},
 \qquad E=\mathbb Q\mathcal O.
\end{aligned}                                                    \tag{F.IC1}\label{maps:eq:F.IC1}
\]
Both lattices are saturated integer kernels of explicitly displayed
linear equations in thirty-six unknowns. Thus Hermite or Smith
reduction computes full bases, rather than sublattices of finite index.

\item We first decide rational isomorphism. For a basis $h_1,\ldots,h_s$ of
$H$, form the polynomial $f(c)=\det(\sum c_jh_j)$. If it is zero,
there is no rational isomorphism. Otherwise successively substitute
one of $0,1,\ldots,6$ for each variable while keeping the remaining
polynomial nonzero. At each stage one of these seven values works,
because the degree in that variable is at most six. The final matrix
$D\in H$ is invertible over $\mathbb Q$. This is a terminating exact
test, including the negative case, and not a search for a lucky
invertible intertwiner.

\item The full lattice $L=D^{-1}H\subset E$ is a right $\mathcal O$-module.
Let $J$ be the Jacobson radical of $E$, let $B=E/J$, and use bars
for images under the quotient map. We next verify the splitting
hypothesis needed for the principal-ideal algorithm. In a rational
$N$-adapted basis,
\[
 N=\begin{pmatrix}0&I_3\\0&0\end{pmatrix},\qquad
 E\subseteq
 \left\{\begin{pmatrix}A&T\\0&A\end{pmatrix}:A,T\in
                                      \operatorname{Mat}_3(\mathbb Q)\right\}.
                                                               \tag{F.IC2}\label{maps:eq:F.IC2}
\]
Projection to $A$ has square-zero kernel and image an algebra
$S\subseteq\operatorname{Mat}_3(\mathbb Q)$. A nilpotent ideal is
contained in the radical, so $E/J\cong S/J(S)$.

\item Choose a composition series of the faithful $S$-module $\mathbb Q^3$.
The kernel of its action on the associated graded module lowers this
filtration and is therefore nilpotent. Conversely the radical kills
every simple factor. Thus this kernel is exactly $J(S)$ and every
simple factor of $S/J(S)$ acts nontrivially on a composition factor.
For a simple algebra $\operatorname{Mat}_n(D)$, where $D$ has centre
$K$ and degree $d$ over that centre, every nonzero simple module has
$\mathbb Q$-dimension $n d^2[K:\mathbb Q]$. If $D$ were
noncommutative this would be at least four, which is impossible in
$\mathbb Q^3$. Consequently
\[
 B\cong\prod_j\operatorname{Mat}_{n_j}(K_j),\qquad
 n_j[K_j:\mathbb Q]\le3.                                    \tag{F.IC3}\label{maps:eq:F.IC3}
\]
Every noncommutative factor is therefore
$\operatorname{Mat}_2(\mathbb Q)$ or
$\operatorname{Mat}_3(\mathbb Q)$; the remaining factors are number
fields of degree at most three. This proves hypothesis (H) of
\cite[Algorithm 8.3, p.~21]{rootBHJ22} for every tuple under consideration.

\item For clarity this decomposition is effective without a split-algebra
oracle. Rational linear algebra computes $J$ as the radical of the
regular trace pairing and computes the quotient multiplication table.
Its simple factors are obtained from its centre. A commutative factor
is a number field. In a noncommutative factor of dimension $n^2$, we
enumerate integral combinations $e$ of a rational basis until the
left ideal $Be$ has dimension $n$. Such an $e$ exists because the
factor is $\operatorname{Mat}_n(\mathbb Q)$, so its rank-one matrices
are rational points; clearing denominators expresses one among the
enumerated combinations. Left multiplication on $Be$ then supplies
an explicit isomorphism to $\operatorname{Mat}_n(\mathbb Q)$.
The program verifies multiplication on every pair of basis elements.

\item The images $\overline{\mathcal O}$ and $\overline L$ are a full order
and a full right lattice in $B$. If the right order of
$\overline L$ differs from $\overline{\mathcal O}$, this lattice
cannot be principal over that order. Otherwise Algorithm 8.3 of
\cite[Algorithm 8.3 and its correctness proof, p.~21]{rootBHJ22}, in its right-ideal form, decides principality and
returns a generator $\overline\alpha$ when one exists. The cited
algorithm passes through a maximal order and its central conductor,
tests local freeness and principality over the maximal order, and
solves the remaining reduced-norm congruence in a finite quotient.
Its correctness proof uses its Theorem 4.5(b) and Theorem 7.9; its
Theorem 8.4 specifies the number-field, finite-field and integer
factorization subproblems. These have unconditional exact algorithms.
In the recorded implementation the class and unit groups of every
number-field component are certified with \texttt{GRH=false} before
the principal-ideal decision is used. The negative decision thus
does not presume the generalized Riemann hypothesis.

\item If $\overline L$ is not principal, $L$ cannot be principal. If it is
principal, solve the integer linear equations expressing
$\overline\alpha$ as a combination of the images of the chosen basis
of $L$. This produces a lift $\alpha\in L$ and an integral
intertwiner $P=D\alpha\in H$. We claim that a simultaneous integral
conjugacy exists if and only if $|\det P|=1$.

\item The sufficient direction is immediate. For the necessary direction,
suppose $P_0\in H$ is unimodular. Every $h\in H$ has
$P_0^{-1}h\in\mathcal O$, so
$L=\alpha_0\mathcal O$, where $\alpha_0=D^{-1}P_0$.
Write $\alpha=\alpha_0u$ with $u\in\mathcal O$.
Because both image elements generate $\overline L$,
$\overline u$ is a unit of $\overline{\mathcal O}$.
Lift its inverse to $v\in\mathcal O$. Then $uv=1+j$ with
$j\in J\cap\mathcal O$, and the finite geometric series for
$(1+j)^{-1}$ shows that $u^{-1}=v(1+j)^{-1}\in\mathcal O$.
Thus $u$ is an integral unimodular matrix and
$P=P_0u$ is unimodular. This also proves that any chosen lift of the
image generator gives the same yes-or-no answer. It is the
right-module form of \cite[Theorem 3.4, p.~6; Propositions 8.1--8.2, p.~21]{rootBHJ22},
with the needed determinant condition written out.

\item Every branch either proves nonexistence or returns a matrix whose
integrality, determinant and all intertwining equations are verified
exactly. This completes the decision procedure.

\end{enumerate}
\end{proof}




\begin{theorem}[a complete canonical-representative algorithm]\label{maps:src:F.2.37}
\StatementRecord{maps:src:F.2.37}{proved}{Historical events: Claude Fable 5.2 (Anthropic), effort high, 16 November 2026: carried as background without fault found. GPT-6 Astra (OpenAI), ChatGPT Pro mode, effort Pro, S6\_09\_PAPERS, 17 November 2026: targeted proof check. Also checked in this writing. Pre-continuation writing events: GPT-6 Astra (OpenAI), ChatGPT Pro mode, effort Pro, S6\_10\_PAPERS, 18 November 2026 02:02:02.394029+00:00: targeted proof check. Continuation mathematical and source checks: GPT-6 Astra (OpenAI), ChatGPT Pro mode, effort Pro, S6\_10\_PAPERS continuation, 19 November 2026: computation or formal certificate; targeted proof check; used and rechecked in this run. Continuation metadata events: GPT-6 Astra (OpenAI), ChatGPT Pro mode, effort Pro, S6\_10\_PAPERS continuation, 19 November 2026: used and rechecked in this run; status-vocabulary reconciliation; metadata reconciliation only; no new mathematical check. \emph{Continuation recorded proof scope.} Every step of the current proof has renewed prereading coverage: the canonical enumeration and actual conjugator argument by the Paper V reviewer, and the free affine lift and marked identity-slot maps by the root reviewer. The finite-order list is the already checked in-paper input. Newly authored exact programs verify the auxiliary computations and the actual word trace; the source and dependency scope was renewed. Former routes are preserved without retrospective prereading credit.}{\cref{maps:src:F.2.2}; \cref{maps:src:F.2.36}; \cref{maps:src:F.2.38}; \cref{maps:src:F.2.16}; the explicit affine lift and marked identity-slot maps displayed here}


Every marked linear tuple in F.2, for either $k=2$ or $k=3$,
has a computable canonical integral representative and a
returned simultaneous unimodular conjugator. Two inputs have
the same representative if and only if they are simultaneously
$\mathrm{GL}_6(\mathbb Z)$ conjugate. The algorithm uses a proved
finite bound depending on the input. No fixed global entry bound
can represent every class.

This is an effective classification, not a claim that its
infinite output has been printed in a finite table. The
ordered involution sector has the separate complete finite
normal form \ref{maps:src:F.2.38}, and all of its $56$ classes have been
computed. The general bounded fallback is potentially very
expensive; its finiteness is a theorem, not a runtime estimate.

\end{theorem}


\begin{proof}
\begin{enumerate}[label=\arabic*.]
\item We remove the existential free-lift condition by an explicit
construction. Let $d>1$ be the order of $M$ and choose a primitive
$u\in\ker(M-I)\cap\mathbb Z^6$. Such a vector is obtained from a
nonzero rational kernel vector by clearing denominators and dividing
by the greatest common divisor of its integer coordinates; both
operations retain the equation $(M-I)u=0$. Put
$g(x)=Mx+u/d^2$ and $S_d=\sum_{a=0}^{d-1}M^a$.
Then $g^j(x)=M^jx+ju/d^2$, $S_dM^j=S_d$ and $S_du=du$.
A fixed point of $g^j$ gives
$(M^j-I)x+ju/d^2=\lambda\in\mathbb Z^6$.
Applying $S_d$ gives $ju/d=S_d\lambda\in\mathbb Z^6$.
A Bezout row pairing to one with $u$ therefore gives $d\mid j$.
Writing $j=dr$, the original equation becomes $ru/d=\lambda$,
so the same row gives $d\mid r$. Thus only multiples of $d^2$
can have fixed points. Since $g^{d^2}$ is translation by the
integer vector $u$, its torus action is the identity. This proves
both its actual order $d^2$ and freeness of every nonidentity power.
For $M=I$, translation by $u/2$ for primitive $u$ has actual order
two and is free. Conversely, when $I-M$ is invertible over
$\mathbb R$, every affine map $Mx+a$ fixes the torus class of
$(I-M)^{-1}a$. The free-lift test is consequently exactly
$\ker(M-I)\ne0$.

\item The finite linear orders here are the already proved list of
\cref{maps:src:F.2.2}:
\[
 d\in\{1,2,3,4,5,6,8,10,12\}.
 \tag{F.CN1}\label{maps:eq:F.CN1}
\]
Taking these finitely many powers, checking the fixed kernel,
the product relation and the cusp Smith factors therefore gives
an exact admissibility test. The actual affine order is the
separate retained gluing datum; the construction in item1 gives
its exact relation to the chosen finite lift.

\item For an admissible full tuple $T=(M_1,\ldots,M_k,M_0)$ let
$H(T)$ be its maximum absolute entry. There is at most one identity
finite slot in a three-point tuple: two such identities would make
$M_0$ finite order, whereas $M_0=I+N$, $N^2=0$ and
$\operatorname{rank}N=3$ give $(I+N)^d=I+dN\ne I$ for every
positive integer $d$. At a fixed identity slot, delete that matrix
and reinsert $I$ at the same slot after classifying the remaining
ordered tuple. These two maps are inverse, and both commute with
$T\mapsto (QM_1Q^{-1},\ldots,QM_0Q^{-1})$ for the same
$Q\in\mathrm{GL}(6,\mathbb Z)$. They therefore retain the actual
returned simultaneous conjugator. The remaining invertible integral
matrices have height at least one, so inserting $I$ leaves height
unchanged. Every entry in the inserted block is constant; the first
unequal row-major entry of two tuples is therefore the same before
and after insertion. This proves preservation of the specified
lexicographic comparison as well. The actual affine order and
numerator at that slot are retained as gluing data and are transported
by this same $Q$. In the two-involution sector use the canonical
representative of \ref{maps:src:F.2.38}, with its own specified
convention. In all remaining sectors, order full integer tuples
first by $H$, then lexicographically by their row-major entries in
the displayed marked order, and take the first admissible member
of the simultaneous conjugacy class.

\item To compute the representative, enumerate heights
$1,\ldots,H(T)$ and all integer entries in the stated order.
There are at most
\[
 (2H(T)+1)^{36(k+1)}
 \tag{F.CN2}\label{maps:eq:F.CN2}
\]
integer lists in this box. Retain the invertible tuples satisfying
the ordered product relation and the stated maximal square-zero
primitive cusp condition, which are finite integer tests, and
apply the decision-and-witness algorithm of
\cref{maps:src:F.2.36} to each candidate and the supplied tuple.
Its general rank-three square-zero input permits precisely this
call. A positive answer transfers all finite-order, orientation
and admissibility conditions through its actual integral
conjugator. The supplied tuple itself occurs, with the identity
as a possible witness, so a positive answer is reached in
finitely many tests. The first one is the required representative.
Record the returned convention as
$P_T^{-1}T_iP_T=\operatorname{Can}(T)_i$.
The displayed candidate bound makes no assumption about entries
of $P_T$; the terminating witness construction is the stated
input from \cref{maps:src:F.2.36}.

\item Every simultaneous orbit has a least element in this order.
Its height is at most the height of each supplied member, so
\eqref{maps:eq:F.CN2} includes that global least element.
Equivalent inputs have exactly the same ordered orbit and
therefore the same representative. Conversely, if $T,T'$ have
the same representative, the actual matrix
$Q=P_TP_{T'}^{-1}$ satisfies $Q^{-1}T_iQ=T'_i$ for every slot.
The separate involution convention has the same property by
\cref{maps:src:F.2.38}, with the product of its actual returned
witnesses giving the corresponding map. Conjugation preserves
the ordered finite orders, so the involution sector is preserved
by those same maps. This proves the classification assertion
and its finite stopping rule.

\item Finally, \ref{maps:src:F.2.16} already gives an unbounded admissible integral
family. In its homology convention direct symbolic
multiplication gives
\[
             \operatorname{tr}(M_1M_2^3)
                       =-v^4+v^3-v^2+2.                 \tag{F.CN3}\label{maps:eq:F.CN3}
\]
This is unbounded even on the previously constructed free
affine subfamily $v\in20\mathbb Z$. It is invariant under
simultaneous conjugacy. A fixed finite entry box contains
only finitely many tuples and hence only finitely many such
traces; it cannot contain a representative of every class.
The input-dependent nature of \eqref{maps:eq:F.CN2} is therefore necessary
for this particular exhaustive procedure.

\end{enumerate}
\end{proof}




\begin{theorem}[the complete ordered involution stratum]\label{maps:src:F.2.38}
\StatementRecord{maps:src:F.2.38}{proved}{Historical events: Claude Fable 5.2 (Anthropic), effort high, 16 November 2026: carried as background without fault found. GPT-6 Astra (OpenAI), ChatGPT Pro mode, effort Pro, S6\_09\_PAPERS, 17 November 2026: targeted proof check. Also checked in this writing. Pre-continuation writing events: GPT-6 Astra (OpenAI), ChatGPT Pro mode, effort Pro, S6\_10\_PAPERS, 18 November 2026 02:43:58.020152+00:00: computation or formal certificate; targeted proof check. Continuation mathematical and source checks: GPT-6 Astra (OpenAI), ChatGPT Pro mode, effort Pro, S6\_10\_PAPERS continuation, 19 November 2026: computation or formal certificate; independent re-proof; used and rechecked in this run. Continuation metadata events: GPT-6 Astra (OpenAI), ChatGPT Pro mode, effort Pro, S6\_10\_PAPERS continuation, 19 November 2026: used and rechecked in this run; status-vocabulary reconciliation; metadata reconciliation only; no new mathematical check. \emph{Pre-continuation recorded proof scope.} GPT-6 Astra (OpenAI, ChatGPT Pro mode), effort Pro, 2026-11-17 UTC; GPT-6 Astra (OpenAI, ChatGPT Pro mode, effort Pro),17November2026,/root; Claude Fable 5.2 (Anthropic), effort high, 16 November 2026; also checked in this writing by GPT-6 Astra (OpenAI), ChatGPT Pro mode, effort Pro, 18 November 2026}{the hypotheses, notation and definitions stated here; \cref{maps:eq:F.IV1}; \cref{maps:eq:F.IV3}; \cref{maps:eq:F.IV4}; \cref{maps:eq:F.IV2}; \cref{maps:src:F.2.31}}


There are exactly 56 simultaneous $\mathrm{GL}_6(\mathbb Z)$
conjugacy classes of ordered tuples $(M_1,M_2,M_0)$ satisfying
the primitive maximal-cusp conditions of F.2 with both finite
linear orders equal to two.  An explicit integral normalizer
returns a canonical representative and a checked unimodular
conjugating matrix for every such input, without any bound on
its entries.  Every representative has free affine lifts of
actual order four at both finite points.  None can give a
simply connected filling, for any choices of free affine
orders or twists.

\end{theorem}


\begin{proof}
\begin{enumerate}[label=\arabic*.]
\item First choose an integral cusp basis, so that
\[
 M_0=I_6+N,\qquad N=\begin{pmatrix}0&I_3\\0&0\end{pmatrix}.
                                                               \tag{F.IV1}\label{maps:eq:F.IV1}
\]
This is effective by Smith reduction.  In fact choose vectors
$\beta_i$ whose images $\alpha_i=N\beta_i$ form a primitive
basis of $\operatorname{im}N=\ker N$.
The index-one hypothesis says that their classes form a basis
of the quotient, so $[\alpha_1,\alpha_2,\alpha_3,
\beta_1,\beta_2,\beta_3]$ is unimodular.

\item Since $M_2=M_1(I-N)$ and $M_1^2=M_2^2=I$, we have
$M_1NM_1=-N$.  Solving this block equation gives precisely
\[
 M_1=\begin{pmatrix}A&B\\0&-A\end{pmatrix},\qquad
 M_2=\begin{pmatrix}A&B-A\\0&-A\end{pmatrix},\qquad
 A^2=I_3,\quad AB=BA.                                     \tag{F.IV2}\label{maps:eq:F.IV2}
\]
Every integral conjugator preserving \eqref{maps:eq:F.IV1} is
$\begin{pmatrix}C&D\\0&C\end{pmatrix}$, with
$C\in\mathrm{GL}_3(\mathbb Z)$ and $D$ arbitrary integral.
Thus one first classifies $A$ integrally; after fixing it,
$B$ is considered modulo $AD+DA$ and modulo conjugation by
the integral centralizer of $A$.

\item We include the integral involution classification needed here.
For the lattice $L=\mathbb Z^3$, let $L_+$ and $L_-$ be the
primitive kernels of $A-I$ and $A+I$.  They satisfy
$2L\subseteq L_+\oplus L_-\subseteq L$.
The quotient $L/(L_+\oplus L_-)$ is a vector space over
$\mathbb F_2$.  Its projections into $L_+/2L_+$ and
$L_-/2L_-$ are injective: a glue vector with only one
nonzero projection would give a new integral vector in the
corresponding eigensublattice, contradicting its definition.
It is consequently the graph of an isomorphism between two
subspaces of equal dimension $r$.
Unimodular changes of bases in $L_\pm$ put this graph in
coordinate form.  For each glued pair $e_+,e_-$ the vectors
$(e_++e_-)/2$ and $(e_+-e_-)/2$ form a swap block.
The unglued vectors give trivial and sign blocks.  Hence
there is a unique type
\[
 A\cong I_p\oplus(-I_q)\oplus
 \begin{pmatrix}0&1\\1&0\end{pmatrix}^{\oplus r},
 \qquad p+q+2r=3.                                         \tag{F.IV3}\label{maps:eq:F.IV3}
\]
The dimensions of the two rational eigenspaces and the
index $[L:L_+\oplus L_-]=2^r$ prove uniqueness.
Here $r$ is zero or one, so this normalization only requires
integer kernels, one Smith reduction and, in the glued case,
one explicit half-sum pair whose integrality is checked.

\item For fixed \eqref{maps:eq:F.IV3}, compute the quotient of the commuting
matrices by the subgroup $\{AD+DA\}$.
On the trivial and sign diagonal blocks the norm is
respectively $2D$ and $-2D$.  The commuting off-diagonal
blocks between the trivial and sign summands are zero.
Every commuting block involving a swap summand is a norm.
For example a trivial-to-swap block is a row $(a,a)$ and
is obtained as $(d,e)+(e,d)$; a sign-to-swap block has the
same calculation with a difference.  For two swap summands,
the commuting block is
$\begin{pmatrix}a&b\\b&a\end{pmatrix}$, and applying the
norm to $\begin{pmatrix}d&e\\f&g\end{pmatrix}$ gives
$\begin{pmatrix}e+f&d+g\\d+g&e+f\end{pmatrix}$.
The reverse mixed blocks have the transposed calculations.
Therefore the full integral quotient is exactly
\[
 \operatorname{End}_{\mathbb Z[A]}(L)/\{AD+DA\}
 \cong \operatorname{Mat}_p(\mathbb F_2)
       \oplus\operatorname{Mat}_q(\mathbb F_2).             \tag{F.IV4}\label{maps:eq:F.IV4}
\]
No finite-index sublattice has been substituted for a norm
image; the displayed block calculations are surjective
integer formulas.

\item The centralizer action on \eqref{maps:eq:F.IV4} is precisely independent
similarity in the two matrix factors.  Indeed, in a product
of commuting block matrices, every contribution factoring
through a swap block is a norm.  The remaining diagonal
blocks are just the reductions of the trivial and sign
blocks of the centralizer matrix.  If that matrix is an
integral unit, its inverse shows that both reductions are
invertible.  Conversely block-diagonal unimodular changes
on the trivial and sign summands induce every element of
$\mathrm{GL}_p(\mathbb F_2)\times
\mathrm{GL}_q(\mathbb F_2)$; elementary matrices give
unimodular lifts.  There is no additional identification
between the two factors, since an intertwiner of a trivial
and a sign lattice is zero over $\mathbb Z$.

\item It remains a genuinely finite classification.  Enumerate
all $2^{n^2}$ matrices of size $n=0,1,2,3$ over $\mathbb F_2$
and conjugate by every invertible matrix of that size.
Their respective numbers of similarity classes are
$1,2,6,14$.  Taking the smallest bit-coded representative
in each orbit gives a fixed canonical choice.
The six integral involution types and their counts are
\[
\begin{array}{c|rrrrrr}
 (p,q,r)&(0,3,0)&(1,2,0)&(2,1,0)&(3,0,0)&(0,1,1)&(1,0,1)\\\hline
 \text{classes}&14&12&12&14&2&2
\end{array}
\]
and their sum is 56.  A representative is \eqref{maps:eq:F.IV2} with
$B$ equal to the chosen $0$--$1$ representatives in its
trivial and sign blocks, and zero in every other block.

\item This also gives an effective normalizer for an arbitrary
input, not merely a list of representatives.  Normalize its
cusp and then its involution by the preceding Smith and
glue constructions.  Use the finite-field orbit lookup to
normalize its two reductions.  In dimensions at most three
the binary lift of an invertible binary matrix is already
unimodular: its determinant is odd and has absolute value
at most two.  The latter bound follows by adjoining a row
and column of ones and applying the $4\times4$ Hadamard
bound.  Finally solve the full integral linear equation
$AD+DA=B-B_{\rm can}$ by Smith reduction.  The preceding
norm calculation proves solvability.  Conjugating by
$\begin{pmatrix}I&D\\0&I\end{pmatrix}$ subtracts this norm
from the upper-right block.  The product of these explicitly
recorded conjugators is integral unimodular and carries all
three marked matrices to their canonical representative.

\item Both matrices in \eqref{maps:eq:F.IV2} have a three-dimensional fixed
space.  Choose a primitive integral fixed vector $v_i$.
Then $g_i(x)=M_ix+v_i/4$ has actual order four:
$g_i^2$ is the nontrivial translation $v_i/2$.
This order-two subgroup is free, and every nontrivial
subgroup of a cyclic group of order four contains it, so
the full action is free.  For any other finite affine lift
the actual order is still even, since its linear part has
order two.  The two-point orbifold quotient of \ref{maps:src:F.2.31} is
therefore the nontrivial cyclic group of order
$\gcd(m_1,m_2)$.  Simple connectivity is impossible for
every choice of twists and actual orders.

\end{enumerate}
\end{proof}




\begin{theorem}[the full effective classification and all marked gluings]\label{maps:src:F.2.39}
\StatementRecord{maps:src:F.2.39}{proved}{Historical events: Claude Fable 5.2 (Anthropic), effort high, 16 November 2026: carried as background without fault found. GPT-6 Astra (OpenAI), ChatGPT Pro mode, effort Pro, S6\_09\_PAPERS, 17 November 2026: targeted proof check. Also checked in this writing. Pre-continuation writing events: GPT-6 Astra (OpenAI), ChatGPT Pro mode, effort Pro, S6\_10\_PAPERS, 18 November 2026 02:02:02.394029+00:00: targeted proof check. Continuation mathematical and source checks: GPT-6 Astra (OpenAI), ChatGPT Pro mode, effort Pro, S6\_10\_PAPERS continuation, 19 November 2026: targeted proof check; used and rechecked in this run. Continuation metadata events: GPT-6 Astra (OpenAI), ChatGPT Pro mode, effort Pro, S6\_10\_PAPERS continuation, 19 November 2026: used and rechecked in this run; status-vocabulary reconciliation; metadata reconciliation only; no new mathematical check. \emph{Pre-continuation recorded proof scope.} GPT-6 Astra (OpenAI, ChatGPT Pro mode), effort Pro, 2026-11-17 UTC; GPT-6 Astra (OpenAI, ChatGPT Pro mode, effort Pro),17November2026,/root; Claude Fable 5.2 (Anthropic), effort high, 16 November 2026; also checked in this writing by GPT-6 Astra (OpenAI), ChatGPT Pro mode, effort Pro, 18 November 2026}{the hypotheses, notation and definitions stated here; \cref{maps:src:F.2.37}; \cref{maps:src:F.2.38}; \cref{maps:src:F.2.31}; \cref{maps:src:F.2.35}; \cref{maps:src:F.2.3}; \cref{maps:src:F.2.30}; \cref{maps:src:F.2.1}; \cref{maps:src:F.2.2}; \cref{maps:src:F.2.36}; \cref{maps:eq:F.CN3}; \cref{maps:src:F.2.33}; \cref{maps:src:F.2.7}}


The linear classification and all-twist topological computation
asked for in F.2 have the following complete effective answer.
For $k=2$ and $k=3$, every admissible ordered integral tuple
has the canonical representative of \ref{maps:src:F.2.37}, with an explicit
simultaneous unimodular conjugator.  Equality of these
representatives is necessary and sufficient for simultaneous
integral conjugacy.  There is no fixed global height cutoff;
the finite stopping bound for an input of height $H$ is
$(2H+1)^{36(k+1)}$ candidates before the proved rejection
filters.  The separate involution normal form has exactly
56 classes and is computed completely in \ref{maps:src:F.2.38}.

For every choice of actual free affine orders and integer
numerators, integral clutch vector, and specified primitive
Kuhn cusp marking, \ref{maps:src:F.2.31} and \ref{maps:src:F.2.35} compute the full
fundamental group and every integral homology group of the
gluing.  This is the fully marked correction to \ref{maps:src:F.2.3};
peripheral matrices are not substituted for a choice of
filling.  The algorithms have no cutoff on affine orders,
twist entries or lattice indices.  The complete list of data
with $\pi_1=1$ and $H_2=H_3=0$ is the constructive
parametrization of \ref{maps:src:F.2.30}, with its exact lattice enumeration
and necessary-and-sufficient finite tests.  Each retained
filling has the homology and oriented smooth type stated
there.

\end{theorem}


\begin{proof}
\begin{enumerate}[label=\arabic*.]
\item The input is finite, so the prime-subgroup test in \ref{maps:src:F.2.31}
decides whether every claimed affine cyclic action is free
and has the stated actual order.  The integral cusp test and
its adapted basis are \ref{maps:src:F.2.1}.  The complete linear admissibility
filter and the canonical representative, including the proved
finite stopping bound, are \ref{maps:src:F.2.2} and \ref{maps:src:F.2.36}--38.  The
characteristic-polynomial filters used by that algorithm are
conjugacy invariants, so none can remove a possible
representative.  The unbounded trace family \eqref{maps:eq:F.CN3} proves
that no fixed finite box could replace the input-dependent
bound.

\item We spell out how the linear classification is used with a
specified gluing.  If $R$ is its returned conjugator and
$P$ is the supplied cusp marking, change the data by
\[
 M_i'=RM_iR^{-1},\qquad v_i'=Rv_i,\qquad
 \kappa'=R\kappa,\qquad P'=RP.                           \tag{F.SY1}\label{maps:eq:F.SY1}
\]
The actual orders are unchanged.  The coordinate map
$x\mapsto Rx$ on the torus intertwines the finite quotient
charts, since it sends $M_ix+v_i/m_i$ to
$M_i'x'+v_i'/m_i$.  At the cusp, its coordinates in the
marked Kuhn model are unchanged:
$(P')^{-1}x'=P^{-1}x$.  Hence it intertwines every toric
chart, translated face identification and collapse map of
that fixed filling.  It also carries the clutch path to
the path with vector $\kappa'$.  These maps agree on the
collars and give a diffeomorphism of the two fully marked
gluings.  Thus homology can be computed in the original
coordinates, while the canonical data and their group
presentation are also returned.  Independently replacing
$P'$ by another Smith marking is not this comparison.

\item An actual order-one site is ordinary, even if its integer
numerator has been recorded.  Delete it by the explicit
comparison of F.2.31: if it is the $i$th site, replace
$\kappa$ by
$\kappa+(M_{i+1}\cdots M_k)^{-1}v_i$ before deletion.
The cusp matrix and its supplied marking are unchanged.
There cannot be an order-one site in an admissible two-point
tuple: its other finite matrix would be the inverse of a
nontrivial unipotent matrix.  Nor can a three-point tuple
have two such sites.  Thus this reduction leaves precisely
the genuine two- or three-point inputs of \ref{maps:src:F.2.35} and loses
no case from the specified input class.

\item For each remaining input, the exact unsaturated normal
closure, cyclic or triangle-group extension and exceptional
spherical Euler relation of \ref{maps:src:F.2.31} give the full fundamental
group.  The finite cap, complement and marked Kuhn chain
maps of \ref{maps:src:F.2.33}--35 give a finite integer Mayer--Vietoris cone
through degree nine.  Integer kernels and Smith reduction
therefore give all $H_0,\ldots,H_8$ with their complete
torsion factors.  This is a calculation for arbitrary
supplied twists, rather than a restriction to the affine
orders of one displayed example.  Its independent
abelianization check compares the resulting $H_1$ with the
van Kampen answer.

\item Finally \ref{maps:src:F.2.7}--30 prove both directions of the survivor
classification.  Genuine three-point orbifold quotients
exclude simple connectivity.  The two-point filters reduce
every survivor to the minimal order-$(5,4)$ spectrum; the
rational normal form, finite stable-lattice and extension
algorithms include every integral representation above it.
The universal affine normalization includes every twist,
and the six exact integral tests decide retention.
Conversely every retained datum satisfies all three
topological conditions by those tests.  \ref{maps:src:F.2.30} computes
its common integral homology and applies the almost-complex
and smooth-classification arguments to identify its oriented
type.  There is no residual untested rational-height bound
in this parametrization.  The canonical-representative
algorithm can remove duplicates from it if desired, but
no finite global enumeration is asserted or needed for
these complete per-input decisions.

\end{enumerate}
\end{proof}




\begin{theorem}\label{night:P1B.1}
\StatementRecord{night:P1B.1}{proved}{Historical events: GPT-6 Astra (OpenAI), ChatGPT Pro mode, effort Pro, S6\_06\_BLIND, 16--18 November 2026: targeted proof check. Also checked in this writing. Pre-continuation writing events: GPT-6 Astra (OpenAI), ChatGPT Pro mode, effort Pro, S6\_10\_PAPERS, 18 November 2026 02:02:02.394029+00:00: targeted proof check. Continuation mathematical and source checks: GPT-6 Astra (OpenAI), ChatGPT Pro mode, effort Pro, S6\_10\_PAPERS continuation, 19 November 2026: used and rechecked in this run. Continuation metadata events: GPT-6 Astra (OpenAI), ChatGPT Pro mode, effort Pro, S6\_10\_PAPERS continuation, 19 November 2026: used and rechecked in this run; status-vocabulary reconciliation; metadata reconciliation only; no new mathematical check. }{Definitions and the displayed constructions}
Let $N$ be an endomorphism of $\Z^6$ with $N^2=0$, rank three,
and primitive image. There is a unimodular basis in which
\[
 I+N=\begin{pmatrix}I_3&I_3\\0&I_3\end{pmatrix}.
\]
Its centralizer in $\operatorname{GL}_6(\Z)$ consists exactly of
$\left(\begin{smallmatrix}P&Q\\0&P\end{smallmatrix}\right)$,
with $P\in\operatorname{GL}_3(\Z)$ and $Q\in\operatorname{Mat}_3(\Z)$.
\end{theorem}
\begin{proof}
\begin{enumerate}[label=\arabic*.]
\item The kernel of an integral homomorphism to a free group is
primitive. Both $\ker N$ and $\operatorname{im}N$ have rank three,
and the latter is contained in the former. Since the image is
primitive in $\Z^6$, their finite quotient is zero. Choose a basis
$e_1,e_2,e_3$ of this common group and choose $f_i$ with $Nf_i=e_i$.
The induced map $\Z^6/\ker N\to\operatorname{im}N$ is an
isomorphism, so the images of the $f_i$ form a quotient basis.
Thus $(e_1,e_2,e_3,f_1,f_2,f_3)$ is a unimodular basis and gives
the displayed matrix. This is also an algorithm: a Smith
decomposition gives the primitive kernel and solves the three
integral equations $Nf_i=e_i$.
Writing a potential centralizer element as
$\left(\begin{smallmatrix}P&Q\\R&S\end{smallmatrix}\right)$,
the equality $CN=NC$ gives $R=0$ and $S=P$.
Its determinant is $(\det P)^2$, so integral invertibility
is equivalent to $P\in\operatorname{GL}_3(\Z)$.

\end{enumerate}
\end{proof}




\begin{theorem}\label{night:P1B.2}
\StatementRecord{night:P1B.2}{proved}{Historical events: GPT-6 Astra (OpenAI), ChatGPT Pro mode, effort Pro, S6\_06\_BLIND, 16--18 November 2026: targeted proof check. Also checked in this writing. Pre-continuation writing events: GPT-6 Astra (OpenAI), ChatGPT Pro mode, effort Pro, S6\_10\_PAPERS, 18 November 2026 02:02:02.394029+00:00: targeted proof check. Continuation mathematical and source checks: GPT-6 Astra (OpenAI), ChatGPT Pro mode, effort Pro, S6\_10\_PAPERS continuation, 19 November 2026: used and rechecked in this run. Continuation metadata events: GPT-6 Astra (OpenAI), ChatGPT Pro mode, effort Pro, S6\_10\_PAPERS continuation, 19 November 2026: used and rechecked in this run; status-vocabulary reconciliation; metadata reconciliation only; no new mathematical check. }{Definitions and the displayed constructions}
A finite-order integral linear action on $\R^6$ admitting a
free finite affine torus lift has characteristic polynomial
among exactly thirty-eight cyclotomic products. Their linear
orders are
\[
 1,2,3,4,5,6,8,10,12.
\]
Exactly nineteen of these polynomials admit a commuting complex
structure on $\R^6$. These are a preliminary spectral list;
no assertion about global gluings is inferred from the list.
\end{theorem}
\begin{proof}
\begin{enumerate}[label=\arabic*.]
\item The linear action has eigenvalue one. Otherwise $I-A$ is
invertible over $\R$, and the affine lift $x\mapsto Ax+b$
has a fixed point, obtained by solving $(I-A)x=b$ in $\R^6$
and projecting to the torus. A finite-order integral matrix is
semisimple and its characteristic polynomial is a product of
cyclotomic polynomials with nonnegative integer exponents.
The one-eigenvalue leaves at most five dimensions for the other
factors. A cyclotomic degree is even for every index greater
than two. The possible degrees are therefore one, two and four,
with indices respectively $1,2$; $3,4,6$; and $5,8,10,12$.
To verify that no index has been omitted, the formula
$\varphi(n)=\prod_{p^a\parallel n}p^{a-1}(p-1)$ with
$\varphi(n)\le4$ restricts $p$ to $2,3,5$, the powers to
$2^a$ with $a\le3$, $3^a$ with $a\le1$, and $5^a$ with
$a\le1$, and gives exactly those indices after multiplication.

\item The executed program \texttt{code/part1/spectral\_filter.py}
enumerates every exponent vector satisfying total degree six
and a positive exponent at $\Phi_1$, with bounds supplied by
these degrees. It writes all thirty-eight exponent vectors,
polynomials and least common multiples of their active indices.
Companion matrices in block diagonal form realize every listed
polynomial and order. A translation by a nonzero rational
fraction in a trivial eigenline makes a finite affine action
free: choose its denominator to be a multiple of the linear
order, and every nonidentity power has nonintegral translation
in that eigenline.

\item A commuting real complex structure requires even multiplicity
of each real eigenvalue $1$ and $-1$. Conversely, on these real
eigenspaces pair basis vectors into oriented complex lines;
on every conjugate nonreal eigenpair use its rotation complex
structure. This constructs a commuting complex structure.
The program checks these two parity conditions and obtains
nineteen products. Since the remaining cyclotomic factors have
even degrees and the total degree is six, the two real
multiplicities have the same parity. Thus the nineteen
complex-compatible products are exactly the orientation-preserving
ones in this preliminary list.

\end{enumerate}
\end{proof}


\section{Source obstructions and the punctured-base period decision}



\begin{proposition}[the previously open period decision]\label{maps:src:F.3.3}
\StatementRecord{maps:src:F.3.3}{proved}{Historical events: Claude Fable 5.2 (Anthropic), effort high, 16 November 2026: carried as background without fault found. GPT-6 Astra (OpenAI), ChatGPT Pro mode, effort Pro, S6\_09\_PAPERS, 17 November 2026: targeted proof check. Also checked in this writing. Pre-continuation writing events: GPT-6 Astra (OpenAI), ChatGPT Pro mode, effort Pro, S6\_10\_PAPERS, 18 November 2026 02:43:58.020152+00:00: targeted proof check. Continuation mathematical and source checks: GPT-6 Astra (OpenAI), ChatGPT Pro mode, effort Pro, S6\_10\_PAPERS continuation, 19 November 2026: used and rechecked in this run. Continuation metadata events: GPT-6 Astra (OpenAI), ChatGPT Pro mode, effort Pro, S6\_10\_PAPERS continuation, 19 November 2026: used and rechecked in this run; status-vocabulary reconciliation; metadata reconciliation only; no new mathematical check. \emph{Pre-continuation recorded proof scope.} Claude Fable 5.2 (Anthropic), effort high, 16 November 2026; also checked in this writing by GPT-6 Astra (OpenAI), ChatGPT Pro mode, effort Pro, 18 November 2026}{\cref{maps:src:F.3.4,maps:src:F.3.7,maps:src:F.3.8,maps:src:F.3.9,maps:src:F.3.10,maps:src:F.3.11}}


Theorem \ref{maps:src:F.3.11} now gives a positive answer for every topological
survivor of F.2: there is a holomorphic period family on the
punctured sphere, with nondegenerate imaginary part in every
local normalized period chart and a bounded holomorphic
correction at its primitive maximal cusp.  The proof is universal
and does not require that the remaining integral classification
be finite.  The proposed polarized abelian-surface source is
excluded by Theorem \ref{maps:src:F.3.4}; the universal construction instead
uses the explicit unpolarized LSAP domain and the relative
section principle in \ref{maps:src:F.3.8}--\ref{maps:src:F.3.11}.  Simultaneous good reduction
at every finite puncture is a stronger compactification condition
than the literal F.3 statement and is not inferred from this
argument.  It is already proved for the explicit compact
fourfold of F.4.

\end{proposition}
\begin{proof}
\begin{enumerate}[label=\arabic*.]
\item Every retained topological datum has determinant-one transport
and the marked primitive maximal cusp. Apply the explicit cusp graph
and relative holomorphic section construction of
\cref{maps:src:F.3.10,maps:src:F.3.11}. Its target plane is transverse
to its conjugate, and its difference from the specified cusp model is
bounded in the displayed cusp coordinate. Quotienting the resulting
holomorphic rank-three vector bundle by its integral period lattice
gives the required punctured-base family, as in \cref{maps:src:F.3.9}.
\item The displayed map from a polarized source would supply the
semisimple exterior-square local system excluded in
\cref{maps:src:F.3.4}. The actual section construction instead uses
the nine spray directions and the relative section maps proved in
\cref{maps:src:F.3.8,maps:src:F.3.10}. Its finite punctures remain
punctures; the theorem asserts precisely the bounded cusp conclusion
just obtained. The separately constructed explicit family has the
finite root-cap extensions proved in \cref{maps:src:F.3.7}.
\end{enumerate}
\end{proof}





\begin{theorem}[exclusion of the abelian-surface source]\label{maps:src:F.3.4}
\StatementRecord{maps:src:F.3.4}{proved}{Historical events: Claude Fable 5.2 (Anthropic), effort high, 16 November 2026: targeted proof check. GPT-6 Astra (OpenAI), ChatGPT Pro mode, effort Pro, S6\_09\_PAPERS, 17 November 2026: targeted proof check. Also checked in this writing. Pre-continuation writing events: GPT-6 Astra (OpenAI), ChatGPT Pro mode, effort Pro, S6\_10\_PAPERS, 18 November 2026 02:02:02.394029+00:00: targeted proof check. Continuation mathematical and source checks: GPT-6 Astra (OpenAI), ChatGPT Pro mode, effort Pro, S6\_10\_PAPERS continuation, 19 November 2026: independent re-proof; used and rechecked in this run. Continuation metadata events: GPT-6 Astra (OpenAI), ChatGPT Pro mode, effort Pro, S6\_10\_PAPERS continuation, 19 November 2026: used and rechecked in this run; status-vocabulary reconciliation; metadata reconciliation only; no new mathematical check. \emph{Pre-continuation recorded proof scope.} GPT-6 Astra (OpenAI, ChatGPT Pro mode), effort Pro, 2026-11-17; Claude Fable 5.2 (Anthropic), effort high, 2026-11-16; also checked in this writing by GPT-6 Astra (OpenAI), ChatGPT Pro mode, effort Pro, 18 November 2026}{the hypotheses, notation and definitions stated here; \cite[Definition 4.2.4 and Proposition 4.2.5, p.~44; Theorem 4.2.6, p.~45]{ffDeligne}; \cref{maps:src:F.3.1}}


The Engel construction proposed in F.3, starting with a polarized
algebraic family of abelian surfaces with a toric-rank-two cusp,
taking a fibrewise degree-zero line bundle, and quotienting its
punctured bundle by a lifted translation, cannot have $b_1=b_2=0$
with the stipulated rank-three square-zero Kuhn cusp.
The two sections need not be proportional.

\end{theorem}


\begin{proof}
\begin{enumerate}[label=\arabic*.]
\item Let $W$ be the rank-four rational homology local system of the
abelian surfaces.  It is semisimple by
\cite[Definition 4.2.4 and Proposition 4.2.5, p.~44; Theorem 4.2.6, p.~45]{ffDeligne}.
Its toric-rank-two cusp has
$N_W^2=0$ and $\ker N_W=\operatorname{im}N_W$ of dimension two.
Semisimplicity splits $W=W^\Gamma\oplus W'$.
The image of $N_W$ lies in $W'$ whereas $W^\Gamma$ lies in its
kernel.  The equality of image and kernel forces $W^\Gamma=0$.
The coinvariants vanish as well.

\item The rank-six rational homology system $T$ after the two Engel
operations has a natural filtration
\[
 0\subset L\subset E\subset T,\qquad
 L\simeq\mathbb Q,\quad E/L\simeq W,\quad T/E\simeq\mathbb Q.
\]
The lower line is the circle in the punctured line bundle and
the upper quotient is the lifted-translation generator.
No polarization on $T$ is assumed.
The triviality of $T/E$ gives $\operatorname{im}N_T\subset E$.
Maximal square-zero cusp monodromy gives
$\ker N_T=\operatorname{im}N_T$; thus every globally fixed vector
lies in $E$.  Since $W^\Gamma=0$, the fixed vectors of $E$ are
exactly $L$.  Hence $\dim T^\Gamma=1$.
The dual filtration, whose middle term is $W^*$, gives
$\dim(T^*)^\Gamma=1$ by the identical argument.

\item Consequently the two ranks $a,b$ in Theorem \ref{maps:src:F.3.1} are both one.
Its Euler-parity calculation applies without requiring self-duality
of $T$, and contradicts $b_1=b_2=0$.
This excludes the proposed algebraic abelian-surface source and,
more generally, every such filtration with semisimple middle term.
It does not exclude an arbitrary unpolarized torus period family.

\end{enumerate}
\end{proof}




\begin{theorem}[the unpolarized period domain has LSAP]\label{maps:src:F.3.8}
\StatementRecord{maps:src:F.3.8}{proved}{Historical events: Claude Fable 5.2 (Anthropic), effort high, 16 November 2026: carried as background without fault found. GPT-6 Astra (OpenAI), ChatGPT Pro mode, effort Pro, S6\_09\_PAPERS, 17 November 2026: targeted proof check. Also checked in this writing. Pre-continuation writing events: GPT-6 Astra (OpenAI), ChatGPT Pro mode, effort Pro, S6\_10\_PAPERS, 18 November 2026 02:02:02.394029+00:00: targeted proof check. Continuation mathematical and source checks: GPT-6 Astra (OpenAI), ChatGPT Pro mode, effort Pro, S6\_10\_PAPERS continuation, 19 November 2026: computation or formal certificate; independent re-proof; targeted proof check; used and rechecked in this run. Continuation metadata events: GPT-6 Astra (OpenAI), ChatGPT Pro mode, effort Pro, S6\_10\_PAPERS continuation, 19 November 2026: used and rechecked in this run; status-vocabulary reconciliation; metadata reconciliation only; no new mathematical check. \emph{Continuation recorded proof scope.} Complete written derivation before reading the target proof: compactified projective lines, a uniform normal neighbourhood, nine spanning directions and the exact graph-chart morphism. The fresh program \path{code/continuation_root/D3_tube_fresh.py} checks both the independently chosen directions and the printed matrices. The Oka-1 spanning-tube and LSAP inputs were opened at their stated locators.}{the hypotheses, notation and definitions stated here; \cite[Definition 2.1, p.~4]{IVVOkaOneV5}; \cite[Proposition 7.13 and preceding paragraph, pp.~23--24]{IVVOkaOneV5}; \cref{maps:eq:F.O3}; \cref{maps:eq:F.O4}}


Let $D_3$ be the component of
\[
 \{F\in\operatorname{Gr}(3,\mathbb C^6):F\cap\overline F=0\}
\]
containing the standard complex structure, so that
$D_3\cong GL^+(6,\mathbb R)/GL(3,\mathbb C)$.
At every point, $D_3$
is dominated by a spanning tube of nine complex lines.  In particular,
it has the local spray approximation property (LSAP) and is Oka--1.
These assertions concern the unpolarized domain.

\end{theorem}


\begin{proof}
\begin{enumerate}[label=\arabic*.]
\item Use the graph chart $F_Z=\operatorname{im}\binom Z{I_3}$.
The open set
\[
 \Omega=\{Z\in M_3(\mathbb C):\det(\operatorname{Im}Z)>0\}
 \tag{F.O1}\label{maps:eq:F.O1}
\]
lies in $D_3$ and contains $iI_3$.  For a real column $w$, write
$K_w x=w\times x$, put $q_w=w^tw$, and define
\[
 A_w=q_wI_3-ww^t+K_w.
 \tag{F.O2}\label{maps:eq:F.O2}
\]
Take the six vectors
\[
 e_1,\ e_2,\ e_3,\ e_1+e_2,\ e_1+e_3,\ e_2+e_3
\]
and the corresponding six matrices $A_w$, followed by the three
matrices $K_{e_1},K_{e_2},K_{e_3}$.  Call these nine matrices
$A_1,\ldots,A_9$.  Their symmetric parts are a basis of the
six-dimensional space of real symmetric matrices: the first three
give all diagonal matrices and the last three add the three distinct
off-diagonal entries.  The final three skew matrices form a basis
of the skew matrices.  Thus the nine directions are linearly
independent over $\mathbb R$ and also over $\mathbb C$.
For reference, the determinant of their nine row-major coordinate
columns is $-16$.

\item Each $A_w$ is real normal, is zero on $\mathbb Rw$, and has
eigenvalues $q_w\pm i\sqrt{q_w}$ on $w^\perp\otimes\mathbb C$.
For every real $s$ its two transverse squared singular values are
\[
 (1+q_ws)^2+q_ws^2
 =q_w(q_w+1)\left(s+\frac1{q_w+1}\right)^2
       +\frac1{q_w+1}\ge\frac13,
 \tag{F.O3}\label{maps:eq:F.O3}
\]
since $q_w\in\{1,2\}$; the remaining squared singular value is one.
For a pure skew direction $K_{e_j}$ all singular values of
$I_3+sK_{e_j}$ are at least one.  Their determinants are positive
at $s=0$ and never vanish, so they are positive for all real $s$.

\item Let $\Lambda=\bigcup_{j=1}^9(iI_3+\mathbb C A_j)$ and let
$B_{1/4}$ be the open Frobenius ball of radius $1/4$ in
$M_3(\mathbb C)$.  Then
\[
 T=\Lambda+B_{1/4}\subset\Omega.
 \tag{F.O4}\label{maps:eq:F.O4}
\]
Indeed, the imaginary part of $iI_3+tA_j+E$ is
$I_3+(\operatorname{Im}t)A_j+\operatorname{Im}E$.
Its unperturbed least singular value is at least $1/\sqrt3$,
whereas $\|\operatorname{Im}E\|_{\rm op}<1/4$.
It therefore stays invertible throughout the straight perturbation
from $E=0$, preserving the positive determinant.
The set $T$ is open and connected and is a union of affine
translates of the spanning tree $\Lambda$.  Thus it is a spanning
tube in the precise sense of \cite[Definition 2.1, p.~4]{IVVOkaOneV5},
not merely a neighborhood of one entire line.  Its inclusion in
$D_3$ is a local biholomorphism at $iI_3$.

\item The real group $GL^+(6,\mathbb R)$ acts holomorphically and
transitively on $D_3$.  Transporting this tube gives the same
domination at every point.  The observation preceding
\cite[Proposition 7.13 and preceding paragraph, pp.~23--24]{IVVOkaOneV5} gives LSAP from
this domination, and its part (b), or Theorem 2.2 of that paper,
gives Oka--1.  The exact program
\texttt{code/fourfolds\_oka\_tubes.py} checks all nine integer
matrices, their linear independence and the symbolic identity
\eqref{maps:eq:F.O3}; the displayed uniform bound proves \eqref{maps:eq:F.O4} for all complex
parameters and is not a finite positivity test.

\end{enumerate}
\end{proof}




\begin{theorem}[holomorphic periods on every open base]\label{maps:src:F.3.9}
\StatementRecord{maps:src:F.3.9}{proved}{Historical events: Claude Fable 5.2 (Anthropic), effort high, 16 November 2026: carried as background without fault found. GPT-6 Astra (OpenAI), ChatGPT Pro mode, effort Pro, S6\_09\_PAPERS, 17 November 2026: targeted proof check. Also checked in this writing. Pre-continuation writing events: GPT-6 Astra (OpenAI), ChatGPT Pro mode, effort Pro, S6\_10\_PAPERS, 18 November 2026 02:43:58.020152+00:00: targeted proof check. Continuation mathematical and source checks: GPT-6 Astra (OpenAI), ChatGPT Pro mode, effort Pro, S6\_10\_PAPERS continuation, 19 November 2026: independent re-proof; used and rechecked in this run. Continuation metadata events: GPT-6 Astra (OpenAI), ChatGPT Pro mode, effort Pro, S6\_10\_PAPERS continuation, 19 November 2026: used and rechecked in this run; status-vocabulary reconciliation; metadata reconciliation only; no new mathematical check. \emph{Pre-continuation recorded proof scope.} GPT-6 Astra (OpenAI, ChatGPT Pro mode), effort Pro, 2026-11-17; Claude Fable 5.2 (Anthropic), effort high, 16 November 2026; also checked in this writing by GPT-6 Astra (OpenAI), ChatGPT Pro mode, effort Pro, 18 November 2026}{the hypotheses, notation and definitions stated here; \cref{maps:src:F.3.8}; \cite[Proposition 7.13(d), p.~24]{IVVOkaOneV5}}


Let $B$ be an open connected Riemann surface and
$\rho:\pi_1(B)\to GL^+(6,\mathbb R)$ any representation.
Its associated flat $D_3$-bundle admits a holomorphic section.
More generally, every continuous section which is holomorphic
near a compact Runge set can be deformed through sections to a
holomorphic section, with arbitrary uniform approximation there
and any prescribed finite-order interpolation at finitely many
points of that set.  For an integral representation, this gives
a holomorphic family of complex three-tori with the prescribed
lattice local system and a zero section over $B$.

\end{theorem}


\begin{proof}
\begin{enumerate}[label=\arabic*.]
\item The flat bundle
$\mathscr D=\widetilde B\times_\rho D_3\to B$ is a holomorphic
fibre bundle: its transition maps are constant real linear
transformations, acting holomorphically on the Grassmannian.
Its fibre has LSAP by Theorem \ref{maps:src:F.3.8}, so the bundle map is an
Oka--1 map by \cite[Proposition 7.13(d), p.~24]{IVVOkaOneV5}.
The lifting property in Definition 7.7 of that paper, applied
to the identity map of $B$, gives precisely the stated section
deformation and approximation, including the finite jets.
This uses the section theorem, rather than an inference from
the Oka--1 property for maps to the fibre alone.

\item There is a continuous section to start with.  An open Riemann
surface has the homotopy type of a one-dimensional CW complex;
a bundle with connected fibre admits a section over its vertices
and edges, and the homotopy lifting property transports this
section to $B$.  Connectedness holds for the chosen component
$D_3$.  Taking the Runge set empty therefore proves existence.

\item For an integral representation let $\mathbb V_\mathbb Z$ be its
lattice local system and let $\mathcal F\subset
\mathbb V_\mathbb Z\otimes\mathcal O_B$ be the resulting
holomorphic rank-three subbundle.  Real transversality says that
projection identifies the real vector bundle with the underlying
real bundle of
$\mathcal Q=(\mathbb V_\mathbb Z\otimes\mathcal O_B)/\mathcal F$.
Hence the image of $\mathbb V_\mathbb Z$ is a full lattice in
every fibre of $\mathcal Q$.  Locally it is uniformly discrete
and acts by holomorphic translations.  The quotient
$\mathcal Q/\mathbb V_\mathbb Z\to B$ is consequently a smooth
proper holomorphic family of compact complex three-tori,
with its evident zero section and the required marking.

\end{enumerate}
\end{proof}




\begin{lemma}[retaining isolated extension points]\label{maps:src:F.3.10}
\StatementRecord{maps:src:F.3.10}{proved}{Historical events: Claude Fable 5.2 (Anthropic), effort high, 16 November 2026: carried as background without fault found. GPT-6 Astra (OpenAI), ChatGPT Pro mode, effort Pro, S6\_09\_PAPERS, 17 November 2026: targeted proof check. Also checked in this writing. Pre-continuation writing events: GPT-6 Astra (OpenAI), ChatGPT Pro mode, effort Pro, S6\_10\_PAPERS, 18 November 2026 02:02:02.394029+00:00: targeted proof check. Continuation mathematical and source checks: GPT-6 Astra (OpenAI), ChatGPT Pro mode, effort Pro, S6\_10\_PAPERS continuation, 19 November 2026: targeted proof check. Continuation metadata events: GPT-6 Astra (OpenAI), ChatGPT Pro mode, effort Pro, S6\_10\_PAPERS continuation, 19 November 2026: used and rechecked in this run; status-vocabulary reconciliation; metadata reconciliation only; no new mathematical check. \emph{Pre-continuation recorded proof scope.} GPT-6 Astra (OpenAI, ChatGPT Pro mode), effort Pro, 2026-11-17; Claude Fable 5.2 (Anthropic), effort high, 16 November 2026; also checked in this writing by GPT-6 Astra (OpenAI), ChatGPT Pro mode, effort Pro, 18 November 2026}{the hypotheses, notation and definitions stated here; \cref{maps:eq:F.O5}; \cite[Theorem 1.2, preprint p.~3]{ffForstnericPrezelj02}; \cite[Theorem 2.13, p.~9]{IVVOkaOneV5}}


Let $R$ be an open Riemann surface, $V\to R$ a holomorphic vector
bundle, and $\mathscr U\subset\operatorname{Gr}_r(V)$ an open
subset whose projection onto $R$ is surjective.  Let $S\subset R$
be finite.  Suppose that over $R\setminus S$ this projection is
a holomorphic fibre bundle with an LSAP fibre.  Let $K\subset R$
be a compact Runge set containing $S$ in its interior.
Every continuous section of $\mathscr U$ which is holomorphic
near $K$ can be deformed through sections to a holomorphic
section, with arbitrary uniform approximation on $K$ and
interpolation of any given finite jets at finitely many points
of its interior.  No fibre-bundle or Serre-fibration assertion
is required at the points of $S$.

\end{lemma}


\begin{proof}
\begin{enumerate}[label=\arabic*.]
\item Here are the details of the relative version of the LSAP argument
in \cite[Lemma 5.1, Theorem 6.1 and Proposition 7.13(d)]
{IVVOkaOneV5}.  They also specify why isolated extension
points cause no new analytic step.  Enlarge $K$ slightly inside
the domain of holomorphicity, retaining the Runge property.
Choose a strongly subharmonic Morse exhaustion of $R$ starting
with this compact set.  All subsequent boundary arcs, attached
discs and critical one-handles can and will be chosen disjoint
from $S$.  On them $\mathscr U$ is an ordinary holomorphic
bundle with LSAP fibre.

\item The global spray on the old compact set can be taken to be
vertical even though the fibres at $S$ differ.  Indeed, let $f$
be the current holomorphic section on a Stein neighborhood $W$
of that set.  Put $A$ equal to the finite interpolation set,
including $S$, and choose an ideal sheaf $\mathcal I$ vanishing
to one more than the prescribed jet orders at its points.
Cartan's theorem A supplies finitely many holomorphic sections
$T_1,\ldots,T_N$ of
$\operatorname{End}(V)|_W\otimes\mathcal I$ which generate
$\operatorname{End}(V)$ on the compact attaching arcs, after
shrinking $W$ if needed.  The formula
\[
 F(z,a)=\exp\left(\sum_{j=1}^N a_jT_j(z)\right)f(z)
 \tag{F.O5}\label{maps:eq:F.O5}
\]
defines a holomorphic spray of sections of the Grassmannian
bundle.  Its differential in $a$ surjects onto the vertical
tangent space at every attaching point, since
$\operatorname{End}(V_z)\to
\operatorname{Hom}(f(z),V_z/f(z))$ is onto.
For $a$ sufficiently small it has range in $\mathscr U$ on
a neighborhood of the old compact set, by openness and
compactness.  Every parameter value preserves all the chosen
jets.  Domination at $A$ is neither asserted nor used.

\item For a noncritical extension, subdivide the new annular regions
into small discs lying in bundle trivializations over
$R\setminus S$.  Restrict \eqref{maps:eq:F.O5} to a thin attaching disc.
LSAP of the fibre extends this spray across the attached disc,
approximating the old spray as closely as desired on the
intersection.  The spray gluing used in the proof of
\cite[Lemma 5.1, equations (5.6)--(5.8)]
{IVVOkaOneV5} now applies verbatim to vertical sprays:
the inverse-function step is needed only on the attaching
intersection, where \eqref{maps:eq:F.O5} is dominating.  It produces
reparametrizations of the two sprays and a glued section.
The old part remains in $\mathscr U$ because it is evaluated
in \eqref{maps:eq:F.O5}; the new part remains in $\mathscr U$ because the
LSAP spray takes values in the bundle fibre.  The fixed jets
survive the reparametrizations, being independent of the
spray parameters on the old side.

\item For a critical extension, first extend over the attached real
arc using the given continuous section and the homotopy already
constructed on the old compact set.  Apply Mergelyan
approximation with the same finite jets to a neighborhood of
the union of the old set and this arc, and then apply the
noncritical step.  To see explicitly that this approximation
is by sections, trivialize $V$ holomorphically on the open
Riemann surface: its underlying complex bundle is trivial over
a graph spine, and the frame bundle therefore has a continuous
section.  Its fibre $GL(\operatorname{rank}V,\mathbb C)$ has
the dominating spray $(g,A)\mapsto g\exp A$, so
\cite[Theorem 1.2, preprint p.~3]{ffForstnericPrezelj02} gives a holomorphic
frame.  The section is then a map to the
fixed Grassmannian, to which
\cite[Theorem 2.13, p.~9]{IVVOkaOneV5} applies.  The base
coordinate is kept unchanged.  Close approximation on the
compact arc and old set, followed by shrinking its
neighborhood, preserves membership in the open subset
$\mathscr U$.  This is the critical case of Theorem 6.1 of
that paper, with no approximation performed at a singular
point of the projection: here the projection is a submersion
everywhere, as $\mathscr U$ is open in a Grassmannian bundle.

\item Choose successive errors summably small on the preceding
compact sets and smaller than their distances to the complement
of $\mathscr U$.  The sections converge uniformly on every
compact set to a holomorphic section with range in
$\mathscr U$.  The finite jets persist.  The close homotopies
on successive stages concatenate in the usual exhaustion
construction.  All new attachment regions lie outside $S$;
the old section remains holomorphic on a full neighborhood of
each point of $S$ at every stage.  This proves the lemma.

\end{enumerate}
\end{proof}




\begin{theorem}[the bounded-cusp period decision]\label{maps:src:F.3.11}
\StatementRecord{maps:src:F.3.11}{proved}{Historical events: Claude Fable 5.2 (Anthropic), effort high, 16 November 2026: carried as background without fault found. GPT-6 Astra (OpenAI), ChatGPT Pro mode, effort Pro, S6\_09\_PAPERS, 17 November 2026: targeted proof check. Also checked in this writing. Pre-continuation writing events: GPT-6 Astra (OpenAI), ChatGPT Pro mode, effort Pro, S6\_10\_PAPERS, 18 November 2026 02:02:02.394029+00:00: targeted proof check. Continuation mathematical and source checks: GPT-6 Astra (OpenAI), ChatGPT Pro mode, effort Pro, S6\_10\_PAPERS continuation, 19 November 2026: targeted proof check; used and rechecked in this run. Continuation metadata events: GPT-6 Astra (OpenAI), ChatGPT Pro mode, effort Pro, S6\_10\_PAPERS continuation, 19 November 2026: used and rechecked in this run; status-vocabulary reconciliation; metadata reconciliation only; no new mathematical check. \emph{Pre-continuation recorded proof scope.} GPT-6 Astra (OpenAI, ChatGPT Pro mode), effort Pro, 2026-11-17; Claude Fable 5.2 (Anthropic), effort high, 16 November 2026; also checked in this writing by GPT-6 Astra (OpenAI), ChatGPT Pro mode, effort Pro, 18 November 2026}{the hypotheses, notation and definitions stated here; \cref{maps:src:F.2.1}; \cref{maps:eq:F.O7}; \cref{maps:src:F.3.8}; \cref{maps:src:F.3.10}; \cref{maps:src:F.2.7}; \cref{maps:src:F.2.15}}


Let
$B^*=\mathbb P^1\setminus\{p_0,p_1,\ldots,p_k\}$,
where $k\ge1$, and let
$\rho:\pi_1(B^*)\to GL^+(6,\mathbb R)$ be a representation.
Assume that the monodromy $M_0$ at $p_0$ satisfies
$(M_0-I)^2=0$ and $\operatorname{rank}(M_0-I)=3$.
Then there is a holomorphic period family with this monodromy
and nondegenerate real transversality everywhere on $B^*$,
equivalently a nondegenerate imaginary part in each local
normalized period chart.
In a marking with $M_0=\left(\begin{smallmatrix}I&I\\0&I\end{smallmatrix}\right)$
its cusp period is
\[
 Z(t)=\frac{\log t}{2\pi i}I_3+C(t),
 \qquad C(t)\text{ holomorphic at }t=0.
 \tag{F.O6}\label{maps:eq:F.O6}
\]
Any prescribed finite jet of $C$ at zero may be imposed.
For an integral primitive cusp the marking can be integral.
Consequently every topological survivor of F.2 admits the
period family requested literally in F.3.  No good-reduction
condition at the other, still punctured, points is included
in this assertion.

\end{theorem}


\begin{proof}
\begin{enumerate}[label=\arabic*.]
\item Write $N=M_0-I$ and choose a real basis with
$N=\left(\begin{smallmatrix}0&I\\0&0\end{smallmatrix}\right)$.
Its integral version is Theorem \ref{maps:src:F.2.1}.  Choose the component
of the period domain selected by the positive logarithmic
model in \eqref{maps:eq:F.O6}.  This choice matters only when an orientation
has been fixed independently of the marking; every element
of $GL^+(6,\mathbb R)$ preserves either component.

\item Retain the cusp as a point of the \emph{open} surface
\[
 B=\mathbb P^1\setminus\{p_1,\ldots,p_k\}.
\]
The flat complex vector bundle on $B^*$ has a holomorphic
extension $V_D$ across $p_0$: in a local coordinate $t$ use
the single-valued frame obtained by the change
$\exp(-sN)$, where $s=(\log t)/(2\pi i)$.
Changing the logarithm by $2\pi i$ is cancelled by $M_0$.
This specifies the extension directly, without a growth
assumption on a period section.

\item Over $B^*$ let $\mathscr U$ be the flat bundle of real-transverse
planes in the chosen component.  At $p_0$ adjoin only the
planes transverse to $\ker N_\mathbb C$, an affine space of
matrices $C\in M_3(\mathbb C)$.  This gives an \emph{open}
subset $\mathscr U\subset\operatorname{Gr}_3(V_D)$.
Indeed, in the chart $F_C=\operatorname{im}\binom C{I_3}$
the actual flat period is
$\exp(sN)F_C=F_{sI_3+C}$.  For $t\ne0$ membership is exactly
\[
 \det\bigl(y(t)I_3+\operatorname{Im}C\bigr)>0,
 \qquad y(t)=-\frac{\log|t|}{2\pi}.
 \tag{F.O7}\label{maps:eq:F.O7}
\]
For every fixed $C_0$, and uniformly for $C$ in a sufficiently
small bounded neighborhood of $C_0$, the matrix in \eqref{maps:eq:F.O7} is
invertible with positive determinant when $|t|$ is sufficiently
small.  This follows, for example, from
$y(t)>\|\operatorname{Im}C\|_{\rm op}$ and a straight-line
homotopy to $y(t)I_3$.  Thus a full neighborhood of
$(0,C_0)$ lies in the adjoined set, proving openness.
The projection is consequently a holomorphic submersion,
although its fibre at $p_0$ is $\mathbb C^9$ rather than $D_3$.

\item Start with a polynomial $C_0(t)$ realizing the desired finite
jet.  Formula \eqref{maps:eq:F.O7} makes it a holomorphic section on a small
closed cusp disc $K$, which is Runge in $B$.
This section extends continuously over $B$.  To justify the
relative extension, $D_3$ is connected and simply connected:
it retracts onto $SO(6)/U(3)$, and
$\pi_1(U(3))\to\pi_1(SO(6))$ is onto, since a unitary rotation
in one complex coordinate is a generator of the real
two-plane rotation loop.  A cellular decomposition of the
surface outside the cusp disc now extends its boundary
section over edges by connectedness and over two-cells by
simple connectivity.  The flat bundle is trivial over each
cell.  This argument uses no fibration property at $p_0$.

\item Over $B\setminus\{p_0\}$ the fibre is the LSAP manifold
$D_3$ from Theorem \ref{maps:src:F.3.8}.  Lemma \ref{maps:src:F.3.10}, with $S=\{p_0\}$,
produces a holomorphic section of $\mathscr U$ over all of
$B$, with the chosen jet at $p_0$.  Since its value at $p_0$
lies in the affine chart just adjoined, it has a holomorphic
matrix $C(t)$ there.  Undoing the canonical frame gives
exactly \eqref{maps:eq:F.O6}, with a bounded holomorphic correction.
Everywhere else membership in $\mathscr U$ is precisely
nondegeneracy of the complex-torus periods.
More explicitly, choose three members of a real lattice basis
whose projected periods are complex-linearly independent, and
normalize them to $I_3$.  The remaining three periods form a
matrix $Z$, and real transversality is exactly
$\det\operatorname{Im}Z\ne0$.  This normalization persists
locally.  A single Schubert chart is not imposed globally,
and none is needed to define a marked torus period family.

\item For F.2 survivors, Theorem \ref{maps:src:F.2.7} excludes $k=3$ and
Proposition \ref{maps:src:F.2.15} proves orientation preservation of the
finite matrices for $k=2$.  The cusp determinant is also one.
Thus all survivors satisfy the hypotheses.  This is a
universal existence decision and does not require completion
of their simultaneous integral classification.

\end{enumerate}
\end{proof}

